\documentclass[a4paper, 11pt, twoside]{article}
\usepackage{wrapfig}
% \usepackage{abstract}
\usepackage{graphicx} % Required for inserting images
\usepackage[a4paper, left=1.15in, right=1.15in, top=1.5in, bottom=1.2in]{geometry}
\usepackage{float}
\usepackage{array}
\usepackage{booktabs}
\usepackage[T1]{fontenc}
\usepackage{amsmath}
\usepackage{tabularx}
\usepackage{graphicx} % Required for inserting images
\usepackage{tikz, xcolor}
\usepackage{hyperref}
\usepackage{cite}
\usepackage{pdfpages}
\usepackage{pdflscape}
\usepackage{afterpage}
\usepackage{textpos}
\usepackage{fancyhdr}
\usepackage{caption}
\usepackage{lmodern,textcomp}
\usepackage[catalan]{babel}
\usepackage{lipsum}
\usepackage{newfloat}
\usepackage{listings}
\usepackage{placeins}
\usepackage[ruled]{algorithm2e}
\usepackage{tcolorbox}
\usepackage[round]{natbib}
\usepackage[toc,page]{appendix}
\usepackage{forest}
\usepackage{pdflscape}
\usepackage{graphicx}
\usepackage{geometry}
\usepackage{quoting,xparse}
\usepackage{svg}



% \NewDocumentCommand{\bywhom}{m}{% the Bourbaki trick
%   {\nobreak\hfill\penalty50\hskip1em\null\nobreak
%    \hfill\mbox{\normalfont(#1)}%
%    \parfillskip=0pt \finalhyphendemerits=0 \par}%
% }

% AIXÒ SERVEIX PER A MODIFICAR UN PROBLEMA AMB ELS CARACTERS --> " + PARAULA 

% \NewDocumentEnvironment{pquotation}{m}
%   {\begin{quoting}[
%      indentfirst=true,
%      leftmargin=\parindent,
%      rightmargin=\parindent]\itshape}
%   {\bywhom{#1}
  
%   \end{quoting}}

% AIXÒ PER INSERIR FRAGMENTS DE CODI

\lstnewenvironment{code}[1][]%
{
   \noindent
   \minipage{\linewidth}
   \vspace{0.5\baselineskip}
   \lstset{
     basicstyle=\ttfamily\footnotesize,
     frame=single,
     breaklines=true,    % Permitir que las líneas se rompan automáticamente
     breakatwhitespace=true, % Permitir la ruptura en los espacios en blanco
     #1
   }
}
{\endminipage}

\lstset{
    language=Java,
    basicstyle=\ttfamily\small,  % Define el tipo de letra y tamaño del código
    keywordstyle=\color{blue},   % Define el color de las palabras clave de Java
    stringstyle=\color{red},     % Define el color de las cadenas de texto
    commentstyle=\color{green},  % Define el color de los comentarios
    morecomment=[s][\color{gray}]{/**}{*/},  % Estilo para comentarios de bloque
    breaklines=true,  % Configura el quiebre automático de línea
    numbers=left,  % Posición de los números de línea
    numberstyle=\tiny\color{gray},  % Estilo para los números de línea
    showspaces=false,  % No mostrar espacios con underscores
    showstringspaces=false,  % No mostrar espacios dentro de strings
    frame=single,  % Agregar un marco alrededor del código
    captionpos=b,  % b = bottom, t = top
}

\definecolor{eclipseStrings}{RGB}{42,0.0,255}
\definecolor{eclipseKeywords}{RGB}{127,0,85}
\colorlet{numb}{magenta!60!black}

\lstdefinelanguage{json}{
    basicstyle=\ttfamily\small,
    showstringspaces=false,
    breaklines=true,
    frame=single,
    literate=
     *{0}{{{\color{numb}0}}}{1}
      {1}{{{\color{numb}1}}}{1}
      {2}{{{\color{numb}2}}}{1}
      {3}{{{\color{numb}3}}}{1}
      {4}{{{\color{numb}4}}}{1}
      {5}{{{\color{numb}5}}}{1}
      {6}{{{\color{numb}6}}}{1}
      {7}{{{\color{numb}7}}}{1}
      {8}{{{\color{numb}8}}}{1}
      {9}{{{\color{numb}9}}}{1}
      {:}{{{\color{punct}{:}}}}{1}
      {,}{{{\color{punct}{,}}}}{1}
      {\{}{{{\color{delim}{\{}}}}{1}
      {\}}{{{\color{delim}{\}}}}}{1}
      {[}{{{\color{delim}{[}}}}{1}
      {]}{{{\color{delim}{]}}}}{1},
    morestring=[b]"
}

\lstdefinelanguage{JavaScript}{
  morekeywords=[1]{break, continue, delete, else, for, function, if, in,
    new, return, this, typeof, var, void, while, with},
  % Literals, primitive types, and reference types.
  morekeywords=[2]{false, null, true, boolean, number, undefined,
    Array, Boolean, Date, Math, Number, String, Object},
  % Built-ins.
  morekeywords=[3]{eval, parseInt, parseFloat, escape, unescape},
  sensitive,
  morecomment=[s]{/*}{*/},
  morecomment=[l]//,
  morecomment=[s]{/**}{*/}, % JavaDoc style comments
  morestring=[b]',
  morestring=[b]"
}[keywords, comments, strings]

\lstdefinelanguage{svelte}{
    morekeywords={async, await, then, if, each, in, on, let, const, function, import, export, from, with, typeof, new, var, return},
    sensitive=true,
    morecomment=[l]{//},
    morecomment=[s]{/*}{*/},
    morecomment=[s]{<!--}{-->},
    morestring=[b]',
    morestring=[b]",
    morestring=[b]`,
    ndkeywords={import, export, return, if, for, while, function, switch, case, break, var, let, const, async, await, then, catch, try, typeof, new, class, extends, implements, interface, constructor, public, private, protected, static, get, set},
    alsoletter={:},
    moredelim=**[is][\color{orange}]{\@\{}{\@\}},
    moredelim=**[is][\color{cyan}]{\$\{}{\}},
    keywordstyle=\color{blue}\bfseries,
    commentstyle=\color{gray}\ttfamily,
    stringstyle=\color{red}\ttfamily,
    identifierstyle=\color{black},
    breaklines=true,
    basicstyle=\ttfamily\small,
    columns=fullflexible,
    showstringspaces=false,
    tabsize=2,
    frame=single,
    captionpos=b,
}


\lstdefinestyle{java}{
    language=Java,
    basicstyle=\ttfamily\small,
    keywordstyle=\color{blue},
    stringstyle=\color{red},
    commentstyle=\color{gray},
    morecomment=[s][\color{gray}]{/**}{*/},
    breaklines=true,
    numbers=left,
    numberstyle=\tiny\color{gray},
    showspaces=false,
    showstringspaces=false,
    frame=single,
    captionpos=b
}

\usepackage{color}
\definecolor{numb}{rgb}{0.58,0,0.82}
\definecolor{punct}{rgb}{0.3,0.3,0.3}
\definecolor{delim}{rgb}{0.5,0.5,0.5}
\definecolor{string}{rgb}{0.6,0.1,0.1}

\lstdefinestyle{json}{
    language=json,
    basicstyle=\ttfamily\small,
    keywordstyle=\color{blue},
    stringstyle=\color{string},
    commentstyle=\color{green},
    morecomment=[s][\color{gray}]{/**}{*/},
    breaklines=true,
    numbers=left,
    numberstyle=\tiny\color{gray},
    showspaces=false,
    showstringspaces=false,
    frame=single,
    captionpos=b
}

\lstdefinelanguage{drl}{
    basicstyle=\ttfamily\small,
    keywordstyle=\color{blue}\bfseries,
    stringstyle=\color{red},
    commentstyle=\color{green},
    morecomment=[s][\color{green}]{/*}{*/},
    morecomment=[l][\color{green}]{//},
    morestring=[b]",
    keywords={rule, when, then, end, salience, eval, update, retract},
    sensitive=true,
    breaklines=true,
    numbers=left,
    numberstyle=\tiny\color{gray},
    showspaces=false,
    showstringspaces=false,
    frame=single,
    captionpos=b
}

\lstdefinestyle{drl}{
    language=drl,
    basicstyle=\ttfamily\small,
    keywordstyle=\color{blue}\bfseries,
    stringstyle=\color{red},
    commentstyle=\color{green},
    breaklines=true,
    numbers=left,
    numberstyle=\tiny\color{gray},
    showspaces=false,
    showstringspaces=false,
    frame=single,
    captionpos=b
}



\renewcommand{\lstlistingname}{Codi}
\renewcommand{\lstlistlistingname}{Índex de Fragments de Codis}
\renewcommand{\appendixpagename}{Annexos}
\renewcommand{\appendixtocname}{Annexos}

\renewcommand*{\listalgorithmcfname}{Llistat de pseudocodi}
\renewcommand*{\algorithmcfname}{Pseudocodi}
\renewcommand*{\algorithmautorefname}{pseudocodi}

\definecolor{darkgreen}{rgb}{0.0, 0.3, 0.0}
\definecolor{lightgreen}{rgb}{0.9, 0.95, 0.9}

\newcommand{\todo}[1]{
    \begin{tcolorbox}[colback=lightgreen, colframe=darkgreen, title=TODO]
        #1
    \end{tcolorbox}
}


% \usepackage{indentfirst} 
% \setlength{\parindent}{1cm}


\usepackage{parskip}
% \setlength{\parindent}{10pt}

% \setlength{\headheight}{15.6pt}
\setlength{\headheight}{16.05263pt}




\vspace{5mm}

\newcolumntype{u}{>{\centering\arraybackslash\hsize=.22\hsize}X}
\newcolumntype{s}{>{\hsize=.45\hsize}X}


%%%%%%%%%%%%%%%%%%%%%%%%%%%%%%%%
% INI creacio subsubsubsecion
%%%%%%%%%%%%%%%%%%%%%%%%%%%%%%%%
\usepackage{titlesec}

\let\oldparagraph\paragraph
\renewcommand{\paragraph}[1]{\oldparagraph{\hspace{30pt}#1}}


\titleclass{\subsubsubsection}{straight}[\subsection]

\newcounter{subsubsubsection}[subsubsection]
\renewcommand\thesubsubsubsection{\thesubsubsection.\arabic{subsubsubsection}}
\renewcommand\theparagraph{\thesubsubsubsection.\arabic{paragraph}} % optional; useful if paragraphs are to be numbered

\titleformat{\subsubsubsection}
  {\normalfont\normalsize\bfseries}{\thesubsubsubsection}{1em}{}
\titlespacing*{\subsubsubsection}
{0pt}{3.25ex plus 1ex minus .2ex}{1.5ex plus .2ex}


\makeatletter
\renewcommand\paragraph{\@startsection{paragraph}{5}{10pt}%
  {3.25ex \@plus1ex \@minus.2ex}%
  {-1em}%
  {\normalfont\normalsize\bfseries}}
\renewcommand\subparagraph{\@startsection{subparagraph}{6}{\parindent}%
  {3.25ex \@plus1ex \@minus .2ex}%
  {-1em}%
  {\normalfont\normalsize\bfseries}}
\def\toclevel@subsubsubsection{4}
\def\toclevel@paragraph{5}
\def\toclevel@paragraph{6}
\def\l@subsubsubsection{\@dottedtocline{4}{7em}{4em}}
\def\l@paragraph{\@dottedtocline{5}{10em}{5em}}
\def\l@subparagraph{\@dottedtocline{6}{14em}{6em}}
\makeatother

\setcounter{secnumdepth}{4}
\setcounter{tocdepth}{3}
%%%%%%%%%%%%%%%%%%%%%%%%%%%%%%%%
% FIN creacio subsubsubtitle
%%%%%%%%%%%%%%%%%%%%%%%%%%%%%%%%


%%%%%%%%%%%%%%%%%%%%%%%%%%%%%%%%
% INI creacio equation
%%%%%%%%%%%%%%%%%%%%%%%%%%%%%%%%
\DeclareFloatingEnvironment[fileext=loeq,placement={!ht},name=Equació]{equationfloat}
%%%%%%%%%%%%%%%%%%%%%%%%%%%%%%%%
% FIN creacio equation
%%%%%%%%%%%%%%%%%%%%%%%%%%%%%%%%

\titleformat{\section}[block]
  {\clearpage\normalfont\Large\bfseries}
  {\thesection}{1em}{}


\renewcommand\refname{Bibliografia}


\linespread{1.25}


\makeatletter
\renewenvironment{abstract}{%
    \if@twocolumn
      \section*{\abstractname}%
    \else 
      \begin{center}%
        {\bfseries \Large\itshape\abstractname\vspace{\z@}}%  % Título en negrita, 14pt y cursiva
      \end{center}%
      \quotation
    \fi}
    {\if@twocolumn\else\endquotation\fi}
\makeatother

% \fancypagestyle{main}{
%     \fancyhf{}
%     \lhead{\small Aplicació Web per a la Visualització i Comparació
% de Trajectòries d’Algorismes d’Optimització}
%     % \rhead{\includegraphics[width=1.2cm]{resources/Icons/logo-fib.png}}
%     \cfoot{\thepage}
%     \setlength{\footskip}{35pt}
% }

\newcommand{\q}[1]{``#1''}

\begin{document}

\hypersetup{
    pdfborder=0 0 0,
}

%% PORTADA I TABLE OF CONTENTS %%

%%%%                                   %%%%
%         INDEX DE LA MEMORIA             %
%%%%                                   %%%%

% \includepdf[pages=-]{portada_fib.pdf}
%%% Portada
\begin{titlepage}
\begin{center}
    \includegraphics[width=1\linewidth]{extra/figures/logo-upc-fib.png}
    \vspace{30mm}

    \LARGE{\textsc{\textbf{{INTEGRACIÓN Y AUTOMATIZACIÓN DE COSTES DE INFRAESTRUCTURA HPC EN ENTORNOS DE INVESTIGACIÓN}}}} \\
    \vspace{25mm}
      
        
    \textbf{\Large{\textsc{PAU NAVARRO ÁLVAREZ}}} \\
    \vspace{20mm}
        
    
    
    \large{\textbf{Director/a} ROBERTO RIVEIRO INSUA (FUNDACIO INSTITUT DE RECERCA BIOMEDICA (IRB BARCELONA))} \\
    \large{\textbf{Ponente:} ALEX PAJUELO GONZÁLEZ (Departamento de Arquitectura de Computadores)} \\
    \large{\textbf{Titulación:} Grado en Ingeniería Informática (Tecnologías de la información)} \\
    
    \large{\textbf{Lliurament 2:} }Planificación temporal \\
    \large{\textbf{Tutor GEP:} CAROLINA CONSOLIDACIÓN SEGURA (Departament d'Organització d'Empreses)} \\
    \large{\textbf{Fecha:}} 29 de Septiembre de 2025\\

    \textbf{\large{Memoria del trabajo de fin de grado}} \\
    \vspace{3mm}
    \textbf{\large{Facultat d'Informàtica de Barcelona (FIB)}} \\
    \vspace{3mm}
    \textbf{\large{Universitat Politècnica de Catalunya (UPC) - BarcelonaTech}} \\

\end{center}
\end{titlepage}
%\input{Extra/portada}
% \cleardoublepage

% \pagestyle{main}
% \setcounter{page}{11}
\renewcommand{\contentsname}{Continguts}

\tableofcontents
\clearpage
\listoffigures
\begingroup
    \let\clearpage\relax
    \listoftables
\endgroup
\clearpage

\newpage
% \cleardoublepage
\documentclass[a4paper, 11pt, twoside]{article}
\usepackage{wrapfig}
% \usepackage{abstract}
\usepackage{graphicx} % Required for inserting images
\usepackage[a4paper, left=1.15in, right=1.15in, top=1.5in, bottom=1.2in]{geometry}
\usepackage{float}
\usepackage{array}
\usepackage{booktabs}
\usepackage[T1]{fontenc}
\usepackage{amsmath}
\usepackage{tabularx}
\usepackage{graphicx} % Required for inserting images
\usepackage{tikz, xcolor}
\usepackage{hyperref}
\usepackage{cite}
\usepackage{pdfpages}
\usepackage{pdflscape}
\usepackage{afterpage}
\usepackage{textpos}
\usepackage{fancyhdr}
\usepackage{caption}
\usepackage{lmodern,textcomp}
\usepackage[catalan]{babel}
\usepackage{lipsum}
\usepackage{newfloat}
\usepackage{listings}
\usepackage{placeins}
\usepackage[ruled]{algorithm2e}
\usepackage{tcolorbox}
% Estilo de citas IEEE sencillo: usar paquete cite para comprimir y ordenar referencias numéricas
\usepackage{cite}
\usepackage[toc,page]{appendix}
\usepackage{forest}
\usepackage{pdflscape}
\usepackage{graphicx}
\usepackage{geometry}
\usepackage{quoting,xparse}
\usepackage{svg}



% \NewDocumentCommand{\bywhom}{m}{% the Bourbaki trick
%   {\nobreak\hfill\penalty50\hskip1em\null\nobreak
%    \hfill\mbox{\normalfont(#1)}%
%    \parfillskip=0pt \finalhyphendemerits=0 \par}%
% }

% AIXÒ SERVEIX PER A MODIFICAR UN PROBLEMA AMB ELS CARACTERS --> " + PARAULA 

% \NewDocumentEnvironment{pquotation}{m}
%   {\begin{quoting}[
%      indentfirst=true,
%      leftmargin=\parindent,
%      rightmargin=\parindent]\itshape}
%   {\bywhom{#1}
  
%   \end{quoting}}

% AIXÒ PER INSERIR FRAGMENTS DE CODI

\lstnewenvironment{code}[1][]%
{
   \noindent
   \minipage{\linewidth}
   \vspace{0.5\baselineskip}
   \lstset{
     basicstyle=\ttfamily\footnotesize,
     frame=single,
     breaklines=true,    % Permitir que las líneas se rompan automáticamente
     breakatwhitespace=true, % Permitir la ruptura en los espacios en blanco
     #1
   }
}
{\endminipage}

\lstset{
    language=Java,
    basicstyle=\ttfamily\small,  % Define el tipo de letra y tamaño del código
    keywordstyle=\color{blue},   % Define el color de las palabras clave de Java
    stringstyle=\color{red},     % Define el color de las cadenas de texto
    commentstyle=\color{green},  % Define el color de los comentarios
    morecomment=[s][\color{gray}]{/**}{*/},  % Estilo para comentarios de bloque
    breaklines=true,  % Configura el quiebre automático de línea
    numbers=left,  % Posición de los números de línea
    numberstyle=\tiny\color{gray},  % Estilo para los números de línea
    showspaces=false,  % No mostrar espacios con underscores
    showstringspaces=false,  % No mostrar espacios dentro de strings
    frame=single,  % Agregar un marco alrededor del código
    captionpos=b,  % b = bottom, t = top
}

\definecolor{eclipseStrings}{RGB}{42,0.0,255}
\definecolor{eclipseKeywords}{RGB}{127,0,85}
\colorlet{numb}{magenta!60!black}

\lstdefinelanguage{json}{
    basicstyle=\ttfamily\small,
    showstringspaces=false,
    breaklines=true,
    frame=single,
    literate=
     *{0}{{{\color{numb}0}}}{1}
      {1}{{{\color{numb}1}}}{1}
      {2}{{{\color{numb}2}}}{1}
      {3}{{{\color{numb}3}}}{1}
      {4}{{{\color{numb}4}}}{1}
      {5}{{{\color{numb}5}}}{1}
      {6}{{{\color{numb}6}}}{1}
      {7}{{{\color{numb}7}}}{1}
      {8}{{{\color{numb}8}}}{1}
      {9}{{{\color{numb}9}}}{1}
      {:}{{{\color{punct}{:}}}}{1}
      {,}{{{\color{punct}{,}}}}{1}
      {\{}{{{\color{delim}{\{}}}}{1}
      {\}}{{{\color{delim}{\}}}}}{1}
      {[}{{{\color{delim}{[}}}}{1}
      {]}{{{\color{delim}{]}}}}{1},
    morestring=[b]"
}

\lstdefinelanguage{JavaScript}{
  morekeywords=[1]{break, continue, delete, else, for, function, if, in,
    new, return, this, typeof, var, void, while, with},
  % Literals, primitive types, and reference types.
  morekeywords=[2]{false, null, true, boolean, number, undefined,
    Array, Boolean, Date, Math, Number, String, Object},
  % Built-ins.
  morekeywords=[3]{eval, parseInt, parseFloat, escape, unescape},
  sensitive,
  morecomment=[s]{/*}{*/},
  morecomment=[l]//,
  morecomment=[s]{/**}{*/}, % JavaDoc style comments
  morestring=[b]',
  morestring=[b]"
}[keywords, comments, strings]

\lstdefinelanguage{svelte}{
    morekeywords={async, await, then, if, each, in, on, let, const, function, import, export, from, with, typeof, new, var, return},
    sensitive=true,
    morecomment=[l]{//},
    morecomment=[s]{/*}{*/},
    morecomment=[s]{<!--}{-->},
    morestring=[b]',
    morestring=[b]",
    morestring=[b]`,
    ndkeywords={import, export, return, if, for, while, function, switch, case, break, var, let, const, async, await, then, catch, try, typeof, new, class, extends, implements, interface, constructor, public, private, protected, static, get, set},
    alsoletter={:},
    moredelim=**[is][\color{orange}]{\@\{}{\@\}},
    moredelim=**[is][\color{cyan}]{\$\{}{\}},
    keywordstyle=\color{blue}\bfseries,
    commentstyle=\color{gray}\ttfamily,
    stringstyle=\color{red}\ttfamily,
    identifierstyle=\color{black},
    breaklines=true,
    basicstyle=\ttfamily\small,
    columns=fullflexible,
    showstringspaces=false,
    tabsize=2,
    frame=single,
    captionpos=b,
}


\lstdefinestyle{java}{
    language=Java,
    basicstyle=\ttfamily\small,
    keywordstyle=\color{blue},
    stringstyle=\color{red},
    commentstyle=\color{gray},
    morecomment=[s][\color{gray}]{/**}{*/},
    breaklines=true,
    numbers=left,
    numberstyle=\tiny\color{gray},
    showspaces=false,
    showstringspaces=false,
    frame=single,
    captionpos=b
}

\usepackage{color}
\definecolor{numb}{rgb}{0.58,0,0.82}
\definecolor{punct}{rgb}{0.3,0.3,0.3}
\definecolor{delim}{rgb}{0.5,0.5,0.5}
\definecolor{string}{rgb}{0.6,0.1,0.1}

\lstdefinestyle{json}{
    language=json,
    basicstyle=\ttfamily\small,
    keywordstyle=\color{blue},
    stringstyle=\color{string},
    commentstyle=\color{green},
    morecomment=[s][\color{gray}]{/**}{*/},
    breaklines=true,
    numbers=left,
    numberstyle=\tiny\color{gray},
    showspaces=false,
    showstringspaces=false,
    frame=single,
    captionpos=b
}

\lstdefinelanguage{drl}{
    basicstyle=\ttfamily\small,
    keywordstyle=\color{blue}\bfseries,
    stringstyle=\color{red},
    commentstyle=\color{green},
    morecomment=[s][\color{green}]{/*}{*/},
    morecomment=[l][\color{green}]{//},
    morestring=[b]",
    keywords={rule, when, then, end, salience, eval, update, retract},
    sensitive=true,
    breaklines=true,
    numbers=left,
    numberstyle=\tiny\color{gray},
    showspaces=false,
    showstringspaces=false,
    frame=single,
    captionpos=b
}

\lstdefinestyle{drl}{
    language=drl,
    basicstyle=\ttfamily\small,
    keywordstyle=\color{blue}\bfseries,
    stringstyle=\color{red},
    commentstyle=\color{green},
    breaklines=true,
    numbers=left,
    numberstyle=\tiny\color{gray},
    showspaces=false,
    showstringspaces=false,
    frame=single,
    captionpos=b
}



\renewcommand{\lstlistingname}{Codi}
\renewcommand{\lstlistlistingname}{Índex de Fragments de Codis}
\renewcommand{\appendixpagename}{Annexos}
\renewcommand{\appendixtocname}{Annexos}

\renewcommand*{\listalgorithmcfname}{Llistat de pseudocodi}
\renewcommand*{\algorithmcfname}{Pseudocodi}
\renewcommand*{\algorithmautorefname}{pseudocodi}

\definecolor{darkgreen}{rgb}{0.0, 0.3, 0.0}
\definecolor{lightgreen}{rgb}{0.9, 0.95, 0.9}

\newcommand{\todo}[1]{
    \begin{tcolorbox}[colback=lightgreen, colframe=darkgreen, title=TODO]
        #1
    \end{tcolorbox}
}


% \usepackage{indentfirst} 
% \setlength{\parindent}{1cm}


\usepackage{parskip}
% \setlength{\parindent}{10pt}

% \setlength{\headheight}{15.6pt}
\setlength{\headheight}{16.05263pt}




\vspace{5mm}

\newcolumntype{u}{>{\centering\arraybackslash\hsize=.22\hsize}X}
\newcolumntype{s}{>{\hsize=.45\hsize}X}


%%%%%%%%%%%%%%%%%%%%%%%%%%%%%%%%
% INI creacio subsubsubsecion
%%%%%%%%%%%%%%%%%%%%%%%%%%%%%%%%
\usepackage{titlesec}

\let\oldparagraph\paragraph
\renewcommand{\paragraph}[1]{\oldparagraph{\hspace{30pt}#1}}


\titleclass{\subsubsubsection}{straight}[\subsection]

\newcounter{subsubsubsection}[subsubsection]
\renewcommand\thesubsubsubsection{\thesubsubsection.\arabic{subsubsubsection}}
\renewcommand\theparagraph{\thesubsubsubsection.\arabic{paragraph}} % optional; useful if paragraphs are to be numbered

\titleformat{\subsubsubsection}
  {\normalfont\normalsize\bfseries}{\thesubsubsubsection}{1em}{}
\titlespacing*{\subsubsubsection}
{0pt}{3.25ex plus 1ex minus .2ex}{1.5ex plus .2ex}


\makeatletter
\renewcommand\paragraph{\@startsection{paragraph}{5}{10pt}%
  {3.25ex \@plus1ex \@minus.2ex}%
  {-1em}%
  {\normalfont\normalsize\bfseries}}
\renewcommand\subparagraph{\@startsection{subparagraph}{6}{\parindent}%
  {3.25ex \@plus1ex \@minus .2ex}%
  {-1em}%
  {\normalfont\normalsize\bfseries}}
\def\toclevel@subsubsubsection{4}
\def\toclevel@paragraph{5}
\def\toclevel@paragraph{6}
\def\l@subsubsubsection{\@dottedtocline{4}{7em}{4em}}
\def\l@paragraph{\@dottedtocline{5}{10em}{5em}}
\def\l@subparagraph{\@dottedtocline{6}{14em}{6em}}
\makeatother

\setcounter{secnumdepth}{4}
\setcounter{tocdepth}{3}
%%%%%%%%%%%%%%%%%%%%%%%%%%%%%%%%
% FIN creacio subsubsubtitle
%%%%%%%%%%%%%%%%%%%%%%%%%%%%%%%%


%%%%%%%%%%%%%%%%%%%%%%%%%%%%%%%%
% INI creacio equation
%%%%%%%%%%%%%%%%%%%%%%%%%%%%%%%%
\DeclareFloatingEnvironment[fileext=loeq,placement={!ht},name=Equació]{equationfloat}
%%%%%%%%%%%%%%%%%%%%%%%%%%%%%%%%
% FIN creacio equation
%%%%%%%%%%%%%%%%%%%%%%%%%%%%%%%%

\titleformat{\section}[block]
  {\clearpage\normalfont\Large\bfseries}
  {\thesection}{1em}{}


\renewcommand\refname{Bibliografia}


\linespread{1.25}


\makeatletter
\renewenvironment{abstract}{%
    \if@twocolumn
      \section*{\abstractname}%
    \else 
      \begin{center}%
        {\bfseries \Large\itshape\abstractname\vspace{\z@}}%  % Título en negrita, 14pt y cursiva
      \end{center}%
      \quotation
    \fi}
    {\if@twocolumn\else\endquotation\fi}
\makeatother

% \fancypagestyle{main}{
%     \fancyhf{}
%     \lhead{\small Aplicació Web per a la Visualització i Comparació
% de Trajectòries d’Algorismes d’Optimització}
%     % \rhead{\includegraphics[width=1.2cm]{resources/Icons/logo-fib.png}}
%     \cfoot{\thepage}
%     \setlength{\footskip}{35pt}
% }

\newcommand{\q}[1]{``#1''}

\begin{document}

\hypersetup{
    pdfborder=0 0 0,
}

%% PORTADA I TABLE OF CONTENTS %%

%%%%                                   %%%%
%         INDEX DE LA MEMORIA             %
%%%%                                   %%%%

% \includepdf[pages=-]{portada_fib.pdf}
%%% Portada
\begin{titlepage}
\begin{center}
    \includegraphics[width=1\linewidth]{extra/figures/logo-upc-fib.png}
    \vspace{30mm}

    \LARGE{\textsc{\textbf{{INTEGRACIÓN Y AUTOMATIZACIÓN DE COSTES DE INFRAESTRUCTURA HPC EN ENTORNOS DE INVESTIGACIÓN}}}} \\
    \vspace{25mm}
      
        
    \textbf{\Large{\textsc{PAU NAVARRO ÁLVAREZ}}} \\
    \vspace{20mm}
        
    
    
    \large{\textbf{Director/a} ROBERTO RIVEIRO INSUA (FUNDACIO INSTITUT DE RECERCA BIOMEDICA (IRB BARCELONA))} \\
    \large{\textbf{Ponente:} ALEX PAJUELO GONZÁLEZ (Departamento de Arquitectura de Computadores)} \\
    \large{\textbf{Titulación:} Grado en Ingeniería Informática (Tecnologías de la información)} \\
    
    \large{\textbf{Lliurament 2:} }Planificación temporal \\
    \large{\textbf{Tutor GEP:} CAROLINA CONSOLIDACIÓN SEGURA (Departament d'Organització d'Empreses)} \\
    \large{\textbf{Fecha:}} 29 de Septiembre de 2025\\

    \textbf{\large{Memoria del trabajo de fin de grado}} \\
    \vspace{3mm}
    \textbf{\large{Facultat d'Informàtica de Barcelona (FIB)}} \\
    \vspace{3mm}
    \textbf{\large{Universitat Politècnica de Catalunya (UPC) - BarcelonaTech}} \\

\end{center}
\end{titlepage}
%\input{Extra/portada}
% \cleardoublepage

% \pagestyle{main}
% \setcounter{page}{11}
\renewcommand{\contentsname}{Continguts}

\tableofcontents
\clearpage
\listoffigures
\begingroup
    \let\clearpage\relax
    \listoftables
\endgroup
\clearpage

\newpage
% \cleardoublepage
\input{Contexto y alcance/Contextualización/main}
\input{Contexto y alcance/Justificación/main}
\input{Contexto y alcance/Alcance/main}
\input{Contexto y alcance/Metodología y herramientas de seguimiento/main}

\documentclass[a4paper, 11pt, twoside]{article}
\usepackage{wrapfig}
% \usepackage{abstract}
\usepackage{graphicx} % Required for inserting images
\usepackage[a4paper, left=1.15in, right=1.15in, top=1.5in, bottom=1.2in]{geometry}
\usepackage{float}
\usepackage{array}
\usepackage{booktabs}
\usepackage[T1]{fontenc}
\usepackage{amsmath}
\usepackage{tabularx}
\usepackage{graphicx} % Required for inserting images
\usepackage{tikz, xcolor}
\usepackage{hyperref}
\usepackage{cite}
\usepackage{pdfpages}
\usepackage{pdflscape}
\usepackage{afterpage}
\usepackage{textpos}
\usepackage{fancyhdr}
\usepackage{caption}
\usepackage{lmodern,textcomp}
\usepackage[catalan]{babel}
\usepackage{lipsum}
\usepackage{newfloat}
\usepackage{listings}
\usepackage{placeins}
\usepackage[ruled]{algorithm2e}
\usepackage{tcolorbox}
% Estilo de citas IEEE sencillo: usar paquete cite para comprimir y ordenar referencias numéricas
\usepackage{cite}
\usepackage[toc,page]{appendix}
\usepackage{forest}
\usepackage{pdflscape}
\usepackage{graphicx}
\usepackage{geometry}
\usepackage{quoting,xparse}
\usepackage{svg}



% \NewDocumentCommand{\bywhom}{m}{% the Bourbaki trick
%   {\nobreak\hfill\penalty50\hskip1em\null\nobreak
%    \hfill\mbox{\normalfont(#1)}%
%    \parfillskip=0pt \finalhyphendemerits=0 \par}%
% }

% AIXÒ SERVEIX PER A MODIFICAR UN PROBLEMA AMB ELS CARACTERS --> " + PARAULA 

% \NewDocumentEnvironment{pquotation}{m}
%   {\begin{quoting}[
%      indentfirst=true,
%      leftmargin=\parindent,
%      rightmargin=\parindent]\itshape}
%   {\bywhom{#1}
  
%   \end{quoting}}

% AIXÒ PER INSERIR FRAGMENTS DE CODI

\lstnewenvironment{code}[1][]%
{
   \noindent
   \minipage{\linewidth}
   \vspace{0.5\baselineskip}
   \lstset{
     basicstyle=\ttfamily\footnotesize,
     frame=single,
     breaklines=true,    % Permitir que las líneas se rompan automáticamente
     breakatwhitespace=true, % Permitir la ruptura en los espacios en blanco
     #1
   }
}
{\endminipage}

\lstset{
    language=Java,
    basicstyle=\ttfamily\small,  % Define el tipo de letra y tamaño del código
    keywordstyle=\color{blue},   % Define el color de las palabras clave de Java
    stringstyle=\color{red},     % Define el color de las cadenas de texto
    commentstyle=\color{green},  % Define el color de los comentarios
    morecomment=[s][\color{gray}]{/**}{*/},  % Estilo para comentarios de bloque
    breaklines=true,  % Configura el quiebre automático de línea
    numbers=left,  % Posición de los números de línea
    numberstyle=\tiny\color{gray},  % Estilo para los números de línea
    showspaces=false,  % No mostrar espacios con underscores
    showstringspaces=false,  % No mostrar espacios dentro de strings
    frame=single,  % Agregar un marco alrededor del código
    captionpos=b,  % b = bottom, t = top
}

\definecolor{eclipseStrings}{RGB}{42,0.0,255}
\definecolor{eclipseKeywords}{RGB}{127,0,85}
\colorlet{numb}{magenta!60!black}

\lstdefinelanguage{json}{
    basicstyle=\ttfamily\small,
    showstringspaces=false,
    breaklines=true,
    frame=single,
    literate=
     *{0}{{{\color{numb}0}}}{1}
      {1}{{{\color{numb}1}}}{1}
      {2}{{{\color{numb}2}}}{1}
      {3}{{{\color{numb}3}}}{1}
      {4}{{{\color{numb}4}}}{1}
      {5}{{{\color{numb}5}}}{1}
      {6}{{{\color{numb}6}}}{1}
      {7}{{{\color{numb}7}}}{1}
      {8}{{{\color{numb}8}}}{1}
      {9}{{{\color{numb}9}}}{1}
      {:}{{{\color{punct}{:}}}}{1}
      {,}{{{\color{punct}{,}}}}{1}
      {\{}{{{\color{delim}{\{}}}}{1}
      {\}}{{{\color{delim}{\}}}}}{1}
      {[}{{{\color{delim}{[}}}}{1}
      {]}{{{\color{delim}{]}}}}{1},
    morestring=[b]"
}

\lstdefinelanguage{JavaScript}{
  morekeywords=[1]{break, continue, delete, else, for, function, if, in,
    new, return, this, typeof, var, void, while, with},
  % Literals, primitive types, and reference types.
  morekeywords=[2]{false, null, true, boolean, number, undefined,
    Array, Boolean, Date, Math, Number, String, Object},
  % Built-ins.
  morekeywords=[3]{eval, parseInt, parseFloat, escape, unescape},
  sensitive,
  morecomment=[s]{/*}{*/},
  morecomment=[l]//,
  morecomment=[s]{/**}{*/}, % JavaDoc style comments
  morestring=[b]',
  morestring=[b]"
}[keywords, comments, strings]

\lstdefinelanguage{svelte}{
    morekeywords={async, await, then, if, each, in, on, let, const, function, import, export, from, with, typeof, new, var, return},
    sensitive=true,
    morecomment=[l]{//},
    morecomment=[s]{/*}{*/},
    morecomment=[s]{<!--}{-->},
    morestring=[b]',
    morestring=[b]",
    morestring=[b]`,
    ndkeywords={import, export, return, if, for, while, function, switch, case, break, var, let, const, async, await, then, catch, try, typeof, new, class, extends, implements, interface, constructor, public, private, protected, static, get, set},
    alsoletter={:},
    moredelim=**[is][\color{orange}]{\@\{}{\@\}},
    moredelim=**[is][\color{cyan}]{\$\{}{\}},
    keywordstyle=\color{blue}\bfseries,
    commentstyle=\color{gray}\ttfamily,
    stringstyle=\color{red}\ttfamily,
    identifierstyle=\color{black},
    breaklines=true,
    basicstyle=\ttfamily\small,
    columns=fullflexible,
    showstringspaces=false,
    tabsize=2,
    frame=single,
    captionpos=b,
}


\lstdefinestyle{java}{
    language=Java,
    basicstyle=\ttfamily\small,
    keywordstyle=\color{blue},
    stringstyle=\color{red},
    commentstyle=\color{gray},
    morecomment=[s][\color{gray}]{/**}{*/},
    breaklines=true,
    numbers=left,
    numberstyle=\tiny\color{gray},
    showspaces=false,
    showstringspaces=false,
    frame=single,
    captionpos=b
}

\usepackage{color}
\definecolor{numb}{rgb}{0.58,0,0.82}
\definecolor{punct}{rgb}{0.3,0.3,0.3}
\definecolor{delim}{rgb}{0.5,0.5,0.5}
\definecolor{string}{rgb}{0.6,0.1,0.1}

\lstdefinestyle{json}{
    language=json,
    basicstyle=\ttfamily\small,
    keywordstyle=\color{blue},
    stringstyle=\color{string},
    commentstyle=\color{green},
    morecomment=[s][\color{gray}]{/**}{*/},
    breaklines=true,
    numbers=left,
    numberstyle=\tiny\color{gray},
    showspaces=false,
    showstringspaces=false,
    frame=single,
    captionpos=b
}

\lstdefinelanguage{drl}{
    basicstyle=\ttfamily\small,
    keywordstyle=\color{blue}\bfseries,
    stringstyle=\color{red},
    commentstyle=\color{green},
    morecomment=[s][\color{green}]{/*}{*/},
    morecomment=[l][\color{green}]{//},
    morestring=[b]",
    keywords={rule, when, then, end, salience, eval, update, retract},
    sensitive=true,
    breaklines=true,
    numbers=left,
    numberstyle=\tiny\color{gray},
    showspaces=false,
    showstringspaces=false,
    frame=single,
    captionpos=b
}

\lstdefinestyle{drl}{
    language=drl,
    basicstyle=\ttfamily\small,
    keywordstyle=\color{blue}\bfseries,
    stringstyle=\color{red},
    commentstyle=\color{green},
    breaklines=true,
    numbers=left,
    numberstyle=\tiny\color{gray},
    showspaces=false,
    showstringspaces=false,
    frame=single,
    captionpos=b
}



\renewcommand{\lstlistingname}{Codi}
\renewcommand{\lstlistlistingname}{Índex de Fragments de Codis}
\renewcommand{\appendixpagename}{Annexos}
\renewcommand{\appendixtocname}{Annexos}

\renewcommand*{\listalgorithmcfname}{Llistat de pseudocodi}
\renewcommand*{\algorithmcfname}{Pseudocodi}
\renewcommand*{\algorithmautorefname}{pseudocodi}

\definecolor{darkgreen}{rgb}{0.0, 0.3, 0.0}
\definecolor{lightgreen}{rgb}{0.9, 0.95, 0.9}

\newcommand{\todo}[1]{
    \begin{tcolorbox}[colback=lightgreen, colframe=darkgreen, title=TODO]
        #1
    \end{tcolorbox}
}


% \usepackage{indentfirst} 
% \setlength{\parindent}{1cm}


\usepackage{parskip}
% \setlength{\parindent}{10pt}

% \setlength{\headheight}{15.6pt}
\setlength{\headheight}{16.05263pt}




\vspace{5mm}

\newcolumntype{u}{>{\centering\arraybackslash\hsize=.22\hsize}X}
\newcolumntype{s}{>{\hsize=.45\hsize}X}


%%%%%%%%%%%%%%%%%%%%%%%%%%%%%%%%
% INI creacio subsubsubsecion
%%%%%%%%%%%%%%%%%%%%%%%%%%%%%%%%
\usepackage{titlesec}

\let\oldparagraph\paragraph
\renewcommand{\paragraph}[1]{\oldparagraph{\hspace{30pt}#1}}


\titleclass{\subsubsubsection}{straight}[\subsection]

\newcounter{subsubsubsection}[subsubsection]
\renewcommand\thesubsubsubsection{\thesubsubsection.\arabic{subsubsubsection}}
\renewcommand\theparagraph{\thesubsubsubsection.\arabic{paragraph}} % optional; useful if paragraphs are to be numbered

\titleformat{\subsubsubsection}
  {\normalfont\normalsize\bfseries}{\thesubsubsubsection}{1em}{}
\titlespacing*{\subsubsubsection}
{0pt}{3.25ex plus 1ex minus .2ex}{1.5ex plus .2ex}


\makeatletter
\renewcommand\paragraph{\@startsection{paragraph}{5}{10pt}%
  {3.25ex \@plus1ex \@minus.2ex}%
  {-1em}%
  {\normalfont\normalsize\bfseries}}
\renewcommand\subparagraph{\@startsection{subparagraph}{6}{\parindent}%
  {3.25ex \@plus1ex \@minus .2ex}%
  {-1em}%
  {\normalfont\normalsize\bfseries}}
\def\toclevel@subsubsubsection{4}
\def\toclevel@paragraph{5}
\def\toclevel@paragraph{6}
\def\l@subsubsubsection{\@dottedtocline{4}{7em}{4em}}
\def\l@paragraph{\@dottedtocline{5}{10em}{5em}}
\def\l@subparagraph{\@dottedtocline{6}{14em}{6em}}
\makeatother

\setcounter{secnumdepth}{4}
\setcounter{tocdepth}{3}
%%%%%%%%%%%%%%%%%%%%%%%%%%%%%%%%
% FIN creacio subsubsubtitle
%%%%%%%%%%%%%%%%%%%%%%%%%%%%%%%%


%%%%%%%%%%%%%%%%%%%%%%%%%%%%%%%%
% INI creacio equation
%%%%%%%%%%%%%%%%%%%%%%%%%%%%%%%%
\DeclareFloatingEnvironment[fileext=loeq,placement={!ht},name=Equació]{equationfloat}
%%%%%%%%%%%%%%%%%%%%%%%%%%%%%%%%
% FIN creacio equation
%%%%%%%%%%%%%%%%%%%%%%%%%%%%%%%%

\titleformat{\section}[block]
  {\clearpage\normalfont\Large\bfseries}
  {\thesection}{1em}{}


\renewcommand\refname{Bibliografia}


\linespread{1.25}


\makeatletter
\renewenvironment{abstract}{%
    \if@twocolumn
      \section*{\abstractname}%
    \else 
      \begin{center}%
        {\bfseries \Large\itshape\abstractname\vspace{\z@}}%  % Título en negrita, 14pt y cursiva
      \end{center}%
      \quotation
    \fi}
    {\if@twocolumn\else\endquotation\fi}
\makeatother

% \fancypagestyle{main}{
%     \fancyhf{}
%     \lhead{\small Aplicació Web per a la Visualització i Comparació
% de Trajectòries d’Algorismes d’Optimització}
%     % \rhead{\includegraphics[width=1.2cm]{resources/Icons/logo-fib.png}}
%     \cfoot{\thepage}
%     \setlength{\footskip}{35pt}
% }

\newcommand{\q}[1]{``#1''}

\begin{document}

\hypersetup{
    pdfborder=0 0 0,
}

%% PORTADA I TABLE OF CONTENTS %%

%%%%                                   %%%%
%         INDEX DE LA MEMORIA             %
%%%%                                   %%%%

% \includepdf[pages=-]{portada_fib.pdf}
\input{Extra/GEP}
%\input{Extra/portada}
% \cleardoublepage

% \pagestyle{main}
% \setcounter{page}{11}
\include{toc}
% \cleardoublepage
\input{Contexto y alcance/main}

\input{Planificación temporal/main}

\input{Gestión económica/main}

\input{Informe de sostenibilidad/main}

% \cleardoublepage

% \cleardoublepage


% \input{Annex/main}


\clearpage
% Estilo bibliográfico IEEE
\bibliographystyle{IEEEtran}
\renewcommand{\refname}{Referencias}
\bibliography{references.bib}

%%%%                                   %%%%
%         FINAL DE LA MEMORIA             %
%%%%                                   %%%%

\end{document}


%% VERSIÓ INCIAL MEMORIA %%



\section{Gestión económica}

La gestión económica de un proyecto es fundamental para asegurar su viabilidad y exito a largo plazo. Ya hemos visto en secciones anteriores, los otros dos vertices 
del triángulo de la gestión de proyectos: el alcance y el tiempo. En esta sección, nos centraremos en cómo gestionar eficazmente los recursos financieros para garantizar
que el proyecto se mantenga dentro del presupuesto y se logren los objetivos establecidos. 

\input{Gestión económica/Recursos humanos.tex}

\input{Gestión económica/Recursos Hardware.tex}

\input{Gestión económica/Recursos Software.tex}

\input{Gestión económica/Contingencias.tex}

\section{Informe de sostenibilidad}

\input{Informe de sostenibilidad/Autoevaluación.tex}


\input{Informe de sostenibilidad/Dimensión ambiental.tex}

\input{Informe de sostenibilidad/Dimensión económica.tex}

\input{Informe de sostenibilidad/Dimensión social.tex}


% \cleardoublepage

% \cleardoublepage


% \documentclass[a4paper, 11pt, twoside]{article}
\usepackage{wrapfig}
% \usepackage{abstract}
\usepackage{graphicx} % Required for inserting images
\usepackage[a4paper, left=1.15in, right=1.15in, top=1.5in, bottom=1.2in]{geometry}
\usepackage{float}
\usepackage{array}
\usepackage{booktabs}
\usepackage[T1]{fontenc}
\usepackage{amsmath}
\usepackage{tabularx}
\usepackage{graphicx} % Required for inserting images
\usepackage{tikz, xcolor}
\usepackage{hyperref}
\usepackage{cite}
\usepackage{pdfpages}
\usepackage{pdflscape}
\usepackage{afterpage}
\usepackage{textpos}
\usepackage{fancyhdr}
\usepackage{caption}
\usepackage{lmodern,textcomp}
\usepackage[catalan]{babel}
\usepackage{lipsum}
\usepackage{newfloat}
\usepackage{listings}
\usepackage{placeins}
\usepackage[ruled]{algorithm2e}
\usepackage{tcolorbox}
% Estilo de citas IEEE sencillo: usar paquete cite para comprimir y ordenar referencias numéricas
\usepackage{cite}
\usepackage[toc,page]{appendix}
\usepackage{forest}
\usepackage{pdflscape}
\usepackage{graphicx}
\usepackage{geometry}
\usepackage{quoting,xparse}
\usepackage{svg}



% \NewDocumentCommand{\bywhom}{m}{% the Bourbaki trick
%   {\nobreak\hfill\penalty50\hskip1em\null\nobreak
%    \hfill\mbox{\normalfont(#1)}%
%    \parfillskip=0pt \finalhyphendemerits=0 \par}%
% }

% AIXÒ SERVEIX PER A MODIFICAR UN PROBLEMA AMB ELS CARACTERS --> " + PARAULA 

% \NewDocumentEnvironment{pquotation}{m}
%   {\begin{quoting}[
%      indentfirst=true,
%      leftmargin=\parindent,
%      rightmargin=\parindent]\itshape}
%   {\bywhom{#1}
  
%   \end{quoting}}

% AIXÒ PER INSERIR FRAGMENTS DE CODI

\lstnewenvironment{code}[1][]%
{
   \noindent
   \minipage{\linewidth}
   \vspace{0.5\baselineskip}
   \lstset{
     basicstyle=\ttfamily\footnotesize,
     frame=single,
     breaklines=true,    % Permitir que las líneas se rompan automáticamente
     breakatwhitespace=true, % Permitir la ruptura en los espacios en blanco
     #1
   }
}
{\endminipage}

\lstset{
    language=Java,
    basicstyle=\ttfamily\small,  % Define el tipo de letra y tamaño del código
    keywordstyle=\color{blue},   % Define el color de las palabras clave de Java
    stringstyle=\color{red},     % Define el color de las cadenas de texto
    commentstyle=\color{green},  % Define el color de los comentarios
    morecomment=[s][\color{gray}]{/**}{*/},  % Estilo para comentarios de bloque
    breaklines=true,  % Configura el quiebre automático de línea
    numbers=left,  % Posición de los números de línea
    numberstyle=\tiny\color{gray},  % Estilo para los números de línea
    showspaces=false,  % No mostrar espacios con underscores
    showstringspaces=false,  % No mostrar espacios dentro de strings
    frame=single,  % Agregar un marco alrededor del código
    captionpos=b,  % b = bottom, t = top
}

\definecolor{eclipseStrings}{RGB}{42,0.0,255}
\definecolor{eclipseKeywords}{RGB}{127,0,85}
\colorlet{numb}{magenta!60!black}

\lstdefinelanguage{json}{
    basicstyle=\ttfamily\small,
    showstringspaces=false,
    breaklines=true,
    frame=single,
    literate=
     *{0}{{{\color{numb}0}}}{1}
      {1}{{{\color{numb}1}}}{1}
      {2}{{{\color{numb}2}}}{1}
      {3}{{{\color{numb}3}}}{1}
      {4}{{{\color{numb}4}}}{1}
      {5}{{{\color{numb}5}}}{1}
      {6}{{{\color{numb}6}}}{1}
      {7}{{{\color{numb}7}}}{1}
      {8}{{{\color{numb}8}}}{1}
      {9}{{{\color{numb}9}}}{1}
      {:}{{{\color{punct}{:}}}}{1}
      {,}{{{\color{punct}{,}}}}{1}
      {\{}{{{\color{delim}{\{}}}}{1}
      {\}}{{{\color{delim}{\}}}}}{1}
      {[}{{{\color{delim}{[}}}}{1}
      {]}{{{\color{delim}{]}}}}{1},
    morestring=[b]"
}

\lstdefinelanguage{JavaScript}{
  morekeywords=[1]{break, continue, delete, else, for, function, if, in,
    new, return, this, typeof, var, void, while, with},
  % Literals, primitive types, and reference types.
  morekeywords=[2]{false, null, true, boolean, number, undefined,
    Array, Boolean, Date, Math, Number, String, Object},
  % Built-ins.
  morekeywords=[3]{eval, parseInt, parseFloat, escape, unescape},
  sensitive,
  morecomment=[s]{/*}{*/},
  morecomment=[l]//,
  morecomment=[s]{/**}{*/}, % JavaDoc style comments
  morestring=[b]',
  morestring=[b]"
}[keywords, comments, strings]

\lstdefinelanguage{svelte}{
    morekeywords={async, await, then, if, each, in, on, let, const, function, import, export, from, with, typeof, new, var, return},
    sensitive=true,
    morecomment=[l]{//},
    morecomment=[s]{/*}{*/},
    morecomment=[s]{<!--}{-->},
    morestring=[b]',
    morestring=[b]",
    morestring=[b]`,
    ndkeywords={import, export, return, if, for, while, function, switch, case, break, var, let, const, async, await, then, catch, try, typeof, new, class, extends, implements, interface, constructor, public, private, protected, static, get, set},
    alsoletter={:},
    moredelim=**[is][\color{orange}]{\@\{}{\@\}},
    moredelim=**[is][\color{cyan}]{\$\{}{\}},
    keywordstyle=\color{blue}\bfseries,
    commentstyle=\color{gray}\ttfamily,
    stringstyle=\color{red}\ttfamily,
    identifierstyle=\color{black},
    breaklines=true,
    basicstyle=\ttfamily\small,
    columns=fullflexible,
    showstringspaces=false,
    tabsize=2,
    frame=single,
    captionpos=b,
}


\lstdefinestyle{java}{
    language=Java,
    basicstyle=\ttfamily\small,
    keywordstyle=\color{blue},
    stringstyle=\color{red},
    commentstyle=\color{gray},
    morecomment=[s][\color{gray}]{/**}{*/},
    breaklines=true,
    numbers=left,
    numberstyle=\tiny\color{gray},
    showspaces=false,
    showstringspaces=false,
    frame=single,
    captionpos=b
}

\usepackage{color}
\definecolor{numb}{rgb}{0.58,0,0.82}
\definecolor{punct}{rgb}{0.3,0.3,0.3}
\definecolor{delim}{rgb}{0.5,0.5,0.5}
\definecolor{string}{rgb}{0.6,0.1,0.1}

\lstdefinestyle{json}{
    language=json,
    basicstyle=\ttfamily\small,
    keywordstyle=\color{blue},
    stringstyle=\color{string},
    commentstyle=\color{green},
    morecomment=[s][\color{gray}]{/**}{*/},
    breaklines=true,
    numbers=left,
    numberstyle=\tiny\color{gray},
    showspaces=false,
    showstringspaces=false,
    frame=single,
    captionpos=b
}

\lstdefinelanguage{drl}{
    basicstyle=\ttfamily\small,
    keywordstyle=\color{blue}\bfseries,
    stringstyle=\color{red},
    commentstyle=\color{green},
    morecomment=[s][\color{green}]{/*}{*/},
    morecomment=[l][\color{green}]{//},
    morestring=[b]",
    keywords={rule, when, then, end, salience, eval, update, retract},
    sensitive=true,
    breaklines=true,
    numbers=left,
    numberstyle=\tiny\color{gray},
    showspaces=false,
    showstringspaces=false,
    frame=single,
    captionpos=b
}

\lstdefinestyle{drl}{
    language=drl,
    basicstyle=\ttfamily\small,
    keywordstyle=\color{blue}\bfseries,
    stringstyle=\color{red},
    commentstyle=\color{green},
    breaklines=true,
    numbers=left,
    numberstyle=\tiny\color{gray},
    showspaces=false,
    showstringspaces=false,
    frame=single,
    captionpos=b
}



\renewcommand{\lstlistingname}{Codi}
\renewcommand{\lstlistlistingname}{Índex de Fragments de Codis}
\renewcommand{\appendixpagename}{Annexos}
\renewcommand{\appendixtocname}{Annexos}

\renewcommand*{\listalgorithmcfname}{Llistat de pseudocodi}
\renewcommand*{\algorithmcfname}{Pseudocodi}
\renewcommand*{\algorithmautorefname}{pseudocodi}

\definecolor{darkgreen}{rgb}{0.0, 0.3, 0.0}
\definecolor{lightgreen}{rgb}{0.9, 0.95, 0.9}

\newcommand{\todo}[1]{
    \begin{tcolorbox}[colback=lightgreen, colframe=darkgreen, title=TODO]
        #1
    \end{tcolorbox}
}


% \usepackage{indentfirst} 
% \setlength{\parindent}{1cm}


\usepackage{parskip}
% \setlength{\parindent}{10pt}

% \setlength{\headheight}{15.6pt}
\setlength{\headheight}{16.05263pt}




\vspace{5mm}

\newcolumntype{u}{>{\centering\arraybackslash\hsize=.22\hsize}X}
\newcolumntype{s}{>{\hsize=.45\hsize}X}


%%%%%%%%%%%%%%%%%%%%%%%%%%%%%%%%
% INI creacio subsubsubsecion
%%%%%%%%%%%%%%%%%%%%%%%%%%%%%%%%
\usepackage{titlesec}

\let\oldparagraph\paragraph
\renewcommand{\paragraph}[1]{\oldparagraph{\hspace{30pt}#1}}


\titleclass{\subsubsubsection}{straight}[\subsection]

\newcounter{subsubsubsection}[subsubsection]
\renewcommand\thesubsubsubsection{\thesubsubsection.\arabic{subsubsubsection}}
\renewcommand\theparagraph{\thesubsubsubsection.\arabic{paragraph}} % optional; useful if paragraphs are to be numbered

\titleformat{\subsubsubsection}
  {\normalfont\normalsize\bfseries}{\thesubsubsubsection}{1em}{}
\titlespacing*{\subsubsubsection}
{0pt}{3.25ex plus 1ex minus .2ex}{1.5ex plus .2ex}


\makeatletter
\renewcommand\paragraph{\@startsection{paragraph}{5}{10pt}%
  {3.25ex \@plus1ex \@minus.2ex}%
  {-1em}%
  {\normalfont\normalsize\bfseries}}
\renewcommand\subparagraph{\@startsection{subparagraph}{6}{\parindent}%
  {3.25ex \@plus1ex \@minus .2ex}%
  {-1em}%
  {\normalfont\normalsize\bfseries}}
\def\toclevel@subsubsubsection{4}
\def\toclevel@paragraph{5}
\def\toclevel@paragraph{6}
\def\l@subsubsubsection{\@dottedtocline{4}{7em}{4em}}
\def\l@paragraph{\@dottedtocline{5}{10em}{5em}}
\def\l@subparagraph{\@dottedtocline{6}{14em}{6em}}
\makeatother

\setcounter{secnumdepth}{4}
\setcounter{tocdepth}{3}
%%%%%%%%%%%%%%%%%%%%%%%%%%%%%%%%
% FIN creacio subsubsubtitle
%%%%%%%%%%%%%%%%%%%%%%%%%%%%%%%%


%%%%%%%%%%%%%%%%%%%%%%%%%%%%%%%%
% INI creacio equation
%%%%%%%%%%%%%%%%%%%%%%%%%%%%%%%%
\DeclareFloatingEnvironment[fileext=loeq,placement={!ht},name=Equació]{equationfloat}
%%%%%%%%%%%%%%%%%%%%%%%%%%%%%%%%
% FIN creacio equation
%%%%%%%%%%%%%%%%%%%%%%%%%%%%%%%%

\titleformat{\section}[block]
  {\clearpage\normalfont\Large\bfseries}
  {\thesection}{1em}{}


\renewcommand\refname{Bibliografia}


\linespread{1.25}


\makeatletter
\renewenvironment{abstract}{%
    \if@twocolumn
      \section*{\abstractname}%
    \else 
      \begin{center}%
        {\bfseries \Large\itshape\abstractname\vspace{\z@}}%  % Título en negrita, 14pt y cursiva
      \end{center}%
      \quotation
    \fi}
    {\if@twocolumn\else\endquotation\fi}
\makeatother

% \fancypagestyle{main}{
%     \fancyhf{}
%     \lhead{\small Aplicació Web per a la Visualització i Comparació
% de Trajectòries d’Algorismes d’Optimització}
%     % \rhead{\includegraphics[width=1.2cm]{resources/Icons/logo-fib.png}}
%     \cfoot{\thepage}
%     \setlength{\footskip}{35pt}
% }

\newcommand{\q}[1]{``#1''}

\begin{document}

\hypersetup{
    pdfborder=0 0 0,
}

%% PORTADA I TABLE OF CONTENTS %%

%%%%                                   %%%%
%         INDEX DE LA MEMORIA             %
%%%%                                   %%%%

% \includepdf[pages=-]{portada_fib.pdf}
\input{Extra/GEP}
%\input{Extra/portada}
% \cleardoublepage

% \pagestyle{main}
% \setcounter{page}{11}
\include{toc}
% \cleardoublepage
\input{Contexto y alcance/main}

\input{Planificación temporal/main}

\input{Gestión económica/main}

\input{Informe de sostenibilidad/main}

% \cleardoublepage

% \cleardoublepage


% \input{Annex/main}


\clearpage
% Estilo bibliográfico IEEE
\bibliographystyle{IEEEtran}
\renewcommand{\refname}{Referencias}
\bibliography{references.bib}

%%%%                                   %%%%
%         FINAL DE LA MEMORIA             %
%%%%                                   %%%%

\end{document}


%% VERSIÓ INCIAL MEMORIA %%




\clearpage
% Estilo bibliográfico IEEE
\bibliographystyle{IEEEtran}
\renewcommand{\refname}{Referencias}
\bibliography{references.bib}

%%%%                                   %%%%
%         FINAL DE LA MEMORIA             %
%%%%                                   %%%%

\end{document}


%% VERSIÓ INCIAL MEMORIA %%


\documentclass[a4paper, 11pt, twoside]{article}
\usepackage{wrapfig}
% \usepackage{abstract}
\usepackage{graphicx} % Required for inserting images
\usepackage[a4paper, left=1.15in, right=1.15in, top=1.5in, bottom=1.2in]{geometry}
\usepackage{float}
\usepackage{array}
\usepackage{booktabs}
\usepackage[T1]{fontenc}
\usepackage{amsmath}
\usepackage{tabularx}
\usepackage{graphicx} % Required for inserting images
\usepackage{tikz, xcolor}
\usepackage{hyperref}
\usepackage{cite}
\usepackage{pdfpages}
\usepackage{pdflscape}
\usepackage{afterpage}
\usepackage{textpos}
\usepackage{fancyhdr}
\usepackage{caption}
\usepackage{lmodern,textcomp}
\usepackage[catalan]{babel}
\usepackage{lipsum}
\usepackage{newfloat}
\usepackage{listings}
\usepackage{placeins}
\usepackage[ruled]{algorithm2e}
\usepackage{tcolorbox}
% Estilo de citas IEEE sencillo: usar paquete cite para comprimir y ordenar referencias numéricas
\usepackage{cite}
\usepackage[toc,page]{appendix}
\usepackage{forest}
\usepackage{pdflscape}
\usepackage{graphicx}
\usepackage{geometry}
\usepackage{quoting,xparse}
\usepackage{svg}



% \NewDocumentCommand{\bywhom}{m}{% the Bourbaki trick
%   {\nobreak\hfill\penalty50\hskip1em\null\nobreak
%    \hfill\mbox{\normalfont(#1)}%
%    \parfillskip=0pt \finalhyphendemerits=0 \par}%
% }

% AIXÒ SERVEIX PER A MODIFICAR UN PROBLEMA AMB ELS CARACTERS --> " + PARAULA 

% \NewDocumentEnvironment{pquotation}{m}
%   {\begin{quoting}[
%      indentfirst=true,
%      leftmargin=\parindent,
%      rightmargin=\parindent]\itshape}
%   {\bywhom{#1}
  
%   \end{quoting}}

% AIXÒ PER INSERIR FRAGMENTS DE CODI

\lstnewenvironment{code}[1][]%
{
   \noindent
   \minipage{\linewidth}
   \vspace{0.5\baselineskip}
   \lstset{
     basicstyle=\ttfamily\footnotesize,
     frame=single,
     breaklines=true,    % Permitir que las líneas se rompan automáticamente
     breakatwhitespace=true, % Permitir la ruptura en los espacios en blanco
     #1
   }
}
{\endminipage}

\lstset{
    language=Java,
    basicstyle=\ttfamily\small,  % Define el tipo de letra y tamaño del código
    keywordstyle=\color{blue},   % Define el color de las palabras clave de Java
    stringstyle=\color{red},     % Define el color de las cadenas de texto
    commentstyle=\color{green},  % Define el color de los comentarios
    morecomment=[s][\color{gray}]{/**}{*/},  % Estilo para comentarios de bloque
    breaklines=true,  % Configura el quiebre automático de línea
    numbers=left,  % Posición de los números de línea
    numberstyle=\tiny\color{gray},  % Estilo para los números de línea
    showspaces=false,  % No mostrar espacios con underscores
    showstringspaces=false,  % No mostrar espacios dentro de strings
    frame=single,  % Agregar un marco alrededor del código
    captionpos=b,  % b = bottom, t = top
}

\definecolor{eclipseStrings}{RGB}{42,0.0,255}
\definecolor{eclipseKeywords}{RGB}{127,0,85}
\colorlet{numb}{magenta!60!black}

\lstdefinelanguage{json}{
    basicstyle=\ttfamily\small,
    showstringspaces=false,
    breaklines=true,
    frame=single,
    literate=
     *{0}{{{\color{numb}0}}}{1}
      {1}{{{\color{numb}1}}}{1}
      {2}{{{\color{numb}2}}}{1}
      {3}{{{\color{numb}3}}}{1}
      {4}{{{\color{numb}4}}}{1}
      {5}{{{\color{numb}5}}}{1}
      {6}{{{\color{numb}6}}}{1}
      {7}{{{\color{numb}7}}}{1}
      {8}{{{\color{numb}8}}}{1}
      {9}{{{\color{numb}9}}}{1}
      {:}{{{\color{punct}{:}}}}{1}
      {,}{{{\color{punct}{,}}}}{1}
      {\{}{{{\color{delim}{\{}}}}{1}
      {\}}{{{\color{delim}{\}}}}}{1}
      {[}{{{\color{delim}{[}}}}{1}
      {]}{{{\color{delim}{]}}}}{1},
    morestring=[b]"
}

\lstdefinelanguage{JavaScript}{
  morekeywords=[1]{break, continue, delete, else, for, function, if, in,
    new, return, this, typeof, var, void, while, with},
  % Literals, primitive types, and reference types.
  morekeywords=[2]{false, null, true, boolean, number, undefined,
    Array, Boolean, Date, Math, Number, String, Object},
  % Built-ins.
  morekeywords=[3]{eval, parseInt, parseFloat, escape, unescape},
  sensitive,
  morecomment=[s]{/*}{*/},
  morecomment=[l]//,
  morecomment=[s]{/**}{*/}, % JavaDoc style comments
  morestring=[b]',
  morestring=[b]"
}[keywords, comments, strings]

\lstdefinelanguage{svelte}{
    morekeywords={async, await, then, if, each, in, on, let, const, function, import, export, from, with, typeof, new, var, return},
    sensitive=true,
    morecomment=[l]{//},
    morecomment=[s]{/*}{*/},
    morecomment=[s]{<!--}{-->},
    morestring=[b]',
    morestring=[b]",
    morestring=[b]`,
    ndkeywords={import, export, return, if, for, while, function, switch, case, break, var, let, const, async, await, then, catch, try, typeof, new, class, extends, implements, interface, constructor, public, private, protected, static, get, set},
    alsoletter={:},
    moredelim=**[is][\color{orange}]{\@\{}{\@\}},
    moredelim=**[is][\color{cyan}]{\$\{}{\}},
    keywordstyle=\color{blue}\bfseries,
    commentstyle=\color{gray}\ttfamily,
    stringstyle=\color{red}\ttfamily,
    identifierstyle=\color{black},
    breaklines=true,
    basicstyle=\ttfamily\small,
    columns=fullflexible,
    showstringspaces=false,
    tabsize=2,
    frame=single,
    captionpos=b,
}


\lstdefinestyle{java}{
    language=Java,
    basicstyle=\ttfamily\small,
    keywordstyle=\color{blue},
    stringstyle=\color{red},
    commentstyle=\color{gray},
    morecomment=[s][\color{gray}]{/**}{*/},
    breaklines=true,
    numbers=left,
    numberstyle=\tiny\color{gray},
    showspaces=false,
    showstringspaces=false,
    frame=single,
    captionpos=b
}

\usepackage{color}
\definecolor{numb}{rgb}{0.58,0,0.82}
\definecolor{punct}{rgb}{0.3,0.3,0.3}
\definecolor{delim}{rgb}{0.5,0.5,0.5}
\definecolor{string}{rgb}{0.6,0.1,0.1}

\lstdefinestyle{json}{
    language=json,
    basicstyle=\ttfamily\small,
    keywordstyle=\color{blue},
    stringstyle=\color{string},
    commentstyle=\color{green},
    morecomment=[s][\color{gray}]{/**}{*/},
    breaklines=true,
    numbers=left,
    numberstyle=\tiny\color{gray},
    showspaces=false,
    showstringspaces=false,
    frame=single,
    captionpos=b
}

\lstdefinelanguage{drl}{
    basicstyle=\ttfamily\small,
    keywordstyle=\color{blue}\bfseries,
    stringstyle=\color{red},
    commentstyle=\color{green},
    morecomment=[s][\color{green}]{/*}{*/},
    morecomment=[l][\color{green}]{//},
    morestring=[b]",
    keywords={rule, when, then, end, salience, eval, update, retract},
    sensitive=true,
    breaklines=true,
    numbers=left,
    numberstyle=\tiny\color{gray},
    showspaces=false,
    showstringspaces=false,
    frame=single,
    captionpos=b
}

\lstdefinestyle{drl}{
    language=drl,
    basicstyle=\ttfamily\small,
    keywordstyle=\color{blue}\bfseries,
    stringstyle=\color{red},
    commentstyle=\color{green},
    breaklines=true,
    numbers=left,
    numberstyle=\tiny\color{gray},
    showspaces=false,
    showstringspaces=false,
    frame=single,
    captionpos=b
}



\renewcommand{\lstlistingname}{Codi}
\renewcommand{\lstlistlistingname}{Índex de Fragments de Codis}
\renewcommand{\appendixpagename}{Annexos}
\renewcommand{\appendixtocname}{Annexos}

\renewcommand*{\listalgorithmcfname}{Llistat de pseudocodi}
\renewcommand*{\algorithmcfname}{Pseudocodi}
\renewcommand*{\algorithmautorefname}{pseudocodi}

\definecolor{darkgreen}{rgb}{0.0, 0.3, 0.0}
\definecolor{lightgreen}{rgb}{0.9, 0.95, 0.9}

\newcommand{\todo}[1]{
    \begin{tcolorbox}[colback=lightgreen, colframe=darkgreen, title=TODO]
        #1
    \end{tcolorbox}
}


% \usepackage{indentfirst} 
% \setlength{\parindent}{1cm}


\usepackage{parskip}
% \setlength{\parindent}{10pt}

% \setlength{\headheight}{15.6pt}
\setlength{\headheight}{16.05263pt}




\vspace{5mm}

\newcolumntype{u}{>{\centering\arraybackslash\hsize=.22\hsize}X}
\newcolumntype{s}{>{\hsize=.45\hsize}X}


%%%%%%%%%%%%%%%%%%%%%%%%%%%%%%%%
% INI creacio subsubsubsecion
%%%%%%%%%%%%%%%%%%%%%%%%%%%%%%%%
\usepackage{titlesec}

\let\oldparagraph\paragraph
\renewcommand{\paragraph}[1]{\oldparagraph{\hspace{30pt}#1}}


\titleclass{\subsubsubsection}{straight}[\subsection]

\newcounter{subsubsubsection}[subsubsection]
\renewcommand\thesubsubsubsection{\thesubsubsection.\arabic{subsubsubsection}}
\renewcommand\theparagraph{\thesubsubsubsection.\arabic{paragraph}} % optional; useful if paragraphs are to be numbered

\titleformat{\subsubsubsection}
  {\normalfont\normalsize\bfseries}{\thesubsubsubsection}{1em}{}
\titlespacing*{\subsubsubsection}
{0pt}{3.25ex plus 1ex minus .2ex}{1.5ex plus .2ex}


\makeatletter
\renewcommand\paragraph{\@startsection{paragraph}{5}{10pt}%
  {3.25ex \@plus1ex \@minus.2ex}%
  {-1em}%
  {\normalfont\normalsize\bfseries}}
\renewcommand\subparagraph{\@startsection{subparagraph}{6}{\parindent}%
  {3.25ex \@plus1ex \@minus .2ex}%
  {-1em}%
  {\normalfont\normalsize\bfseries}}
\def\toclevel@subsubsubsection{4}
\def\toclevel@paragraph{5}
\def\toclevel@paragraph{6}
\def\l@subsubsubsection{\@dottedtocline{4}{7em}{4em}}
\def\l@paragraph{\@dottedtocline{5}{10em}{5em}}
\def\l@subparagraph{\@dottedtocline{6}{14em}{6em}}
\makeatother

\setcounter{secnumdepth}{4}
\setcounter{tocdepth}{3}
%%%%%%%%%%%%%%%%%%%%%%%%%%%%%%%%
% FIN creacio subsubsubtitle
%%%%%%%%%%%%%%%%%%%%%%%%%%%%%%%%


%%%%%%%%%%%%%%%%%%%%%%%%%%%%%%%%
% INI creacio equation
%%%%%%%%%%%%%%%%%%%%%%%%%%%%%%%%
\DeclareFloatingEnvironment[fileext=loeq,placement={!ht},name=Equació]{equationfloat}
%%%%%%%%%%%%%%%%%%%%%%%%%%%%%%%%
% FIN creacio equation
%%%%%%%%%%%%%%%%%%%%%%%%%%%%%%%%

\titleformat{\section}[block]
  {\clearpage\normalfont\Large\bfseries}
  {\thesection}{1em}{}


\renewcommand\refname{Bibliografia}


\linespread{1.25}


\makeatletter
\renewenvironment{abstract}{%
    \if@twocolumn
      \section*{\abstractname}%
    \else 
      \begin{center}%
        {\bfseries \Large\itshape\abstractname\vspace{\z@}}%  % Título en negrita, 14pt y cursiva
      \end{center}%
      \quotation
    \fi}
    {\if@twocolumn\else\endquotation\fi}
\makeatother

% \fancypagestyle{main}{
%     \fancyhf{}
%     \lhead{\small Aplicació Web per a la Visualització i Comparació
% de Trajectòries d’Algorismes d’Optimització}
%     % \rhead{\includegraphics[width=1.2cm]{resources/Icons/logo-fib.png}}
%     \cfoot{\thepage}
%     \setlength{\footskip}{35pt}
% }

\newcommand{\q}[1]{``#1''}

\begin{document}

\hypersetup{
    pdfborder=0 0 0,
}

%% PORTADA I TABLE OF CONTENTS %%

%%%%                                   %%%%
%         INDEX DE LA MEMORIA             %
%%%%                                   %%%%

% \includepdf[pages=-]{portada_fib.pdf}
%%% Portada
\begin{titlepage}
\begin{center}
    \includegraphics[width=1\linewidth]{extra/figures/logo-upc-fib.png}
    \vspace{30mm}

    \LARGE{\textsc{\textbf{{INTEGRACIÓN Y AUTOMATIZACIÓN DE COSTES DE INFRAESTRUCTURA HPC EN ENTORNOS DE INVESTIGACIÓN}}}} \\
    \vspace{25mm}
      
        
    \textbf{\Large{\textsc{PAU NAVARRO ÁLVAREZ}}} \\
    \vspace{20mm}
        
    
    
    \large{\textbf{Director/a} ROBERTO RIVEIRO INSUA (FUNDACIO INSTITUT DE RECERCA BIOMEDICA (IRB BARCELONA))} \\
    \large{\textbf{Ponente:} ALEX PAJUELO GONZÁLEZ (Departamento de Arquitectura de Computadores)} \\
    \large{\textbf{Titulación:} Grado en Ingeniería Informática (Tecnologías de la información)} \\
    
    \large{\textbf{Lliurament 2:} }Planificación temporal \\
    \large{\textbf{Tutor GEP:} CAROLINA CONSOLIDACIÓN SEGURA (Departament d'Organització d'Empreses)} \\
    \large{\textbf{Fecha:}} 29 de Septiembre de 2025\\

    \textbf{\large{Memoria del trabajo de fin de grado}} \\
    \vspace{3mm}
    \textbf{\large{Facultat d'Informàtica de Barcelona (FIB)}} \\
    \vspace{3mm}
    \textbf{\large{Universitat Politècnica de Catalunya (UPC) - BarcelonaTech}} \\

\end{center}
\end{titlepage}
%\input{Extra/portada}
% \cleardoublepage

% \pagestyle{main}
% \setcounter{page}{11}
\renewcommand{\contentsname}{Continguts}

\tableofcontents
\clearpage
\listoffigures
\begingroup
    \let\clearpage\relax
    \listoftables
\endgroup
\clearpage

\newpage
% \cleardoublepage
\input{Contexto y alcance/Contextualización/main}
\input{Contexto y alcance/Justificación/main}
\input{Contexto y alcance/Alcance/main}
\input{Contexto y alcance/Metodología y herramientas de seguimiento/main}

\documentclass[a4paper, 11pt, twoside]{article}
\usepackage{wrapfig}
% \usepackage{abstract}
\usepackage{graphicx} % Required for inserting images
\usepackage[a4paper, left=1.15in, right=1.15in, top=1.5in, bottom=1.2in]{geometry}
\usepackage{float}
\usepackage{array}
\usepackage{booktabs}
\usepackage[T1]{fontenc}
\usepackage{amsmath}
\usepackage{tabularx}
\usepackage{graphicx} % Required for inserting images
\usepackage{tikz, xcolor}
\usepackage{hyperref}
\usepackage{cite}
\usepackage{pdfpages}
\usepackage{pdflscape}
\usepackage{afterpage}
\usepackage{textpos}
\usepackage{fancyhdr}
\usepackage{caption}
\usepackage{lmodern,textcomp}
\usepackage[catalan]{babel}
\usepackage{lipsum}
\usepackage{newfloat}
\usepackage{listings}
\usepackage{placeins}
\usepackage[ruled]{algorithm2e}
\usepackage{tcolorbox}
% Estilo de citas IEEE sencillo: usar paquete cite para comprimir y ordenar referencias numéricas
\usepackage{cite}
\usepackage[toc,page]{appendix}
\usepackage{forest}
\usepackage{pdflscape}
\usepackage{graphicx}
\usepackage{geometry}
\usepackage{quoting,xparse}
\usepackage{svg}



% \NewDocumentCommand{\bywhom}{m}{% the Bourbaki trick
%   {\nobreak\hfill\penalty50\hskip1em\null\nobreak
%    \hfill\mbox{\normalfont(#1)}%
%    \parfillskip=0pt \finalhyphendemerits=0 \par}%
% }

% AIXÒ SERVEIX PER A MODIFICAR UN PROBLEMA AMB ELS CARACTERS --> " + PARAULA 

% \NewDocumentEnvironment{pquotation}{m}
%   {\begin{quoting}[
%      indentfirst=true,
%      leftmargin=\parindent,
%      rightmargin=\parindent]\itshape}
%   {\bywhom{#1}
  
%   \end{quoting}}

% AIXÒ PER INSERIR FRAGMENTS DE CODI

\lstnewenvironment{code}[1][]%
{
   \noindent
   \minipage{\linewidth}
   \vspace{0.5\baselineskip}
   \lstset{
     basicstyle=\ttfamily\footnotesize,
     frame=single,
     breaklines=true,    % Permitir que las líneas se rompan automáticamente
     breakatwhitespace=true, % Permitir la ruptura en los espacios en blanco
     #1
   }
}
{\endminipage}

\lstset{
    language=Java,
    basicstyle=\ttfamily\small,  % Define el tipo de letra y tamaño del código
    keywordstyle=\color{blue},   % Define el color de las palabras clave de Java
    stringstyle=\color{red},     % Define el color de las cadenas de texto
    commentstyle=\color{green},  % Define el color de los comentarios
    morecomment=[s][\color{gray}]{/**}{*/},  % Estilo para comentarios de bloque
    breaklines=true,  % Configura el quiebre automático de línea
    numbers=left,  % Posición de los números de línea
    numberstyle=\tiny\color{gray},  % Estilo para los números de línea
    showspaces=false,  % No mostrar espacios con underscores
    showstringspaces=false,  % No mostrar espacios dentro de strings
    frame=single,  % Agregar un marco alrededor del código
    captionpos=b,  % b = bottom, t = top
}

\definecolor{eclipseStrings}{RGB}{42,0.0,255}
\definecolor{eclipseKeywords}{RGB}{127,0,85}
\colorlet{numb}{magenta!60!black}

\lstdefinelanguage{json}{
    basicstyle=\ttfamily\small,
    showstringspaces=false,
    breaklines=true,
    frame=single,
    literate=
     *{0}{{{\color{numb}0}}}{1}
      {1}{{{\color{numb}1}}}{1}
      {2}{{{\color{numb}2}}}{1}
      {3}{{{\color{numb}3}}}{1}
      {4}{{{\color{numb}4}}}{1}
      {5}{{{\color{numb}5}}}{1}
      {6}{{{\color{numb}6}}}{1}
      {7}{{{\color{numb}7}}}{1}
      {8}{{{\color{numb}8}}}{1}
      {9}{{{\color{numb}9}}}{1}
      {:}{{{\color{punct}{:}}}}{1}
      {,}{{{\color{punct}{,}}}}{1}
      {\{}{{{\color{delim}{\{}}}}{1}
      {\}}{{{\color{delim}{\}}}}}{1}
      {[}{{{\color{delim}{[}}}}{1}
      {]}{{{\color{delim}{]}}}}{1},
    morestring=[b]"
}

\lstdefinelanguage{JavaScript}{
  morekeywords=[1]{break, continue, delete, else, for, function, if, in,
    new, return, this, typeof, var, void, while, with},
  % Literals, primitive types, and reference types.
  morekeywords=[2]{false, null, true, boolean, number, undefined,
    Array, Boolean, Date, Math, Number, String, Object},
  % Built-ins.
  morekeywords=[3]{eval, parseInt, parseFloat, escape, unescape},
  sensitive,
  morecomment=[s]{/*}{*/},
  morecomment=[l]//,
  morecomment=[s]{/**}{*/}, % JavaDoc style comments
  morestring=[b]',
  morestring=[b]"
}[keywords, comments, strings]

\lstdefinelanguage{svelte}{
    morekeywords={async, await, then, if, each, in, on, let, const, function, import, export, from, with, typeof, new, var, return},
    sensitive=true,
    morecomment=[l]{//},
    morecomment=[s]{/*}{*/},
    morecomment=[s]{<!--}{-->},
    morestring=[b]',
    morestring=[b]",
    morestring=[b]`,
    ndkeywords={import, export, return, if, for, while, function, switch, case, break, var, let, const, async, await, then, catch, try, typeof, new, class, extends, implements, interface, constructor, public, private, protected, static, get, set},
    alsoletter={:},
    moredelim=**[is][\color{orange}]{\@\{}{\@\}},
    moredelim=**[is][\color{cyan}]{\$\{}{\}},
    keywordstyle=\color{blue}\bfseries,
    commentstyle=\color{gray}\ttfamily,
    stringstyle=\color{red}\ttfamily,
    identifierstyle=\color{black},
    breaklines=true,
    basicstyle=\ttfamily\small,
    columns=fullflexible,
    showstringspaces=false,
    tabsize=2,
    frame=single,
    captionpos=b,
}


\lstdefinestyle{java}{
    language=Java,
    basicstyle=\ttfamily\small,
    keywordstyle=\color{blue},
    stringstyle=\color{red},
    commentstyle=\color{gray},
    morecomment=[s][\color{gray}]{/**}{*/},
    breaklines=true,
    numbers=left,
    numberstyle=\tiny\color{gray},
    showspaces=false,
    showstringspaces=false,
    frame=single,
    captionpos=b
}

\usepackage{color}
\definecolor{numb}{rgb}{0.58,0,0.82}
\definecolor{punct}{rgb}{0.3,0.3,0.3}
\definecolor{delim}{rgb}{0.5,0.5,0.5}
\definecolor{string}{rgb}{0.6,0.1,0.1}

\lstdefinestyle{json}{
    language=json,
    basicstyle=\ttfamily\small,
    keywordstyle=\color{blue},
    stringstyle=\color{string},
    commentstyle=\color{green},
    morecomment=[s][\color{gray}]{/**}{*/},
    breaklines=true,
    numbers=left,
    numberstyle=\tiny\color{gray},
    showspaces=false,
    showstringspaces=false,
    frame=single,
    captionpos=b
}

\lstdefinelanguage{drl}{
    basicstyle=\ttfamily\small,
    keywordstyle=\color{blue}\bfseries,
    stringstyle=\color{red},
    commentstyle=\color{green},
    morecomment=[s][\color{green}]{/*}{*/},
    morecomment=[l][\color{green}]{//},
    morestring=[b]",
    keywords={rule, when, then, end, salience, eval, update, retract},
    sensitive=true,
    breaklines=true,
    numbers=left,
    numberstyle=\tiny\color{gray},
    showspaces=false,
    showstringspaces=false,
    frame=single,
    captionpos=b
}

\lstdefinestyle{drl}{
    language=drl,
    basicstyle=\ttfamily\small,
    keywordstyle=\color{blue}\bfseries,
    stringstyle=\color{red},
    commentstyle=\color{green},
    breaklines=true,
    numbers=left,
    numberstyle=\tiny\color{gray},
    showspaces=false,
    showstringspaces=false,
    frame=single,
    captionpos=b
}



\renewcommand{\lstlistingname}{Codi}
\renewcommand{\lstlistlistingname}{Índex de Fragments de Codis}
\renewcommand{\appendixpagename}{Annexos}
\renewcommand{\appendixtocname}{Annexos}

\renewcommand*{\listalgorithmcfname}{Llistat de pseudocodi}
\renewcommand*{\algorithmcfname}{Pseudocodi}
\renewcommand*{\algorithmautorefname}{pseudocodi}

\definecolor{darkgreen}{rgb}{0.0, 0.3, 0.0}
\definecolor{lightgreen}{rgb}{0.9, 0.95, 0.9}

\newcommand{\todo}[1]{
    \begin{tcolorbox}[colback=lightgreen, colframe=darkgreen, title=TODO]
        #1
    \end{tcolorbox}
}


% \usepackage{indentfirst} 
% \setlength{\parindent}{1cm}


\usepackage{parskip}
% \setlength{\parindent}{10pt}

% \setlength{\headheight}{15.6pt}
\setlength{\headheight}{16.05263pt}




\vspace{5mm}

\newcolumntype{u}{>{\centering\arraybackslash\hsize=.22\hsize}X}
\newcolumntype{s}{>{\hsize=.45\hsize}X}


%%%%%%%%%%%%%%%%%%%%%%%%%%%%%%%%
% INI creacio subsubsubsecion
%%%%%%%%%%%%%%%%%%%%%%%%%%%%%%%%
\usepackage{titlesec}

\let\oldparagraph\paragraph
\renewcommand{\paragraph}[1]{\oldparagraph{\hspace{30pt}#1}}


\titleclass{\subsubsubsection}{straight}[\subsection]

\newcounter{subsubsubsection}[subsubsection]
\renewcommand\thesubsubsubsection{\thesubsubsection.\arabic{subsubsubsection}}
\renewcommand\theparagraph{\thesubsubsubsection.\arabic{paragraph}} % optional; useful if paragraphs are to be numbered

\titleformat{\subsubsubsection}
  {\normalfont\normalsize\bfseries}{\thesubsubsubsection}{1em}{}
\titlespacing*{\subsubsubsection}
{0pt}{3.25ex plus 1ex minus .2ex}{1.5ex plus .2ex}


\makeatletter
\renewcommand\paragraph{\@startsection{paragraph}{5}{10pt}%
  {3.25ex \@plus1ex \@minus.2ex}%
  {-1em}%
  {\normalfont\normalsize\bfseries}}
\renewcommand\subparagraph{\@startsection{subparagraph}{6}{\parindent}%
  {3.25ex \@plus1ex \@minus .2ex}%
  {-1em}%
  {\normalfont\normalsize\bfseries}}
\def\toclevel@subsubsubsection{4}
\def\toclevel@paragraph{5}
\def\toclevel@paragraph{6}
\def\l@subsubsubsection{\@dottedtocline{4}{7em}{4em}}
\def\l@paragraph{\@dottedtocline{5}{10em}{5em}}
\def\l@subparagraph{\@dottedtocline{6}{14em}{6em}}
\makeatother

\setcounter{secnumdepth}{4}
\setcounter{tocdepth}{3}
%%%%%%%%%%%%%%%%%%%%%%%%%%%%%%%%
% FIN creacio subsubsubtitle
%%%%%%%%%%%%%%%%%%%%%%%%%%%%%%%%


%%%%%%%%%%%%%%%%%%%%%%%%%%%%%%%%
% INI creacio equation
%%%%%%%%%%%%%%%%%%%%%%%%%%%%%%%%
\DeclareFloatingEnvironment[fileext=loeq,placement={!ht},name=Equació]{equationfloat}
%%%%%%%%%%%%%%%%%%%%%%%%%%%%%%%%
% FIN creacio equation
%%%%%%%%%%%%%%%%%%%%%%%%%%%%%%%%

\titleformat{\section}[block]
  {\clearpage\normalfont\Large\bfseries}
  {\thesection}{1em}{}


\renewcommand\refname{Bibliografia}


\linespread{1.25}


\makeatletter
\renewenvironment{abstract}{%
    \if@twocolumn
      \section*{\abstractname}%
    \else 
      \begin{center}%
        {\bfseries \Large\itshape\abstractname\vspace{\z@}}%  % Título en negrita, 14pt y cursiva
      \end{center}%
      \quotation
    \fi}
    {\if@twocolumn\else\endquotation\fi}
\makeatother

% \fancypagestyle{main}{
%     \fancyhf{}
%     \lhead{\small Aplicació Web per a la Visualització i Comparació
% de Trajectòries d’Algorismes d’Optimització}
%     % \rhead{\includegraphics[width=1.2cm]{resources/Icons/logo-fib.png}}
%     \cfoot{\thepage}
%     \setlength{\footskip}{35pt}
% }

\newcommand{\q}[1]{``#1''}

\begin{document}

\hypersetup{
    pdfborder=0 0 0,
}

%% PORTADA I TABLE OF CONTENTS %%

%%%%                                   %%%%
%         INDEX DE LA MEMORIA             %
%%%%                                   %%%%

% \includepdf[pages=-]{portada_fib.pdf}
\input{Extra/GEP}
%\input{Extra/portada}
% \cleardoublepage

% \pagestyle{main}
% \setcounter{page}{11}
\include{toc}
% \cleardoublepage
\input{Contexto y alcance/main}

\input{Planificación temporal/main}

\input{Gestión económica/main}

\input{Informe de sostenibilidad/main}

% \cleardoublepage

% \cleardoublepage


% \input{Annex/main}


\clearpage
% Estilo bibliográfico IEEE
\bibliographystyle{IEEEtran}
\renewcommand{\refname}{Referencias}
\bibliography{references.bib}

%%%%                                   %%%%
%         FINAL DE LA MEMORIA             %
%%%%                                   %%%%

\end{document}


%% VERSIÓ INCIAL MEMORIA %%



\section{Gestión económica}

La gestión económica de un proyecto es fundamental para asegurar su viabilidad y exito a largo plazo. Ya hemos visto en secciones anteriores, los otros dos vertices 
del triángulo de la gestión de proyectos: el alcance y el tiempo. En esta sección, nos centraremos en cómo gestionar eficazmente los recursos financieros para garantizar
que el proyecto se mantenga dentro del presupuesto y se logren los objetivos establecidos. 

\input{Gestión económica/Recursos humanos.tex}

\input{Gestión económica/Recursos Hardware.tex}

\input{Gestión económica/Recursos Software.tex}

\input{Gestión económica/Contingencias.tex}

\section{Informe de sostenibilidad}

\input{Informe de sostenibilidad/Autoevaluación.tex}


\input{Informe de sostenibilidad/Dimensión ambiental.tex}

\input{Informe de sostenibilidad/Dimensión económica.tex}

\input{Informe de sostenibilidad/Dimensión social.tex}


% \cleardoublepage

% \cleardoublepage


% \documentclass[a4paper, 11pt, twoside]{article}
\usepackage{wrapfig}
% \usepackage{abstract}
\usepackage{graphicx} % Required for inserting images
\usepackage[a4paper, left=1.15in, right=1.15in, top=1.5in, bottom=1.2in]{geometry}
\usepackage{float}
\usepackage{array}
\usepackage{booktabs}
\usepackage[T1]{fontenc}
\usepackage{amsmath}
\usepackage{tabularx}
\usepackage{graphicx} % Required for inserting images
\usepackage{tikz, xcolor}
\usepackage{hyperref}
\usepackage{cite}
\usepackage{pdfpages}
\usepackage{pdflscape}
\usepackage{afterpage}
\usepackage{textpos}
\usepackage{fancyhdr}
\usepackage{caption}
\usepackage{lmodern,textcomp}
\usepackage[catalan]{babel}
\usepackage{lipsum}
\usepackage{newfloat}
\usepackage{listings}
\usepackage{placeins}
\usepackage[ruled]{algorithm2e}
\usepackage{tcolorbox}
% Estilo de citas IEEE sencillo: usar paquete cite para comprimir y ordenar referencias numéricas
\usepackage{cite}
\usepackage[toc,page]{appendix}
\usepackage{forest}
\usepackage{pdflscape}
\usepackage{graphicx}
\usepackage{geometry}
\usepackage{quoting,xparse}
\usepackage{svg}



% \NewDocumentCommand{\bywhom}{m}{% the Bourbaki trick
%   {\nobreak\hfill\penalty50\hskip1em\null\nobreak
%    \hfill\mbox{\normalfont(#1)}%
%    \parfillskip=0pt \finalhyphendemerits=0 \par}%
% }

% AIXÒ SERVEIX PER A MODIFICAR UN PROBLEMA AMB ELS CARACTERS --> " + PARAULA 

% \NewDocumentEnvironment{pquotation}{m}
%   {\begin{quoting}[
%      indentfirst=true,
%      leftmargin=\parindent,
%      rightmargin=\parindent]\itshape}
%   {\bywhom{#1}
  
%   \end{quoting}}

% AIXÒ PER INSERIR FRAGMENTS DE CODI

\lstnewenvironment{code}[1][]%
{
   \noindent
   \minipage{\linewidth}
   \vspace{0.5\baselineskip}
   \lstset{
     basicstyle=\ttfamily\footnotesize,
     frame=single,
     breaklines=true,    % Permitir que las líneas se rompan automáticamente
     breakatwhitespace=true, % Permitir la ruptura en los espacios en blanco
     #1
   }
}
{\endminipage}

\lstset{
    language=Java,
    basicstyle=\ttfamily\small,  % Define el tipo de letra y tamaño del código
    keywordstyle=\color{blue},   % Define el color de las palabras clave de Java
    stringstyle=\color{red},     % Define el color de las cadenas de texto
    commentstyle=\color{green},  % Define el color de los comentarios
    morecomment=[s][\color{gray}]{/**}{*/},  % Estilo para comentarios de bloque
    breaklines=true,  % Configura el quiebre automático de línea
    numbers=left,  % Posición de los números de línea
    numberstyle=\tiny\color{gray},  % Estilo para los números de línea
    showspaces=false,  % No mostrar espacios con underscores
    showstringspaces=false,  % No mostrar espacios dentro de strings
    frame=single,  % Agregar un marco alrededor del código
    captionpos=b,  % b = bottom, t = top
}

\definecolor{eclipseStrings}{RGB}{42,0.0,255}
\definecolor{eclipseKeywords}{RGB}{127,0,85}
\colorlet{numb}{magenta!60!black}

\lstdefinelanguage{json}{
    basicstyle=\ttfamily\small,
    showstringspaces=false,
    breaklines=true,
    frame=single,
    literate=
     *{0}{{{\color{numb}0}}}{1}
      {1}{{{\color{numb}1}}}{1}
      {2}{{{\color{numb}2}}}{1}
      {3}{{{\color{numb}3}}}{1}
      {4}{{{\color{numb}4}}}{1}
      {5}{{{\color{numb}5}}}{1}
      {6}{{{\color{numb}6}}}{1}
      {7}{{{\color{numb}7}}}{1}
      {8}{{{\color{numb}8}}}{1}
      {9}{{{\color{numb}9}}}{1}
      {:}{{{\color{punct}{:}}}}{1}
      {,}{{{\color{punct}{,}}}}{1}
      {\{}{{{\color{delim}{\{}}}}{1}
      {\}}{{{\color{delim}{\}}}}}{1}
      {[}{{{\color{delim}{[}}}}{1}
      {]}{{{\color{delim}{]}}}}{1},
    morestring=[b]"
}

\lstdefinelanguage{JavaScript}{
  morekeywords=[1]{break, continue, delete, else, for, function, if, in,
    new, return, this, typeof, var, void, while, with},
  % Literals, primitive types, and reference types.
  morekeywords=[2]{false, null, true, boolean, number, undefined,
    Array, Boolean, Date, Math, Number, String, Object},
  % Built-ins.
  morekeywords=[3]{eval, parseInt, parseFloat, escape, unescape},
  sensitive,
  morecomment=[s]{/*}{*/},
  morecomment=[l]//,
  morecomment=[s]{/**}{*/}, % JavaDoc style comments
  morestring=[b]',
  morestring=[b]"
}[keywords, comments, strings]

\lstdefinelanguage{svelte}{
    morekeywords={async, await, then, if, each, in, on, let, const, function, import, export, from, with, typeof, new, var, return},
    sensitive=true,
    morecomment=[l]{//},
    morecomment=[s]{/*}{*/},
    morecomment=[s]{<!--}{-->},
    morestring=[b]',
    morestring=[b]",
    morestring=[b]`,
    ndkeywords={import, export, return, if, for, while, function, switch, case, break, var, let, const, async, await, then, catch, try, typeof, new, class, extends, implements, interface, constructor, public, private, protected, static, get, set},
    alsoletter={:},
    moredelim=**[is][\color{orange}]{\@\{}{\@\}},
    moredelim=**[is][\color{cyan}]{\$\{}{\}},
    keywordstyle=\color{blue}\bfseries,
    commentstyle=\color{gray}\ttfamily,
    stringstyle=\color{red}\ttfamily,
    identifierstyle=\color{black},
    breaklines=true,
    basicstyle=\ttfamily\small,
    columns=fullflexible,
    showstringspaces=false,
    tabsize=2,
    frame=single,
    captionpos=b,
}


\lstdefinestyle{java}{
    language=Java,
    basicstyle=\ttfamily\small,
    keywordstyle=\color{blue},
    stringstyle=\color{red},
    commentstyle=\color{gray},
    morecomment=[s][\color{gray}]{/**}{*/},
    breaklines=true,
    numbers=left,
    numberstyle=\tiny\color{gray},
    showspaces=false,
    showstringspaces=false,
    frame=single,
    captionpos=b
}

\usepackage{color}
\definecolor{numb}{rgb}{0.58,0,0.82}
\definecolor{punct}{rgb}{0.3,0.3,0.3}
\definecolor{delim}{rgb}{0.5,0.5,0.5}
\definecolor{string}{rgb}{0.6,0.1,0.1}

\lstdefinestyle{json}{
    language=json,
    basicstyle=\ttfamily\small,
    keywordstyle=\color{blue},
    stringstyle=\color{string},
    commentstyle=\color{green},
    morecomment=[s][\color{gray}]{/**}{*/},
    breaklines=true,
    numbers=left,
    numberstyle=\tiny\color{gray},
    showspaces=false,
    showstringspaces=false,
    frame=single,
    captionpos=b
}

\lstdefinelanguage{drl}{
    basicstyle=\ttfamily\small,
    keywordstyle=\color{blue}\bfseries,
    stringstyle=\color{red},
    commentstyle=\color{green},
    morecomment=[s][\color{green}]{/*}{*/},
    morecomment=[l][\color{green}]{//},
    morestring=[b]",
    keywords={rule, when, then, end, salience, eval, update, retract},
    sensitive=true,
    breaklines=true,
    numbers=left,
    numberstyle=\tiny\color{gray},
    showspaces=false,
    showstringspaces=false,
    frame=single,
    captionpos=b
}

\lstdefinestyle{drl}{
    language=drl,
    basicstyle=\ttfamily\small,
    keywordstyle=\color{blue}\bfseries,
    stringstyle=\color{red},
    commentstyle=\color{green},
    breaklines=true,
    numbers=left,
    numberstyle=\tiny\color{gray},
    showspaces=false,
    showstringspaces=false,
    frame=single,
    captionpos=b
}



\renewcommand{\lstlistingname}{Codi}
\renewcommand{\lstlistlistingname}{Índex de Fragments de Codis}
\renewcommand{\appendixpagename}{Annexos}
\renewcommand{\appendixtocname}{Annexos}

\renewcommand*{\listalgorithmcfname}{Llistat de pseudocodi}
\renewcommand*{\algorithmcfname}{Pseudocodi}
\renewcommand*{\algorithmautorefname}{pseudocodi}

\definecolor{darkgreen}{rgb}{0.0, 0.3, 0.0}
\definecolor{lightgreen}{rgb}{0.9, 0.95, 0.9}

\newcommand{\todo}[1]{
    \begin{tcolorbox}[colback=lightgreen, colframe=darkgreen, title=TODO]
        #1
    \end{tcolorbox}
}


% \usepackage{indentfirst} 
% \setlength{\parindent}{1cm}


\usepackage{parskip}
% \setlength{\parindent}{10pt}

% \setlength{\headheight}{15.6pt}
\setlength{\headheight}{16.05263pt}




\vspace{5mm}

\newcolumntype{u}{>{\centering\arraybackslash\hsize=.22\hsize}X}
\newcolumntype{s}{>{\hsize=.45\hsize}X}


%%%%%%%%%%%%%%%%%%%%%%%%%%%%%%%%
% INI creacio subsubsubsecion
%%%%%%%%%%%%%%%%%%%%%%%%%%%%%%%%
\usepackage{titlesec}

\let\oldparagraph\paragraph
\renewcommand{\paragraph}[1]{\oldparagraph{\hspace{30pt}#1}}


\titleclass{\subsubsubsection}{straight}[\subsection]

\newcounter{subsubsubsection}[subsubsection]
\renewcommand\thesubsubsubsection{\thesubsubsection.\arabic{subsubsubsection}}
\renewcommand\theparagraph{\thesubsubsubsection.\arabic{paragraph}} % optional; useful if paragraphs are to be numbered

\titleformat{\subsubsubsection}
  {\normalfont\normalsize\bfseries}{\thesubsubsubsection}{1em}{}
\titlespacing*{\subsubsubsection}
{0pt}{3.25ex plus 1ex minus .2ex}{1.5ex plus .2ex}


\makeatletter
\renewcommand\paragraph{\@startsection{paragraph}{5}{10pt}%
  {3.25ex \@plus1ex \@minus.2ex}%
  {-1em}%
  {\normalfont\normalsize\bfseries}}
\renewcommand\subparagraph{\@startsection{subparagraph}{6}{\parindent}%
  {3.25ex \@plus1ex \@minus .2ex}%
  {-1em}%
  {\normalfont\normalsize\bfseries}}
\def\toclevel@subsubsubsection{4}
\def\toclevel@paragraph{5}
\def\toclevel@paragraph{6}
\def\l@subsubsubsection{\@dottedtocline{4}{7em}{4em}}
\def\l@paragraph{\@dottedtocline{5}{10em}{5em}}
\def\l@subparagraph{\@dottedtocline{6}{14em}{6em}}
\makeatother

\setcounter{secnumdepth}{4}
\setcounter{tocdepth}{3}
%%%%%%%%%%%%%%%%%%%%%%%%%%%%%%%%
% FIN creacio subsubsubtitle
%%%%%%%%%%%%%%%%%%%%%%%%%%%%%%%%


%%%%%%%%%%%%%%%%%%%%%%%%%%%%%%%%
% INI creacio equation
%%%%%%%%%%%%%%%%%%%%%%%%%%%%%%%%
\DeclareFloatingEnvironment[fileext=loeq,placement={!ht},name=Equació]{equationfloat}
%%%%%%%%%%%%%%%%%%%%%%%%%%%%%%%%
% FIN creacio equation
%%%%%%%%%%%%%%%%%%%%%%%%%%%%%%%%

\titleformat{\section}[block]
  {\clearpage\normalfont\Large\bfseries}
  {\thesection}{1em}{}


\renewcommand\refname{Bibliografia}


\linespread{1.25}


\makeatletter
\renewenvironment{abstract}{%
    \if@twocolumn
      \section*{\abstractname}%
    \else 
      \begin{center}%
        {\bfseries \Large\itshape\abstractname\vspace{\z@}}%  % Título en negrita, 14pt y cursiva
      \end{center}%
      \quotation
    \fi}
    {\if@twocolumn\else\endquotation\fi}
\makeatother

% \fancypagestyle{main}{
%     \fancyhf{}
%     \lhead{\small Aplicació Web per a la Visualització i Comparació
% de Trajectòries d’Algorismes d’Optimització}
%     % \rhead{\includegraphics[width=1.2cm]{resources/Icons/logo-fib.png}}
%     \cfoot{\thepage}
%     \setlength{\footskip}{35pt}
% }

\newcommand{\q}[1]{``#1''}

\begin{document}

\hypersetup{
    pdfborder=0 0 0,
}

%% PORTADA I TABLE OF CONTENTS %%

%%%%                                   %%%%
%         INDEX DE LA MEMORIA             %
%%%%                                   %%%%

% \includepdf[pages=-]{portada_fib.pdf}
\input{Extra/GEP}
%\input{Extra/portada}
% \cleardoublepage

% \pagestyle{main}
% \setcounter{page}{11}
\include{toc}
% \cleardoublepage
\input{Contexto y alcance/main}

\input{Planificación temporal/main}

\input{Gestión económica/main}

\input{Informe de sostenibilidad/main}

% \cleardoublepage

% \cleardoublepage


% \input{Annex/main}


\clearpage
% Estilo bibliográfico IEEE
\bibliographystyle{IEEEtran}
\renewcommand{\refname}{Referencias}
\bibliography{references.bib}

%%%%                                   %%%%
%         FINAL DE LA MEMORIA             %
%%%%                                   %%%%

\end{document}


%% VERSIÓ INCIAL MEMORIA %%




\clearpage
% Estilo bibliográfico IEEE
\bibliographystyle{IEEEtran}
\renewcommand{\refname}{Referencias}
\bibliography{references.bib}

%%%%                                   %%%%
%         FINAL DE LA MEMORIA             %
%%%%                                   %%%%

\end{document}


%% VERSIÓ INCIAL MEMORIA %%


\section{Alcance}

\subsection{Objetivo principal}

\subsection{Objetivos secundarios}
Otros objetivos que deben cumplirse, pero sin ser críticos para el éxito del proyecto son las siguientes:
\begin{itemize}
    \item Desarrollar un front-end (Open OnDemand) que permita a los usuarios consultar de manera sencilla las cuotas de uso y los costes asociados.
    \item Diseñar el sistema para que sea portable, en previsión de la futura migración de los clústeres a un servicio externo de housing\footnote{El servicio de \textit{housing} en un Centro de Procesamiento de Datos (CPD), también llamado \textit{colocation}, consiste en alquilar espacio en el CPD para instalar y operar servidores propios. El proveedor gestiona la infraestructura (electricidad, refrigeración, seguridad y conectividad), mientras el cliente mantiene la propiedad y gestión del hardware.}.
\end{itemize}

\subsection{Requerimientos funcionales}
Los requerimientos funcionales de este proyecto son los siguientes:

\begin{itemize}
    \item El sistema deberá obtener diariamente los datos de consumo de cómputo y almacenamiento de los laboratorios y Core Facilities.
    \item Ampliación del código fuente de Slurm mediante plugins para recopilar información adicional sobre la heterogeneidad de arquitecturas de CPU y la disponibilidad de nodos on spot.
    \item Los datos deberán almacenarse en Google BigQuery como base central de análisis.
    \item La aplicación deberá procesar la información y calcular la imputación proporcional de costes según el uso de hardware y recursos.
    \item El sistema deberá integrar los resultados en el sistema financiero SAP, utilizando el Service Layer (API de SAP Business One).
    \item El proceso deberá ejecutarse de manera automática y periódica en una máquina virtual Linux.
    \item La solución deberá generar registros (logs) de ejecución para garantizar trazabilidad y control de errores.
\end{itemize}


\subsection{Requerimientos no funcionales}
A continuación, se detallan los requerimientos no funcionales para el correcto funcionamiento del proyecto: 
\begin{itemize}
    \item La aplicación deberá estar programada en Python, siguiendo buenas prácticas de desarrollo y con control de versiones en GitHub.
    \item El sistema deberá ser escalable, permitiendo la adaptación a un mayor número de usuarios o recursos en el futuro.
    \item El tiempo de procesamiento diario deberá ser razonable, evitando retrasos en la imputación de costes.
    \item La solución deberá ser segura, garantizando la integridad de los datos transferidos entre BigQuery y SAP.
\end{itemize}


\subsection{Posibles obstáculos y riesgos}
\documentclass[a4paper, 11pt, twoside]{article}
\usepackage{wrapfig}
% \usepackage{abstract}
\usepackage{graphicx} % Required for inserting images
\usepackage[a4paper, left=1.15in, right=1.15in, top=1.5in, bottom=1.2in]{geometry}
\usepackage{float}
\usepackage{array}
\usepackage{booktabs}
\usepackage[T1]{fontenc}
\usepackage{amsmath}
\usepackage{tabularx}
\usepackage{graphicx} % Required for inserting images
\usepackage{tikz, xcolor}
\usepackage{hyperref}
\usepackage{cite}
\usepackage{pdfpages}
\usepackage{pdflscape}
\usepackage{afterpage}
\usepackage{textpos}
\usepackage{fancyhdr}
\usepackage{caption}
\usepackage{lmodern,textcomp}
\usepackage[catalan]{babel}
\usepackage{lipsum}
\usepackage{newfloat}
\usepackage{listings}
\usepackage{placeins}
\usepackage[ruled]{algorithm2e}
\usepackage{tcolorbox}
% Estilo de citas IEEE sencillo: usar paquete cite para comprimir y ordenar referencias numéricas
\usepackage{cite}
\usepackage[toc,page]{appendix}
\usepackage{forest}
\usepackage{pdflscape}
\usepackage{graphicx}
\usepackage{geometry}
\usepackage{quoting,xparse}
\usepackage{svg}



% \NewDocumentCommand{\bywhom}{m}{% the Bourbaki trick
%   {\nobreak\hfill\penalty50\hskip1em\null\nobreak
%    \hfill\mbox{\normalfont(#1)}%
%    \parfillskip=0pt \finalhyphendemerits=0 \par}%
% }

% AIXÒ SERVEIX PER A MODIFICAR UN PROBLEMA AMB ELS CARACTERS --> " + PARAULA 

% \NewDocumentEnvironment{pquotation}{m}
%   {\begin{quoting}[
%      indentfirst=true,
%      leftmargin=\parindent,
%      rightmargin=\parindent]\itshape}
%   {\bywhom{#1}
  
%   \end{quoting}}

% AIXÒ PER INSERIR FRAGMENTS DE CODI

\lstnewenvironment{code}[1][]%
{
   \noindent
   \minipage{\linewidth}
   \vspace{0.5\baselineskip}
   \lstset{
     basicstyle=\ttfamily\footnotesize,
     frame=single,
     breaklines=true,    % Permitir que las líneas se rompan automáticamente
     breakatwhitespace=true, % Permitir la ruptura en los espacios en blanco
     #1
   }
}
{\endminipage}

\lstset{
    language=Java,
    basicstyle=\ttfamily\small,  % Define el tipo de letra y tamaño del código
    keywordstyle=\color{blue},   % Define el color de las palabras clave de Java
    stringstyle=\color{red},     % Define el color de las cadenas de texto
    commentstyle=\color{green},  % Define el color de los comentarios
    morecomment=[s][\color{gray}]{/**}{*/},  % Estilo para comentarios de bloque
    breaklines=true,  % Configura el quiebre automático de línea
    numbers=left,  % Posición de los números de línea
    numberstyle=\tiny\color{gray},  % Estilo para los números de línea
    showspaces=false,  % No mostrar espacios con underscores
    showstringspaces=false,  % No mostrar espacios dentro de strings
    frame=single,  % Agregar un marco alrededor del código
    captionpos=b,  % b = bottom, t = top
}

\definecolor{eclipseStrings}{RGB}{42,0.0,255}
\definecolor{eclipseKeywords}{RGB}{127,0,85}
\colorlet{numb}{magenta!60!black}

\lstdefinelanguage{json}{
    basicstyle=\ttfamily\small,
    showstringspaces=false,
    breaklines=true,
    frame=single,
    literate=
     *{0}{{{\color{numb}0}}}{1}
      {1}{{{\color{numb}1}}}{1}
      {2}{{{\color{numb}2}}}{1}
      {3}{{{\color{numb}3}}}{1}
      {4}{{{\color{numb}4}}}{1}
      {5}{{{\color{numb}5}}}{1}
      {6}{{{\color{numb}6}}}{1}
      {7}{{{\color{numb}7}}}{1}
      {8}{{{\color{numb}8}}}{1}
      {9}{{{\color{numb}9}}}{1}
      {:}{{{\color{punct}{:}}}}{1}
      {,}{{{\color{punct}{,}}}}{1}
      {\{}{{{\color{delim}{\{}}}}{1}
      {\}}{{{\color{delim}{\}}}}}{1}
      {[}{{{\color{delim}{[}}}}{1}
      {]}{{{\color{delim}{]}}}}{1},
    morestring=[b]"
}

\lstdefinelanguage{JavaScript}{
  morekeywords=[1]{break, continue, delete, else, for, function, if, in,
    new, return, this, typeof, var, void, while, with},
  % Literals, primitive types, and reference types.
  morekeywords=[2]{false, null, true, boolean, number, undefined,
    Array, Boolean, Date, Math, Number, String, Object},
  % Built-ins.
  morekeywords=[3]{eval, parseInt, parseFloat, escape, unescape},
  sensitive,
  morecomment=[s]{/*}{*/},
  morecomment=[l]//,
  morecomment=[s]{/**}{*/}, % JavaDoc style comments
  morestring=[b]',
  morestring=[b]"
}[keywords, comments, strings]

\lstdefinelanguage{svelte}{
    morekeywords={async, await, then, if, each, in, on, let, const, function, import, export, from, with, typeof, new, var, return},
    sensitive=true,
    morecomment=[l]{//},
    morecomment=[s]{/*}{*/},
    morecomment=[s]{<!--}{-->},
    morestring=[b]',
    morestring=[b]",
    morestring=[b]`,
    ndkeywords={import, export, return, if, for, while, function, switch, case, break, var, let, const, async, await, then, catch, try, typeof, new, class, extends, implements, interface, constructor, public, private, protected, static, get, set},
    alsoletter={:},
    moredelim=**[is][\color{orange}]{\@\{}{\@\}},
    moredelim=**[is][\color{cyan}]{\$\{}{\}},
    keywordstyle=\color{blue}\bfseries,
    commentstyle=\color{gray}\ttfamily,
    stringstyle=\color{red}\ttfamily,
    identifierstyle=\color{black},
    breaklines=true,
    basicstyle=\ttfamily\small,
    columns=fullflexible,
    showstringspaces=false,
    tabsize=2,
    frame=single,
    captionpos=b,
}


\lstdefinestyle{java}{
    language=Java,
    basicstyle=\ttfamily\small,
    keywordstyle=\color{blue},
    stringstyle=\color{red},
    commentstyle=\color{gray},
    morecomment=[s][\color{gray}]{/**}{*/},
    breaklines=true,
    numbers=left,
    numberstyle=\tiny\color{gray},
    showspaces=false,
    showstringspaces=false,
    frame=single,
    captionpos=b
}

\usepackage{color}
\definecolor{numb}{rgb}{0.58,0,0.82}
\definecolor{punct}{rgb}{0.3,0.3,0.3}
\definecolor{delim}{rgb}{0.5,0.5,0.5}
\definecolor{string}{rgb}{0.6,0.1,0.1}

\lstdefinestyle{json}{
    language=json,
    basicstyle=\ttfamily\small,
    keywordstyle=\color{blue},
    stringstyle=\color{string},
    commentstyle=\color{green},
    morecomment=[s][\color{gray}]{/**}{*/},
    breaklines=true,
    numbers=left,
    numberstyle=\tiny\color{gray},
    showspaces=false,
    showstringspaces=false,
    frame=single,
    captionpos=b
}

\lstdefinelanguage{drl}{
    basicstyle=\ttfamily\small,
    keywordstyle=\color{blue}\bfseries,
    stringstyle=\color{red},
    commentstyle=\color{green},
    morecomment=[s][\color{green}]{/*}{*/},
    morecomment=[l][\color{green}]{//},
    morestring=[b]",
    keywords={rule, when, then, end, salience, eval, update, retract},
    sensitive=true,
    breaklines=true,
    numbers=left,
    numberstyle=\tiny\color{gray},
    showspaces=false,
    showstringspaces=false,
    frame=single,
    captionpos=b
}

\lstdefinestyle{drl}{
    language=drl,
    basicstyle=\ttfamily\small,
    keywordstyle=\color{blue}\bfseries,
    stringstyle=\color{red},
    commentstyle=\color{green},
    breaklines=true,
    numbers=left,
    numberstyle=\tiny\color{gray},
    showspaces=false,
    showstringspaces=false,
    frame=single,
    captionpos=b
}



\renewcommand{\lstlistingname}{Codi}
\renewcommand{\lstlistlistingname}{Índex de Fragments de Codis}
\renewcommand{\appendixpagename}{Annexos}
\renewcommand{\appendixtocname}{Annexos}

\renewcommand*{\listalgorithmcfname}{Llistat de pseudocodi}
\renewcommand*{\algorithmcfname}{Pseudocodi}
\renewcommand*{\algorithmautorefname}{pseudocodi}

\definecolor{darkgreen}{rgb}{0.0, 0.3, 0.0}
\definecolor{lightgreen}{rgb}{0.9, 0.95, 0.9}

\newcommand{\todo}[1]{
    \begin{tcolorbox}[colback=lightgreen, colframe=darkgreen, title=TODO]
        #1
    \end{tcolorbox}
}


% \usepackage{indentfirst} 
% \setlength{\parindent}{1cm}


\usepackage{parskip}
% \setlength{\parindent}{10pt}

% \setlength{\headheight}{15.6pt}
\setlength{\headheight}{16.05263pt}




\vspace{5mm}

\newcolumntype{u}{>{\centering\arraybackslash\hsize=.22\hsize}X}
\newcolumntype{s}{>{\hsize=.45\hsize}X}


%%%%%%%%%%%%%%%%%%%%%%%%%%%%%%%%
% INI creacio subsubsubsecion
%%%%%%%%%%%%%%%%%%%%%%%%%%%%%%%%
\usepackage{titlesec}

\let\oldparagraph\paragraph
\renewcommand{\paragraph}[1]{\oldparagraph{\hspace{30pt}#1}}


\titleclass{\subsubsubsection}{straight}[\subsection]

\newcounter{subsubsubsection}[subsubsection]
\renewcommand\thesubsubsubsection{\thesubsubsection.\arabic{subsubsubsection}}
\renewcommand\theparagraph{\thesubsubsubsection.\arabic{paragraph}} % optional; useful if paragraphs are to be numbered

\titleformat{\subsubsubsection}
  {\normalfont\normalsize\bfseries}{\thesubsubsubsection}{1em}{}
\titlespacing*{\subsubsubsection}
{0pt}{3.25ex plus 1ex minus .2ex}{1.5ex plus .2ex}


\makeatletter
\renewcommand\paragraph{\@startsection{paragraph}{5}{10pt}%
  {3.25ex \@plus1ex \@minus.2ex}%
  {-1em}%
  {\normalfont\normalsize\bfseries}}
\renewcommand\subparagraph{\@startsection{subparagraph}{6}{\parindent}%
  {3.25ex \@plus1ex \@minus .2ex}%
  {-1em}%
  {\normalfont\normalsize\bfseries}}
\def\toclevel@subsubsubsection{4}
\def\toclevel@paragraph{5}
\def\toclevel@paragraph{6}
\def\l@subsubsubsection{\@dottedtocline{4}{7em}{4em}}
\def\l@paragraph{\@dottedtocline{5}{10em}{5em}}
\def\l@subparagraph{\@dottedtocline{6}{14em}{6em}}
\makeatother

\setcounter{secnumdepth}{4}
\setcounter{tocdepth}{3}
%%%%%%%%%%%%%%%%%%%%%%%%%%%%%%%%
% FIN creacio subsubsubtitle
%%%%%%%%%%%%%%%%%%%%%%%%%%%%%%%%


%%%%%%%%%%%%%%%%%%%%%%%%%%%%%%%%
% INI creacio equation
%%%%%%%%%%%%%%%%%%%%%%%%%%%%%%%%
\DeclareFloatingEnvironment[fileext=loeq,placement={!ht},name=Equació]{equationfloat}
%%%%%%%%%%%%%%%%%%%%%%%%%%%%%%%%
% FIN creacio equation
%%%%%%%%%%%%%%%%%%%%%%%%%%%%%%%%

\titleformat{\section}[block]
  {\clearpage\normalfont\Large\bfseries}
  {\thesection}{1em}{}


\renewcommand\refname{Bibliografia}


\linespread{1.25}


\makeatletter
\renewenvironment{abstract}{%
    \if@twocolumn
      \section*{\abstractname}%
    \else 
      \begin{center}%
        {\bfseries \Large\itshape\abstractname\vspace{\z@}}%  % Título en negrita, 14pt y cursiva
      \end{center}%
      \quotation
    \fi}
    {\if@twocolumn\else\endquotation\fi}
\makeatother

% \fancypagestyle{main}{
%     \fancyhf{}
%     \lhead{\small Aplicació Web per a la Visualització i Comparació
% de Trajectòries d’Algorismes d’Optimització}
%     % \rhead{\includegraphics[width=1.2cm]{resources/Icons/logo-fib.png}}
%     \cfoot{\thepage}
%     \setlength{\footskip}{35pt}
% }

\newcommand{\q}[1]{``#1''}

\begin{document}

\hypersetup{
    pdfborder=0 0 0,
}

%% PORTADA I TABLE OF CONTENTS %%

%%%%                                   %%%%
%         INDEX DE LA MEMORIA             %
%%%%                                   %%%%

% \includepdf[pages=-]{portada_fib.pdf}
%%% Portada
\begin{titlepage}
\begin{center}
    \includegraphics[width=1\linewidth]{extra/figures/logo-upc-fib.png}
    \vspace{30mm}

    \LARGE{\textsc{\textbf{{INTEGRACIÓN Y AUTOMATIZACIÓN DE COSTES DE INFRAESTRUCTURA HPC EN ENTORNOS DE INVESTIGACIÓN}}}} \\
    \vspace{25mm}
      
        
    \textbf{\Large{\textsc{PAU NAVARRO ÁLVAREZ}}} \\
    \vspace{20mm}
        
    
    
    \large{\textbf{Director/a} ROBERTO RIVEIRO INSUA (FUNDACIO INSTITUT DE RECERCA BIOMEDICA (IRB BARCELONA))} \\
    \large{\textbf{Ponente:} ALEX PAJUELO GONZÁLEZ (Departamento de Arquitectura de Computadores)} \\
    \large{\textbf{Titulación:} Grado en Ingeniería Informática (Tecnologías de la información)} \\
    
    \large{\textbf{Lliurament 2:} }Planificación temporal \\
    \large{\textbf{Tutor GEP:} CAROLINA CONSOLIDACIÓN SEGURA (Departament d'Organització d'Empreses)} \\
    \large{\textbf{Fecha:}} 29 de Septiembre de 2025\\

    \textbf{\large{Memoria del trabajo de fin de grado}} \\
    \vspace{3mm}
    \textbf{\large{Facultat d'Informàtica de Barcelona (FIB)}} \\
    \vspace{3mm}
    \textbf{\large{Universitat Politècnica de Catalunya (UPC) - BarcelonaTech}} \\

\end{center}
\end{titlepage}
%\input{Extra/portada}
% \cleardoublepage

% \pagestyle{main}
% \setcounter{page}{11}
\renewcommand{\contentsname}{Continguts}

\tableofcontents
\clearpage
\listoffigures
\begingroup
    \let\clearpage\relax
    \listoftables
\endgroup
\clearpage

\newpage
% \cleardoublepage
\input{Contexto y alcance/Contextualización/main}
\input{Contexto y alcance/Justificación/main}
\input{Contexto y alcance/Alcance/main}
\input{Contexto y alcance/Metodología y herramientas de seguimiento/main}

\documentclass[a4paper, 11pt, twoside]{article}
\usepackage{wrapfig}
% \usepackage{abstract}
\usepackage{graphicx} % Required for inserting images
\usepackage[a4paper, left=1.15in, right=1.15in, top=1.5in, bottom=1.2in]{geometry}
\usepackage{float}
\usepackage{array}
\usepackage{booktabs}
\usepackage[T1]{fontenc}
\usepackage{amsmath}
\usepackage{tabularx}
\usepackage{graphicx} % Required for inserting images
\usepackage{tikz, xcolor}
\usepackage{hyperref}
\usepackage{cite}
\usepackage{pdfpages}
\usepackage{pdflscape}
\usepackage{afterpage}
\usepackage{textpos}
\usepackage{fancyhdr}
\usepackage{caption}
\usepackage{lmodern,textcomp}
\usepackage[catalan]{babel}
\usepackage{lipsum}
\usepackage{newfloat}
\usepackage{listings}
\usepackage{placeins}
\usepackage[ruled]{algorithm2e}
\usepackage{tcolorbox}
% Estilo de citas IEEE sencillo: usar paquete cite para comprimir y ordenar referencias numéricas
\usepackage{cite}
\usepackage[toc,page]{appendix}
\usepackage{forest}
\usepackage{pdflscape}
\usepackage{graphicx}
\usepackage{geometry}
\usepackage{quoting,xparse}
\usepackage{svg}



% \NewDocumentCommand{\bywhom}{m}{% the Bourbaki trick
%   {\nobreak\hfill\penalty50\hskip1em\null\nobreak
%    \hfill\mbox{\normalfont(#1)}%
%    \parfillskip=0pt \finalhyphendemerits=0 \par}%
% }

% AIXÒ SERVEIX PER A MODIFICAR UN PROBLEMA AMB ELS CARACTERS --> " + PARAULA 

% \NewDocumentEnvironment{pquotation}{m}
%   {\begin{quoting}[
%      indentfirst=true,
%      leftmargin=\parindent,
%      rightmargin=\parindent]\itshape}
%   {\bywhom{#1}
  
%   \end{quoting}}

% AIXÒ PER INSERIR FRAGMENTS DE CODI

\lstnewenvironment{code}[1][]%
{
   \noindent
   \minipage{\linewidth}
   \vspace{0.5\baselineskip}
   \lstset{
     basicstyle=\ttfamily\footnotesize,
     frame=single,
     breaklines=true,    % Permitir que las líneas se rompan automáticamente
     breakatwhitespace=true, % Permitir la ruptura en los espacios en blanco
     #1
   }
}
{\endminipage}

\lstset{
    language=Java,
    basicstyle=\ttfamily\small,  % Define el tipo de letra y tamaño del código
    keywordstyle=\color{blue},   % Define el color de las palabras clave de Java
    stringstyle=\color{red},     % Define el color de las cadenas de texto
    commentstyle=\color{green},  % Define el color de los comentarios
    morecomment=[s][\color{gray}]{/**}{*/},  % Estilo para comentarios de bloque
    breaklines=true,  % Configura el quiebre automático de línea
    numbers=left,  % Posición de los números de línea
    numberstyle=\tiny\color{gray},  % Estilo para los números de línea
    showspaces=false,  % No mostrar espacios con underscores
    showstringspaces=false,  % No mostrar espacios dentro de strings
    frame=single,  % Agregar un marco alrededor del código
    captionpos=b,  % b = bottom, t = top
}

\definecolor{eclipseStrings}{RGB}{42,0.0,255}
\definecolor{eclipseKeywords}{RGB}{127,0,85}
\colorlet{numb}{magenta!60!black}

\lstdefinelanguage{json}{
    basicstyle=\ttfamily\small,
    showstringspaces=false,
    breaklines=true,
    frame=single,
    literate=
     *{0}{{{\color{numb}0}}}{1}
      {1}{{{\color{numb}1}}}{1}
      {2}{{{\color{numb}2}}}{1}
      {3}{{{\color{numb}3}}}{1}
      {4}{{{\color{numb}4}}}{1}
      {5}{{{\color{numb}5}}}{1}
      {6}{{{\color{numb}6}}}{1}
      {7}{{{\color{numb}7}}}{1}
      {8}{{{\color{numb}8}}}{1}
      {9}{{{\color{numb}9}}}{1}
      {:}{{{\color{punct}{:}}}}{1}
      {,}{{{\color{punct}{,}}}}{1}
      {\{}{{{\color{delim}{\{}}}}{1}
      {\}}{{{\color{delim}{\}}}}}{1}
      {[}{{{\color{delim}{[}}}}{1}
      {]}{{{\color{delim}{]}}}}{1},
    morestring=[b]"
}

\lstdefinelanguage{JavaScript}{
  morekeywords=[1]{break, continue, delete, else, for, function, if, in,
    new, return, this, typeof, var, void, while, with},
  % Literals, primitive types, and reference types.
  morekeywords=[2]{false, null, true, boolean, number, undefined,
    Array, Boolean, Date, Math, Number, String, Object},
  % Built-ins.
  morekeywords=[3]{eval, parseInt, parseFloat, escape, unescape},
  sensitive,
  morecomment=[s]{/*}{*/},
  morecomment=[l]//,
  morecomment=[s]{/**}{*/}, % JavaDoc style comments
  morestring=[b]',
  morestring=[b]"
}[keywords, comments, strings]

\lstdefinelanguage{svelte}{
    morekeywords={async, await, then, if, each, in, on, let, const, function, import, export, from, with, typeof, new, var, return},
    sensitive=true,
    morecomment=[l]{//},
    morecomment=[s]{/*}{*/},
    morecomment=[s]{<!--}{-->},
    morestring=[b]',
    morestring=[b]",
    morestring=[b]`,
    ndkeywords={import, export, return, if, for, while, function, switch, case, break, var, let, const, async, await, then, catch, try, typeof, new, class, extends, implements, interface, constructor, public, private, protected, static, get, set},
    alsoletter={:},
    moredelim=**[is][\color{orange}]{\@\{}{\@\}},
    moredelim=**[is][\color{cyan}]{\$\{}{\}},
    keywordstyle=\color{blue}\bfseries,
    commentstyle=\color{gray}\ttfamily,
    stringstyle=\color{red}\ttfamily,
    identifierstyle=\color{black},
    breaklines=true,
    basicstyle=\ttfamily\small,
    columns=fullflexible,
    showstringspaces=false,
    tabsize=2,
    frame=single,
    captionpos=b,
}


\lstdefinestyle{java}{
    language=Java,
    basicstyle=\ttfamily\small,
    keywordstyle=\color{blue},
    stringstyle=\color{red},
    commentstyle=\color{gray},
    morecomment=[s][\color{gray}]{/**}{*/},
    breaklines=true,
    numbers=left,
    numberstyle=\tiny\color{gray},
    showspaces=false,
    showstringspaces=false,
    frame=single,
    captionpos=b
}

\usepackage{color}
\definecolor{numb}{rgb}{0.58,0,0.82}
\definecolor{punct}{rgb}{0.3,0.3,0.3}
\definecolor{delim}{rgb}{0.5,0.5,0.5}
\definecolor{string}{rgb}{0.6,0.1,0.1}

\lstdefinestyle{json}{
    language=json,
    basicstyle=\ttfamily\small,
    keywordstyle=\color{blue},
    stringstyle=\color{string},
    commentstyle=\color{green},
    morecomment=[s][\color{gray}]{/**}{*/},
    breaklines=true,
    numbers=left,
    numberstyle=\tiny\color{gray},
    showspaces=false,
    showstringspaces=false,
    frame=single,
    captionpos=b
}

\lstdefinelanguage{drl}{
    basicstyle=\ttfamily\small,
    keywordstyle=\color{blue}\bfseries,
    stringstyle=\color{red},
    commentstyle=\color{green},
    morecomment=[s][\color{green}]{/*}{*/},
    morecomment=[l][\color{green}]{//},
    morestring=[b]",
    keywords={rule, when, then, end, salience, eval, update, retract},
    sensitive=true,
    breaklines=true,
    numbers=left,
    numberstyle=\tiny\color{gray},
    showspaces=false,
    showstringspaces=false,
    frame=single,
    captionpos=b
}

\lstdefinestyle{drl}{
    language=drl,
    basicstyle=\ttfamily\small,
    keywordstyle=\color{blue}\bfseries,
    stringstyle=\color{red},
    commentstyle=\color{green},
    breaklines=true,
    numbers=left,
    numberstyle=\tiny\color{gray},
    showspaces=false,
    showstringspaces=false,
    frame=single,
    captionpos=b
}



\renewcommand{\lstlistingname}{Codi}
\renewcommand{\lstlistlistingname}{Índex de Fragments de Codis}
\renewcommand{\appendixpagename}{Annexos}
\renewcommand{\appendixtocname}{Annexos}

\renewcommand*{\listalgorithmcfname}{Llistat de pseudocodi}
\renewcommand*{\algorithmcfname}{Pseudocodi}
\renewcommand*{\algorithmautorefname}{pseudocodi}

\definecolor{darkgreen}{rgb}{0.0, 0.3, 0.0}
\definecolor{lightgreen}{rgb}{0.9, 0.95, 0.9}

\newcommand{\todo}[1]{
    \begin{tcolorbox}[colback=lightgreen, colframe=darkgreen, title=TODO]
        #1
    \end{tcolorbox}
}


% \usepackage{indentfirst} 
% \setlength{\parindent}{1cm}


\usepackage{parskip}
% \setlength{\parindent}{10pt}

% \setlength{\headheight}{15.6pt}
\setlength{\headheight}{16.05263pt}




\vspace{5mm}

\newcolumntype{u}{>{\centering\arraybackslash\hsize=.22\hsize}X}
\newcolumntype{s}{>{\hsize=.45\hsize}X}


%%%%%%%%%%%%%%%%%%%%%%%%%%%%%%%%
% INI creacio subsubsubsecion
%%%%%%%%%%%%%%%%%%%%%%%%%%%%%%%%
\usepackage{titlesec}

\let\oldparagraph\paragraph
\renewcommand{\paragraph}[1]{\oldparagraph{\hspace{30pt}#1}}


\titleclass{\subsubsubsection}{straight}[\subsection]

\newcounter{subsubsubsection}[subsubsection]
\renewcommand\thesubsubsubsection{\thesubsubsection.\arabic{subsubsubsection}}
\renewcommand\theparagraph{\thesubsubsubsection.\arabic{paragraph}} % optional; useful if paragraphs are to be numbered

\titleformat{\subsubsubsection}
  {\normalfont\normalsize\bfseries}{\thesubsubsubsection}{1em}{}
\titlespacing*{\subsubsubsection}
{0pt}{3.25ex plus 1ex minus .2ex}{1.5ex plus .2ex}


\makeatletter
\renewcommand\paragraph{\@startsection{paragraph}{5}{10pt}%
  {3.25ex \@plus1ex \@minus.2ex}%
  {-1em}%
  {\normalfont\normalsize\bfseries}}
\renewcommand\subparagraph{\@startsection{subparagraph}{6}{\parindent}%
  {3.25ex \@plus1ex \@minus .2ex}%
  {-1em}%
  {\normalfont\normalsize\bfseries}}
\def\toclevel@subsubsubsection{4}
\def\toclevel@paragraph{5}
\def\toclevel@paragraph{6}
\def\l@subsubsubsection{\@dottedtocline{4}{7em}{4em}}
\def\l@paragraph{\@dottedtocline{5}{10em}{5em}}
\def\l@subparagraph{\@dottedtocline{6}{14em}{6em}}
\makeatother

\setcounter{secnumdepth}{4}
\setcounter{tocdepth}{3}
%%%%%%%%%%%%%%%%%%%%%%%%%%%%%%%%
% FIN creacio subsubsubtitle
%%%%%%%%%%%%%%%%%%%%%%%%%%%%%%%%


%%%%%%%%%%%%%%%%%%%%%%%%%%%%%%%%
% INI creacio equation
%%%%%%%%%%%%%%%%%%%%%%%%%%%%%%%%
\DeclareFloatingEnvironment[fileext=loeq,placement={!ht},name=Equació]{equationfloat}
%%%%%%%%%%%%%%%%%%%%%%%%%%%%%%%%
% FIN creacio equation
%%%%%%%%%%%%%%%%%%%%%%%%%%%%%%%%

\titleformat{\section}[block]
  {\clearpage\normalfont\Large\bfseries}
  {\thesection}{1em}{}


\renewcommand\refname{Bibliografia}


\linespread{1.25}


\makeatletter
\renewenvironment{abstract}{%
    \if@twocolumn
      \section*{\abstractname}%
    \else 
      \begin{center}%
        {\bfseries \Large\itshape\abstractname\vspace{\z@}}%  % Título en negrita, 14pt y cursiva
      \end{center}%
      \quotation
    \fi}
    {\if@twocolumn\else\endquotation\fi}
\makeatother

% \fancypagestyle{main}{
%     \fancyhf{}
%     \lhead{\small Aplicació Web per a la Visualització i Comparació
% de Trajectòries d’Algorismes d’Optimització}
%     % \rhead{\includegraphics[width=1.2cm]{resources/Icons/logo-fib.png}}
%     \cfoot{\thepage}
%     \setlength{\footskip}{35pt}
% }

\newcommand{\q}[1]{``#1''}

\begin{document}

\hypersetup{
    pdfborder=0 0 0,
}

%% PORTADA I TABLE OF CONTENTS %%

%%%%                                   %%%%
%         INDEX DE LA MEMORIA             %
%%%%                                   %%%%

% \includepdf[pages=-]{portada_fib.pdf}
\input{Extra/GEP}
%\input{Extra/portada}
% \cleardoublepage

% \pagestyle{main}
% \setcounter{page}{11}
\include{toc}
% \cleardoublepage
\input{Contexto y alcance/main}

\input{Planificación temporal/main}

\input{Gestión económica/main}

\input{Informe de sostenibilidad/main}

% \cleardoublepage

% \cleardoublepage


% \input{Annex/main}


\clearpage
% Estilo bibliográfico IEEE
\bibliographystyle{IEEEtran}
\renewcommand{\refname}{Referencias}
\bibliography{references.bib}

%%%%                                   %%%%
%         FINAL DE LA MEMORIA             %
%%%%                                   %%%%

\end{document}


%% VERSIÓ INCIAL MEMORIA %%



\section{Gestión económica}

La gestión económica de un proyecto es fundamental para asegurar su viabilidad y exito a largo plazo. Ya hemos visto en secciones anteriores, los otros dos vertices 
del triángulo de la gestión de proyectos: el alcance y el tiempo. En esta sección, nos centraremos en cómo gestionar eficazmente los recursos financieros para garantizar
que el proyecto se mantenga dentro del presupuesto y se logren los objetivos establecidos. 

\input{Gestión económica/Recursos humanos.tex}

\input{Gestión económica/Recursos Hardware.tex}

\input{Gestión económica/Recursos Software.tex}

\input{Gestión económica/Contingencias.tex}

\section{Informe de sostenibilidad}

\input{Informe de sostenibilidad/Autoevaluación.tex}


\input{Informe de sostenibilidad/Dimensión ambiental.tex}

\input{Informe de sostenibilidad/Dimensión económica.tex}

\input{Informe de sostenibilidad/Dimensión social.tex}


% \cleardoublepage

% \cleardoublepage


% \documentclass[a4paper, 11pt, twoside]{article}
\usepackage{wrapfig}
% \usepackage{abstract}
\usepackage{graphicx} % Required for inserting images
\usepackage[a4paper, left=1.15in, right=1.15in, top=1.5in, bottom=1.2in]{geometry}
\usepackage{float}
\usepackage{array}
\usepackage{booktabs}
\usepackage[T1]{fontenc}
\usepackage{amsmath}
\usepackage{tabularx}
\usepackage{graphicx} % Required for inserting images
\usepackage{tikz, xcolor}
\usepackage{hyperref}
\usepackage{cite}
\usepackage{pdfpages}
\usepackage{pdflscape}
\usepackage{afterpage}
\usepackage{textpos}
\usepackage{fancyhdr}
\usepackage{caption}
\usepackage{lmodern,textcomp}
\usepackage[catalan]{babel}
\usepackage{lipsum}
\usepackage{newfloat}
\usepackage{listings}
\usepackage{placeins}
\usepackage[ruled]{algorithm2e}
\usepackage{tcolorbox}
% Estilo de citas IEEE sencillo: usar paquete cite para comprimir y ordenar referencias numéricas
\usepackage{cite}
\usepackage[toc,page]{appendix}
\usepackage{forest}
\usepackage{pdflscape}
\usepackage{graphicx}
\usepackage{geometry}
\usepackage{quoting,xparse}
\usepackage{svg}



% \NewDocumentCommand{\bywhom}{m}{% the Bourbaki trick
%   {\nobreak\hfill\penalty50\hskip1em\null\nobreak
%    \hfill\mbox{\normalfont(#1)}%
%    \parfillskip=0pt \finalhyphendemerits=0 \par}%
% }

% AIXÒ SERVEIX PER A MODIFICAR UN PROBLEMA AMB ELS CARACTERS --> " + PARAULA 

% \NewDocumentEnvironment{pquotation}{m}
%   {\begin{quoting}[
%      indentfirst=true,
%      leftmargin=\parindent,
%      rightmargin=\parindent]\itshape}
%   {\bywhom{#1}
  
%   \end{quoting}}

% AIXÒ PER INSERIR FRAGMENTS DE CODI

\lstnewenvironment{code}[1][]%
{
   \noindent
   \minipage{\linewidth}
   \vspace{0.5\baselineskip}
   \lstset{
     basicstyle=\ttfamily\footnotesize,
     frame=single,
     breaklines=true,    % Permitir que las líneas se rompan automáticamente
     breakatwhitespace=true, % Permitir la ruptura en los espacios en blanco
     #1
   }
}
{\endminipage}

\lstset{
    language=Java,
    basicstyle=\ttfamily\small,  % Define el tipo de letra y tamaño del código
    keywordstyle=\color{blue},   % Define el color de las palabras clave de Java
    stringstyle=\color{red},     % Define el color de las cadenas de texto
    commentstyle=\color{green},  % Define el color de los comentarios
    morecomment=[s][\color{gray}]{/**}{*/},  % Estilo para comentarios de bloque
    breaklines=true,  % Configura el quiebre automático de línea
    numbers=left,  % Posición de los números de línea
    numberstyle=\tiny\color{gray},  % Estilo para los números de línea
    showspaces=false,  % No mostrar espacios con underscores
    showstringspaces=false,  % No mostrar espacios dentro de strings
    frame=single,  % Agregar un marco alrededor del código
    captionpos=b,  % b = bottom, t = top
}

\definecolor{eclipseStrings}{RGB}{42,0.0,255}
\definecolor{eclipseKeywords}{RGB}{127,0,85}
\colorlet{numb}{magenta!60!black}

\lstdefinelanguage{json}{
    basicstyle=\ttfamily\small,
    showstringspaces=false,
    breaklines=true,
    frame=single,
    literate=
     *{0}{{{\color{numb}0}}}{1}
      {1}{{{\color{numb}1}}}{1}
      {2}{{{\color{numb}2}}}{1}
      {3}{{{\color{numb}3}}}{1}
      {4}{{{\color{numb}4}}}{1}
      {5}{{{\color{numb}5}}}{1}
      {6}{{{\color{numb}6}}}{1}
      {7}{{{\color{numb}7}}}{1}
      {8}{{{\color{numb}8}}}{1}
      {9}{{{\color{numb}9}}}{1}
      {:}{{{\color{punct}{:}}}}{1}
      {,}{{{\color{punct}{,}}}}{1}
      {\{}{{{\color{delim}{\{}}}}{1}
      {\}}{{{\color{delim}{\}}}}}{1}
      {[}{{{\color{delim}{[}}}}{1}
      {]}{{{\color{delim}{]}}}}{1},
    morestring=[b]"
}

\lstdefinelanguage{JavaScript}{
  morekeywords=[1]{break, continue, delete, else, for, function, if, in,
    new, return, this, typeof, var, void, while, with},
  % Literals, primitive types, and reference types.
  morekeywords=[2]{false, null, true, boolean, number, undefined,
    Array, Boolean, Date, Math, Number, String, Object},
  % Built-ins.
  morekeywords=[3]{eval, parseInt, parseFloat, escape, unescape},
  sensitive,
  morecomment=[s]{/*}{*/},
  morecomment=[l]//,
  morecomment=[s]{/**}{*/}, % JavaDoc style comments
  morestring=[b]',
  morestring=[b]"
}[keywords, comments, strings]

\lstdefinelanguage{svelte}{
    morekeywords={async, await, then, if, each, in, on, let, const, function, import, export, from, with, typeof, new, var, return},
    sensitive=true,
    morecomment=[l]{//},
    morecomment=[s]{/*}{*/},
    morecomment=[s]{<!--}{-->},
    morestring=[b]',
    morestring=[b]",
    morestring=[b]`,
    ndkeywords={import, export, return, if, for, while, function, switch, case, break, var, let, const, async, await, then, catch, try, typeof, new, class, extends, implements, interface, constructor, public, private, protected, static, get, set},
    alsoletter={:},
    moredelim=**[is][\color{orange}]{\@\{}{\@\}},
    moredelim=**[is][\color{cyan}]{\$\{}{\}},
    keywordstyle=\color{blue}\bfseries,
    commentstyle=\color{gray}\ttfamily,
    stringstyle=\color{red}\ttfamily,
    identifierstyle=\color{black},
    breaklines=true,
    basicstyle=\ttfamily\small,
    columns=fullflexible,
    showstringspaces=false,
    tabsize=2,
    frame=single,
    captionpos=b,
}


\lstdefinestyle{java}{
    language=Java,
    basicstyle=\ttfamily\small,
    keywordstyle=\color{blue},
    stringstyle=\color{red},
    commentstyle=\color{gray},
    morecomment=[s][\color{gray}]{/**}{*/},
    breaklines=true,
    numbers=left,
    numberstyle=\tiny\color{gray},
    showspaces=false,
    showstringspaces=false,
    frame=single,
    captionpos=b
}

\usepackage{color}
\definecolor{numb}{rgb}{0.58,0,0.82}
\definecolor{punct}{rgb}{0.3,0.3,0.3}
\definecolor{delim}{rgb}{0.5,0.5,0.5}
\definecolor{string}{rgb}{0.6,0.1,0.1}

\lstdefinestyle{json}{
    language=json,
    basicstyle=\ttfamily\small,
    keywordstyle=\color{blue},
    stringstyle=\color{string},
    commentstyle=\color{green},
    morecomment=[s][\color{gray}]{/**}{*/},
    breaklines=true,
    numbers=left,
    numberstyle=\tiny\color{gray},
    showspaces=false,
    showstringspaces=false,
    frame=single,
    captionpos=b
}

\lstdefinelanguage{drl}{
    basicstyle=\ttfamily\small,
    keywordstyle=\color{blue}\bfseries,
    stringstyle=\color{red},
    commentstyle=\color{green},
    morecomment=[s][\color{green}]{/*}{*/},
    morecomment=[l][\color{green}]{//},
    morestring=[b]",
    keywords={rule, when, then, end, salience, eval, update, retract},
    sensitive=true,
    breaklines=true,
    numbers=left,
    numberstyle=\tiny\color{gray},
    showspaces=false,
    showstringspaces=false,
    frame=single,
    captionpos=b
}

\lstdefinestyle{drl}{
    language=drl,
    basicstyle=\ttfamily\small,
    keywordstyle=\color{blue}\bfseries,
    stringstyle=\color{red},
    commentstyle=\color{green},
    breaklines=true,
    numbers=left,
    numberstyle=\tiny\color{gray},
    showspaces=false,
    showstringspaces=false,
    frame=single,
    captionpos=b
}



\renewcommand{\lstlistingname}{Codi}
\renewcommand{\lstlistlistingname}{Índex de Fragments de Codis}
\renewcommand{\appendixpagename}{Annexos}
\renewcommand{\appendixtocname}{Annexos}

\renewcommand*{\listalgorithmcfname}{Llistat de pseudocodi}
\renewcommand*{\algorithmcfname}{Pseudocodi}
\renewcommand*{\algorithmautorefname}{pseudocodi}

\definecolor{darkgreen}{rgb}{0.0, 0.3, 0.0}
\definecolor{lightgreen}{rgb}{0.9, 0.95, 0.9}

\newcommand{\todo}[1]{
    \begin{tcolorbox}[colback=lightgreen, colframe=darkgreen, title=TODO]
        #1
    \end{tcolorbox}
}


% \usepackage{indentfirst} 
% \setlength{\parindent}{1cm}


\usepackage{parskip}
% \setlength{\parindent}{10pt}

% \setlength{\headheight}{15.6pt}
\setlength{\headheight}{16.05263pt}




\vspace{5mm}

\newcolumntype{u}{>{\centering\arraybackslash\hsize=.22\hsize}X}
\newcolumntype{s}{>{\hsize=.45\hsize}X}


%%%%%%%%%%%%%%%%%%%%%%%%%%%%%%%%
% INI creacio subsubsubsecion
%%%%%%%%%%%%%%%%%%%%%%%%%%%%%%%%
\usepackage{titlesec}

\let\oldparagraph\paragraph
\renewcommand{\paragraph}[1]{\oldparagraph{\hspace{30pt}#1}}


\titleclass{\subsubsubsection}{straight}[\subsection]

\newcounter{subsubsubsection}[subsubsection]
\renewcommand\thesubsubsubsection{\thesubsubsection.\arabic{subsubsubsection}}
\renewcommand\theparagraph{\thesubsubsubsection.\arabic{paragraph}} % optional; useful if paragraphs are to be numbered

\titleformat{\subsubsubsection}
  {\normalfont\normalsize\bfseries}{\thesubsubsubsection}{1em}{}
\titlespacing*{\subsubsubsection}
{0pt}{3.25ex plus 1ex minus .2ex}{1.5ex plus .2ex}


\makeatletter
\renewcommand\paragraph{\@startsection{paragraph}{5}{10pt}%
  {3.25ex \@plus1ex \@minus.2ex}%
  {-1em}%
  {\normalfont\normalsize\bfseries}}
\renewcommand\subparagraph{\@startsection{subparagraph}{6}{\parindent}%
  {3.25ex \@plus1ex \@minus .2ex}%
  {-1em}%
  {\normalfont\normalsize\bfseries}}
\def\toclevel@subsubsubsection{4}
\def\toclevel@paragraph{5}
\def\toclevel@paragraph{6}
\def\l@subsubsubsection{\@dottedtocline{4}{7em}{4em}}
\def\l@paragraph{\@dottedtocline{5}{10em}{5em}}
\def\l@subparagraph{\@dottedtocline{6}{14em}{6em}}
\makeatother

\setcounter{secnumdepth}{4}
\setcounter{tocdepth}{3}
%%%%%%%%%%%%%%%%%%%%%%%%%%%%%%%%
% FIN creacio subsubsubtitle
%%%%%%%%%%%%%%%%%%%%%%%%%%%%%%%%


%%%%%%%%%%%%%%%%%%%%%%%%%%%%%%%%
% INI creacio equation
%%%%%%%%%%%%%%%%%%%%%%%%%%%%%%%%
\DeclareFloatingEnvironment[fileext=loeq,placement={!ht},name=Equació]{equationfloat}
%%%%%%%%%%%%%%%%%%%%%%%%%%%%%%%%
% FIN creacio equation
%%%%%%%%%%%%%%%%%%%%%%%%%%%%%%%%

\titleformat{\section}[block]
  {\clearpage\normalfont\Large\bfseries}
  {\thesection}{1em}{}


\renewcommand\refname{Bibliografia}


\linespread{1.25}


\makeatletter
\renewenvironment{abstract}{%
    \if@twocolumn
      \section*{\abstractname}%
    \else 
      \begin{center}%
        {\bfseries \Large\itshape\abstractname\vspace{\z@}}%  % Título en negrita, 14pt y cursiva
      \end{center}%
      \quotation
    \fi}
    {\if@twocolumn\else\endquotation\fi}
\makeatother

% \fancypagestyle{main}{
%     \fancyhf{}
%     \lhead{\small Aplicació Web per a la Visualització i Comparació
% de Trajectòries d’Algorismes d’Optimització}
%     % \rhead{\includegraphics[width=1.2cm]{resources/Icons/logo-fib.png}}
%     \cfoot{\thepage}
%     \setlength{\footskip}{35pt}
% }

\newcommand{\q}[1]{``#1''}

\begin{document}

\hypersetup{
    pdfborder=0 0 0,
}

%% PORTADA I TABLE OF CONTENTS %%

%%%%                                   %%%%
%         INDEX DE LA MEMORIA             %
%%%%                                   %%%%

% \includepdf[pages=-]{portada_fib.pdf}
\input{Extra/GEP}
%\input{Extra/portada}
% \cleardoublepage

% \pagestyle{main}
% \setcounter{page}{11}
\include{toc}
% \cleardoublepage
\input{Contexto y alcance/main}

\input{Planificación temporal/main}

\input{Gestión económica/main}

\input{Informe de sostenibilidad/main}

% \cleardoublepage

% \cleardoublepage


% \input{Annex/main}


\clearpage
% Estilo bibliográfico IEEE
\bibliographystyle{IEEEtran}
\renewcommand{\refname}{Referencias}
\bibliography{references.bib}

%%%%                                   %%%%
%         FINAL DE LA MEMORIA             %
%%%%                                   %%%%

\end{document}


%% VERSIÓ INCIAL MEMORIA %%




\clearpage
% Estilo bibliográfico IEEE
\bibliographystyle{IEEEtran}
\renewcommand{\refname}{Referencias}
\bibliography{references.bib}

%%%%                                   %%%%
%         FINAL DE LA MEMORIA             %
%%%%                                   %%%%

\end{document}


%% VERSIÓ INCIAL MEMORIA %%



\documentclass[a4paper, 11pt, twoside]{article}
\usepackage{wrapfig}
% \usepackage{abstract}
\usepackage{graphicx} % Required for inserting images
\usepackage[a4paper, left=1.15in, right=1.15in, top=1.5in, bottom=1.2in]{geometry}
\usepackage{float}
\usepackage{array}
\usepackage{booktabs}
\usepackage[T1]{fontenc}
\usepackage{amsmath}
\usepackage{tabularx}
\usepackage{graphicx} % Required for inserting images
\usepackage{tikz, xcolor}
\usepackage{hyperref}
\usepackage{cite}
\usepackage{pdfpages}
\usepackage{pdflscape}
\usepackage{afterpage}
\usepackage{textpos}
\usepackage{fancyhdr}
\usepackage{caption}
\usepackage{lmodern,textcomp}
\usepackage[catalan]{babel}
\usepackage{lipsum}
\usepackage{newfloat}
\usepackage{listings}
\usepackage{placeins}
\usepackage[ruled]{algorithm2e}
\usepackage{tcolorbox}
% Estilo de citas IEEE sencillo: usar paquete cite para comprimir y ordenar referencias numéricas
\usepackage{cite}
\usepackage[toc,page]{appendix}
\usepackage{forest}
\usepackage{pdflscape}
\usepackage{graphicx}
\usepackage{geometry}
\usepackage{quoting,xparse}
\usepackage{svg}



% \NewDocumentCommand{\bywhom}{m}{% the Bourbaki trick
%   {\nobreak\hfill\penalty50\hskip1em\null\nobreak
%    \hfill\mbox{\normalfont(#1)}%
%    \parfillskip=0pt \finalhyphendemerits=0 \par}%
% }

% AIXÒ SERVEIX PER A MODIFICAR UN PROBLEMA AMB ELS CARACTERS --> " + PARAULA 

% \NewDocumentEnvironment{pquotation}{m}
%   {\begin{quoting}[
%      indentfirst=true,
%      leftmargin=\parindent,
%      rightmargin=\parindent]\itshape}
%   {\bywhom{#1}
  
%   \end{quoting}}

% AIXÒ PER INSERIR FRAGMENTS DE CODI

\lstnewenvironment{code}[1][]%
{
   \noindent
   \minipage{\linewidth}
   \vspace{0.5\baselineskip}
   \lstset{
     basicstyle=\ttfamily\footnotesize,
     frame=single,
     breaklines=true,    % Permitir que las líneas se rompan automáticamente
     breakatwhitespace=true, % Permitir la ruptura en los espacios en blanco
     #1
   }
}
{\endminipage}

\lstset{
    language=Java,
    basicstyle=\ttfamily\small,  % Define el tipo de letra y tamaño del código
    keywordstyle=\color{blue},   % Define el color de las palabras clave de Java
    stringstyle=\color{red},     % Define el color de las cadenas de texto
    commentstyle=\color{green},  % Define el color de los comentarios
    morecomment=[s][\color{gray}]{/**}{*/},  % Estilo para comentarios de bloque
    breaklines=true,  % Configura el quiebre automático de línea
    numbers=left,  % Posición de los números de línea
    numberstyle=\tiny\color{gray},  % Estilo para los números de línea
    showspaces=false,  % No mostrar espacios con underscores
    showstringspaces=false,  % No mostrar espacios dentro de strings
    frame=single,  % Agregar un marco alrededor del código
    captionpos=b,  % b = bottom, t = top
}

\definecolor{eclipseStrings}{RGB}{42,0.0,255}
\definecolor{eclipseKeywords}{RGB}{127,0,85}
\colorlet{numb}{magenta!60!black}

\lstdefinelanguage{json}{
    basicstyle=\ttfamily\small,
    showstringspaces=false,
    breaklines=true,
    frame=single,
    literate=
     *{0}{{{\color{numb}0}}}{1}
      {1}{{{\color{numb}1}}}{1}
      {2}{{{\color{numb}2}}}{1}
      {3}{{{\color{numb}3}}}{1}
      {4}{{{\color{numb}4}}}{1}
      {5}{{{\color{numb}5}}}{1}
      {6}{{{\color{numb}6}}}{1}
      {7}{{{\color{numb}7}}}{1}
      {8}{{{\color{numb}8}}}{1}
      {9}{{{\color{numb}9}}}{1}
      {:}{{{\color{punct}{:}}}}{1}
      {,}{{{\color{punct}{,}}}}{1}
      {\{}{{{\color{delim}{\{}}}}{1}
      {\}}{{{\color{delim}{\}}}}}{1}
      {[}{{{\color{delim}{[}}}}{1}
      {]}{{{\color{delim}{]}}}}{1},
    morestring=[b]"
}

\lstdefinelanguage{JavaScript}{
  morekeywords=[1]{break, continue, delete, else, for, function, if, in,
    new, return, this, typeof, var, void, while, with},
  % Literals, primitive types, and reference types.
  morekeywords=[2]{false, null, true, boolean, number, undefined,
    Array, Boolean, Date, Math, Number, String, Object},
  % Built-ins.
  morekeywords=[3]{eval, parseInt, parseFloat, escape, unescape},
  sensitive,
  morecomment=[s]{/*}{*/},
  morecomment=[l]//,
  morecomment=[s]{/**}{*/}, % JavaDoc style comments
  morestring=[b]',
  morestring=[b]"
}[keywords, comments, strings]

\lstdefinelanguage{svelte}{
    morekeywords={async, await, then, if, each, in, on, let, const, function, import, export, from, with, typeof, new, var, return},
    sensitive=true,
    morecomment=[l]{//},
    morecomment=[s]{/*}{*/},
    morecomment=[s]{<!--}{-->},
    morestring=[b]',
    morestring=[b]",
    morestring=[b]`,
    ndkeywords={import, export, return, if, for, while, function, switch, case, break, var, let, const, async, await, then, catch, try, typeof, new, class, extends, implements, interface, constructor, public, private, protected, static, get, set},
    alsoletter={:},
    moredelim=**[is][\color{orange}]{\@\{}{\@\}},
    moredelim=**[is][\color{cyan}]{\$\{}{\}},
    keywordstyle=\color{blue}\bfseries,
    commentstyle=\color{gray}\ttfamily,
    stringstyle=\color{red}\ttfamily,
    identifierstyle=\color{black},
    breaklines=true,
    basicstyle=\ttfamily\small,
    columns=fullflexible,
    showstringspaces=false,
    tabsize=2,
    frame=single,
    captionpos=b,
}


\lstdefinestyle{java}{
    language=Java,
    basicstyle=\ttfamily\small,
    keywordstyle=\color{blue},
    stringstyle=\color{red},
    commentstyle=\color{gray},
    morecomment=[s][\color{gray}]{/**}{*/},
    breaklines=true,
    numbers=left,
    numberstyle=\tiny\color{gray},
    showspaces=false,
    showstringspaces=false,
    frame=single,
    captionpos=b
}

\usepackage{color}
\definecolor{numb}{rgb}{0.58,0,0.82}
\definecolor{punct}{rgb}{0.3,0.3,0.3}
\definecolor{delim}{rgb}{0.5,0.5,0.5}
\definecolor{string}{rgb}{0.6,0.1,0.1}

\lstdefinestyle{json}{
    language=json,
    basicstyle=\ttfamily\small,
    keywordstyle=\color{blue},
    stringstyle=\color{string},
    commentstyle=\color{green},
    morecomment=[s][\color{gray}]{/**}{*/},
    breaklines=true,
    numbers=left,
    numberstyle=\tiny\color{gray},
    showspaces=false,
    showstringspaces=false,
    frame=single,
    captionpos=b
}

\lstdefinelanguage{drl}{
    basicstyle=\ttfamily\small,
    keywordstyle=\color{blue}\bfseries,
    stringstyle=\color{red},
    commentstyle=\color{green},
    morecomment=[s][\color{green}]{/*}{*/},
    morecomment=[l][\color{green}]{//},
    morestring=[b]",
    keywords={rule, when, then, end, salience, eval, update, retract},
    sensitive=true,
    breaklines=true,
    numbers=left,
    numberstyle=\tiny\color{gray},
    showspaces=false,
    showstringspaces=false,
    frame=single,
    captionpos=b
}

\lstdefinestyle{drl}{
    language=drl,
    basicstyle=\ttfamily\small,
    keywordstyle=\color{blue}\bfseries,
    stringstyle=\color{red},
    commentstyle=\color{green},
    breaklines=true,
    numbers=left,
    numberstyle=\tiny\color{gray},
    showspaces=false,
    showstringspaces=false,
    frame=single,
    captionpos=b
}



\renewcommand{\lstlistingname}{Codi}
\renewcommand{\lstlistlistingname}{Índex de Fragments de Codis}
\renewcommand{\appendixpagename}{Annexos}
\renewcommand{\appendixtocname}{Annexos}

\renewcommand*{\listalgorithmcfname}{Llistat de pseudocodi}
\renewcommand*{\algorithmcfname}{Pseudocodi}
\renewcommand*{\algorithmautorefname}{pseudocodi}

\definecolor{darkgreen}{rgb}{0.0, 0.3, 0.0}
\definecolor{lightgreen}{rgb}{0.9, 0.95, 0.9}

\newcommand{\todo}[1]{
    \begin{tcolorbox}[colback=lightgreen, colframe=darkgreen, title=TODO]
        #1
    \end{tcolorbox}
}


% \usepackage{indentfirst} 
% \setlength{\parindent}{1cm}


\usepackage{parskip}
% \setlength{\parindent}{10pt}

% \setlength{\headheight}{15.6pt}
\setlength{\headheight}{16.05263pt}




\vspace{5mm}

\newcolumntype{u}{>{\centering\arraybackslash\hsize=.22\hsize}X}
\newcolumntype{s}{>{\hsize=.45\hsize}X}


%%%%%%%%%%%%%%%%%%%%%%%%%%%%%%%%
% INI creacio subsubsubsecion
%%%%%%%%%%%%%%%%%%%%%%%%%%%%%%%%
\usepackage{titlesec}

\let\oldparagraph\paragraph
\renewcommand{\paragraph}[1]{\oldparagraph{\hspace{30pt}#1}}


\titleclass{\subsubsubsection}{straight}[\subsection]

\newcounter{subsubsubsection}[subsubsection]
\renewcommand\thesubsubsubsection{\thesubsubsection.\arabic{subsubsubsection}}
\renewcommand\theparagraph{\thesubsubsubsection.\arabic{paragraph}} % optional; useful if paragraphs are to be numbered

\titleformat{\subsubsubsection}
  {\normalfont\normalsize\bfseries}{\thesubsubsubsection}{1em}{}
\titlespacing*{\subsubsubsection}
{0pt}{3.25ex plus 1ex minus .2ex}{1.5ex plus .2ex}


\makeatletter
\renewcommand\paragraph{\@startsection{paragraph}{5}{10pt}%
  {3.25ex \@plus1ex \@minus.2ex}%
  {-1em}%
  {\normalfont\normalsize\bfseries}}
\renewcommand\subparagraph{\@startsection{subparagraph}{6}{\parindent}%
  {3.25ex \@plus1ex \@minus .2ex}%
  {-1em}%
  {\normalfont\normalsize\bfseries}}
\def\toclevel@subsubsubsection{4}
\def\toclevel@paragraph{5}
\def\toclevel@paragraph{6}
\def\l@subsubsubsection{\@dottedtocline{4}{7em}{4em}}
\def\l@paragraph{\@dottedtocline{5}{10em}{5em}}
\def\l@subparagraph{\@dottedtocline{6}{14em}{6em}}
\makeatother

\setcounter{secnumdepth}{4}
\setcounter{tocdepth}{3}
%%%%%%%%%%%%%%%%%%%%%%%%%%%%%%%%
% FIN creacio subsubsubtitle
%%%%%%%%%%%%%%%%%%%%%%%%%%%%%%%%


%%%%%%%%%%%%%%%%%%%%%%%%%%%%%%%%
% INI creacio equation
%%%%%%%%%%%%%%%%%%%%%%%%%%%%%%%%
\DeclareFloatingEnvironment[fileext=loeq,placement={!ht},name=Equació]{equationfloat}
%%%%%%%%%%%%%%%%%%%%%%%%%%%%%%%%
% FIN creacio equation
%%%%%%%%%%%%%%%%%%%%%%%%%%%%%%%%

\titleformat{\section}[block]
  {\clearpage\normalfont\Large\bfseries}
  {\thesection}{1em}{}


\renewcommand\refname{Bibliografia}


\linespread{1.25}


\makeatletter
\renewenvironment{abstract}{%
    \if@twocolumn
      \section*{\abstractname}%
    \else 
      \begin{center}%
        {\bfseries \Large\itshape\abstractname\vspace{\z@}}%  % Título en negrita, 14pt y cursiva
      \end{center}%
      \quotation
    \fi}
    {\if@twocolumn\else\endquotation\fi}
\makeatother

% \fancypagestyle{main}{
%     \fancyhf{}
%     \lhead{\small Aplicació Web per a la Visualització i Comparació
% de Trajectòries d’Algorismes d’Optimització}
%     % \rhead{\includegraphics[width=1.2cm]{resources/Icons/logo-fib.png}}
%     \cfoot{\thepage}
%     \setlength{\footskip}{35pt}
% }

\newcommand{\q}[1]{``#1''}

\begin{document}

\hypersetup{
    pdfborder=0 0 0,
}

%% PORTADA I TABLE OF CONTENTS %%

%%%%                                   %%%%
%         INDEX DE LA MEMORIA             %
%%%%                                   %%%%

% \includepdf[pages=-]{portada_fib.pdf}
%%% Portada
\begin{titlepage}
\begin{center}
    \includegraphics[width=1\linewidth]{extra/figures/logo-upc-fib.png}
    \vspace{30mm}

    \LARGE{\textsc{\textbf{{INTEGRACIÓN Y AUTOMATIZACIÓN DE COSTES DE INFRAESTRUCTURA HPC EN ENTORNOS DE INVESTIGACIÓN}}}} \\
    \vspace{25mm}
      
        
    \textbf{\Large{\textsc{PAU NAVARRO ÁLVAREZ}}} \\
    \vspace{20mm}
        
    
    
    \large{\textbf{Director/a} ROBERTO RIVEIRO INSUA (FUNDACIO INSTITUT DE RECERCA BIOMEDICA (IRB BARCELONA))} \\
    \large{\textbf{Ponente:} ALEX PAJUELO GONZÁLEZ (Departamento de Arquitectura de Computadores)} \\
    \large{\textbf{Titulación:} Grado en Ingeniería Informática (Tecnologías de la información)} \\
    
    \large{\textbf{Lliurament 2:} }Planificación temporal \\
    \large{\textbf{Tutor GEP:} CAROLINA CONSOLIDACIÓN SEGURA (Departament d'Organització d'Empreses)} \\
    \large{\textbf{Fecha:}} 29 de Septiembre de 2025\\

    \textbf{\large{Memoria del trabajo de fin de grado}} \\
    \vspace{3mm}
    \textbf{\large{Facultat d'Informàtica de Barcelona (FIB)}} \\
    \vspace{3mm}
    \textbf{\large{Universitat Politècnica de Catalunya (UPC) - BarcelonaTech}} \\

\end{center}
\end{titlepage}
%\input{Extra/portada}
% \cleardoublepage

% \pagestyle{main}
% \setcounter{page}{11}
\renewcommand{\contentsname}{Continguts}

\tableofcontents
\clearpage
\listoffigures
\begingroup
    \let\clearpage\relax
    \listoftables
\endgroup
\clearpage

\newpage
% \cleardoublepage
\documentclass[a4paper, 11pt, twoside]{article}
\usepackage{wrapfig}
% \usepackage{abstract}
\usepackage{graphicx} % Required for inserting images
\usepackage[a4paper, left=1.15in, right=1.15in, top=1.5in, bottom=1.2in]{geometry}
\usepackage{float}
\usepackage{array}
\usepackage{booktabs}
\usepackage[T1]{fontenc}
\usepackage{amsmath}
\usepackage{tabularx}
\usepackage{graphicx} % Required for inserting images
\usepackage{tikz, xcolor}
\usepackage{hyperref}
\usepackage{cite}
\usepackage{pdfpages}
\usepackage{pdflscape}
\usepackage{afterpage}
\usepackage{textpos}
\usepackage{fancyhdr}
\usepackage{caption}
\usepackage{lmodern,textcomp}
\usepackage[catalan]{babel}
\usepackage{lipsum}
\usepackage{newfloat}
\usepackage{listings}
\usepackage{placeins}
\usepackage[ruled]{algorithm2e}
\usepackage{tcolorbox}
% Estilo de citas IEEE sencillo: usar paquete cite para comprimir y ordenar referencias numéricas
\usepackage{cite}
\usepackage[toc,page]{appendix}
\usepackage{forest}
\usepackage{pdflscape}
\usepackage{graphicx}
\usepackage{geometry}
\usepackage{quoting,xparse}
\usepackage{svg}



% \NewDocumentCommand{\bywhom}{m}{% the Bourbaki trick
%   {\nobreak\hfill\penalty50\hskip1em\null\nobreak
%    \hfill\mbox{\normalfont(#1)}%
%    \parfillskip=0pt \finalhyphendemerits=0 \par}%
% }

% AIXÒ SERVEIX PER A MODIFICAR UN PROBLEMA AMB ELS CARACTERS --> " + PARAULA 

% \NewDocumentEnvironment{pquotation}{m}
%   {\begin{quoting}[
%      indentfirst=true,
%      leftmargin=\parindent,
%      rightmargin=\parindent]\itshape}
%   {\bywhom{#1}
  
%   \end{quoting}}

% AIXÒ PER INSERIR FRAGMENTS DE CODI

\lstnewenvironment{code}[1][]%
{
   \noindent
   \minipage{\linewidth}
   \vspace{0.5\baselineskip}
   \lstset{
     basicstyle=\ttfamily\footnotesize,
     frame=single,
     breaklines=true,    % Permitir que las líneas se rompan automáticamente
     breakatwhitespace=true, % Permitir la ruptura en los espacios en blanco
     #1
   }
}
{\endminipage}

\lstset{
    language=Java,
    basicstyle=\ttfamily\small,  % Define el tipo de letra y tamaño del código
    keywordstyle=\color{blue},   % Define el color de las palabras clave de Java
    stringstyle=\color{red},     % Define el color de las cadenas de texto
    commentstyle=\color{green},  % Define el color de los comentarios
    morecomment=[s][\color{gray}]{/**}{*/},  % Estilo para comentarios de bloque
    breaklines=true,  % Configura el quiebre automático de línea
    numbers=left,  % Posición de los números de línea
    numberstyle=\tiny\color{gray},  % Estilo para los números de línea
    showspaces=false,  % No mostrar espacios con underscores
    showstringspaces=false,  % No mostrar espacios dentro de strings
    frame=single,  % Agregar un marco alrededor del código
    captionpos=b,  % b = bottom, t = top
}

\definecolor{eclipseStrings}{RGB}{42,0.0,255}
\definecolor{eclipseKeywords}{RGB}{127,0,85}
\colorlet{numb}{magenta!60!black}

\lstdefinelanguage{json}{
    basicstyle=\ttfamily\small,
    showstringspaces=false,
    breaklines=true,
    frame=single,
    literate=
     *{0}{{{\color{numb}0}}}{1}
      {1}{{{\color{numb}1}}}{1}
      {2}{{{\color{numb}2}}}{1}
      {3}{{{\color{numb}3}}}{1}
      {4}{{{\color{numb}4}}}{1}
      {5}{{{\color{numb}5}}}{1}
      {6}{{{\color{numb}6}}}{1}
      {7}{{{\color{numb}7}}}{1}
      {8}{{{\color{numb}8}}}{1}
      {9}{{{\color{numb}9}}}{1}
      {:}{{{\color{punct}{:}}}}{1}
      {,}{{{\color{punct}{,}}}}{1}
      {\{}{{{\color{delim}{\{}}}}{1}
      {\}}{{{\color{delim}{\}}}}}{1}
      {[}{{{\color{delim}{[}}}}{1}
      {]}{{{\color{delim}{]}}}}{1},
    morestring=[b]"
}

\lstdefinelanguage{JavaScript}{
  morekeywords=[1]{break, continue, delete, else, for, function, if, in,
    new, return, this, typeof, var, void, while, with},
  % Literals, primitive types, and reference types.
  morekeywords=[2]{false, null, true, boolean, number, undefined,
    Array, Boolean, Date, Math, Number, String, Object},
  % Built-ins.
  morekeywords=[3]{eval, parseInt, parseFloat, escape, unescape},
  sensitive,
  morecomment=[s]{/*}{*/},
  morecomment=[l]//,
  morecomment=[s]{/**}{*/}, % JavaDoc style comments
  morestring=[b]',
  morestring=[b]"
}[keywords, comments, strings]

\lstdefinelanguage{svelte}{
    morekeywords={async, await, then, if, each, in, on, let, const, function, import, export, from, with, typeof, new, var, return},
    sensitive=true,
    morecomment=[l]{//},
    morecomment=[s]{/*}{*/},
    morecomment=[s]{<!--}{-->},
    morestring=[b]',
    morestring=[b]",
    morestring=[b]`,
    ndkeywords={import, export, return, if, for, while, function, switch, case, break, var, let, const, async, await, then, catch, try, typeof, new, class, extends, implements, interface, constructor, public, private, protected, static, get, set},
    alsoletter={:},
    moredelim=**[is][\color{orange}]{\@\{}{\@\}},
    moredelim=**[is][\color{cyan}]{\$\{}{\}},
    keywordstyle=\color{blue}\bfseries,
    commentstyle=\color{gray}\ttfamily,
    stringstyle=\color{red}\ttfamily,
    identifierstyle=\color{black},
    breaklines=true,
    basicstyle=\ttfamily\small,
    columns=fullflexible,
    showstringspaces=false,
    tabsize=2,
    frame=single,
    captionpos=b,
}


\lstdefinestyle{java}{
    language=Java,
    basicstyle=\ttfamily\small,
    keywordstyle=\color{blue},
    stringstyle=\color{red},
    commentstyle=\color{gray},
    morecomment=[s][\color{gray}]{/**}{*/},
    breaklines=true,
    numbers=left,
    numberstyle=\tiny\color{gray},
    showspaces=false,
    showstringspaces=false,
    frame=single,
    captionpos=b
}

\usepackage{color}
\definecolor{numb}{rgb}{0.58,0,0.82}
\definecolor{punct}{rgb}{0.3,0.3,0.3}
\definecolor{delim}{rgb}{0.5,0.5,0.5}
\definecolor{string}{rgb}{0.6,0.1,0.1}

\lstdefinestyle{json}{
    language=json,
    basicstyle=\ttfamily\small,
    keywordstyle=\color{blue},
    stringstyle=\color{string},
    commentstyle=\color{green},
    morecomment=[s][\color{gray}]{/**}{*/},
    breaklines=true,
    numbers=left,
    numberstyle=\tiny\color{gray},
    showspaces=false,
    showstringspaces=false,
    frame=single,
    captionpos=b
}

\lstdefinelanguage{drl}{
    basicstyle=\ttfamily\small,
    keywordstyle=\color{blue}\bfseries,
    stringstyle=\color{red},
    commentstyle=\color{green},
    morecomment=[s][\color{green}]{/*}{*/},
    morecomment=[l][\color{green}]{//},
    morestring=[b]",
    keywords={rule, when, then, end, salience, eval, update, retract},
    sensitive=true,
    breaklines=true,
    numbers=left,
    numberstyle=\tiny\color{gray},
    showspaces=false,
    showstringspaces=false,
    frame=single,
    captionpos=b
}

\lstdefinestyle{drl}{
    language=drl,
    basicstyle=\ttfamily\small,
    keywordstyle=\color{blue}\bfseries,
    stringstyle=\color{red},
    commentstyle=\color{green},
    breaklines=true,
    numbers=left,
    numberstyle=\tiny\color{gray},
    showspaces=false,
    showstringspaces=false,
    frame=single,
    captionpos=b
}



\renewcommand{\lstlistingname}{Codi}
\renewcommand{\lstlistlistingname}{Índex de Fragments de Codis}
\renewcommand{\appendixpagename}{Annexos}
\renewcommand{\appendixtocname}{Annexos}

\renewcommand*{\listalgorithmcfname}{Llistat de pseudocodi}
\renewcommand*{\algorithmcfname}{Pseudocodi}
\renewcommand*{\algorithmautorefname}{pseudocodi}

\definecolor{darkgreen}{rgb}{0.0, 0.3, 0.0}
\definecolor{lightgreen}{rgb}{0.9, 0.95, 0.9}

\newcommand{\todo}[1]{
    \begin{tcolorbox}[colback=lightgreen, colframe=darkgreen, title=TODO]
        #1
    \end{tcolorbox}
}


% \usepackage{indentfirst} 
% \setlength{\parindent}{1cm}


\usepackage{parskip}
% \setlength{\parindent}{10pt}

% \setlength{\headheight}{15.6pt}
\setlength{\headheight}{16.05263pt}




\vspace{5mm}

\newcolumntype{u}{>{\centering\arraybackslash\hsize=.22\hsize}X}
\newcolumntype{s}{>{\hsize=.45\hsize}X}


%%%%%%%%%%%%%%%%%%%%%%%%%%%%%%%%
% INI creacio subsubsubsecion
%%%%%%%%%%%%%%%%%%%%%%%%%%%%%%%%
\usepackage{titlesec}

\let\oldparagraph\paragraph
\renewcommand{\paragraph}[1]{\oldparagraph{\hspace{30pt}#1}}


\titleclass{\subsubsubsection}{straight}[\subsection]

\newcounter{subsubsubsection}[subsubsection]
\renewcommand\thesubsubsubsection{\thesubsubsection.\arabic{subsubsubsection}}
\renewcommand\theparagraph{\thesubsubsubsection.\arabic{paragraph}} % optional; useful if paragraphs are to be numbered

\titleformat{\subsubsubsection}
  {\normalfont\normalsize\bfseries}{\thesubsubsubsection}{1em}{}
\titlespacing*{\subsubsubsection}
{0pt}{3.25ex plus 1ex minus .2ex}{1.5ex plus .2ex}


\makeatletter
\renewcommand\paragraph{\@startsection{paragraph}{5}{10pt}%
  {3.25ex \@plus1ex \@minus.2ex}%
  {-1em}%
  {\normalfont\normalsize\bfseries}}
\renewcommand\subparagraph{\@startsection{subparagraph}{6}{\parindent}%
  {3.25ex \@plus1ex \@minus .2ex}%
  {-1em}%
  {\normalfont\normalsize\bfseries}}
\def\toclevel@subsubsubsection{4}
\def\toclevel@paragraph{5}
\def\toclevel@paragraph{6}
\def\l@subsubsubsection{\@dottedtocline{4}{7em}{4em}}
\def\l@paragraph{\@dottedtocline{5}{10em}{5em}}
\def\l@subparagraph{\@dottedtocline{6}{14em}{6em}}
\makeatother

\setcounter{secnumdepth}{4}
\setcounter{tocdepth}{3}
%%%%%%%%%%%%%%%%%%%%%%%%%%%%%%%%
% FIN creacio subsubsubtitle
%%%%%%%%%%%%%%%%%%%%%%%%%%%%%%%%


%%%%%%%%%%%%%%%%%%%%%%%%%%%%%%%%
% INI creacio equation
%%%%%%%%%%%%%%%%%%%%%%%%%%%%%%%%
\DeclareFloatingEnvironment[fileext=loeq,placement={!ht},name=Equació]{equationfloat}
%%%%%%%%%%%%%%%%%%%%%%%%%%%%%%%%
% FIN creacio equation
%%%%%%%%%%%%%%%%%%%%%%%%%%%%%%%%

\titleformat{\section}[block]
  {\clearpage\normalfont\Large\bfseries}
  {\thesection}{1em}{}


\renewcommand\refname{Bibliografia}


\linespread{1.25}


\makeatletter
\renewenvironment{abstract}{%
    \if@twocolumn
      \section*{\abstractname}%
    \else 
      \begin{center}%
        {\bfseries \Large\itshape\abstractname\vspace{\z@}}%  % Título en negrita, 14pt y cursiva
      \end{center}%
      \quotation
    \fi}
    {\if@twocolumn\else\endquotation\fi}
\makeatother

% \fancypagestyle{main}{
%     \fancyhf{}
%     \lhead{\small Aplicació Web per a la Visualització i Comparació
% de Trajectòries d’Algorismes d’Optimització}
%     % \rhead{\includegraphics[width=1.2cm]{resources/Icons/logo-fib.png}}
%     \cfoot{\thepage}
%     \setlength{\footskip}{35pt}
% }

\newcommand{\q}[1]{``#1''}

\begin{document}

\hypersetup{
    pdfborder=0 0 0,
}

%% PORTADA I TABLE OF CONTENTS %%

%%%%                                   %%%%
%         INDEX DE LA MEMORIA             %
%%%%                                   %%%%

% \includepdf[pages=-]{portada_fib.pdf}
\input{Extra/GEP}
%\input{Extra/portada}
% \cleardoublepage

% \pagestyle{main}
% \setcounter{page}{11}
\include{toc}
% \cleardoublepage
\input{Contexto y alcance/main}

\input{Planificación temporal/main}

\input{Gestión económica/main}

\input{Informe de sostenibilidad/main}

% \cleardoublepage

% \cleardoublepage


% \input{Annex/main}


\clearpage
% Estilo bibliográfico IEEE
\bibliographystyle{IEEEtran}
\renewcommand{\refname}{Referencias}
\bibliography{references.bib}

%%%%                                   %%%%
%         FINAL DE LA MEMORIA             %
%%%%                                   %%%%

\end{document}


%% VERSIÓ INCIAL MEMORIA %%


\documentclass[a4paper, 11pt, twoside]{article}
\usepackage{wrapfig}
% \usepackage{abstract}
\usepackage{graphicx} % Required for inserting images
\usepackage[a4paper, left=1.15in, right=1.15in, top=1.5in, bottom=1.2in]{geometry}
\usepackage{float}
\usepackage{array}
\usepackage{booktabs}
\usepackage[T1]{fontenc}
\usepackage{amsmath}
\usepackage{tabularx}
\usepackage{graphicx} % Required for inserting images
\usepackage{tikz, xcolor}
\usepackage{hyperref}
\usepackage{cite}
\usepackage{pdfpages}
\usepackage{pdflscape}
\usepackage{afterpage}
\usepackage{textpos}
\usepackage{fancyhdr}
\usepackage{caption}
\usepackage{lmodern,textcomp}
\usepackage[catalan]{babel}
\usepackage{lipsum}
\usepackage{newfloat}
\usepackage{listings}
\usepackage{placeins}
\usepackage[ruled]{algorithm2e}
\usepackage{tcolorbox}
% Estilo de citas IEEE sencillo: usar paquete cite para comprimir y ordenar referencias numéricas
\usepackage{cite}
\usepackage[toc,page]{appendix}
\usepackage{forest}
\usepackage{pdflscape}
\usepackage{graphicx}
\usepackage{geometry}
\usepackage{quoting,xparse}
\usepackage{svg}



% \NewDocumentCommand{\bywhom}{m}{% the Bourbaki trick
%   {\nobreak\hfill\penalty50\hskip1em\null\nobreak
%    \hfill\mbox{\normalfont(#1)}%
%    \parfillskip=0pt \finalhyphendemerits=0 \par}%
% }

% AIXÒ SERVEIX PER A MODIFICAR UN PROBLEMA AMB ELS CARACTERS --> " + PARAULA 

% \NewDocumentEnvironment{pquotation}{m}
%   {\begin{quoting}[
%      indentfirst=true,
%      leftmargin=\parindent,
%      rightmargin=\parindent]\itshape}
%   {\bywhom{#1}
  
%   \end{quoting}}

% AIXÒ PER INSERIR FRAGMENTS DE CODI

\lstnewenvironment{code}[1][]%
{
   \noindent
   \minipage{\linewidth}
   \vspace{0.5\baselineskip}
   \lstset{
     basicstyle=\ttfamily\footnotesize,
     frame=single,
     breaklines=true,    % Permitir que las líneas se rompan automáticamente
     breakatwhitespace=true, % Permitir la ruptura en los espacios en blanco
     #1
   }
}
{\endminipage}

\lstset{
    language=Java,
    basicstyle=\ttfamily\small,  % Define el tipo de letra y tamaño del código
    keywordstyle=\color{blue},   % Define el color de las palabras clave de Java
    stringstyle=\color{red},     % Define el color de las cadenas de texto
    commentstyle=\color{green},  % Define el color de los comentarios
    morecomment=[s][\color{gray}]{/**}{*/},  % Estilo para comentarios de bloque
    breaklines=true,  % Configura el quiebre automático de línea
    numbers=left,  % Posición de los números de línea
    numberstyle=\tiny\color{gray},  % Estilo para los números de línea
    showspaces=false,  % No mostrar espacios con underscores
    showstringspaces=false,  % No mostrar espacios dentro de strings
    frame=single,  % Agregar un marco alrededor del código
    captionpos=b,  % b = bottom, t = top
}

\definecolor{eclipseStrings}{RGB}{42,0.0,255}
\definecolor{eclipseKeywords}{RGB}{127,0,85}
\colorlet{numb}{magenta!60!black}

\lstdefinelanguage{json}{
    basicstyle=\ttfamily\small,
    showstringspaces=false,
    breaklines=true,
    frame=single,
    literate=
     *{0}{{{\color{numb}0}}}{1}
      {1}{{{\color{numb}1}}}{1}
      {2}{{{\color{numb}2}}}{1}
      {3}{{{\color{numb}3}}}{1}
      {4}{{{\color{numb}4}}}{1}
      {5}{{{\color{numb}5}}}{1}
      {6}{{{\color{numb}6}}}{1}
      {7}{{{\color{numb}7}}}{1}
      {8}{{{\color{numb}8}}}{1}
      {9}{{{\color{numb}9}}}{1}
      {:}{{{\color{punct}{:}}}}{1}
      {,}{{{\color{punct}{,}}}}{1}
      {\{}{{{\color{delim}{\{}}}}{1}
      {\}}{{{\color{delim}{\}}}}}{1}
      {[}{{{\color{delim}{[}}}}{1}
      {]}{{{\color{delim}{]}}}}{1},
    morestring=[b]"
}

\lstdefinelanguage{JavaScript}{
  morekeywords=[1]{break, continue, delete, else, for, function, if, in,
    new, return, this, typeof, var, void, while, with},
  % Literals, primitive types, and reference types.
  morekeywords=[2]{false, null, true, boolean, number, undefined,
    Array, Boolean, Date, Math, Number, String, Object},
  % Built-ins.
  morekeywords=[3]{eval, parseInt, parseFloat, escape, unescape},
  sensitive,
  morecomment=[s]{/*}{*/},
  morecomment=[l]//,
  morecomment=[s]{/**}{*/}, % JavaDoc style comments
  morestring=[b]',
  morestring=[b]"
}[keywords, comments, strings]

\lstdefinelanguage{svelte}{
    morekeywords={async, await, then, if, each, in, on, let, const, function, import, export, from, with, typeof, new, var, return},
    sensitive=true,
    morecomment=[l]{//},
    morecomment=[s]{/*}{*/},
    morecomment=[s]{<!--}{-->},
    morestring=[b]',
    morestring=[b]",
    morestring=[b]`,
    ndkeywords={import, export, return, if, for, while, function, switch, case, break, var, let, const, async, await, then, catch, try, typeof, new, class, extends, implements, interface, constructor, public, private, protected, static, get, set},
    alsoletter={:},
    moredelim=**[is][\color{orange}]{\@\{}{\@\}},
    moredelim=**[is][\color{cyan}]{\$\{}{\}},
    keywordstyle=\color{blue}\bfseries,
    commentstyle=\color{gray}\ttfamily,
    stringstyle=\color{red}\ttfamily,
    identifierstyle=\color{black},
    breaklines=true,
    basicstyle=\ttfamily\small,
    columns=fullflexible,
    showstringspaces=false,
    tabsize=2,
    frame=single,
    captionpos=b,
}


\lstdefinestyle{java}{
    language=Java,
    basicstyle=\ttfamily\small,
    keywordstyle=\color{blue},
    stringstyle=\color{red},
    commentstyle=\color{gray},
    morecomment=[s][\color{gray}]{/**}{*/},
    breaklines=true,
    numbers=left,
    numberstyle=\tiny\color{gray},
    showspaces=false,
    showstringspaces=false,
    frame=single,
    captionpos=b
}

\usepackage{color}
\definecolor{numb}{rgb}{0.58,0,0.82}
\definecolor{punct}{rgb}{0.3,0.3,0.3}
\definecolor{delim}{rgb}{0.5,0.5,0.5}
\definecolor{string}{rgb}{0.6,0.1,0.1}

\lstdefinestyle{json}{
    language=json,
    basicstyle=\ttfamily\small,
    keywordstyle=\color{blue},
    stringstyle=\color{string},
    commentstyle=\color{green},
    morecomment=[s][\color{gray}]{/**}{*/},
    breaklines=true,
    numbers=left,
    numberstyle=\tiny\color{gray},
    showspaces=false,
    showstringspaces=false,
    frame=single,
    captionpos=b
}

\lstdefinelanguage{drl}{
    basicstyle=\ttfamily\small,
    keywordstyle=\color{blue}\bfseries,
    stringstyle=\color{red},
    commentstyle=\color{green},
    morecomment=[s][\color{green}]{/*}{*/},
    morecomment=[l][\color{green}]{//},
    morestring=[b]",
    keywords={rule, when, then, end, salience, eval, update, retract},
    sensitive=true,
    breaklines=true,
    numbers=left,
    numberstyle=\tiny\color{gray},
    showspaces=false,
    showstringspaces=false,
    frame=single,
    captionpos=b
}

\lstdefinestyle{drl}{
    language=drl,
    basicstyle=\ttfamily\small,
    keywordstyle=\color{blue}\bfseries,
    stringstyle=\color{red},
    commentstyle=\color{green},
    breaklines=true,
    numbers=left,
    numberstyle=\tiny\color{gray},
    showspaces=false,
    showstringspaces=false,
    frame=single,
    captionpos=b
}



\renewcommand{\lstlistingname}{Codi}
\renewcommand{\lstlistlistingname}{Índex de Fragments de Codis}
\renewcommand{\appendixpagename}{Annexos}
\renewcommand{\appendixtocname}{Annexos}

\renewcommand*{\listalgorithmcfname}{Llistat de pseudocodi}
\renewcommand*{\algorithmcfname}{Pseudocodi}
\renewcommand*{\algorithmautorefname}{pseudocodi}

\definecolor{darkgreen}{rgb}{0.0, 0.3, 0.0}
\definecolor{lightgreen}{rgb}{0.9, 0.95, 0.9}

\newcommand{\todo}[1]{
    \begin{tcolorbox}[colback=lightgreen, colframe=darkgreen, title=TODO]
        #1
    \end{tcolorbox}
}


% \usepackage{indentfirst} 
% \setlength{\parindent}{1cm}


\usepackage{parskip}
% \setlength{\parindent}{10pt}

% \setlength{\headheight}{15.6pt}
\setlength{\headheight}{16.05263pt}




\vspace{5mm}

\newcolumntype{u}{>{\centering\arraybackslash\hsize=.22\hsize}X}
\newcolumntype{s}{>{\hsize=.45\hsize}X}


%%%%%%%%%%%%%%%%%%%%%%%%%%%%%%%%
% INI creacio subsubsubsecion
%%%%%%%%%%%%%%%%%%%%%%%%%%%%%%%%
\usepackage{titlesec}

\let\oldparagraph\paragraph
\renewcommand{\paragraph}[1]{\oldparagraph{\hspace{30pt}#1}}


\titleclass{\subsubsubsection}{straight}[\subsection]

\newcounter{subsubsubsection}[subsubsection]
\renewcommand\thesubsubsubsection{\thesubsubsection.\arabic{subsubsubsection}}
\renewcommand\theparagraph{\thesubsubsubsection.\arabic{paragraph}} % optional; useful if paragraphs are to be numbered

\titleformat{\subsubsubsection}
  {\normalfont\normalsize\bfseries}{\thesubsubsubsection}{1em}{}
\titlespacing*{\subsubsubsection}
{0pt}{3.25ex plus 1ex minus .2ex}{1.5ex plus .2ex}


\makeatletter
\renewcommand\paragraph{\@startsection{paragraph}{5}{10pt}%
  {3.25ex \@plus1ex \@minus.2ex}%
  {-1em}%
  {\normalfont\normalsize\bfseries}}
\renewcommand\subparagraph{\@startsection{subparagraph}{6}{\parindent}%
  {3.25ex \@plus1ex \@minus .2ex}%
  {-1em}%
  {\normalfont\normalsize\bfseries}}
\def\toclevel@subsubsubsection{4}
\def\toclevel@paragraph{5}
\def\toclevel@paragraph{6}
\def\l@subsubsubsection{\@dottedtocline{4}{7em}{4em}}
\def\l@paragraph{\@dottedtocline{5}{10em}{5em}}
\def\l@subparagraph{\@dottedtocline{6}{14em}{6em}}
\makeatother

\setcounter{secnumdepth}{4}
\setcounter{tocdepth}{3}
%%%%%%%%%%%%%%%%%%%%%%%%%%%%%%%%
% FIN creacio subsubsubtitle
%%%%%%%%%%%%%%%%%%%%%%%%%%%%%%%%


%%%%%%%%%%%%%%%%%%%%%%%%%%%%%%%%
% INI creacio equation
%%%%%%%%%%%%%%%%%%%%%%%%%%%%%%%%
\DeclareFloatingEnvironment[fileext=loeq,placement={!ht},name=Equació]{equationfloat}
%%%%%%%%%%%%%%%%%%%%%%%%%%%%%%%%
% FIN creacio equation
%%%%%%%%%%%%%%%%%%%%%%%%%%%%%%%%

\titleformat{\section}[block]
  {\clearpage\normalfont\Large\bfseries}
  {\thesection}{1em}{}


\renewcommand\refname{Bibliografia}


\linespread{1.25}


\makeatletter
\renewenvironment{abstract}{%
    \if@twocolumn
      \section*{\abstractname}%
    \else 
      \begin{center}%
        {\bfseries \Large\itshape\abstractname\vspace{\z@}}%  % Título en negrita, 14pt y cursiva
      \end{center}%
      \quotation
    \fi}
    {\if@twocolumn\else\endquotation\fi}
\makeatother

% \fancypagestyle{main}{
%     \fancyhf{}
%     \lhead{\small Aplicació Web per a la Visualització i Comparació
% de Trajectòries d’Algorismes d’Optimització}
%     % \rhead{\includegraphics[width=1.2cm]{resources/Icons/logo-fib.png}}
%     \cfoot{\thepage}
%     \setlength{\footskip}{35pt}
% }

\newcommand{\q}[1]{``#1''}

\begin{document}

\hypersetup{
    pdfborder=0 0 0,
}

%% PORTADA I TABLE OF CONTENTS %%

%%%%                                   %%%%
%         INDEX DE LA MEMORIA             %
%%%%                                   %%%%

% \includepdf[pages=-]{portada_fib.pdf}
\input{Extra/GEP}
%\input{Extra/portada}
% \cleardoublepage

% \pagestyle{main}
% \setcounter{page}{11}
\include{toc}
% \cleardoublepage
\input{Contexto y alcance/main}

\input{Planificación temporal/main}

\input{Gestión económica/main}

\input{Informe de sostenibilidad/main}

% \cleardoublepage

% \cleardoublepage


% \input{Annex/main}


\clearpage
% Estilo bibliográfico IEEE
\bibliographystyle{IEEEtran}
\renewcommand{\refname}{Referencias}
\bibliography{references.bib}

%%%%                                   %%%%
%         FINAL DE LA MEMORIA             %
%%%%                                   %%%%

\end{document}


%% VERSIÓ INCIAL MEMORIA %%


\section{Alcance}

\input{Contexto y alcance/Alcance/Objetivo principal}

\input{Contexto y alcance/Alcance/Objetivos secundarios}

\input{Contexto y alcance/Alcance/Requerimientos funcionales}

\input{Contexto y alcance/Alcance/Requerimientos no funcionales}

\input{Contexto y alcance/Alcance/Posibles obstaculos y riesgos}
\documentclass[a4paper, 11pt, twoside]{article}
\usepackage{wrapfig}
% \usepackage{abstract}
\usepackage{graphicx} % Required for inserting images
\usepackage[a4paper, left=1.15in, right=1.15in, top=1.5in, bottom=1.2in]{geometry}
\usepackage{float}
\usepackage{array}
\usepackage{booktabs}
\usepackage[T1]{fontenc}
\usepackage{amsmath}
\usepackage{tabularx}
\usepackage{graphicx} % Required for inserting images
\usepackage{tikz, xcolor}
\usepackage{hyperref}
\usepackage{cite}
\usepackage{pdfpages}
\usepackage{pdflscape}
\usepackage{afterpage}
\usepackage{textpos}
\usepackage{fancyhdr}
\usepackage{caption}
\usepackage{lmodern,textcomp}
\usepackage[catalan]{babel}
\usepackage{lipsum}
\usepackage{newfloat}
\usepackage{listings}
\usepackage{placeins}
\usepackage[ruled]{algorithm2e}
\usepackage{tcolorbox}
% Estilo de citas IEEE sencillo: usar paquete cite para comprimir y ordenar referencias numéricas
\usepackage{cite}
\usepackage[toc,page]{appendix}
\usepackage{forest}
\usepackage{pdflscape}
\usepackage{graphicx}
\usepackage{geometry}
\usepackage{quoting,xparse}
\usepackage{svg}



% \NewDocumentCommand{\bywhom}{m}{% the Bourbaki trick
%   {\nobreak\hfill\penalty50\hskip1em\null\nobreak
%    \hfill\mbox{\normalfont(#1)}%
%    \parfillskip=0pt \finalhyphendemerits=0 \par}%
% }

% AIXÒ SERVEIX PER A MODIFICAR UN PROBLEMA AMB ELS CARACTERS --> " + PARAULA 

% \NewDocumentEnvironment{pquotation}{m}
%   {\begin{quoting}[
%      indentfirst=true,
%      leftmargin=\parindent,
%      rightmargin=\parindent]\itshape}
%   {\bywhom{#1}
  
%   \end{quoting}}

% AIXÒ PER INSERIR FRAGMENTS DE CODI

\lstnewenvironment{code}[1][]%
{
   \noindent
   \minipage{\linewidth}
   \vspace{0.5\baselineskip}
   \lstset{
     basicstyle=\ttfamily\footnotesize,
     frame=single,
     breaklines=true,    % Permitir que las líneas se rompan automáticamente
     breakatwhitespace=true, % Permitir la ruptura en los espacios en blanco
     #1
   }
}
{\endminipage}

\lstset{
    language=Java,
    basicstyle=\ttfamily\small,  % Define el tipo de letra y tamaño del código
    keywordstyle=\color{blue},   % Define el color de las palabras clave de Java
    stringstyle=\color{red},     % Define el color de las cadenas de texto
    commentstyle=\color{green},  % Define el color de los comentarios
    morecomment=[s][\color{gray}]{/**}{*/},  % Estilo para comentarios de bloque
    breaklines=true,  % Configura el quiebre automático de línea
    numbers=left,  % Posición de los números de línea
    numberstyle=\tiny\color{gray},  % Estilo para los números de línea
    showspaces=false,  % No mostrar espacios con underscores
    showstringspaces=false,  % No mostrar espacios dentro de strings
    frame=single,  % Agregar un marco alrededor del código
    captionpos=b,  % b = bottom, t = top
}

\definecolor{eclipseStrings}{RGB}{42,0.0,255}
\definecolor{eclipseKeywords}{RGB}{127,0,85}
\colorlet{numb}{magenta!60!black}

\lstdefinelanguage{json}{
    basicstyle=\ttfamily\small,
    showstringspaces=false,
    breaklines=true,
    frame=single,
    literate=
     *{0}{{{\color{numb}0}}}{1}
      {1}{{{\color{numb}1}}}{1}
      {2}{{{\color{numb}2}}}{1}
      {3}{{{\color{numb}3}}}{1}
      {4}{{{\color{numb}4}}}{1}
      {5}{{{\color{numb}5}}}{1}
      {6}{{{\color{numb}6}}}{1}
      {7}{{{\color{numb}7}}}{1}
      {8}{{{\color{numb}8}}}{1}
      {9}{{{\color{numb}9}}}{1}
      {:}{{{\color{punct}{:}}}}{1}
      {,}{{{\color{punct}{,}}}}{1}
      {\{}{{{\color{delim}{\{}}}}{1}
      {\}}{{{\color{delim}{\}}}}}{1}
      {[}{{{\color{delim}{[}}}}{1}
      {]}{{{\color{delim}{]}}}}{1},
    morestring=[b]"
}

\lstdefinelanguage{JavaScript}{
  morekeywords=[1]{break, continue, delete, else, for, function, if, in,
    new, return, this, typeof, var, void, while, with},
  % Literals, primitive types, and reference types.
  morekeywords=[2]{false, null, true, boolean, number, undefined,
    Array, Boolean, Date, Math, Number, String, Object},
  % Built-ins.
  morekeywords=[3]{eval, parseInt, parseFloat, escape, unescape},
  sensitive,
  morecomment=[s]{/*}{*/},
  morecomment=[l]//,
  morecomment=[s]{/**}{*/}, % JavaDoc style comments
  morestring=[b]',
  morestring=[b]"
}[keywords, comments, strings]

\lstdefinelanguage{svelte}{
    morekeywords={async, await, then, if, each, in, on, let, const, function, import, export, from, with, typeof, new, var, return},
    sensitive=true,
    morecomment=[l]{//},
    morecomment=[s]{/*}{*/},
    morecomment=[s]{<!--}{-->},
    morestring=[b]',
    morestring=[b]",
    morestring=[b]`,
    ndkeywords={import, export, return, if, for, while, function, switch, case, break, var, let, const, async, await, then, catch, try, typeof, new, class, extends, implements, interface, constructor, public, private, protected, static, get, set},
    alsoletter={:},
    moredelim=**[is][\color{orange}]{\@\{}{\@\}},
    moredelim=**[is][\color{cyan}]{\$\{}{\}},
    keywordstyle=\color{blue}\bfseries,
    commentstyle=\color{gray}\ttfamily,
    stringstyle=\color{red}\ttfamily,
    identifierstyle=\color{black},
    breaklines=true,
    basicstyle=\ttfamily\small,
    columns=fullflexible,
    showstringspaces=false,
    tabsize=2,
    frame=single,
    captionpos=b,
}


\lstdefinestyle{java}{
    language=Java,
    basicstyle=\ttfamily\small,
    keywordstyle=\color{blue},
    stringstyle=\color{red},
    commentstyle=\color{gray},
    morecomment=[s][\color{gray}]{/**}{*/},
    breaklines=true,
    numbers=left,
    numberstyle=\tiny\color{gray},
    showspaces=false,
    showstringspaces=false,
    frame=single,
    captionpos=b
}

\usepackage{color}
\definecolor{numb}{rgb}{0.58,0,0.82}
\definecolor{punct}{rgb}{0.3,0.3,0.3}
\definecolor{delim}{rgb}{0.5,0.5,0.5}
\definecolor{string}{rgb}{0.6,0.1,0.1}

\lstdefinestyle{json}{
    language=json,
    basicstyle=\ttfamily\small,
    keywordstyle=\color{blue},
    stringstyle=\color{string},
    commentstyle=\color{green},
    morecomment=[s][\color{gray}]{/**}{*/},
    breaklines=true,
    numbers=left,
    numberstyle=\tiny\color{gray},
    showspaces=false,
    showstringspaces=false,
    frame=single,
    captionpos=b
}

\lstdefinelanguage{drl}{
    basicstyle=\ttfamily\small,
    keywordstyle=\color{blue}\bfseries,
    stringstyle=\color{red},
    commentstyle=\color{green},
    morecomment=[s][\color{green}]{/*}{*/},
    morecomment=[l][\color{green}]{//},
    morestring=[b]",
    keywords={rule, when, then, end, salience, eval, update, retract},
    sensitive=true,
    breaklines=true,
    numbers=left,
    numberstyle=\tiny\color{gray},
    showspaces=false,
    showstringspaces=false,
    frame=single,
    captionpos=b
}

\lstdefinestyle{drl}{
    language=drl,
    basicstyle=\ttfamily\small,
    keywordstyle=\color{blue}\bfseries,
    stringstyle=\color{red},
    commentstyle=\color{green},
    breaklines=true,
    numbers=left,
    numberstyle=\tiny\color{gray},
    showspaces=false,
    showstringspaces=false,
    frame=single,
    captionpos=b
}



\renewcommand{\lstlistingname}{Codi}
\renewcommand{\lstlistlistingname}{Índex de Fragments de Codis}
\renewcommand{\appendixpagename}{Annexos}
\renewcommand{\appendixtocname}{Annexos}

\renewcommand*{\listalgorithmcfname}{Llistat de pseudocodi}
\renewcommand*{\algorithmcfname}{Pseudocodi}
\renewcommand*{\algorithmautorefname}{pseudocodi}

\definecolor{darkgreen}{rgb}{0.0, 0.3, 0.0}
\definecolor{lightgreen}{rgb}{0.9, 0.95, 0.9}

\newcommand{\todo}[1]{
    \begin{tcolorbox}[colback=lightgreen, colframe=darkgreen, title=TODO]
        #1
    \end{tcolorbox}
}


% \usepackage{indentfirst} 
% \setlength{\parindent}{1cm}


\usepackage{parskip}
% \setlength{\parindent}{10pt}

% \setlength{\headheight}{15.6pt}
\setlength{\headheight}{16.05263pt}




\vspace{5mm}

\newcolumntype{u}{>{\centering\arraybackslash\hsize=.22\hsize}X}
\newcolumntype{s}{>{\hsize=.45\hsize}X}


%%%%%%%%%%%%%%%%%%%%%%%%%%%%%%%%
% INI creacio subsubsubsecion
%%%%%%%%%%%%%%%%%%%%%%%%%%%%%%%%
\usepackage{titlesec}

\let\oldparagraph\paragraph
\renewcommand{\paragraph}[1]{\oldparagraph{\hspace{30pt}#1}}


\titleclass{\subsubsubsection}{straight}[\subsection]

\newcounter{subsubsubsection}[subsubsection]
\renewcommand\thesubsubsubsection{\thesubsubsection.\arabic{subsubsubsection}}
\renewcommand\theparagraph{\thesubsubsubsection.\arabic{paragraph}} % optional; useful if paragraphs are to be numbered

\titleformat{\subsubsubsection}
  {\normalfont\normalsize\bfseries}{\thesubsubsubsection}{1em}{}
\titlespacing*{\subsubsubsection}
{0pt}{3.25ex plus 1ex minus .2ex}{1.5ex plus .2ex}


\makeatletter
\renewcommand\paragraph{\@startsection{paragraph}{5}{10pt}%
  {3.25ex \@plus1ex \@minus.2ex}%
  {-1em}%
  {\normalfont\normalsize\bfseries}}
\renewcommand\subparagraph{\@startsection{subparagraph}{6}{\parindent}%
  {3.25ex \@plus1ex \@minus .2ex}%
  {-1em}%
  {\normalfont\normalsize\bfseries}}
\def\toclevel@subsubsubsection{4}
\def\toclevel@paragraph{5}
\def\toclevel@paragraph{6}
\def\l@subsubsubsection{\@dottedtocline{4}{7em}{4em}}
\def\l@paragraph{\@dottedtocline{5}{10em}{5em}}
\def\l@subparagraph{\@dottedtocline{6}{14em}{6em}}
\makeatother

\setcounter{secnumdepth}{4}
\setcounter{tocdepth}{3}
%%%%%%%%%%%%%%%%%%%%%%%%%%%%%%%%
% FIN creacio subsubsubtitle
%%%%%%%%%%%%%%%%%%%%%%%%%%%%%%%%


%%%%%%%%%%%%%%%%%%%%%%%%%%%%%%%%
% INI creacio equation
%%%%%%%%%%%%%%%%%%%%%%%%%%%%%%%%
\DeclareFloatingEnvironment[fileext=loeq,placement={!ht},name=Equació]{equationfloat}
%%%%%%%%%%%%%%%%%%%%%%%%%%%%%%%%
% FIN creacio equation
%%%%%%%%%%%%%%%%%%%%%%%%%%%%%%%%

\titleformat{\section}[block]
  {\clearpage\normalfont\Large\bfseries}
  {\thesection}{1em}{}


\renewcommand\refname{Bibliografia}


\linespread{1.25}


\makeatletter
\renewenvironment{abstract}{%
    \if@twocolumn
      \section*{\abstractname}%
    \else 
      \begin{center}%
        {\bfseries \Large\itshape\abstractname\vspace{\z@}}%  % Título en negrita, 14pt y cursiva
      \end{center}%
      \quotation
    \fi}
    {\if@twocolumn\else\endquotation\fi}
\makeatother

% \fancypagestyle{main}{
%     \fancyhf{}
%     \lhead{\small Aplicació Web per a la Visualització i Comparació
% de Trajectòries d’Algorismes d’Optimització}
%     % \rhead{\includegraphics[width=1.2cm]{resources/Icons/logo-fib.png}}
%     \cfoot{\thepage}
%     \setlength{\footskip}{35pt}
% }

\newcommand{\q}[1]{``#1''}

\begin{document}

\hypersetup{
    pdfborder=0 0 0,
}

%% PORTADA I TABLE OF CONTENTS %%

%%%%                                   %%%%
%         INDEX DE LA MEMORIA             %
%%%%                                   %%%%

% \includepdf[pages=-]{portada_fib.pdf}
\input{Extra/GEP}
%\input{Extra/portada}
% \cleardoublepage

% \pagestyle{main}
% \setcounter{page}{11}
\include{toc}
% \cleardoublepage
\input{Contexto y alcance/main}

\input{Planificación temporal/main}

\input{Gestión económica/main}

\input{Informe de sostenibilidad/main}

% \cleardoublepage

% \cleardoublepage


% \input{Annex/main}


\clearpage
% Estilo bibliográfico IEEE
\bibliographystyle{IEEEtran}
\renewcommand{\refname}{Referencias}
\bibliography{references.bib}

%%%%                                   %%%%
%         FINAL DE LA MEMORIA             %
%%%%                                   %%%%

\end{document}


%% VERSIÓ INCIAL MEMORIA %%



\documentclass[a4paper, 11pt, twoside]{article}
\usepackage{wrapfig}
% \usepackage{abstract}
\usepackage{graphicx} % Required for inserting images
\usepackage[a4paper, left=1.15in, right=1.15in, top=1.5in, bottom=1.2in]{geometry}
\usepackage{float}
\usepackage{array}
\usepackage{booktabs}
\usepackage[T1]{fontenc}
\usepackage{amsmath}
\usepackage{tabularx}
\usepackage{graphicx} % Required for inserting images
\usepackage{tikz, xcolor}
\usepackage{hyperref}
\usepackage{cite}
\usepackage{pdfpages}
\usepackage{pdflscape}
\usepackage{afterpage}
\usepackage{textpos}
\usepackage{fancyhdr}
\usepackage{caption}
\usepackage{lmodern,textcomp}
\usepackage[catalan]{babel}
\usepackage{lipsum}
\usepackage{newfloat}
\usepackage{listings}
\usepackage{placeins}
\usepackage[ruled]{algorithm2e}
\usepackage{tcolorbox}
% Estilo de citas IEEE sencillo: usar paquete cite para comprimir y ordenar referencias numéricas
\usepackage{cite}
\usepackage[toc,page]{appendix}
\usepackage{forest}
\usepackage{pdflscape}
\usepackage{graphicx}
\usepackage{geometry}
\usepackage{quoting,xparse}
\usepackage{svg}



% \NewDocumentCommand{\bywhom}{m}{% the Bourbaki trick
%   {\nobreak\hfill\penalty50\hskip1em\null\nobreak
%    \hfill\mbox{\normalfont(#1)}%
%    \parfillskip=0pt \finalhyphendemerits=0 \par}%
% }

% AIXÒ SERVEIX PER A MODIFICAR UN PROBLEMA AMB ELS CARACTERS --> " + PARAULA 

% \NewDocumentEnvironment{pquotation}{m}
%   {\begin{quoting}[
%      indentfirst=true,
%      leftmargin=\parindent,
%      rightmargin=\parindent]\itshape}
%   {\bywhom{#1}
  
%   \end{quoting}}

% AIXÒ PER INSERIR FRAGMENTS DE CODI

\lstnewenvironment{code}[1][]%
{
   \noindent
   \minipage{\linewidth}
   \vspace{0.5\baselineskip}
   \lstset{
     basicstyle=\ttfamily\footnotesize,
     frame=single,
     breaklines=true,    % Permitir que las líneas se rompan automáticamente
     breakatwhitespace=true, % Permitir la ruptura en los espacios en blanco
     #1
   }
}
{\endminipage}

\lstset{
    language=Java,
    basicstyle=\ttfamily\small,  % Define el tipo de letra y tamaño del código
    keywordstyle=\color{blue},   % Define el color de las palabras clave de Java
    stringstyle=\color{red},     % Define el color de las cadenas de texto
    commentstyle=\color{green},  % Define el color de los comentarios
    morecomment=[s][\color{gray}]{/**}{*/},  % Estilo para comentarios de bloque
    breaklines=true,  % Configura el quiebre automático de línea
    numbers=left,  % Posición de los números de línea
    numberstyle=\tiny\color{gray},  % Estilo para los números de línea
    showspaces=false,  % No mostrar espacios con underscores
    showstringspaces=false,  % No mostrar espacios dentro de strings
    frame=single,  % Agregar un marco alrededor del código
    captionpos=b,  % b = bottom, t = top
}

\definecolor{eclipseStrings}{RGB}{42,0.0,255}
\definecolor{eclipseKeywords}{RGB}{127,0,85}
\colorlet{numb}{magenta!60!black}

\lstdefinelanguage{json}{
    basicstyle=\ttfamily\small,
    showstringspaces=false,
    breaklines=true,
    frame=single,
    literate=
     *{0}{{{\color{numb}0}}}{1}
      {1}{{{\color{numb}1}}}{1}
      {2}{{{\color{numb}2}}}{1}
      {3}{{{\color{numb}3}}}{1}
      {4}{{{\color{numb}4}}}{1}
      {5}{{{\color{numb}5}}}{1}
      {6}{{{\color{numb}6}}}{1}
      {7}{{{\color{numb}7}}}{1}
      {8}{{{\color{numb}8}}}{1}
      {9}{{{\color{numb}9}}}{1}
      {:}{{{\color{punct}{:}}}}{1}
      {,}{{{\color{punct}{,}}}}{1}
      {\{}{{{\color{delim}{\{}}}}{1}
      {\}}{{{\color{delim}{\}}}}}{1}
      {[}{{{\color{delim}{[}}}}{1}
      {]}{{{\color{delim}{]}}}}{1},
    morestring=[b]"
}

\lstdefinelanguage{JavaScript}{
  morekeywords=[1]{break, continue, delete, else, for, function, if, in,
    new, return, this, typeof, var, void, while, with},
  % Literals, primitive types, and reference types.
  morekeywords=[2]{false, null, true, boolean, number, undefined,
    Array, Boolean, Date, Math, Number, String, Object},
  % Built-ins.
  morekeywords=[3]{eval, parseInt, parseFloat, escape, unescape},
  sensitive,
  morecomment=[s]{/*}{*/},
  morecomment=[l]//,
  morecomment=[s]{/**}{*/}, % JavaDoc style comments
  morestring=[b]',
  morestring=[b]"
}[keywords, comments, strings]

\lstdefinelanguage{svelte}{
    morekeywords={async, await, then, if, each, in, on, let, const, function, import, export, from, with, typeof, new, var, return},
    sensitive=true,
    morecomment=[l]{//},
    morecomment=[s]{/*}{*/},
    morecomment=[s]{<!--}{-->},
    morestring=[b]',
    morestring=[b]",
    morestring=[b]`,
    ndkeywords={import, export, return, if, for, while, function, switch, case, break, var, let, const, async, await, then, catch, try, typeof, new, class, extends, implements, interface, constructor, public, private, protected, static, get, set},
    alsoletter={:},
    moredelim=**[is][\color{orange}]{\@\{}{\@\}},
    moredelim=**[is][\color{cyan}]{\$\{}{\}},
    keywordstyle=\color{blue}\bfseries,
    commentstyle=\color{gray}\ttfamily,
    stringstyle=\color{red}\ttfamily,
    identifierstyle=\color{black},
    breaklines=true,
    basicstyle=\ttfamily\small,
    columns=fullflexible,
    showstringspaces=false,
    tabsize=2,
    frame=single,
    captionpos=b,
}


\lstdefinestyle{java}{
    language=Java,
    basicstyle=\ttfamily\small,
    keywordstyle=\color{blue},
    stringstyle=\color{red},
    commentstyle=\color{gray},
    morecomment=[s][\color{gray}]{/**}{*/},
    breaklines=true,
    numbers=left,
    numberstyle=\tiny\color{gray},
    showspaces=false,
    showstringspaces=false,
    frame=single,
    captionpos=b
}

\usepackage{color}
\definecolor{numb}{rgb}{0.58,0,0.82}
\definecolor{punct}{rgb}{0.3,0.3,0.3}
\definecolor{delim}{rgb}{0.5,0.5,0.5}
\definecolor{string}{rgb}{0.6,0.1,0.1}

\lstdefinestyle{json}{
    language=json,
    basicstyle=\ttfamily\small,
    keywordstyle=\color{blue},
    stringstyle=\color{string},
    commentstyle=\color{green},
    morecomment=[s][\color{gray}]{/**}{*/},
    breaklines=true,
    numbers=left,
    numberstyle=\tiny\color{gray},
    showspaces=false,
    showstringspaces=false,
    frame=single,
    captionpos=b
}

\lstdefinelanguage{drl}{
    basicstyle=\ttfamily\small,
    keywordstyle=\color{blue}\bfseries,
    stringstyle=\color{red},
    commentstyle=\color{green},
    morecomment=[s][\color{green}]{/*}{*/},
    morecomment=[l][\color{green}]{//},
    morestring=[b]",
    keywords={rule, when, then, end, salience, eval, update, retract},
    sensitive=true,
    breaklines=true,
    numbers=left,
    numberstyle=\tiny\color{gray},
    showspaces=false,
    showstringspaces=false,
    frame=single,
    captionpos=b
}

\lstdefinestyle{drl}{
    language=drl,
    basicstyle=\ttfamily\small,
    keywordstyle=\color{blue}\bfseries,
    stringstyle=\color{red},
    commentstyle=\color{green},
    breaklines=true,
    numbers=left,
    numberstyle=\tiny\color{gray},
    showspaces=false,
    showstringspaces=false,
    frame=single,
    captionpos=b
}



\renewcommand{\lstlistingname}{Codi}
\renewcommand{\lstlistlistingname}{Índex de Fragments de Codis}
\renewcommand{\appendixpagename}{Annexos}
\renewcommand{\appendixtocname}{Annexos}

\renewcommand*{\listalgorithmcfname}{Llistat de pseudocodi}
\renewcommand*{\algorithmcfname}{Pseudocodi}
\renewcommand*{\algorithmautorefname}{pseudocodi}

\definecolor{darkgreen}{rgb}{0.0, 0.3, 0.0}
\definecolor{lightgreen}{rgb}{0.9, 0.95, 0.9}

\newcommand{\todo}[1]{
    \begin{tcolorbox}[colback=lightgreen, colframe=darkgreen, title=TODO]
        #1
    \end{tcolorbox}
}


% \usepackage{indentfirst} 
% \setlength{\parindent}{1cm}


\usepackage{parskip}
% \setlength{\parindent}{10pt}

% \setlength{\headheight}{15.6pt}
\setlength{\headheight}{16.05263pt}




\vspace{5mm}

\newcolumntype{u}{>{\centering\arraybackslash\hsize=.22\hsize}X}
\newcolumntype{s}{>{\hsize=.45\hsize}X}


%%%%%%%%%%%%%%%%%%%%%%%%%%%%%%%%
% INI creacio subsubsubsecion
%%%%%%%%%%%%%%%%%%%%%%%%%%%%%%%%
\usepackage{titlesec}

\let\oldparagraph\paragraph
\renewcommand{\paragraph}[1]{\oldparagraph{\hspace{30pt}#1}}


\titleclass{\subsubsubsection}{straight}[\subsection]

\newcounter{subsubsubsection}[subsubsection]
\renewcommand\thesubsubsubsection{\thesubsubsection.\arabic{subsubsubsection}}
\renewcommand\theparagraph{\thesubsubsubsection.\arabic{paragraph}} % optional; useful if paragraphs are to be numbered

\titleformat{\subsubsubsection}
  {\normalfont\normalsize\bfseries}{\thesubsubsubsection}{1em}{}
\titlespacing*{\subsubsubsection}
{0pt}{3.25ex plus 1ex minus .2ex}{1.5ex plus .2ex}


\makeatletter
\renewcommand\paragraph{\@startsection{paragraph}{5}{10pt}%
  {3.25ex \@plus1ex \@minus.2ex}%
  {-1em}%
  {\normalfont\normalsize\bfseries}}
\renewcommand\subparagraph{\@startsection{subparagraph}{6}{\parindent}%
  {3.25ex \@plus1ex \@minus .2ex}%
  {-1em}%
  {\normalfont\normalsize\bfseries}}
\def\toclevel@subsubsubsection{4}
\def\toclevel@paragraph{5}
\def\toclevel@paragraph{6}
\def\l@subsubsubsection{\@dottedtocline{4}{7em}{4em}}
\def\l@paragraph{\@dottedtocline{5}{10em}{5em}}
\def\l@subparagraph{\@dottedtocline{6}{14em}{6em}}
\makeatother

\setcounter{secnumdepth}{4}
\setcounter{tocdepth}{3}
%%%%%%%%%%%%%%%%%%%%%%%%%%%%%%%%
% FIN creacio subsubsubtitle
%%%%%%%%%%%%%%%%%%%%%%%%%%%%%%%%


%%%%%%%%%%%%%%%%%%%%%%%%%%%%%%%%
% INI creacio equation
%%%%%%%%%%%%%%%%%%%%%%%%%%%%%%%%
\DeclareFloatingEnvironment[fileext=loeq,placement={!ht},name=Equació]{equationfloat}
%%%%%%%%%%%%%%%%%%%%%%%%%%%%%%%%
% FIN creacio equation
%%%%%%%%%%%%%%%%%%%%%%%%%%%%%%%%

\titleformat{\section}[block]
  {\clearpage\normalfont\Large\bfseries}
  {\thesection}{1em}{}


\renewcommand\refname{Bibliografia}


\linespread{1.25}


\makeatletter
\renewenvironment{abstract}{%
    \if@twocolumn
      \section*{\abstractname}%
    \else 
      \begin{center}%
        {\bfseries \Large\itshape\abstractname\vspace{\z@}}%  % Título en negrita, 14pt y cursiva
      \end{center}%
      \quotation
    \fi}
    {\if@twocolumn\else\endquotation\fi}
\makeatother

% \fancypagestyle{main}{
%     \fancyhf{}
%     \lhead{\small Aplicació Web per a la Visualització i Comparació
% de Trajectòries d’Algorismes d’Optimització}
%     % \rhead{\includegraphics[width=1.2cm]{resources/Icons/logo-fib.png}}
%     \cfoot{\thepage}
%     \setlength{\footskip}{35pt}
% }

\newcommand{\q}[1]{``#1''}

\begin{document}

\hypersetup{
    pdfborder=0 0 0,
}

%% PORTADA I TABLE OF CONTENTS %%

%%%%                                   %%%%
%         INDEX DE LA MEMORIA             %
%%%%                                   %%%%

% \includepdf[pages=-]{portada_fib.pdf}
%%% Portada
\begin{titlepage}
\begin{center}
    \includegraphics[width=1\linewidth]{extra/figures/logo-upc-fib.png}
    \vspace{30mm}

    \LARGE{\textsc{\textbf{{INTEGRACIÓN Y AUTOMATIZACIÓN DE COSTES DE INFRAESTRUCTURA HPC EN ENTORNOS DE INVESTIGACIÓN}}}} \\
    \vspace{25mm}
      
        
    \textbf{\Large{\textsc{PAU NAVARRO ÁLVAREZ}}} \\
    \vspace{20mm}
        
    
    
    \large{\textbf{Director/a} ROBERTO RIVEIRO INSUA (FUNDACIO INSTITUT DE RECERCA BIOMEDICA (IRB BARCELONA))} \\
    \large{\textbf{Ponente:} ALEX PAJUELO GONZÁLEZ (Departamento de Arquitectura de Computadores)} \\
    \large{\textbf{Titulación:} Grado en Ingeniería Informática (Tecnologías de la información)} \\
    
    \large{\textbf{Lliurament 2:} }Planificación temporal \\
    \large{\textbf{Tutor GEP:} CAROLINA CONSOLIDACIÓN SEGURA (Departament d'Organització d'Empreses)} \\
    \large{\textbf{Fecha:}} 29 de Septiembre de 2025\\

    \textbf{\large{Memoria del trabajo de fin de grado}} \\
    \vspace{3mm}
    \textbf{\large{Facultat d'Informàtica de Barcelona (FIB)}} \\
    \vspace{3mm}
    \textbf{\large{Universitat Politècnica de Catalunya (UPC) - BarcelonaTech}} \\

\end{center}
\end{titlepage}
%\input{Extra/portada}
% \cleardoublepage

% \pagestyle{main}
% \setcounter{page}{11}
\renewcommand{\contentsname}{Continguts}

\tableofcontents
\clearpage
\listoffigures
\begingroup
    \let\clearpage\relax
    \listoftables
\endgroup
\clearpage

\newpage
% \cleardoublepage
\input{Contexto y alcance/Contextualización/main}
\input{Contexto y alcance/Justificación/main}
\input{Contexto y alcance/Alcance/main}
\input{Contexto y alcance/Metodología y herramientas de seguimiento/main}

\documentclass[a4paper, 11pt, twoside]{article}
\usepackage{wrapfig}
% \usepackage{abstract}
\usepackage{graphicx} % Required for inserting images
\usepackage[a4paper, left=1.15in, right=1.15in, top=1.5in, bottom=1.2in]{geometry}
\usepackage{float}
\usepackage{array}
\usepackage{booktabs}
\usepackage[T1]{fontenc}
\usepackage{amsmath}
\usepackage{tabularx}
\usepackage{graphicx} % Required for inserting images
\usepackage{tikz, xcolor}
\usepackage{hyperref}
\usepackage{cite}
\usepackage{pdfpages}
\usepackage{pdflscape}
\usepackage{afterpage}
\usepackage{textpos}
\usepackage{fancyhdr}
\usepackage{caption}
\usepackage{lmodern,textcomp}
\usepackage[catalan]{babel}
\usepackage{lipsum}
\usepackage{newfloat}
\usepackage{listings}
\usepackage{placeins}
\usepackage[ruled]{algorithm2e}
\usepackage{tcolorbox}
% Estilo de citas IEEE sencillo: usar paquete cite para comprimir y ordenar referencias numéricas
\usepackage{cite}
\usepackage[toc,page]{appendix}
\usepackage{forest}
\usepackage{pdflscape}
\usepackage{graphicx}
\usepackage{geometry}
\usepackage{quoting,xparse}
\usepackage{svg}



% \NewDocumentCommand{\bywhom}{m}{% the Bourbaki trick
%   {\nobreak\hfill\penalty50\hskip1em\null\nobreak
%    \hfill\mbox{\normalfont(#1)}%
%    \parfillskip=0pt \finalhyphendemerits=0 \par}%
% }

% AIXÒ SERVEIX PER A MODIFICAR UN PROBLEMA AMB ELS CARACTERS --> " + PARAULA 

% \NewDocumentEnvironment{pquotation}{m}
%   {\begin{quoting}[
%      indentfirst=true,
%      leftmargin=\parindent,
%      rightmargin=\parindent]\itshape}
%   {\bywhom{#1}
  
%   \end{quoting}}

% AIXÒ PER INSERIR FRAGMENTS DE CODI

\lstnewenvironment{code}[1][]%
{
   \noindent
   \minipage{\linewidth}
   \vspace{0.5\baselineskip}
   \lstset{
     basicstyle=\ttfamily\footnotesize,
     frame=single,
     breaklines=true,    % Permitir que las líneas se rompan automáticamente
     breakatwhitespace=true, % Permitir la ruptura en los espacios en blanco
     #1
   }
}
{\endminipage}

\lstset{
    language=Java,
    basicstyle=\ttfamily\small,  % Define el tipo de letra y tamaño del código
    keywordstyle=\color{blue},   % Define el color de las palabras clave de Java
    stringstyle=\color{red},     % Define el color de las cadenas de texto
    commentstyle=\color{green},  % Define el color de los comentarios
    morecomment=[s][\color{gray}]{/**}{*/},  % Estilo para comentarios de bloque
    breaklines=true,  % Configura el quiebre automático de línea
    numbers=left,  % Posición de los números de línea
    numberstyle=\tiny\color{gray},  % Estilo para los números de línea
    showspaces=false,  % No mostrar espacios con underscores
    showstringspaces=false,  % No mostrar espacios dentro de strings
    frame=single,  % Agregar un marco alrededor del código
    captionpos=b,  % b = bottom, t = top
}

\definecolor{eclipseStrings}{RGB}{42,0.0,255}
\definecolor{eclipseKeywords}{RGB}{127,0,85}
\colorlet{numb}{magenta!60!black}

\lstdefinelanguage{json}{
    basicstyle=\ttfamily\small,
    showstringspaces=false,
    breaklines=true,
    frame=single,
    literate=
     *{0}{{{\color{numb}0}}}{1}
      {1}{{{\color{numb}1}}}{1}
      {2}{{{\color{numb}2}}}{1}
      {3}{{{\color{numb}3}}}{1}
      {4}{{{\color{numb}4}}}{1}
      {5}{{{\color{numb}5}}}{1}
      {6}{{{\color{numb}6}}}{1}
      {7}{{{\color{numb}7}}}{1}
      {8}{{{\color{numb}8}}}{1}
      {9}{{{\color{numb}9}}}{1}
      {:}{{{\color{punct}{:}}}}{1}
      {,}{{{\color{punct}{,}}}}{1}
      {\{}{{{\color{delim}{\{}}}}{1}
      {\}}{{{\color{delim}{\}}}}}{1}
      {[}{{{\color{delim}{[}}}}{1}
      {]}{{{\color{delim}{]}}}}{1},
    morestring=[b]"
}

\lstdefinelanguage{JavaScript}{
  morekeywords=[1]{break, continue, delete, else, for, function, if, in,
    new, return, this, typeof, var, void, while, with},
  % Literals, primitive types, and reference types.
  morekeywords=[2]{false, null, true, boolean, number, undefined,
    Array, Boolean, Date, Math, Number, String, Object},
  % Built-ins.
  morekeywords=[3]{eval, parseInt, parseFloat, escape, unescape},
  sensitive,
  morecomment=[s]{/*}{*/},
  morecomment=[l]//,
  morecomment=[s]{/**}{*/}, % JavaDoc style comments
  morestring=[b]',
  morestring=[b]"
}[keywords, comments, strings]

\lstdefinelanguage{svelte}{
    morekeywords={async, await, then, if, each, in, on, let, const, function, import, export, from, with, typeof, new, var, return},
    sensitive=true,
    morecomment=[l]{//},
    morecomment=[s]{/*}{*/},
    morecomment=[s]{<!--}{-->},
    morestring=[b]',
    morestring=[b]",
    morestring=[b]`,
    ndkeywords={import, export, return, if, for, while, function, switch, case, break, var, let, const, async, await, then, catch, try, typeof, new, class, extends, implements, interface, constructor, public, private, protected, static, get, set},
    alsoletter={:},
    moredelim=**[is][\color{orange}]{\@\{}{\@\}},
    moredelim=**[is][\color{cyan}]{\$\{}{\}},
    keywordstyle=\color{blue}\bfseries,
    commentstyle=\color{gray}\ttfamily,
    stringstyle=\color{red}\ttfamily,
    identifierstyle=\color{black},
    breaklines=true,
    basicstyle=\ttfamily\small,
    columns=fullflexible,
    showstringspaces=false,
    tabsize=2,
    frame=single,
    captionpos=b,
}


\lstdefinestyle{java}{
    language=Java,
    basicstyle=\ttfamily\small,
    keywordstyle=\color{blue},
    stringstyle=\color{red},
    commentstyle=\color{gray},
    morecomment=[s][\color{gray}]{/**}{*/},
    breaklines=true,
    numbers=left,
    numberstyle=\tiny\color{gray},
    showspaces=false,
    showstringspaces=false,
    frame=single,
    captionpos=b
}

\usepackage{color}
\definecolor{numb}{rgb}{0.58,0,0.82}
\definecolor{punct}{rgb}{0.3,0.3,0.3}
\definecolor{delim}{rgb}{0.5,0.5,0.5}
\definecolor{string}{rgb}{0.6,0.1,0.1}

\lstdefinestyle{json}{
    language=json,
    basicstyle=\ttfamily\small,
    keywordstyle=\color{blue},
    stringstyle=\color{string},
    commentstyle=\color{green},
    morecomment=[s][\color{gray}]{/**}{*/},
    breaklines=true,
    numbers=left,
    numberstyle=\tiny\color{gray},
    showspaces=false,
    showstringspaces=false,
    frame=single,
    captionpos=b
}

\lstdefinelanguage{drl}{
    basicstyle=\ttfamily\small,
    keywordstyle=\color{blue}\bfseries,
    stringstyle=\color{red},
    commentstyle=\color{green},
    morecomment=[s][\color{green}]{/*}{*/},
    morecomment=[l][\color{green}]{//},
    morestring=[b]",
    keywords={rule, when, then, end, salience, eval, update, retract},
    sensitive=true,
    breaklines=true,
    numbers=left,
    numberstyle=\tiny\color{gray},
    showspaces=false,
    showstringspaces=false,
    frame=single,
    captionpos=b
}

\lstdefinestyle{drl}{
    language=drl,
    basicstyle=\ttfamily\small,
    keywordstyle=\color{blue}\bfseries,
    stringstyle=\color{red},
    commentstyle=\color{green},
    breaklines=true,
    numbers=left,
    numberstyle=\tiny\color{gray},
    showspaces=false,
    showstringspaces=false,
    frame=single,
    captionpos=b
}



\renewcommand{\lstlistingname}{Codi}
\renewcommand{\lstlistlistingname}{Índex de Fragments de Codis}
\renewcommand{\appendixpagename}{Annexos}
\renewcommand{\appendixtocname}{Annexos}

\renewcommand*{\listalgorithmcfname}{Llistat de pseudocodi}
\renewcommand*{\algorithmcfname}{Pseudocodi}
\renewcommand*{\algorithmautorefname}{pseudocodi}

\definecolor{darkgreen}{rgb}{0.0, 0.3, 0.0}
\definecolor{lightgreen}{rgb}{0.9, 0.95, 0.9}

\newcommand{\todo}[1]{
    \begin{tcolorbox}[colback=lightgreen, colframe=darkgreen, title=TODO]
        #1
    \end{tcolorbox}
}


% \usepackage{indentfirst} 
% \setlength{\parindent}{1cm}


\usepackage{parskip}
% \setlength{\parindent}{10pt}

% \setlength{\headheight}{15.6pt}
\setlength{\headheight}{16.05263pt}




\vspace{5mm}

\newcolumntype{u}{>{\centering\arraybackslash\hsize=.22\hsize}X}
\newcolumntype{s}{>{\hsize=.45\hsize}X}


%%%%%%%%%%%%%%%%%%%%%%%%%%%%%%%%
% INI creacio subsubsubsecion
%%%%%%%%%%%%%%%%%%%%%%%%%%%%%%%%
\usepackage{titlesec}

\let\oldparagraph\paragraph
\renewcommand{\paragraph}[1]{\oldparagraph{\hspace{30pt}#1}}


\titleclass{\subsubsubsection}{straight}[\subsection]

\newcounter{subsubsubsection}[subsubsection]
\renewcommand\thesubsubsubsection{\thesubsubsection.\arabic{subsubsubsection}}
\renewcommand\theparagraph{\thesubsubsubsection.\arabic{paragraph}} % optional; useful if paragraphs are to be numbered

\titleformat{\subsubsubsection}
  {\normalfont\normalsize\bfseries}{\thesubsubsubsection}{1em}{}
\titlespacing*{\subsubsubsection}
{0pt}{3.25ex plus 1ex minus .2ex}{1.5ex plus .2ex}


\makeatletter
\renewcommand\paragraph{\@startsection{paragraph}{5}{10pt}%
  {3.25ex \@plus1ex \@minus.2ex}%
  {-1em}%
  {\normalfont\normalsize\bfseries}}
\renewcommand\subparagraph{\@startsection{subparagraph}{6}{\parindent}%
  {3.25ex \@plus1ex \@minus .2ex}%
  {-1em}%
  {\normalfont\normalsize\bfseries}}
\def\toclevel@subsubsubsection{4}
\def\toclevel@paragraph{5}
\def\toclevel@paragraph{6}
\def\l@subsubsubsection{\@dottedtocline{4}{7em}{4em}}
\def\l@paragraph{\@dottedtocline{5}{10em}{5em}}
\def\l@subparagraph{\@dottedtocline{6}{14em}{6em}}
\makeatother

\setcounter{secnumdepth}{4}
\setcounter{tocdepth}{3}
%%%%%%%%%%%%%%%%%%%%%%%%%%%%%%%%
% FIN creacio subsubsubtitle
%%%%%%%%%%%%%%%%%%%%%%%%%%%%%%%%


%%%%%%%%%%%%%%%%%%%%%%%%%%%%%%%%
% INI creacio equation
%%%%%%%%%%%%%%%%%%%%%%%%%%%%%%%%
\DeclareFloatingEnvironment[fileext=loeq,placement={!ht},name=Equació]{equationfloat}
%%%%%%%%%%%%%%%%%%%%%%%%%%%%%%%%
% FIN creacio equation
%%%%%%%%%%%%%%%%%%%%%%%%%%%%%%%%

\titleformat{\section}[block]
  {\clearpage\normalfont\Large\bfseries}
  {\thesection}{1em}{}


\renewcommand\refname{Bibliografia}


\linespread{1.25}


\makeatletter
\renewenvironment{abstract}{%
    \if@twocolumn
      \section*{\abstractname}%
    \else 
      \begin{center}%
        {\bfseries \Large\itshape\abstractname\vspace{\z@}}%  % Título en negrita, 14pt y cursiva
      \end{center}%
      \quotation
    \fi}
    {\if@twocolumn\else\endquotation\fi}
\makeatother

% \fancypagestyle{main}{
%     \fancyhf{}
%     \lhead{\small Aplicació Web per a la Visualització i Comparació
% de Trajectòries d’Algorismes d’Optimització}
%     % \rhead{\includegraphics[width=1.2cm]{resources/Icons/logo-fib.png}}
%     \cfoot{\thepage}
%     \setlength{\footskip}{35pt}
% }

\newcommand{\q}[1]{``#1''}

\begin{document}

\hypersetup{
    pdfborder=0 0 0,
}

%% PORTADA I TABLE OF CONTENTS %%

%%%%                                   %%%%
%         INDEX DE LA MEMORIA             %
%%%%                                   %%%%

% \includepdf[pages=-]{portada_fib.pdf}
\input{Extra/GEP}
%\input{Extra/portada}
% \cleardoublepage

% \pagestyle{main}
% \setcounter{page}{11}
\include{toc}
% \cleardoublepage
\input{Contexto y alcance/main}

\input{Planificación temporal/main}

\input{Gestión económica/main}

\input{Informe de sostenibilidad/main}

% \cleardoublepage

% \cleardoublepage


% \input{Annex/main}


\clearpage
% Estilo bibliográfico IEEE
\bibliographystyle{IEEEtran}
\renewcommand{\refname}{Referencias}
\bibliography{references.bib}

%%%%                                   %%%%
%         FINAL DE LA MEMORIA             %
%%%%                                   %%%%

\end{document}


%% VERSIÓ INCIAL MEMORIA %%



\section{Gestión económica}

La gestión económica de un proyecto es fundamental para asegurar su viabilidad y exito a largo plazo. Ya hemos visto en secciones anteriores, los otros dos vertices 
del triángulo de la gestión de proyectos: el alcance y el tiempo. En esta sección, nos centraremos en cómo gestionar eficazmente los recursos financieros para garantizar
que el proyecto se mantenga dentro del presupuesto y se logren los objetivos establecidos. 

\input{Gestión económica/Recursos humanos.tex}

\input{Gestión económica/Recursos Hardware.tex}

\input{Gestión económica/Recursos Software.tex}

\input{Gestión económica/Contingencias.tex}

\section{Informe de sostenibilidad}

\input{Informe de sostenibilidad/Autoevaluación.tex}


\input{Informe de sostenibilidad/Dimensión ambiental.tex}

\input{Informe de sostenibilidad/Dimensión económica.tex}

\input{Informe de sostenibilidad/Dimensión social.tex}


% \cleardoublepage

% \cleardoublepage


% \documentclass[a4paper, 11pt, twoside]{article}
\usepackage{wrapfig}
% \usepackage{abstract}
\usepackage{graphicx} % Required for inserting images
\usepackage[a4paper, left=1.15in, right=1.15in, top=1.5in, bottom=1.2in]{geometry}
\usepackage{float}
\usepackage{array}
\usepackage{booktabs}
\usepackage[T1]{fontenc}
\usepackage{amsmath}
\usepackage{tabularx}
\usepackage{graphicx} % Required for inserting images
\usepackage{tikz, xcolor}
\usepackage{hyperref}
\usepackage{cite}
\usepackage{pdfpages}
\usepackage{pdflscape}
\usepackage{afterpage}
\usepackage{textpos}
\usepackage{fancyhdr}
\usepackage{caption}
\usepackage{lmodern,textcomp}
\usepackage[catalan]{babel}
\usepackage{lipsum}
\usepackage{newfloat}
\usepackage{listings}
\usepackage{placeins}
\usepackage[ruled]{algorithm2e}
\usepackage{tcolorbox}
% Estilo de citas IEEE sencillo: usar paquete cite para comprimir y ordenar referencias numéricas
\usepackage{cite}
\usepackage[toc,page]{appendix}
\usepackage{forest}
\usepackage{pdflscape}
\usepackage{graphicx}
\usepackage{geometry}
\usepackage{quoting,xparse}
\usepackage{svg}



% \NewDocumentCommand{\bywhom}{m}{% the Bourbaki trick
%   {\nobreak\hfill\penalty50\hskip1em\null\nobreak
%    \hfill\mbox{\normalfont(#1)}%
%    \parfillskip=0pt \finalhyphendemerits=0 \par}%
% }

% AIXÒ SERVEIX PER A MODIFICAR UN PROBLEMA AMB ELS CARACTERS --> " + PARAULA 

% \NewDocumentEnvironment{pquotation}{m}
%   {\begin{quoting}[
%      indentfirst=true,
%      leftmargin=\parindent,
%      rightmargin=\parindent]\itshape}
%   {\bywhom{#1}
  
%   \end{quoting}}

% AIXÒ PER INSERIR FRAGMENTS DE CODI

\lstnewenvironment{code}[1][]%
{
   \noindent
   \minipage{\linewidth}
   \vspace{0.5\baselineskip}
   \lstset{
     basicstyle=\ttfamily\footnotesize,
     frame=single,
     breaklines=true,    % Permitir que las líneas se rompan automáticamente
     breakatwhitespace=true, % Permitir la ruptura en los espacios en blanco
     #1
   }
}
{\endminipage}

\lstset{
    language=Java,
    basicstyle=\ttfamily\small,  % Define el tipo de letra y tamaño del código
    keywordstyle=\color{blue},   % Define el color de las palabras clave de Java
    stringstyle=\color{red},     % Define el color de las cadenas de texto
    commentstyle=\color{green},  % Define el color de los comentarios
    morecomment=[s][\color{gray}]{/**}{*/},  % Estilo para comentarios de bloque
    breaklines=true,  % Configura el quiebre automático de línea
    numbers=left,  % Posición de los números de línea
    numberstyle=\tiny\color{gray},  % Estilo para los números de línea
    showspaces=false,  % No mostrar espacios con underscores
    showstringspaces=false,  % No mostrar espacios dentro de strings
    frame=single,  % Agregar un marco alrededor del código
    captionpos=b,  % b = bottom, t = top
}

\definecolor{eclipseStrings}{RGB}{42,0.0,255}
\definecolor{eclipseKeywords}{RGB}{127,0,85}
\colorlet{numb}{magenta!60!black}

\lstdefinelanguage{json}{
    basicstyle=\ttfamily\small,
    showstringspaces=false,
    breaklines=true,
    frame=single,
    literate=
     *{0}{{{\color{numb}0}}}{1}
      {1}{{{\color{numb}1}}}{1}
      {2}{{{\color{numb}2}}}{1}
      {3}{{{\color{numb}3}}}{1}
      {4}{{{\color{numb}4}}}{1}
      {5}{{{\color{numb}5}}}{1}
      {6}{{{\color{numb}6}}}{1}
      {7}{{{\color{numb}7}}}{1}
      {8}{{{\color{numb}8}}}{1}
      {9}{{{\color{numb}9}}}{1}
      {:}{{{\color{punct}{:}}}}{1}
      {,}{{{\color{punct}{,}}}}{1}
      {\{}{{{\color{delim}{\{}}}}{1}
      {\}}{{{\color{delim}{\}}}}}{1}
      {[}{{{\color{delim}{[}}}}{1}
      {]}{{{\color{delim}{]}}}}{1},
    morestring=[b]"
}

\lstdefinelanguage{JavaScript}{
  morekeywords=[1]{break, continue, delete, else, for, function, if, in,
    new, return, this, typeof, var, void, while, with},
  % Literals, primitive types, and reference types.
  morekeywords=[2]{false, null, true, boolean, number, undefined,
    Array, Boolean, Date, Math, Number, String, Object},
  % Built-ins.
  morekeywords=[3]{eval, parseInt, parseFloat, escape, unescape},
  sensitive,
  morecomment=[s]{/*}{*/},
  morecomment=[l]//,
  morecomment=[s]{/**}{*/}, % JavaDoc style comments
  morestring=[b]',
  morestring=[b]"
}[keywords, comments, strings]

\lstdefinelanguage{svelte}{
    morekeywords={async, await, then, if, each, in, on, let, const, function, import, export, from, with, typeof, new, var, return},
    sensitive=true,
    morecomment=[l]{//},
    morecomment=[s]{/*}{*/},
    morecomment=[s]{<!--}{-->},
    morestring=[b]',
    morestring=[b]",
    morestring=[b]`,
    ndkeywords={import, export, return, if, for, while, function, switch, case, break, var, let, const, async, await, then, catch, try, typeof, new, class, extends, implements, interface, constructor, public, private, protected, static, get, set},
    alsoletter={:},
    moredelim=**[is][\color{orange}]{\@\{}{\@\}},
    moredelim=**[is][\color{cyan}]{\$\{}{\}},
    keywordstyle=\color{blue}\bfseries,
    commentstyle=\color{gray}\ttfamily,
    stringstyle=\color{red}\ttfamily,
    identifierstyle=\color{black},
    breaklines=true,
    basicstyle=\ttfamily\small,
    columns=fullflexible,
    showstringspaces=false,
    tabsize=2,
    frame=single,
    captionpos=b,
}


\lstdefinestyle{java}{
    language=Java,
    basicstyle=\ttfamily\small,
    keywordstyle=\color{blue},
    stringstyle=\color{red},
    commentstyle=\color{gray},
    morecomment=[s][\color{gray}]{/**}{*/},
    breaklines=true,
    numbers=left,
    numberstyle=\tiny\color{gray},
    showspaces=false,
    showstringspaces=false,
    frame=single,
    captionpos=b
}

\usepackage{color}
\definecolor{numb}{rgb}{0.58,0,0.82}
\definecolor{punct}{rgb}{0.3,0.3,0.3}
\definecolor{delim}{rgb}{0.5,0.5,0.5}
\definecolor{string}{rgb}{0.6,0.1,0.1}

\lstdefinestyle{json}{
    language=json,
    basicstyle=\ttfamily\small,
    keywordstyle=\color{blue},
    stringstyle=\color{string},
    commentstyle=\color{green},
    morecomment=[s][\color{gray}]{/**}{*/},
    breaklines=true,
    numbers=left,
    numberstyle=\tiny\color{gray},
    showspaces=false,
    showstringspaces=false,
    frame=single,
    captionpos=b
}

\lstdefinelanguage{drl}{
    basicstyle=\ttfamily\small,
    keywordstyle=\color{blue}\bfseries,
    stringstyle=\color{red},
    commentstyle=\color{green},
    morecomment=[s][\color{green}]{/*}{*/},
    morecomment=[l][\color{green}]{//},
    morestring=[b]",
    keywords={rule, when, then, end, salience, eval, update, retract},
    sensitive=true,
    breaklines=true,
    numbers=left,
    numberstyle=\tiny\color{gray},
    showspaces=false,
    showstringspaces=false,
    frame=single,
    captionpos=b
}

\lstdefinestyle{drl}{
    language=drl,
    basicstyle=\ttfamily\small,
    keywordstyle=\color{blue}\bfseries,
    stringstyle=\color{red},
    commentstyle=\color{green},
    breaklines=true,
    numbers=left,
    numberstyle=\tiny\color{gray},
    showspaces=false,
    showstringspaces=false,
    frame=single,
    captionpos=b
}



\renewcommand{\lstlistingname}{Codi}
\renewcommand{\lstlistlistingname}{Índex de Fragments de Codis}
\renewcommand{\appendixpagename}{Annexos}
\renewcommand{\appendixtocname}{Annexos}

\renewcommand*{\listalgorithmcfname}{Llistat de pseudocodi}
\renewcommand*{\algorithmcfname}{Pseudocodi}
\renewcommand*{\algorithmautorefname}{pseudocodi}

\definecolor{darkgreen}{rgb}{0.0, 0.3, 0.0}
\definecolor{lightgreen}{rgb}{0.9, 0.95, 0.9}

\newcommand{\todo}[1]{
    \begin{tcolorbox}[colback=lightgreen, colframe=darkgreen, title=TODO]
        #1
    \end{tcolorbox}
}


% \usepackage{indentfirst} 
% \setlength{\parindent}{1cm}


\usepackage{parskip}
% \setlength{\parindent}{10pt}

% \setlength{\headheight}{15.6pt}
\setlength{\headheight}{16.05263pt}




\vspace{5mm}

\newcolumntype{u}{>{\centering\arraybackslash\hsize=.22\hsize}X}
\newcolumntype{s}{>{\hsize=.45\hsize}X}


%%%%%%%%%%%%%%%%%%%%%%%%%%%%%%%%
% INI creacio subsubsubsecion
%%%%%%%%%%%%%%%%%%%%%%%%%%%%%%%%
\usepackage{titlesec}

\let\oldparagraph\paragraph
\renewcommand{\paragraph}[1]{\oldparagraph{\hspace{30pt}#1}}


\titleclass{\subsubsubsection}{straight}[\subsection]

\newcounter{subsubsubsection}[subsubsection]
\renewcommand\thesubsubsubsection{\thesubsubsection.\arabic{subsubsubsection}}
\renewcommand\theparagraph{\thesubsubsubsection.\arabic{paragraph}} % optional; useful if paragraphs are to be numbered

\titleformat{\subsubsubsection}
  {\normalfont\normalsize\bfseries}{\thesubsubsubsection}{1em}{}
\titlespacing*{\subsubsubsection}
{0pt}{3.25ex plus 1ex minus .2ex}{1.5ex plus .2ex}


\makeatletter
\renewcommand\paragraph{\@startsection{paragraph}{5}{10pt}%
  {3.25ex \@plus1ex \@minus.2ex}%
  {-1em}%
  {\normalfont\normalsize\bfseries}}
\renewcommand\subparagraph{\@startsection{subparagraph}{6}{\parindent}%
  {3.25ex \@plus1ex \@minus .2ex}%
  {-1em}%
  {\normalfont\normalsize\bfseries}}
\def\toclevel@subsubsubsection{4}
\def\toclevel@paragraph{5}
\def\toclevel@paragraph{6}
\def\l@subsubsubsection{\@dottedtocline{4}{7em}{4em}}
\def\l@paragraph{\@dottedtocline{5}{10em}{5em}}
\def\l@subparagraph{\@dottedtocline{6}{14em}{6em}}
\makeatother

\setcounter{secnumdepth}{4}
\setcounter{tocdepth}{3}
%%%%%%%%%%%%%%%%%%%%%%%%%%%%%%%%
% FIN creacio subsubsubtitle
%%%%%%%%%%%%%%%%%%%%%%%%%%%%%%%%


%%%%%%%%%%%%%%%%%%%%%%%%%%%%%%%%
% INI creacio equation
%%%%%%%%%%%%%%%%%%%%%%%%%%%%%%%%
\DeclareFloatingEnvironment[fileext=loeq,placement={!ht},name=Equació]{equationfloat}
%%%%%%%%%%%%%%%%%%%%%%%%%%%%%%%%
% FIN creacio equation
%%%%%%%%%%%%%%%%%%%%%%%%%%%%%%%%

\titleformat{\section}[block]
  {\clearpage\normalfont\Large\bfseries}
  {\thesection}{1em}{}


\renewcommand\refname{Bibliografia}


\linespread{1.25}


\makeatletter
\renewenvironment{abstract}{%
    \if@twocolumn
      \section*{\abstractname}%
    \else 
      \begin{center}%
        {\bfseries \Large\itshape\abstractname\vspace{\z@}}%  % Título en negrita, 14pt y cursiva
      \end{center}%
      \quotation
    \fi}
    {\if@twocolumn\else\endquotation\fi}
\makeatother

% \fancypagestyle{main}{
%     \fancyhf{}
%     \lhead{\small Aplicació Web per a la Visualització i Comparació
% de Trajectòries d’Algorismes d’Optimització}
%     % \rhead{\includegraphics[width=1.2cm]{resources/Icons/logo-fib.png}}
%     \cfoot{\thepage}
%     \setlength{\footskip}{35pt}
% }

\newcommand{\q}[1]{``#1''}

\begin{document}

\hypersetup{
    pdfborder=0 0 0,
}

%% PORTADA I TABLE OF CONTENTS %%

%%%%                                   %%%%
%         INDEX DE LA MEMORIA             %
%%%%                                   %%%%

% \includepdf[pages=-]{portada_fib.pdf}
\input{Extra/GEP}
%\input{Extra/portada}
% \cleardoublepage

% \pagestyle{main}
% \setcounter{page}{11}
\include{toc}
% \cleardoublepage
\input{Contexto y alcance/main}

\input{Planificación temporal/main}

\input{Gestión económica/main}

\input{Informe de sostenibilidad/main}

% \cleardoublepage

% \cleardoublepage


% \input{Annex/main}


\clearpage
% Estilo bibliográfico IEEE
\bibliographystyle{IEEEtran}
\renewcommand{\refname}{Referencias}
\bibliography{references.bib}

%%%%                                   %%%%
%         FINAL DE LA MEMORIA             %
%%%%                                   %%%%

\end{document}


%% VERSIÓ INCIAL MEMORIA %%




\clearpage
% Estilo bibliográfico IEEE
\bibliographystyle{IEEEtran}
\renewcommand{\refname}{Referencias}
\bibliography{references.bib}

%%%%                                   %%%%
%         FINAL DE LA MEMORIA             %
%%%%                                   %%%%

\end{document}


%% VERSIÓ INCIAL MEMORIA %%



\section{Gestión económica}

La gestión económica de un proyecto es fundamental para asegurar su viabilidad y exito a largo plazo. Ya hemos visto en secciones anteriores, los otros dos vertices 
del triángulo de la gestión de proyectos: el alcance y el tiempo. En esta sección, nos centraremos en cómo gestionar eficazmente los recursos financieros para garantizar
que el proyecto se mantenga dentro del presupuesto y se logren los objetivos establecidos. 

\input{Gestión económica/Recursos humanos.tex}

\subsection{Recursos hardware}
Los gastos de los recursos hardware hacen referencia al coste de todo los equipos físicos necesarios para el desarrollo del proyecto.
En este caso, se ha considerado el uso del mismo equipo par los diferentes roles desempeñados, que contiene el material 
necesario para que cada uno pueda realizar sus funciones. El equipo utilizado es el siguiente:

\begin{table}[H]
    \centering
    \begin{tabular}{|c|c|c|c|c|}
        \hline
        \textbf{Cantidad} & \textbf{Nombre} & \textbf{Coste}\\
        \hline
        3 & hp elitebook 640 14 inch g11 notebook pc & 859,00 € \\
        3 & hp 125 Wired Keyboard PERP & 27,10 € \\
        3 & hp Wired Mouse 1000 & 5,99 € \\
        \hline
    \end{tabular}
    \caption{Resumen del tiempo aproximado para cada tarea. Elaboración propia.}
    \label{tab:estimaciones}
\end{table}


\subsection{Recursos software}
La naturaleza de la organización como entidad sin ánimo de lucro impulsa una estrategia 
rigurosa de optimización de costes, priorizando el uso de soluciones de código abierto 
(\textit{Open Source}) siempre que sea posible. Este enfoque no solo minimiza el gasto 
financiero, sino que también fomenta la transparencia y la adaptabilidad técnica del proyecto.

Para la gestión de recursos de computación y planificación de tareas, se ha optado por 
Slurm, una solución de código abierto ampliamente reconocida en entornos de alto 
rendimiento. En cuanto a la infraestructura colaborativa, la empresa dispone de un acuerdo 
para el uso de Google Workspace, lo que permite gestionar la comunicación y la 
documentación bajo una filosofía de costes reducidos o nulos que se alinea con la estrategia 
\textit{Open Source}.

Para la infraestructura de datos, se utiliza Google Cloud en su plan Estándar, con 
almacenamiento y análisis en BigQuery On-Demand. Según la estimación de cargas del 
proyecto, se prevé un volumen aproximado de:
\begin{itemize}
    \item \textbf{Almacenamiento}: 100--200 GB de datos históricos generados por Slurm, lo que implica un coste mensual aproximado de 4 a 8€ (0,02 €/GB/mes).
    \item \textbf{Consultas}: unos 2 TB procesadas al mes, lo que corresponde a 10 € mensuales (5 €/TB procesado).
\end{itemize}

En total, se estima un coste global de unos \textbf{20–25 € mensuales}, es decir, alrededor de \textbf{80–100 €} para los 4 meses de duración del proyecto. Se ha imputado un presupuesto fijo de \textbf{100 €} para cubrir holguras en picos de almacenamiento o consultas.

Finalmente, la gestión de tareas se realiza a través de \textbf{Jira}, donde se aprovecha una reducción de precio substancial por el estado non-profit de la empresa.

\begin{table}[H]
    \begin{center}
    \begin{tabular}{|c|c|c|c|}
        \hline
        \textbf{Nombre del recurso} & \textbf{Coste} & \textbf{Vida útil} & \textbf{Coste Horario}\\
        \hline
        Google Cloud (BigQuery On-Demand, 4 meses) & 100,00 € & 539 h & 0,186 € \\
        Jira (3 usuarios, coste reducido) & 50,00 € & 539 h & 0,093 € \\
        Slurm (Open Source) & 0,00 € & 539 h & 0,000 € \\
        Google Workspace (Acuerdo non-profit) & 0,00 € & 539 h & 0,000 € \\
        \hline
        \multicolumn{3}{|l|}{\textbf{Total horario} (asumiendo los 3 roles)} & \textbf{0,279 €} \\
        \hline
        \multicolumn{3}{|l|}{\textbf{Total del proyecto} (asumiendo los 3 roles)} & \textbf{150,00 €} \\
        \hline
    \end{tabular}
    \caption{Coste horario de los recursos software indirectos. Elaboración propia.}
    \label{tab:recursos_software}
    \end{center}
\end{table}


\subsection{Contingencias}
Existe una gran posibilidad de que durante la ejecución del proyecto surjan imprevistos que 
requieran recursos adicionales no contemplados en el presupuesto inicial. Es por eso que 
para mitigar el riesgo que implican estos imprevistos, se ha decidido asignar un porcentaje 
del presupuesto total a contingencias. El porcentaje habitual de contingencia para un proyecto 
de desarrollo de software (o sector tecnológico en general) suele oscilar entre el 10\% y el 20\% 
del presupuesto total. Se ha optado por asignar un \textbf{15\%} del presupuesto total a contingencias, 
lo que nos porporciona los siguientes valores:


\begin{table}[H]
    \begin{center}
    \begin{tabular}{|c|c|c|c|}
        \hline
        \textbf{Recurso} & \textbf{Importe} & \textbf{Contingencias} \\ 
        \hline
        Recursos humanos & 12.176,58 € & 1.826,49 € \\
        Recursos materiales directos & 2.676,00 € & 401,40 € \\
        Recursos materiales indirectos  & 6.447,00 € & 967,05 € \\
        \hline
        \multicolumn{2}{|l|}{\textbf{Total}}  & \textbf{3.194,94 €} \\
        \hline
    \end{tabular}
    \caption{Tabla de contingencias. Elaboración propia.}
    \label{tab:calculo_contingencias}
    \end{center}
\end{table}

\section{Informe de sostenibilidad}

\input{Informe de sostenibilidad/Autoevaluación.tex}


\input{Informe de sostenibilidad/Dimensión ambiental.tex}

\subsection{Dimensión económica}

\textbf{¿Has estimado el coste de la realización del proyecto (recursos humanos
y materiales)?}

Sí, se ha realizado una estimación del coste del proyecto, una parte fundamental del proyecto.
Dicha estimación se encuentra en los apartados 6.x en la Gestión económica del proyecto.

\textbf{¿Cómo se resuelve actualmente el problema que quieres abordar (estado del arte)? 
¿En qué mejorará económicamente tu solución a las existentes?}

Siguiendo con la respuesta del apartado anterorior no existe una solución actual al problema a abordar.
Pero se espera que la solución propuesta mejore económicamente las cuentas de la institución. Al optimizar
los costes de los recursos y asignarlos de manera eficiente, se espera que la institución pueda reaprovechar
recursos y así poder comprar equipamiento nuevo que mejore la calidad del servicio ofrecido.

\input{Informe de sostenibilidad/Dimensión social.tex}


% \cleardoublepage

% \cleardoublepage


% \documentclass[a4paper, 11pt, twoside]{article}
\usepackage{wrapfig}
% \usepackage{abstract}
\usepackage{graphicx} % Required for inserting images
\usepackage[a4paper, left=1.15in, right=1.15in, top=1.5in, bottom=1.2in]{geometry}
\usepackage{float}
\usepackage{array}
\usepackage{booktabs}
\usepackage[T1]{fontenc}
\usepackage{amsmath}
\usepackage{tabularx}
\usepackage{graphicx} % Required for inserting images
\usepackage{tikz, xcolor}
\usepackage{hyperref}
\usepackage{cite}
\usepackage{pdfpages}
\usepackage{pdflscape}
\usepackage{afterpage}
\usepackage{textpos}
\usepackage{fancyhdr}
\usepackage{caption}
\usepackage{lmodern,textcomp}
\usepackage[catalan]{babel}
\usepackage{lipsum}
\usepackage{newfloat}
\usepackage{listings}
\usepackage{placeins}
\usepackage[ruled]{algorithm2e}
\usepackage{tcolorbox}
% Estilo de citas IEEE sencillo: usar paquete cite para comprimir y ordenar referencias numéricas
\usepackage{cite}
\usepackage[toc,page]{appendix}
\usepackage{forest}
\usepackage{pdflscape}
\usepackage{graphicx}
\usepackage{geometry}
\usepackage{quoting,xparse}
\usepackage{svg}



% \NewDocumentCommand{\bywhom}{m}{% the Bourbaki trick
%   {\nobreak\hfill\penalty50\hskip1em\null\nobreak
%    \hfill\mbox{\normalfont(#1)}%
%    \parfillskip=0pt \finalhyphendemerits=0 \par}%
% }

% AIXÒ SERVEIX PER A MODIFICAR UN PROBLEMA AMB ELS CARACTERS --> " + PARAULA 

% \NewDocumentEnvironment{pquotation}{m}
%   {\begin{quoting}[
%      indentfirst=true,
%      leftmargin=\parindent,
%      rightmargin=\parindent]\itshape}
%   {\bywhom{#1}
  
%   \end{quoting}}

% AIXÒ PER INSERIR FRAGMENTS DE CODI

\lstnewenvironment{code}[1][]%
{
   \noindent
   \minipage{\linewidth}
   \vspace{0.5\baselineskip}
   \lstset{
     basicstyle=\ttfamily\footnotesize,
     frame=single,
     breaklines=true,    % Permitir que las líneas se rompan automáticamente
     breakatwhitespace=true, % Permitir la ruptura en los espacios en blanco
     #1
   }
}
{\endminipage}

\lstset{
    language=Java,
    basicstyle=\ttfamily\small,  % Define el tipo de letra y tamaño del código
    keywordstyle=\color{blue},   % Define el color de las palabras clave de Java
    stringstyle=\color{red},     % Define el color de las cadenas de texto
    commentstyle=\color{green},  % Define el color de los comentarios
    morecomment=[s][\color{gray}]{/**}{*/},  % Estilo para comentarios de bloque
    breaklines=true,  % Configura el quiebre automático de línea
    numbers=left,  % Posición de los números de línea
    numberstyle=\tiny\color{gray},  % Estilo para los números de línea
    showspaces=false,  % No mostrar espacios con underscores
    showstringspaces=false,  % No mostrar espacios dentro de strings
    frame=single,  % Agregar un marco alrededor del código
    captionpos=b,  % b = bottom, t = top
}

\definecolor{eclipseStrings}{RGB}{42,0.0,255}
\definecolor{eclipseKeywords}{RGB}{127,0,85}
\colorlet{numb}{magenta!60!black}

\lstdefinelanguage{json}{
    basicstyle=\ttfamily\small,
    showstringspaces=false,
    breaklines=true,
    frame=single,
    literate=
     *{0}{{{\color{numb}0}}}{1}
      {1}{{{\color{numb}1}}}{1}
      {2}{{{\color{numb}2}}}{1}
      {3}{{{\color{numb}3}}}{1}
      {4}{{{\color{numb}4}}}{1}
      {5}{{{\color{numb}5}}}{1}
      {6}{{{\color{numb}6}}}{1}
      {7}{{{\color{numb}7}}}{1}
      {8}{{{\color{numb}8}}}{1}
      {9}{{{\color{numb}9}}}{1}
      {:}{{{\color{punct}{:}}}}{1}
      {,}{{{\color{punct}{,}}}}{1}
      {\{}{{{\color{delim}{\{}}}}{1}
      {\}}{{{\color{delim}{\}}}}}{1}
      {[}{{{\color{delim}{[}}}}{1}
      {]}{{{\color{delim}{]}}}}{1},
    morestring=[b]"
}

\lstdefinelanguage{JavaScript}{
  morekeywords=[1]{break, continue, delete, else, for, function, if, in,
    new, return, this, typeof, var, void, while, with},
  % Literals, primitive types, and reference types.
  morekeywords=[2]{false, null, true, boolean, number, undefined,
    Array, Boolean, Date, Math, Number, String, Object},
  % Built-ins.
  morekeywords=[3]{eval, parseInt, parseFloat, escape, unescape},
  sensitive,
  morecomment=[s]{/*}{*/},
  morecomment=[l]//,
  morecomment=[s]{/**}{*/}, % JavaDoc style comments
  morestring=[b]',
  morestring=[b]"
}[keywords, comments, strings]

\lstdefinelanguage{svelte}{
    morekeywords={async, await, then, if, each, in, on, let, const, function, import, export, from, with, typeof, new, var, return},
    sensitive=true,
    morecomment=[l]{//},
    morecomment=[s]{/*}{*/},
    morecomment=[s]{<!--}{-->},
    morestring=[b]',
    morestring=[b]",
    morestring=[b]`,
    ndkeywords={import, export, return, if, for, while, function, switch, case, break, var, let, const, async, await, then, catch, try, typeof, new, class, extends, implements, interface, constructor, public, private, protected, static, get, set},
    alsoletter={:},
    moredelim=**[is][\color{orange}]{\@\{}{\@\}},
    moredelim=**[is][\color{cyan}]{\$\{}{\}},
    keywordstyle=\color{blue}\bfseries,
    commentstyle=\color{gray}\ttfamily,
    stringstyle=\color{red}\ttfamily,
    identifierstyle=\color{black},
    breaklines=true,
    basicstyle=\ttfamily\small,
    columns=fullflexible,
    showstringspaces=false,
    tabsize=2,
    frame=single,
    captionpos=b,
}


\lstdefinestyle{java}{
    language=Java,
    basicstyle=\ttfamily\small,
    keywordstyle=\color{blue},
    stringstyle=\color{red},
    commentstyle=\color{gray},
    morecomment=[s][\color{gray}]{/**}{*/},
    breaklines=true,
    numbers=left,
    numberstyle=\tiny\color{gray},
    showspaces=false,
    showstringspaces=false,
    frame=single,
    captionpos=b
}

\usepackage{color}
\definecolor{numb}{rgb}{0.58,0,0.82}
\definecolor{punct}{rgb}{0.3,0.3,0.3}
\definecolor{delim}{rgb}{0.5,0.5,0.5}
\definecolor{string}{rgb}{0.6,0.1,0.1}

\lstdefinestyle{json}{
    language=json,
    basicstyle=\ttfamily\small,
    keywordstyle=\color{blue},
    stringstyle=\color{string},
    commentstyle=\color{green},
    morecomment=[s][\color{gray}]{/**}{*/},
    breaklines=true,
    numbers=left,
    numberstyle=\tiny\color{gray},
    showspaces=false,
    showstringspaces=false,
    frame=single,
    captionpos=b
}

\lstdefinelanguage{drl}{
    basicstyle=\ttfamily\small,
    keywordstyle=\color{blue}\bfseries,
    stringstyle=\color{red},
    commentstyle=\color{green},
    morecomment=[s][\color{green}]{/*}{*/},
    morecomment=[l][\color{green}]{//},
    morestring=[b]",
    keywords={rule, when, then, end, salience, eval, update, retract},
    sensitive=true,
    breaklines=true,
    numbers=left,
    numberstyle=\tiny\color{gray},
    showspaces=false,
    showstringspaces=false,
    frame=single,
    captionpos=b
}

\lstdefinestyle{drl}{
    language=drl,
    basicstyle=\ttfamily\small,
    keywordstyle=\color{blue}\bfseries,
    stringstyle=\color{red},
    commentstyle=\color{green},
    breaklines=true,
    numbers=left,
    numberstyle=\tiny\color{gray},
    showspaces=false,
    showstringspaces=false,
    frame=single,
    captionpos=b
}



\renewcommand{\lstlistingname}{Codi}
\renewcommand{\lstlistlistingname}{Índex de Fragments de Codis}
\renewcommand{\appendixpagename}{Annexos}
\renewcommand{\appendixtocname}{Annexos}

\renewcommand*{\listalgorithmcfname}{Llistat de pseudocodi}
\renewcommand*{\algorithmcfname}{Pseudocodi}
\renewcommand*{\algorithmautorefname}{pseudocodi}

\definecolor{darkgreen}{rgb}{0.0, 0.3, 0.0}
\definecolor{lightgreen}{rgb}{0.9, 0.95, 0.9}

\newcommand{\todo}[1]{
    \begin{tcolorbox}[colback=lightgreen, colframe=darkgreen, title=TODO]
        #1
    \end{tcolorbox}
}


% \usepackage{indentfirst} 
% \setlength{\parindent}{1cm}


\usepackage{parskip}
% \setlength{\parindent}{10pt}

% \setlength{\headheight}{15.6pt}
\setlength{\headheight}{16.05263pt}




\vspace{5mm}

\newcolumntype{u}{>{\centering\arraybackslash\hsize=.22\hsize}X}
\newcolumntype{s}{>{\hsize=.45\hsize}X}


%%%%%%%%%%%%%%%%%%%%%%%%%%%%%%%%
% INI creacio subsubsubsecion
%%%%%%%%%%%%%%%%%%%%%%%%%%%%%%%%
\usepackage{titlesec}

\let\oldparagraph\paragraph
\renewcommand{\paragraph}[1]{\oldparagraph{\hspace{30pt}#1}}


\titleclass{\subsubsubsection}{straight}[\subsection]

\newcounter{subsubsubsection}[subsubsection]
\renewcommand\thesubsubsubsection{\thesubsubsection.\arabic{subsubsubsection}}
\renewcommand\theparagraph{\thesubsubsubsection.\arabic{paragraph}} % optional; useful if paragraphs are to be numbered

\titleformat{\subsubsubsection}
  {\normalfont\normalsize\bfseries}{\thesubsubsubsection}{1em}{}
\titlespacing*{\subsubsubsection}
{0pt}{3.25ex plus 1ex minus .2ex}{1.5ex plus .2ex}


\makeatletter
\renewcommand\paragraph{\@startsection{paragraph}{5}{10pt}%
  {3.25ex \@plus1ex \@minus.2ex}%
  {-1em}%
  {\normalfont\normalsize\bfseries}}
\renewcommand\subparagraph{\@startsection{subparagraph}{6}{\parindent}%
  {3.25ex \@plus1ex \@minus .2ex}%
  {-1em}%
  {\normalfont\normalsize\bfseries}}
\def\toclevel@subsubsubsection{4}
\def\toclevel@paragraph{5}
\def\toclevel@paragraph{6}
\def\l@subsubsubsection{\@dottedtocline{4}{7em}{4em}}
\def\l@paragraph{\@dottedtocline{5}{10em}{5em}}
\def\l@subparagraph{\@dottedtocline{6}{14em}{6em}}
\makeatother

\setcounter{secnumdepth}{4}
\setcounter{tocdepth}{3}
%%%%%%%%%%%%%%%%%%%%%%%%%%%%%%%%
% FIN creacio subsubsubtitle
%%%%%%%%%%%%%%%%%%%%%%%%%%%%%%%%


%%%%%%%%%%%%%%%%%%%%%%%%%%%%%%%%
% INI creacio equation
%%%%%%%%%%%%%%%%%%%%%%%%%%%%%%%%
\DeclareFloatingEnvironment[fileext=loeq,placement={!ht},name=Equació]{equationfloat}
%%%%%%%%%%%%%%%%%%%%%%%%%%%%%%%%
% FIN creacio equation
%%%%%%%%%%%%%%%%%%%%%%%%%%%%%%%%

\titleformat{\section}[block]
  {\clearpage\normalfont\Large\bfseries}
  {\thesection}{1em}{}


\renewcommand\refname{Bibliografia}


\linespread{1.25}


\makeatletter
\renewenvironment{abstract}{%
    \if@twocolumn
      \section*{\abstractname}%
    \else 
      \begin{center}%
        {\bfseries \Large\itshape\abstractname\vspace{\z@}}%  % Título en negrita, 14pt y cursiva
      \end{center}%
      \quotation
    \fi}
    {\if@twocolumn\else\endquotation\fi}
\makeatother

% \fancypagestyle{main}{
%     \fancyhf{}
%     \lhead{\small Aplicació Web per a la Visualització i Comparació
% de Trajectòries d’Algorismes d’Optimització}
%     % \rhead{\includegraphics[width=1.2cm]{resources/Icons/logo-fib.png}}
%     \cfoot{\thepage}
%     \setlength{\footskip}{35pt}
% }

\newcommand{\q}[1]{``#1''}

\begin{document}

\hypersetup{
    pdfborder=0 0 0,
}

%% PORTADA I TABLE OF CONTENTS %%

%%%%                                   %%%%
%         INDEX DE LA MEMORIA             %
%%%%                                   %%%%

% \includepdf[pages=-]{portada_fib.pdf}
%%% Portada
\begin{titlepage}
\begin{center}
    \includegraphics[width=1\linewidth]{extra/figures/logo-upc-fib.png}
    \vspace{30mm}

    \LARGE{\textsc{\textbf{{INTEGRACIÓN Y AUTOMATIZACIÓN DE COSTES DE INFRAESTRUCTURA HPC EN ENTORNOS DE INVESTIGACIÓN}}}} \\
    \vspace{25mm}
      
        
    \textbf{\Large{\textsc{PAU NAVARRO ÁLVAREZ}}} \\
    \vspace{20mm}
        
    
    
    \large{\textbf{Director/a} ROBERTO RIVEIRO INSUA (FUNDACIO INSTITUT DE RECERCA BIOMEDICA (IRB BARCELONA))} \\
    \large{\textbf{Ponente:} ALEX PAJUELO GONZÁLEZ (Departamento de Arquitectura de Computadores)} \\
    \large{\textbf{Titulación:} Grado en Ingeniería Informática (Tecnologías de la información)} \\
    
    \large{\textbf{Lliurament 2:} }Planificación temporal \\
    \large{\textbf{Tutor GEP:} CAROLINA CONSOLIDACIÓN SEGURA (Departament d'Organització d'Empreses)} \\
    \large{\textbf{Fecha:}} 29 de Septiembre de 2025\\

    \textbf{\large{Memoria del trabajo de fin de grado}} \\
    \vspace{3mm}
    \textbf{\large{Facultat d'Informàtica de Barcelona (FIB)}} \\
    \vspace{3mm}
    \textbf{\large{Universitat Politècnica de Catalunya (UPC) - BarcelonaTech}} \\

\end{center}
\end{titlepage}
%\input{Extra/portada}
% \cleardoublepage

% \pagestyle{main}
% \setcounter{page}{11}
\renewcommand{\contentsname}{Continguts}

\tableofcontents
\clearpage
\listoffigures
\begingroup
    \let\clearpage\relax
    \listoftables
\endgroup
\clearpage

\newpage
% \cleardoublepage
\input{Contexto y alcance/Contextualización/main}
\input{Contexto y alcance/Justificación/main}
\input{Contexto y alcance/Alcance/main}
\input{Contexto y alcance/Metodología y herramientas de seguimiento/main}

\documentclass[a4paper, 11pt, twoside]{article}
\usepackage{wrapfig}
% \usepackage{abstract}
\usepackage{graphicx} % Required for inserting images
\usepackage[a4paper, left=1.15in, right=1.15in, top=1.5in, bottom=1.2in]{geometry}
\usepackage{float}
\usepackage{array}
\usepackage{booktabs}
\usepackage[T1]{fontenc}
\usepackage{amsmath}
\usepackage{tabularx}
\usepackage{graphicx} % Required for inserting images
\usepackage{tikz, xcolor}
\usepackage{hyperref}
\usepackage{cite}
\usepackage{pdfpages}
\usepackage{pdflscape}
\usepackage{afterpage}
\usepackage{textpos}
\usepackage{fancyhdr}
\usepackage{caption}
\usepackage{lmodern,textcomp}
\usepackage[catalan]{babel}
\usepackage{lipsum}
\usepackage{newfloat}
\usepackage{listings}
\usepackage{placeins}
\usepackage[ruled]{algorithm2e}
\usepackage{tcolorbox}
% Estilo de citas IEEE sencillo: usar paquete cite para comprimir y ordenar referencias numéricas
\usepackage{cite}
\usepackage[toc,page]{appendix}
\usepackage{forest}
\usepackage{pdflscape}
\usepackage{graphicx}
\usepackage{geometry}
\usepackage{quoting,xparse}
\usepackage{svg}



% \NewDocumentCommand{\bywhom}{m}{% the Bourbaki trick
%   {\nobreak\hfill\penalty50\hskip1em\null\nobreak
%    \hfill\mbox{\normalfont(#1)}%
%    \parfillskip=0pt \finalhyphendemerits=0 \par}%
% }

% AIXÒ SERVEIX PER A MODIFICAR UN PROBLEMA AMB ELS CARACTERS --> " + PARAULA 

% \NewDocumentEnvironment{pquotation}{m}
%   {\begin{quoting}[
%      indentfirst=true,
%      leftmargin=\parindent,
%      rightmargin=\parindent]\itshape}
%   {\bywhom{#1}
  
%   \end{quoting}}

% AIXÒ PER INSERIR FRAGMENTS DE CODI

\lstnewenvironment{code}[1][]%
{
   \noindent
   \minipage{\linewidth}
   \vspace{0.5\baselineskip}
   \lstset{
     basicstyle=\ttfamily\footnotesize,
     frame=single,
     breaklines=true,    % Permitir que las líneas se rompan automáticamente
     breakatwhitespace=true, % Permitir la ruptura en los espacios en blanco
     #1
   }
}
{\endminipage}

\lstset{
    language=Java,
    basicstyle=\ttfamily\small,  % Define el tipo de letra y tamaño del código
    keywordstyle=\color{blue},   % Define el color de las palabras clave de Java
    stringstyle=\color{red},     % Define el color de las cadenas de texto
    commentstyle=\color{green},  % Define el color de los comentarios
    morecomment=[s][\color{gray}]{/**}{*/},  % Estilo para comentarios de bloque
    breaklines=true,  % Configura el quiebre automático de línea
    numbers=left,  % Posición de los números de línea
    numberstyle=\tiny\color{gray},  % Estilo para los números de línea
    showspaces=false,  % No mostrar espacios con underscores
    showstringspaces=false,  % No mostrar espacios dentro de strings
    frame=single,  % Agregar un marco alrededor del código
    captionpos=b,  % b = bottom, t = top
}

\definecolor{eclipseStrings}{RGB}{42,0.0,255}
\definecolor{eclipseKeywords}{RGB}{127,0,85}
\colorlet{numb}{magenta!60!black}

\lstdefinelanguage{json}{
    basicstyle=\ttfamily\small,
    showstringspaces=false,
    breaklines=true,
    frame=single,
    literate=
     *{0}{{{\color{numb}0}}}{1}
      {1}{{{\color{numb}1}}}{1}
      {2}{{{\color{numb}2}}}{1}
      {3}{{{\color{numb}3}}}{1}
      {4}{{{\color{numb}4}}}{1}
      {5}{{{\color{numb}5}}}{1}
      {6}{{{\color{numb}6}}}{1}
      {7}{{{\color{numb}7}}}{1}
      {8}{{{\color{numb}8}}}{1}
      {9}{{{\color{numb}9}}}{1}
      {:}{{{\color{punct}{:}}}}{1}
      {,}{{{\color{punct}{,}}}}{1}
      {\{}{{{\color{delim}{\{}}}}{1}
      {\}}{{{\color{delim}{\}}}}}{1}
      {[}{{{\color{delim}{[}}}}{1}
      {]}{{{\color{delim}{]}}}}{1},
    morestring=[b]"
}

\lstdefinelanguage{JavaScript}{
  morekeywords=[1]{break, continue, delete, else, for, function, if, in,
    new, return, this, typeof, var, void, while, with},
  % Literals, primitive types, and reference types.
  morekeywords=[2]{false, null, true, boolean, number, undefined,
    Array, Boolean, Date, Math, Number, String, Object},
  % Built-ins.
  morekeywords=[3]{eval, parseInt, parseFloat, escape, unescape},
  sensitive,
  morecomment=[s]{/*}{*/},
  morecomment=[l]//,
  morecomment=[s]{/**}{*/}, % JavaDoc style comments
  morestring=[b]',
  morestring=[b]"
}[keywords, comments, strings]

\lstdefinelanguage{svelte}{
    morekeywords={async, await, then, if, each, in, on, let, const, function, import, export, from, with, typeof, new, var, return},
    sensitive=true,
    morecomment=[l]{//},
    morecomment=[s]{/*}{*/},
    morecomment=[s]{<!--}{-->},
    morestring=[b]',
    morestring=[b]",
    morestring=[b]`,
    ndkeywords={import, export, return, if, for, while, function, switch, case, break, var, let, const, async, await, then, catch, try, typeof, new, class, extends, implements, interface, constructor, public, private, protected, static, get, set},
    alsoletter={:},
    moredelim=**[is][\color{orange}]{\@\{}{\@\}},
    moredelim=**[is][\color{cyan}]{\$\{}{\}},
    keywordstyle=\color{blue}\bfseries,
    commentstyle=\color{gray}\ttfamily,
    stringstyle=\color{red}\ttfamily,
    identifierstyle=\color{black},
    breaklines=true,
    basicstyle=\ttfamily\small,
    columns=fullflexible,
    showstringspaces=false,
    tabsize=2,
    frame=single,
    captionpos=b,
}


\lstdefinestyle{java}{
    language=Java,
    basicstyle=\ttfamily\small,
    keywordstyle=\color{blue},
    stringstyle=\color{red},
    commentstyle=\color{gray},
    morecomment=[s][\color{gray}]{/**}{*/},
    breaklines=true,
    numbers=left,
    numberstyle=\tiny\color{gray},
    showspaces=false,
    showstringspaces=false,
    frame=single,
    captionpos=b
}

\usepackage{color}
\definecolor{numb}{rgb}{0.58,0,0.82}
\definecolor{punct}{rgb}{0.3,0.3,0.3}
\definecolor{delim}{rgb}{0.5,0.5,0.5}
\definecolor{string}{rgb}{0.6,0.1,0.1}

\lstdefinestyle{json}{
    language=json,
    basicstyle=\ttfamily\small,
    keywordstyle=\color{blue},
    stringstyle=\color{string},
    commentstyle=\color{green},
    morecomment=[s][\color{gray}]{/**}{*/},
    breaklines=true,
    numbers=left,
    numberstyle=\tiny\color{gray},
    showspaces=false,
    showstringspaces=false,
    frame=single,
    captionpos=b
}

\lstdefinelanguage{drl}{
    basicstyle=\ttfamily\small,
    keywordstyle=\color{blue}\bfseries,
    stringstyle=\color{red},
    commentstyle=\color{green},
    morecomment=[s][\color{green}]{/*}{*/},
    morecomment=[l][\color{green}]{//},
    morestring=[b]",
    keywords={rule, when, then, end, salience, eval, update, retract},
    sensitive=true,
    breaklines=true,
    numbers=left,
    numberstyle=\tiny\color{gray},
    showspaces=false,
    showstringspaces=false,
    frame=single,
    captionpos=b
}

\lstdefinestyle{drl}{
    language=drl,
    basicstyle=\ttfamily\small,
    keywordstyle=\color{blue}\bfseries,
    stringstyle=\color{red},
    commentstyle=\color{green},
    breaklines=true,
    numbers=left,
    numberstyle=\tiny\color{gray},
    showspaces=false,
    showstringspaces=false,
    frame=single,
    captionpos=b
}



\renewcommand{\lstlistingname}{Codi}
\renewcommand{\lstlistlistingname}{Índex de Fragments de Codis}
\renewcommand{\appendixpagename}{Annexos}
\renewcommand{\appendixtocname}{Annexos}

\renewcommand*{\listalgorithmcfname}{Llistat de pseudocodi}
\renewcommand*{\algorithmcfname}{Pseudocodi}
\renewcommand*{\algorithmautorefname}{pseudocodi}

\definecolor{darkgreen}{rgb}{0.0, 0.3, 0.0}
\definecolor{lightgreen}{rgb}{0.9, 0.95, 0.9}

\newcommand{\todo}[1]{
    \begin{tcolorbox}[colback=lightgreen, colframe=darkgreen, title=TODO]
        #1
    \end{tcolorbox}
}


% \usepackage{indentfirst} 
% \setlength{\parindent}{1cm}


\usepackage{parskip}
% \setlength{\parindent}{10pt}

% \setlength{\headheight}{15.6pt}
\setlength{\headheight}{16.05263pt}




\vspace{5mm}

\newcolumntype{u}{>{\centering\arraybackslash\hsize=.22\hsize}X}
\newcolumntype{s}{>{\hsize=.45\hsize}X}


%%%%%%%%%%%%%%%%%%%%%%%%%%%%%%%%
% INI creacio subsubsubsecion
%%%%%%%%%%%%%%%%%%%%%%%%%%%%%%%%
\usepackage{titlesec}

\let\oldparagraph\paragraph
\renewcommand{\paragraph}[1]{\oldparagraph{\hspace{30pt}#1}}


\titleclass{\subsubsubsection}{straight}[\subsection]

\newcounter{subsubsubsection}[subsubsection]
\renewcommand\thesubsubsubsection{\thesubsubsection.\arabic{subsubsubsection}}
\renewcommand\theparagraph{\thesubsubsubsection.\arabic{paragraph}} % optional; useful if paragraphs are to be numbered

\titleformat{\subsubsubsection}
  {\normalfont\normalsize\bfseries}{\thesubsubsubsection}{1em}{}
\titlespacing*{\subsubsubsection}
{0pt}{3.25ex plus 1ex minus .2ex}{1.5ex plus .2ex}


\makeatletter
\renewcommand\paragraph{\@startsection{paragraph}{5}{10pt}%
  {3.25ex \@plus1ex \@minus.2ex}%
  {-1em}%
  {\normalfont\normalsize\bfseries}}
\renewcommand\subparagraph{\@startsection{subparagraph}{6}{\parindent}%
  {3.25ex \@plus1ex \@minus .2ex}%
  {-1em}%
  {\normalfont\normalsize\bfseries}}
\def\toclevel@subsubsubsection{4}
\def\toclevel@paragraph{5}
\def\toclevel@paragraph{6}
\def\l@subsubsubsection{\@dottedtocline{4}{7em}{4em}}
\def\l@paragraph{\@dottedtocline{5}{10em}{5em}}
\def\l@subparagraph{\@dottedtocline{6}{14em}{6em}}
\makeatother

\setcounter{secnumdepth}{4}
\setcounter{tocdepth}{3}
%%%%%%%%%%%%%%%%%%%%%%%%%%%%%%%%
% FIN creacio subsubsubtitle
%%%%%%%%%%%%%%%%%%%%%%%%%%%%%%%%


%%%%%%%%%%%%%%%%%%%%%%%%%%%%%%%%
% INI creacio equation
%%%%%%%%%%%%%%%%%%%%%%%%%%%%%%%%
\DeclareFloatingEnvironment[fileext=loeq,placement={!ht},name=Equació]{equationfloat}
%%%%%%%%%%%%%%%%%%%%%%%%%%%%%%%%
% FIN creacio equation
%%%%%%%%%%%%%%%%%%%%%%%%%%%%%%%%

\titleformat{\section}[block]
  {\clearpage\normalfont\Large\bfseries}
  {\thesection}{1em}{}


\renewcommand\refname{Bibliografia}


\linespread{1.25}


\makeatletter
\renewenvironment{abstract}{%
    \if@twocolumn
      \section*{\abstractname}%
    \else 
      \begin{center}%
        {\bfseries \Large\itshape\abstractname\vspace{\z@}}%  % Título en negrita, 14pt y cursiva
      \end{center}%
      \quotation
    \fi}
    {\if@twocolumn\else\endquotation\fi}
\makeatother

% \fancypagestyle{main}{
%     \fancyhf{}
%     \lhead{\small Aplicació Web per a la Visualització i Comparació
% de Trajectòries d’Algorismes d’Optimització}
%     % \rhead{\includegraphics[width=1.2cm]{resources/Icons/logo-fib.png}}
%     \cfoot{\thepage}
%     \setlength{\footskip}{35pt}
% }

\newcommand{\q}[1]{``#1''}

\begin{document}

\hypersetup{
    pdfborder=0 0 0,
}

%% PORTADA I TABLE OF CONTENTS %%

%%%%                                   %%%%
%         INDEX DE LA MEMORIA             %
%%%%                                   %%%%

% \includepdf[pages=-]{portada_fib.pdf}
\input{Extra/GEP}
%\input{Extra/portada}
% \cleardoublepage

% \pagestyle{main}
% \setcounter{page}{11}
\include{toc}
% \cleardoublepage
\input{Contexto y alcance/main}

\input{Planificación temporal/main}

\input{Gestión económica/main}

\input{Informe de sostenibilidad/main}

% \cleardoublepage

% \cleardoublepage


% \input{Annex/main}


\clearpage
% Estilo bibliográfico IEEE
\bibliographystyle{IEEEtran}
\renewcommand{\refname}{Referencias}
\bibliography{references.bib}

%%%%                                   %%%%
%         FINAL DE LA MEMORIA             %
%%%%                                   %%%%

\end{document}


%% VERSIÓ INCIAL MEMORIA %%



\section{Gestión económica}

La gestión económica de un proyecto es fundamental para asegurar su viabilidad y exito a largo plazo. Ya hemos visto en secciones anteriores, los otros dos vertices 
del triángulo de la gestión de proyectos: el alcance y el tiempo. En esta sección, nos centraremos en cómo gestionar eficazmente los recursos financieros para garantizar
que el proyecto se mantenga dentro del presupuesto y se logren los objetivos establecidos. 

\input{Gestión económica/Recursos humanos.tex}

\input{Gestión económica/Recursos Hardware.tex}

\input{Gestión económica/Recursos Software.tex}

\input{Gestión económica/Contingencias.tex}

\section{Informe de sostenibilidad}

\input{Informe de sostenibilidad/Autoevaluación.tex}


\input{Informe de sostenibilidad/Dimensión ambiental.tex}

\input{Informe de sostenibilidad/Dimensión económica.tex}

\input{Informe de sostenibilidad/Dimensión social.tex}


% \cleardoublepage

% \cleardoublepage


% \documentclass[a4paper, 11pt, twoside]{article}
\usepackage{wrapfig}
% \usepackage{abstract}
\usepackage{graphicx} % Required for inserting images
\usepackage[a4paper, left=1.15in, right=1.15in, top=1.5in, bottom=1.2in]{geometry}
\usepackage{float}
\usepackage{array}
\usepackage{booktabs}
\usepackage[T1]{fontenc}
\usepackage{amsmath}
\usepackage{tabularx}
\usepackage{graphicx} % Required for inserting images
\usepackage{tikz, xcolor}
\usepackage{hyperref}
\usepackage{cite}
\usepackage{pdfpages}
\usepackage{pdflscape}
\usepackage{afterpage}
\usepackage{textpos}
\usepackage{fancyhdr}
\usepackage{caption}
\usepackage{lmodern,textcomp}
\usepackage[catalan]{babel}
\usepackage{lipsum}
\usepackage{newfloat}
\usepackage{listings}
\usepackage{placeins}
\usepackage[ruled]{algorithm2e}
\usepackage{tcolorbox}
% Estilo de citas IEEE sencillo: usar paquete cite para comprimir y ordenar referencias numéricas
\usepackage{cite}
\usepackage[toc,page]{appendix}
\usepackage{forest}
\usepackage{pdflscape}
\usepackage{graphicx}
\usepackage{geometry}
\usepackage{quoting,xparse}
\usepackage{svg}



% \NewDocumentCommand{\bywhom}{m}{% the Bourbaki trick
%   {\nobreak\hfill\penalty50\hskip1em\null\nobreak
%    \hfill\mbox{\normalfont(#1)}%
%    \parfillskip=0pt \finalhyphendemerits=0 \par}%
% }

% AIXÒ SERVEIX PER A MODIFICAR UN PROBLEMA AMB ELS CARACTERS --> " + PARAULA 

% \NewDocumentEnvironment{pquotation}{m}
%   {\begin{quoting}[
%      indentfirst=true,
%      leftmargin=\parindent,
%      rightmargin=\parindent]\itshape}
%   {\bywhom{#1}
  
%   \end{quoting}}

% AIXÒ PER INSERIR FRAGMENTS DE CODI

\lstnewenvironment{code}[1][]%
{
   \noindent
   \minipage{\linewidth}
   \vspace{0.5\baselineskip}
   \lstset{
     basicstyle=\ttfamily\footnotesize,
     frame=single,
     breaklines=true,    % Permitir que las líneas se rompan automáticamente
     breakatwhitespace=true, % Permitir la ruptura en los espacios en blanco
     #1
   }
}
{\endminipage}

\lstset{
    language=Java,
    basicstyle=\ttfamily\small,  % Define el tipo de letra y tamaño del código
    keywordstyle=\color{blue},   % Define el color de las palabras clave de Java
    stringstyle=\color{red},     % Define el color de las cadenas de texto
    commentstyle=\color{green},  % Define el color de los comentarios
    morecomment=[s][\color{gray}]{/**}{*/},  % Estilo para comentarios de bloque
    breaklines=true,  % Configura el quiebre automático de línea
    numbers=left,  % Posición de los números de línea
    numberstyle=\tiny\color{gray},  % Estilo para los números de línea
    showspaces=false,  % No mostrar espacios con underscores
    showstringspaces=false,  % No mostrar espacios dentro de strings
    frame=single,  % Agregar un marco alrededor del código
    captionpos=b,  % b = bottom, t = top
}

\definecolor{eclipseStrings}{RGB}{42,0.0,255}
\definecolor{eclipseKeywords}{RGB}{127,0,85}
\colorlet{numb}{magenta!60!black}

\lstdefinelanguage{json}{
    basicstyle=\ttfamily\small,
    showstringspaces=false,
    breaklines=true,
    frame=single,
    literate=
     *{0}{{{\color{numb}0}}}{1}
      {1}{{{\color{numb}1}}}{1}
      {2}{{{\color{numb}2}}}{1}
      {3}{{{\color{numb}3}}}{1}
      {4}{{{\color{numb}4}}}{1}
      {5}{{{\color{numb}5}}}{1}
      {6}{{{\color{numb}6}}}{1}
      {7}{{{\color{numb}7}}}{1}
      {8}{{{\color{numb}8}}}{1}
      {9}{{{\color{numb}9}}}{1}
      {:}{{{\color{punct}{:}}}}{1}
      {,}{{{\color{punct}{,}}}}{1}
      {\{}{{{\color{delim}{\{}}}}{1}
      {\}}{{{\color{delim}{\}}}}}{1}
      {[}{{{\color{delim}{[}}}}{1}
      {]}{{{\color{delim}{]}}}}{1},
    morestring=[b]"
}

\lstdefinelanguage{JavaScript}{
  morekeywords=[1]{break, continue, delete, else, for, function, if, in,
    new, return, this, typeof, var, void, while, with},
  % Literals, primitive types, and reference types.
  morekeywords=[2]{false, null, true, boolean, number, undefined,
    Array, Boolean, Date, Math, Number, String, Object},
  % Built-ins.
  morekeywords=[3]{eval, parseInt, parseFloat, escape, unescape},
  sensitive,
  morecomment=[s]{/*}{*/},
  morecomment=[l]//,
  morecomment=[s]{/**}{*/}, % JavaDoc style comments
  morestring=[b]',
  morestring=[b]"
}[keywords, comments, strings]

\lstdefinelanguage{svelte}{
    morekeywords={async, await, then, if, each, in, on, let, const, function, import, export, from, with, typeof, new, var, return},
    sensitive=true,
    morecomment=[l]{//},
    morecomment=[s]{/*}{*/},
    morecomment=[s]{<!--}{-->},
    morestring=[b]',
    morestring=[b]",
    morestring=[b]`,
    ndkeywords={import, export, return, if, for, while, function, switch, case, break, var, let, const, async, await, then, catch, try, typeof, new, class, extends, implements, interface, constructor, public, private, protected, static, get, set},
    alsoletter={:},
    moredelim=**[is][\color{orange}]{\@\{}{\@\}},
    moredelim=**[is][\color{cyan}]{\$\{}{\}},
    keywordstyle=\color{blue}\bfseries,
    commentstyle=\color{gray}\ttfamily,
    stringstyle=\color{red}\ttfamily,
    identifierstyle=\color{black},
    breaklines=true,
    basicstyle=\ttfamily\small,
    columns=fullflexible,
    showstringspaces=false,
    tabsize=2,
    frame=single,
    captionpos=b,
}


\lstdefinestyle{java}{
    language=Java,
    basicstyle=\ttfamily\small,
    keywordstyle=\color{blue},
    stringstyle=\color{red},
    commentstyle=\color{gray},
    morecomment=[s][\color{gray}]{/**}{*/},
    breaklines=true,
    numbers=left,
    numberstyle=\tiny\color{gray},
    showspaces=false,
    showstringspaces=false,
    frame=single,
    captionpos=b
}

\usepackage{color}
\definecolor{numb}{rgb}{0.58,0,0.82}
\definecolor{punct}{rgb}{0.3,0.3,0.3}
\definecolor{delim}{rgb}{0.5,0.5,0.5}
\definecolor{string}{rgb}{0.6,0.1,0.1}

\lstdefinestyle{json}{
    language=json,
    basicstyle=\ttfamily\small,
    keywordstyle=\color{blue},
    stringstyle=\color{string},
    commentstyle=\color{green},
    morecomment=[s][\color{gray}]{/**}{*/},
    breaklines=true,
    numbers=left,
    numberstyle=\tiny\color{gray},
    showspaces=false,
    showstringspaces=false,
    frame=single,
    captionpos=b
}

\lstdefinelanguage{drl}{
    basicstyle=\ttfamily\small,
    keywordstyle=\color{blue}\bfseries,
    stringstyle=\color{red},
    commentstyle=\color{green},
    morecomment=[s][\color{green}]{/*}{*/},
    morecomment=[l][\color{green}]{//},
    morestring=[b]",
    keywords={rule, when, then, end, salience, eval, update, retract},
    sensitive=true,
    breaklines=true,
    numbers=left,
    numberstyle=\tiny\color{gray},
    showspaces=false,
    showstringspaces=false,
    frame=single,
    captionpos=b
}

\lstdefinestyle{drl}{
    language=drl,
    basicstyle=\ttfamily\small,
    keywordstyle=\color{blue}\bfseries,
    stringstyle=\color{red},
    commentstyle=\color{green},
    breaklines=true,
    numbers=left,
    numberstyle=\tiny\color{gray},
    showspaces=false,
    showstringspaces=false,
    frame=single,
    captionpos=b
}



\renewcommand{\lstlistingname}{Codi}
\renewcommand{\lstlistlistingname}{Índex de Fragments de Codis}
\renewcommand{\appendixpagename}{Annexos}
\renewcommand{\appendixtocname}{Annexos}

\renewcommand*{\listalgorithmcfname}{Llistat de pseudocodi}
\renewcommand*{\algorithmcfname}{Pseudocodi}
\renewcommand*{\algorithmautorefname}{pseudocodi}

\definecolor{darkgreen}{rgb}{0.0, 0.3, 0.0}
\definecolor{lightgreen}{rgb}{0.9, 0.95, 0.9}

\newcommand{\todo}[1]{
    \begin{tcolorbox}[colback=lightgreen, colframe=darkgreen, title=TODO]
        #1
    \end{tcolorbox}
}


% \usepackage{indentfirst} 
% \setlength{\parindent}{1cm}


\usepackage{parskip}
% \setlength{\parindent}{10pt}

% \setlength{\headheight}{15.6pt}
\setlength{\headheight}{16.05263pt}




\vspace{5mm}

\newcolumntype{u}{>{\centering\arraybackslash\hsize=.22\hsize}X}
\newcolumntype{s}{>{\hsize=.45\hsize}X}


%%%%%%%%%%%%%%%%%%%%%%%%%%%%%%%%
% INI creacio subsubsubsecion
%%%%%%%%%%%%%%%%%%%%%%%%%%%%%%%%
\usepackage{titlesec}

\let\oldparagraph\paragraph
\renewcommand{\paragraph}[1]{\oldparagraph{\hspace{30pt}#1}}


\titleclass{\subsubsubsection}{straight}[\subsection]

\newcounter{subsubsubsection}[subsubsection]
\renewcommand\thesubsubsubsection{\thesubsubsection.\arabic{subsubsubsection}}
\renewcommand\theparagraph{\thesubsubsubsection.\arabic{paragraph}} % optional; useful if paragraphs are to be numbered

\titleformat{\subsubsubsection}
  {\normalfont\normalsize\bfseries}{\thesubsubsubsection}{1em}{}
\titlespacing*{\subsubsubsection}
{0pt}{3.25ex plus 1ex minus .2ex}{1.5ex plus .2ex}


\makeatletter
\renewcommand\paragraph{\@startsection{paragraph}{5}{10pt}%
  {3.25ex \@plus1ex \@minus.2ex}%
  {-1em}%
  {\normalfont\normalsize\bfseries}}
\renewcommand\subparagraph{\@startsection{subparagraph}{6}{\parindent}%
  {3.25ex \@plus1ex \@minus .2ex}%
  {-1em}%
  {\normalfont\normalsize\bfseries}}
\def\toclevel@subsubsubsection{4}
\def\toclevel@paragraph{5}
\def\toclevel@paragraph{6}
\def\l@subsubsubsection{\@dottedtocline{4}{7em}{4em}}
\def\l@paragraph{\@dottedtocline{5}{10em}{5em}}
\def\l@subparagraph{\@dottedtocline{6}{14em}{6em}}
\makeatother

\setcounter{secnumdepth}{4}
\setcounter{tocdepth}{3}
%%%%%%%%%%%%%%%%%%%%%%%%%%%%%%%%
% FIN creacio subsubsubtitle
%%%%%%%%%%%%%%%%%%%%%%%%%%%%%%%%


%%%%%%%%%%%%%%%%%%%%%%%%%%%%%%%%
% INI creacio equation
%%%%%%%%%%%%%%%%%%%%%%%%%%%%%%%%
\DeclareFloatingEnvironment[fileext=loeq,placement={!ht},name=Equació]{equationfloat}
%%%%%%%%%%%%%%%%%%%%%%%%%%%%%%%%
% FIN creacio equation
%%%%%%%%%%%%%%%%%%%%%%%%%%%%%%%%

\titleformat{\section}[block]
  {\clearpage\normalfont\Large\bfseries}
  {\thesection}{1em}{}


\renewcommand\refname{Bibliografia}


\linespread{1.25}


\makeatletter
\renewenvironment{abstract}{%
    \if@twocolumn
      \section*{\abstractname}%
    \else 
      \begin{center}%
        {\bfseries \Large\itshape\abstractname\vspace{\z@}}%  % Título en negrita, 14pt y cursiva
      \end{center}%
      \quotation
    \fi}
    {\if@twocolumn\else\endquotation\fi}
\makeatother

% \fancypagestyle{main}{
%     \fancyhf{}
%     \lhead{\small Aplicació Web per a la Visualització i Comparació
% de Trajectòries d’Algorismes d’Optimització}
%     % \rhead{\includegraphics[width=1.2cm]{resources/Icons/logo-fib.png}}
%     \cfoot{\thepage}
%     \setlength{\footskip}{35pt}
% }

\newcommand{\q}[1]{``#1''}

\begin{document}

\hypersetup{
    pdfborder=0 0 0,
}

%% PORTADA I TABLE OF CONTENTS %%

%%%%                                   %%%%
%         INDEX DE LA MEMORIA             %
%%%%                                   %%%%

% \includepdf[pages=-]{portada_fib.pdf}
\input{Extra/GEP}
%\input{Extra/portada}
% \cleardoublepage

% \pagestyle{main}
% \setcounter{page}{11}
\include{toc}
% \cleardoublepage
\input{Contexto y alcance/main}

\input{Planificación temporal/main}

\input{Gestión económica/main}

\input{Informe de sostenibilidad/main}

% \cleardoublepage

% \cleardoublepage


% \input{Annex/main}


\clearpage
% Estilo bibliográfico IEEE
\bibliographystyle{IEEEtran}
\renewcommand{\refname}{Referencias}
\bibliography{references.bib}

%%%%                                   %%%%
%         FINAL DE LA MEMORIA             %
%%%%                                   %%%%

\end{document}


%% VERSIÓ INCIAL MEMORIA %%




\clearpage
% Estilo bibliográfico IEEE
\bibliographystyle{IEEEtran}
\renewcommand{\refname}{Referencias}
\bibliography{references.bib}

%%%%                                   %%%%
%         FINAL DE LA MEMORIA             %
%%%%                                   %%%%

\end{document}


%% VERSIÓ INCIAL MEMORIA %%




\clearpage
% Estilo bibliográfico IEEE
\bibliographystyle{IEEEtran}
\renewcommand{\refname}{Referencias}
\bibliography{references.bib}

%%%%                                   %%%%
%         FINAL DE LA MEMORIA             %
%%%%                                   %%%%

\end{document}


%% VERSIÓ INCIAL MEMORIA %%



\section{Gestión económica}

La gestión económica de un proyecto es fundamental para asegurar su viabilidad y exito a largo plazo. Ya hemos visto en secciones anteriores, los otros dos vertices 
del triángulo de la gestión de proyectos: el alcance y el tiempo. En esta sección, nos centraremos en cómo gestionar eficazmente los recursos financieros para garantizar
que el proyecto se mantenga dentro del presupuesto y se logren los objetivos establecidos. 

\input{Gestión económica/Recursos humanos.tex}

\subsection{Recursos hardware}
Los gastos de los recursos hardware hacen referencia al coste de todo los equipos físicos necesarios para el desarrollo del proyecto.
En este caso, se ha considerado el uso del mismo equipo par los diferentes roles desempeñados, que contiene el material 
necesario para que cada uno pueda realizar sus funciones. El equipo utilizado es el siguiente:

\begin{table}[H]
    \centering
    \begin{tabular}{|c|c|c|c|c|}
        \hline
        \textbf{Cantidad} & \textbf{Nombre} & \textbf{Coste}\\
        \hline
        3 & hp elitebook 640 14 inch g11 notebook pc & 859,00 € \\
        3 & hp 125 Wired Keyboard PERP & 27,10 € \\
        3 & hp Wired Mouse 1000 & 5,99 € \\
        \hline
    \end{tabular}
    \caption{Resumen del tiempo aproximado para cada tarea. Elaboración propia.}
    \label{tab:estimaciones}
\end{table}


\subsection{Recursos software}
La naturaleza de la organización como entidad sin ánimo de lucro impulsa una estrategia 
rigurosa de optimización de costes, priorizando el uso de soluciones de código abierto 
(\textit{Open Source}) siempre que sea posible. Este enfoque no solo minimiza el gasto 
financiero, sino que también fomenta la transparencia y la adaptabilidad técnica del proyecto.

Para la gestión de recursos de computación y planificación de tareas, se ha optado por 
Slurm, una solución de código abierto ampliamente reconocida en entornos de alto 
rendimiento. En cuanto a la infraestructura colaborativa, la empresa dispone de un acuerdo 
para el uso de Google Workspace, lo que permite gestionar la comunicación y la 
documentación bajo una filosofía de costes reducidos o nulos que se alinea con la estrategia 
\textit{Open Source}.

Para la infraestructura de datos, se utiliza Google Cloud en su plan Estándar, con 
almacenamiento y análisis en BigQuery On-Demand. Según la estimación de cargas del 
proyecto, se prevé un volumen aproximado de:
\begin{itemize}
    \item \textbf{Almacenamiento}: 100--200 GB de datos históricos generados por Slurm, lo que implica un coste mensual aproximado de 4 a 8€ (0,02 €/GB/mes).
    \item \textbf{Consultas}: unos 2 TB procesadas al mes, lo que corresponde a 10 € mensuales (5 €/TB procesado).
\end{itemize}

En total, se estima un coste global de unos \textbf{20–25 € mensuales}, es decir, alrededor de \textbf{80–100 €} para los 4 meses de duración del proyecto. Se ha imputado un presupuesto fijo de \textbf{100 €} para cubrir holguras en picos de almacenamiento o consultas.

Finalmente, la gestión de tareas se realiza a través de \textbf{Jira}, donde se aprovecha una reducción de precio substancial por el estado non-profit de la empresa.

\begin{table}[H]
    \begin{center}
    \begin{tabular}{|c|c|c|c|}
        \hline
        \textbf{Nombre del recurso} & \textbf{Coste} & \textbf{Vida útil} & \textbf{Coste Horario}\\
        \hline
        Google Cloud (BigQuery On-Demand, 4 meses) & 100,00 € & 539 h & 0,186 € \\
        Jira (3 usuarios, coste reducido) & 50,00 € & 539 h & 0,093 € \\
        Slurm (Open Source) & 0,00 € & 539 h & 0,000 € \\
        Google Workspace (Acuerdo non-profit) & 0,00 € & 539 h & 0,000 € \\
        \hline
        \multicolumn{3}{|l|}{\textbf{Total horario} (asumiendo los 3 roles)} & \textbf{0,279 €} \\
        \hline
        \multicolumn{3}{|l|}{\textbf{Total del proyecto} (asumiendo los 3 roles)} & \textbf{150,00 €} \\
        \hline
    \end{tabular}
    \caption{Coste horario de los recursos software indirectos. Elaboración propia.}
    \label{tab:recursos_software}
    \end{center}
\end{table}


\subsection{Contingencias}
Existe una gran posibilidad de que durante la ejecución del proyecto surjan imprevistos que 
requieran recursos adicionales no contemplados en el presupuesto inicial. Es por eso que 
para mitigar el riesgo que implican estos imprevistos, se ha decidido asignar un porcentaje 
del presupuesto total a contingencias. El porcentaje habitual de contingencia para un proyecto 
de desarrollo de software (o sector tecnológico en general) suele oscilar entre el 10\% y el 20\% 
del presupuesto total. Se ha optado por asignar un \textbf{15\%} del presupuesto total a contingencias, 
lo que nos porporciona los siguientes valores:


\begin{table}[H]
    \begin{center}
    \begin{tabular}{|c|c|c|c|}
        \hline
        \textbf{Recurso} & \textbf{Importe} & \textbf{Contingencias} \\ 
        \hline
        Recursos humanos & 12.176,58 € & 1.826,49 € \\
        Recursos materiales directos & 2.676,00 € & 401,40 € \\
        Recursos materiales indirectos  & 6.447,00 € & 967,05 € \\
        \hline
        \multicolumn{2}{|l|}{\textbf{Total}}  & \textbf{3.194,94 €} \\
        \hline
    \end{tabular}
    \caption{Tabla de contingencias. Elaboración propia.}
    \label{tab:calculo_contingencias}
    \end{center}
\end{table}

\section{Informe de sostenibilidad}

\input{Informe de sostenibilidad/Autoevaluación.tex}


\input{Informe de sostenibilidad/Dimensión ambiental.tex}

\subsection{Dimensión económica}

\textbf{¿Has estimado el coste de la realización del proyecto (recursos humanos
y materiales)?}

Sí, se ha realizado una estimación del coste del proyecto, una parte fundamental del proyecto.
Dicha estimación se encuentra en los apartados 6.x en la Gestión económica del proyecto.

\textbf{¿Cómo se resuelve actualmente el problema que quieres abordar (estado del arte)? 
¿En qué mejorará económicamente tu solución a las existentes?}

Siguiendo con la respuesta del apartado anterorior no existe una solución actual al problema a abordar.
Pero se espera que la solución propuesta mejore económicamente las cuentas de la institución. Al optimizar
los costes de los recursos y asignarlos de manera eficiente, se espera que la institución pueda reaprovechar
recursos y así poder comprar equipamiento nuevo que mejore la calidad del servicio ofrecido.

\input{Informe de sostenibilidad/Dimensión social.tex}


% \cleardoublepage

% \cleardoublepage


% \documentclass[a4paper, 11pt, twoside]{article}
\usepackage{wrapfig}
% \usepackage{abstract}
\usepackage{graphicx} % Required for inserting images
\usepackage[a4paper, left=1.15in, right=1.15in, top=1.5in, bottom=1.2in]{geometry}
\usepackage{float}
\usepackage{array}
\usepackage{booktabs}
\usepackage[T1]{fontenc}
\usepackage{amsmath}
\usepackage{tabularx}
\usepackage{graphicx} % Required for inserting images
\usepackage{tikz, xcolor}
\usepackage{hyperref}
\usepackage{cite}
\usepackage{pdfpages}
\usepackage{pdflscape}
\usepackage{afterpage}
\usepackage{textpos}
\usepackage{fancyhdr}
\usepackage{caption}
\usepackage{lmodern,textcomp}
\usepackage[catalan]{babel}
\usepackage{lipsum}
\usepackage{newfloat}
\usepackage{listings}
\usepackage{placeins}
\usepackage[ruled]{algorithm2e}
\usepackage{tcolorbox}
% Estilo de citas IEEE sencillo: usar paquete cite para comprimir y ordenar referencias numéricas
\usepackage{cite}
\usepackage[toc,page]{appendix}
\usepackage{forest}
\usepackage{pdflscape}
\usepackage{graphicx}
\usepackage{geometry}
\usepackage{quoting,xparse}
\usepackage{svg}



% \NewDocumentCommand{\bywhom}{m}{% the Bourbaki trick
%   {\nobreak\hfill\penalty50\hskip1em\null\nobreak
%    \hfill\mbox{\normalfont(#1)}%
%    \parfillskip=0pt \finalhyphendemerits=0 \par}%
% }

% AIXÒ SERVEIX PER A MODIFICAR UN PROBLEMA AMB ELS CARACTERS --> " + PARAULA 

% \NewDocumentEnvironment{pquotation}{m}
%   {\begin{quoting}[
%      indentfirst=true,
%      leftmargin=\parindent,
%      rightmargin=\parindent]\itshape}
%   {\bywhom{#1}
  
%   \end{quoting}}

% AIXÒ PER INSERIR FRAGMENTS DE CODI

\lstnewenvironment{code}[1][]%
{
   \noindent
   \minipage{\linewidth}
   \vspace{0.5\baselineskip}
   \lstset{
     basicstyle=\ttfamily\footnotesize,
     frame=single,
     breaklines=true,    % Permitir que las líneas se rompan automáticamente
     breakatwhitespace=true, % Permitir la ruptura en los espacios en blanco
     #1
   }
}
{\endminipage}

\lstset{
    language=Java,
    basicstyle=\ttfamily\small,  % Define el tipo de letra y tamaño del código
    keywordstyle=\color{blue},   % Define el color de las palabras clave de Java
    stringstyle=\color{red},     % Define el color de las cadenas de texto
    commentstyle=\color{green},  % Define el color de los comentarios
    morecomment=[s][\color{gray}]{/**}{*/},  % Estilo para comentarios de bloque
    breaklines=true,  % Configura el quiebre automático de línea
    numbers=left,  % Posición de los números de línea
    numberstyle=\tiny\color{gray},  % Estilo para los números de línea
    showspaces=false,  % No mostrar espacios con underscores
    showstringspaces=false,  % No mostrar espacios dentro de strings
    frame=single,  % Agregar un marco alrededor del código
    captionpos=b,  % b = bottom, t = top
}

\definecolor{eclipseStrings}{RGB}{42,0.0,255}
\definecolor{eclipseKeywords}{RGB}{127,0,85}
\colorlet{numb}{magenta!60!black}

\lstdefinelanguage{json}{
    basicstyle=\ttfamily\small,
    showstringspaces=false,
    breaklines=true,
    frame=single,
    literate=
     *{0}{{{\color{numb}0}}}{1}
      {1}{{{\color{numb}1}}}{1}
      {2}{{{\color{numb}2}}}{1}
      {3}{{{\color{numb}3}}}{1}
      {4}{{{\color{numb}4}}}{1}
      {5}{{{\color{numb}5}}}{1}
      {6}{{{\color{numb}6}}}{1}
      {7}{{{\color{numb}7}}}{1}
      {8}{{{\color{numb}8}}}{1}
      {9}{{{\color{numb}9}}}{1}
      {:}{{{\color{punct}{:}}}}{1}
      {,}{{{\color{punct}{,}}}}{1}
      {\{}{{{\color{delim}{\{}}}}{1}
      {\}}{{{\color{delim}{\}}}}}{1}
      {[}{{{\color{delim}{[}}}}{1}
      {]}{{{\color{delim}{]}}}}{1},
    morestring=[b]"
}

\lstdefinelanguage{JavaScript}{
  morekeywords=[1]{break, continue, delete, else, for, function, if, in,
    new, return, this, typeof, var, void, while, with},
  % Literals, primitive types, and reference types.
  morekeywords=[2]{false, null, true, boolean, number, undefined,
    Array, Boolean, Date, Math, Number, String, Object},
  % Built-ins.
  morekeywords=[3]{eval, parseInt, parseFloat, escape, unescape},
  sensitive,
  morecomment=[s]{/*}{*/},
  morecomment=[l]//,
  morecomment=[s]{/**}{*/}, % JavaDoc style comments
  morestring=[b]',
  morestring=[b]"
}[keywords, comments, strings]

\lstdefinelanguage{svelte}{
    morekeywords={async, await, then, if, each, in, on, let, const, function, import, export, from, with, typeof, new, var, return},
    sensitive=true,
    morecomment=[l]{//},
    morecomment=[s]{/*}{*/},
    morecomment=[s]{<!--}{-->},
    morestring=[b]',
    morestring=[b]",
    morestring=[b]`,
    ndkeywords={import, export, return, if, for, while, function, switch, case, break, var, let, const, async, await, then, catch, try, typeof, new, class, extends, implements, interface, constructor, public, private, protected, static, get, set},
    alsoletter={:},
    moredelim=**[is][\color{orange}]{\@\{}{\@\}},
    moredelim=**[is][\color{cyan}]{\$\{}{\}},
    keywordstyle=\color{blue}\bfseries,
    commentstyle=\color{gray}\ttfamily,
    stringstyle=\color{red}\ttfamily,
    identifierstyle=\color{black},
    breaklines=true,
    basicstyle=\ttfamily\small,
    columns=fullflexible,
    showstringspaces=false,
    tabsize=2,
    frame=single,
    captionpos=b,
}


\lstdefinestyle{java}{
    language=Java,
    basicstyle=\ttfamily\small,
    keywordstyle=\color{blue},
    stringstyle=\color{red},
    commentstyle=\color{gray},
    morecomment=[s][\color{gray}]{/**}{*/},
    breaklines=true,
    numbers=left,
    numberstyle=\tiny\color{gray},
    showspaces=false,
    showstringspaces=false,
    frame=single,
    captionpos=b
}

\usepackage{color}
\definecolor{numb}{rgb}{0.58,0,0.82}
\definecolor{punct}{rgb}{0.3,0.3,0.3}
\definecolor{delim}{rgb}{0.5,0.5,0.5}
\definecolor{string}{rgb}{0.6,0.1,0.1}

\lstdefinestyle{json}{
    language=json,
    basicstyle=\ttfamily\small,
    keywordstyle=\color{blue},
    stringstyle=\color{string},
    commentstyle=\color{green},
    morecomment=[s][\color{gray}]{/**}{*/},
    breaklines=true,
    numbers=left,
    numberstyle=\tiny\color{gray},
    showspaces=false,
    showstringspaces=false,
    frame=single,
    captionpos=b
}

\lstdefinelanguage{drl}{
    basicstyle=\ttfamily\small,
    keywordstyle=\color{blue}\bfseries,
    stringstyle=\color{red},
    commentstyle=\color{green},
    morecomment=[s][\color{green}]{/*}{*/},
    morecomment=[l][\color{green}]{//},
    morestring=[b]",
    keywords={rule, when, then, end, salience, eval, update, retract},
    sensitive=true,
    breaklines=true,
    numbers=left,
    numberstyle=\tiny\color{gray},
    showspaces=false,
    showstringspaces=false,
    frame=single,
    captionpos=b
}

\lstdefinestyle{drl}{
    language=drl,
    basicstyle=\ttfamily\small,
    keywordstyle=\color{blue}\bfseries,
    stringstyle=\color{red},
    commentstyle=\color{green},
    breaklines=true,
    numbers=left,
    numberstyle=\tiny\color{gray},
    showspaces=false,
    showstringspaces=false,
    frame=single,
    captionpos=b
}



\renewcommand{\lstlistingname}{Codi}
\renewcommand{\lstlistlistingname}{Índex de Fragments de Codis}
\renewcommand{\appendixpagename}{Annexos}
\renewcommand{\appendixtocname}{Annexos}

\renewcommand*{\listalgorithmcfname}{Llistat de pseudocodi}
\renewcommand*{\algorithmcfname}{Pseudocodi}
\renewcommand*{\algorithmautorefname}{pseudocodi}

\definecolor{darkgreen}{rgb}{0.0, 0.3, 0.0}
\definecolor{lightgreen}{rgb}{0.9, 0.95, 0.9}

\newcommand{\todo}[1]{
    \begin{tcolorbox}[colback=lightgreen, colframe=darkgreen, title=TODO]
        #1
    \end{tcolorbox}
}


% \usepackage{indentfirst} 
% \setlength{\parindent}{1cm}


\usepackage{parskip}
% \setlength{\parindent}{10pt}

% \setlength{\headheight}{15.6pt}
\setlength{\headheight}{16.05263pt}




\vspace{5mm}

\newcolumntype{u}{>{\centering\arraybackslash\hsize=.22\hsize}X}
\newcolumntype{s}{>{\hsize=.45\hsize}X}


%%%%%%%%%%%%%%%%%%%%%%%%%%%%%%%%
% INI creacio subsubsubsecion
%%%%%%%%%%%%%%%%%%%%%%%%%%%%%%%%
\usepackage{titlesec}

\let\oldparagraph\paragraph
\renewcommand{\paragraph}[1]{\oldparagraph{\hspace{30pt}#1}}


\titleclass{\subsubsubsection}{straight}[\subsection]

\newcounter{subsubsubsection}[subsubsection]
\renewcommand\thesubsubsubsection{\thesubsubsection.\arabic{subsubsubsection}}
\renewcommand\theparagraph{\thesubsubsubsection.\arabic{paragraph}} % optional; useful if paragraphs are to be numbered

\titleformat{\subsubsubsection}
  {\normalfont\normalsize\bfseries}{\thesubsubsubsection}{1em}{}
\titlespacing*{\subsubsubsection}
{0pt}{3.25ex plus 1ex minus .2ex}{1.5ex plus .2ex}


\makeatletter
\renewcommand\paragraph{\@startsection{paragraph}{5}{10pt}%
  {3.25ex \@plus1ex \@minus.2ex}%
  {-1em}%
  {\normalfont\normalsize\bfseries}}
\renewcommand\subparagraph{\@startsection{subparagraph}{6}{\parindent}%
  {3.25ex \@plus1ex \@minus .2ex}%
  {-1em}%
  {\normalfont\normalsize\bfseries}}
\def\toclevel@subsubsubsection{4}
\def\toclevel@paragraph{5}
\def\toclevel@paragraph{6}
\def\l@subsubsubsection{\@dottedtocline{4}{7em}{4em}}
\def\l@paragraph{\@dottedtocline{5}{10em}{5em}}
\def\l@subparagraph{\@dottedtocline{6}{14em}{6em}}
\makeatother

\setcounter{secnumdepth}{4}
\setcounter{tocdepth}{3}
%%%%%%%%%%%%%%%%%%%%%%%%%%%%%%%%
% FIN creacio subsubsubtitle
%%%%%%%%%%%%%%%%%%%%%%%%%%%%%%%%


%%%%%%%%%%%%%%%%%%%%%%%%%%%%%%%%
% INI creacio equation
%%%%%%%%%%%%%%%%%%%%%%%%%%%%%%%%
\DeclareFloatingEnvironment[fileext=loeq,placement={!ht},name=Equació]{equationfloat}
%%%%%%%%%%%%%%%%%%%%%%%%%%%%%%%%
% FIN creacio equation
%%%%%%%%%%%%%%%%%%%%%%%%%%%%%%%%

\titleformat{\section}[block]
  {\clearpage\normalfont\Large\bfseries}
  {\thesection}{1em}{}


\renewcommand\refname{Bibliografia}


\linespread{1.25}


\makeatletter
\renewenvironment{abstract}{%
    \if@twocolumn
      \section*{\abstractname}%
    \else 
      \begin{center}%
        {\bfseries \Large\itshape\abstractname\vspace{\z@}}%  % Título en negrita, 14pt y cursiva
      \end{center}%
      \quotation
    \fi}
    {\if@twocolumn\else\endquotation\fi}
\makeatother

% \fancypagestyle{main}{
%     \fancyhf{}
%     \lhead{\small Aplicació Web per a la Visualització i Comparació
% de Trajectòries d’Algorismes d’Optimització}
%     % \rhead{\includegraphics[width=1.2cm]{resources/Icons/logo-fib.png}}
%     \cfoot{\thepage}
%     \setlength{\footskip}{35pt}
% }

\newcommand{\q}[1]{``#1''}

\begin{document}

\hypersetup{
    pdfborder=0 0 0,
}

%% PORTADA I TABLE OF CONTENTS %%

%%%%                                   %%%%
%         INDEX DE LA MEMORIA             %
%%%%                                   %%%%

% \includepdf[pages=-]{portada_fib.pdf}
%%% Portada
\begin{titlepage}
\begin{center}
    \includegraphics[width=1\linewidth]{extra/figures/logo-upc-fib.png}
    \vspace{30mm}

    \LARGE{\textsc{\textbf{{INTEGRACIÓN Y AUTOMATIZACIÓN DE COSTES DE INFRAESTRUCTURA HPC EN ENTORNOS DE INVESTIGACIÓN}}}} \\
    \vspace{25mm}
      
        
    \textbf{\Large{\textsc{PAU NAVARRO ÁLVAREZ}}} \\
    \vspace{20mm}
        
    
    
    \large{\textbf{Director/a} ROBERTO RIVEIRO INSUA (FUNDACIO INSTITUT DE RECERCA BIOMEDICA (IRB BARCELONA))} \\
    \large{\textbf{Ponente:} ALEX PAJUELO GONZÁLEZ (Departamento de Arquitectura de Computadores)} \\
    \large{\textbf{Titulación:} Grado en Ingeniería Informática (Tecnologías de la información)} \\
    
    \large{\textbf{Lliurament 2:} }Planificación temporal \\
    \large{\textbf{Tutor GEP:} CAROLINA CONSOLIDACIÓN SEGURA (Departament d'Organització d'Empreses)} \\
    \large{\textbf{Fecha:}} 29 de Septiembre de 2025\\

    \textbf{\large{Memoria del trabajo de fin de grado}} \\
    \vspace{3mm}
    \textbf{\large{Facultat d'Informàtica de Barcelona (FIB)}} \\
    \vspace{3mm}
    \textbf{\large{Universitat Politècnica de Catalunya (UPC) - BarcelonaTech}} \\

\end{center}
\end{titlepage}
%\input{Extra/portada}
% \cleardoublepage

% \pagestyle{main}
% \setcounter{page}{11}
\renewcommand{\contentsname}{Continguts}

\tableofcontents
\clearpage
\listoffigures
\begingroup
    \let\clearpage\relax
    \listoftables
\endgroup
\clearpage

\newpage
% \cleardoublepage
\documentclass[a4paper, 11pt, twoside]{article}
\usepackage{wrapfig}
% \usepackage{abstract}
\usepackage{graphicx} % Required for inserting images
\usepackage[a4paper, left=1.15in, right=1.15in, top=1.5in, bottom=1.2in]{geometry}
\usepackage{float}
\usepackage{array}
\usepackage{booktabs}
\usepackage[T1]{fontenc}
\usepackage{amsmath}
\usepackage{tabularx}
\usepackage{graphicx} % Required for inserting images
\usepackage{tikz, xcolor}
\usepackage{hyperref}
\usepackage{cite}
\usepackage{pdfpages}
\usepackage{pdflscape}
\usepackage{afterpage}
\usepackage{textpos}
\usepackage{fancyhdr}
\usepackage{caption}
\usepackage{lmodern,textcomp}
\usepackage[catalan]{babel}
\usepackage{lipsum}
\usepackage{newfloat}
\usepackage{listings}
\usepackage{placeins}
\usepackage[ruled]{algorithm2e}
\usepackage{tcolorbox}
% Estilo de citas IEEE sencillo: usar paquete cite para comprimir y ordenar referencias numéricas
\usepackage{cite}
\usepackage[toc,page]{appendix}
\usepackage{forest}
\usepackage{pdflscape}
\usepackage{graphicx}
\usepackage{geometry}
\usepackage{quoting,xparse}
\usepackage{svg}



% \NewDocumentCommand{\bywhom}{m}{% the Bourbaki trick
%   {\nobreak\hfill\penalty50\hskip1em\null\nobreak
%    \hfill\mbox{\normalfont(#1)}%
%    \parfillskip=0pt \finalhyphendemerits=0 \par}%
% }

% AIXÒ SERVEIX PER A MODIFICAR UN PROBLEMA AMB ELS CARACTERS --> " + PARAULA 

% \NewDocumentEnvironment{pquotation}{m}
%   {\begin{quoting}[
%      indentfirst=true,
%      leftmargin=\parindent,
%      rightmargin=\parindent]\itshape}
%   {\bywhom{#1}
  
%   \end{quoting}}

% AIXÒ PER INSERIR FRAGMENTS DE CODI

\lstnewenvironment{code}[1][]%
{
   \noindent
   \minipage{\linewidth}
   \vspace{0.5\baselineskip}
   \lstset{
     basicstyle=\ttfamily\footnotesize,
     frame=single,
     breaklines=true,    % Permitir que las líneas se rompan automáticamente
     breakatwhitespace=true, % Permitir la ruptura en los espacios en blanco
     #1
   }
}
{\endminipage}

\lstset{
    language=Java,
    basicstyle=\ttfamily\small,  % Define el tipo de letra y tamaño del código
    keywordstyle=\color{blue},   % Define el color de las palabras clave de Java
    stringstyle=\color{red},     % Define el color de las cadenas de texto
    commentstyle=\color{green},  % Define el color de los comentarios
    morecomment=[s][\color{gray}]{/**}{*/},  % Estilo para comentarios de bloque
    breaklines=true,  % Configura el quiebre automático de línea
    numbers=left,  % Posición de los números de línea
    numberstyle=\tiny\color{gray},  % Estilo para los números de línea
    showspaces=false,  % No mostrar espacios con underscores
    showstringspaces=false,  % No mostrar espacios dentro de strings
    frame=single,  % Agregar un marco alrededor del código
    captionpos=b,  % b = bottom, t = top
}

\definecolor{eclipseStrings}{RGB}{42,0.0,255}
\definecolor{eclipseKeywords}{RGB}{127,0,85}
\colorlet{numb}{magenta!60!black}

\lstdefinelanguage{json}{
    basicstyle=\ttfamily\small,
    showstringspaces=false,
    breaklines=true,
    frame=single,
    literate=
     *{0}{{{\color{numb}0}}}{1}
      {1}{{{\color{numb}1}}}{1}
      {2}{{{\color{numb}2}}}{1}
      {3}{{{\color{numb}3}}}{1}
      {4}{{{\color{numb}4}}}{1}
      {5}{{{\color{numb}5}}}{1}
      {6}{{{\color{numb}6}}}{1}
      {7}{{{\color{numb}7}}}{1}
      {8}{{{\color{numb}8}}}{1}
      {9}{{{\color{numb}9}}}{1}
      {:}{{{\color{punct}{:}}}}{1}
      {,}{{{\color{punct}{,}}}}{1}
      {\{}{{{\color{delim}{\{}}}}{1}
      {\}}{{{\color{delim}{\}}}}}{1}
      {[}{{{\color{delim}{[}}}}{1}
      {]}{{{\color{delim}{]}}}}{1},
    morestring=[b]"
}

\lstdefinelanguage{JavaScript}{
  morekeywords=[1]{break, continue, delete, else, for, function, if, in,
    new, return, this, typeof, var, void, while, with},
  % Literals, primitive types, and reference types.
  morekeywords=[2]{false, null, true, boolean, number, undefined,
    Array, Boolean, Date, Math, Number, String, Object},
  % Built-ins.
  morekeywords=[3]{eval, parseInt, parseFloat, escape, unescape},
  sensitive,
  morecomment=[s]{/*}{*/},
  morecomment=[l]//,
  morecomment=[s]{/**}{*/}, % JavaDoc style comments
  morestring=[b]',
  morestring=[b]"
}[keywords, comments, strings]

\lstdefinelanguage{svelte}{
    morekeywords={async, await, then, if, each, in, on, let, const, function, import, export, from, with, typeof, new, var, return},
    sensitive=true,
    morecomment=[l]{//},
    morecomment=[s]{/*}{*/},
    morecomment=[s]{<!--}{-->},
    morestring=[b]',
    morestring=[b]",
    morestring=[b]`,
    ndkeywords={import, export, return, if, for, while, function, switch, case, break, var, let, const, async, await, then, catch, try, typeof, new, class, extends, implements, interface, constructor, public, private, protected, static, get, set},
    alsoletter={:},
    moredelim=**[is][\color{orange}]{\@\{}{\@\}},
    moredelim=**[is][\color{cyan}]{\$\{}{\}},
    keywordstyle=\color{blue}\bfseries,
    commentstyle=\color{gray}\ttfamily,
    stringstyle=\color{red}\ttfamily,
    identifierstyle=\color{black},
    breaklines=true,
    basicstyle=\ttfamily\small,
    columns=fullflexible,
    showstringspaces=false,
    tabsize=2,
    frame=single,
    captionpos=b,
}


\lstdefinestyle{java}{
    language=Java,
    basicstyle=\ttfamily\small,
    keywordstyle=\color{blue},
    stringstyle=\color{red},
    commentstyle=\color{gray},
    morecomment=[s][\color{gray}]{/**}{*/},
    breaklines=true,
    numbers=left,
    numberstyle=\tiny\color{gray},
    showspaces=false,
    showstringspaces=false,
    frame=single,
    captionpos=b
}

\usepackage{color}
\definecolor{numb}{rgb}{0.58,0,0.82}
\definecolor{punct}{rgb}{0.3,0.3,0.3}
\definecolor{delim}{rgb}{0.5,0.5,0.5}
\definecolor{string}{rgb}{0.6,0.1,0.1}

\lstdefinestyle{json}{
    language=json,
    basicstyle=\ttfamily\small,
    keywordstyle=\color{blue},
    stringstyle=\color{string},
    commentstyle=\color{green},
    morecomment=[s][\color{gray}]{/**}{*/},
    breaklines=true,
    numbers=left,
    numberstyle=\tiny\color{gray},
    showspaces=false,
    showstringspaces=false,
    frame=single,
    captionpos=b
}

\lstdefinelanguage{drl}{
    basicstyle=\ttfamily\small,
    keywordstyle=\color{blue}\bfseries,
    stringstyle=\color{red},
    commentstyle=\color{green},
    morecomment=[s][\color{green}]{/*}{*/},
    morecomment=[l][\color{green}]{//},
    morestring=[b]",
    keywords={rule, when, then, end, salience, eval, update, retract},
    sensitive=true,
    breaklines=true,
    numbers=left,
    numberstyle=\tiny\color{gray},
    showspaces=false,
    showstringspaces=false,
    frame=single,
    captionpos=b
}

\lstdefinestyle{drl}{
    language=drl,
    basicstyle=\ttfamily\small,
    keywordstyle=\color{blue}\bfseries,
    stringstyle=\color{red},
    commentstyle=\color{green},
    breaklines=true,
    numbers=left,
    numberstyle=\tiny\color{gray},
    showspaces=false,
    showstringspaces=false,
    frame=single,
    captionpos=b
}



\renewcommand{\lstlistingname}{Codi}
\renewcommand{\lstlistlistingname}{Índex de Fragments de Codis}
\renewcommand{\appendixpagename}{Annexos}
\renewcommand{\appendixtocname}{Annexos}

\renewcommand*{\listalgorithmcfname}{Llistat de pseudocodi}
\renewcommand*{\algorithmcfname}{Pseudocodi}
\renewcommand*{\algorithmautorefname}{pseudocodi}

\definecolor{darkgreen}{rgb}{0.0, 0.3, 0.0}
\definecolor{lightgreen}{rgb}{0.9, 0.95, 0.9}

\newcommand{\todo}[1]{
    \begin{tcolorbox}[colback=lightgreen, colframe=darkgreen, title=TODO]
        #1
    \end{tcolorbox}
}


% \usepackage{indentfirst} 
% \setlength{\parindent}{1cm}


\usepackage{parskip}
% \setlength{\parindent}{10pt}

% \setlength{\headheight}{15.6pt}
\setlength{\headheight}{16.05263pt}




\vspace{5mm}

\newcolumntype{u}{>{\centering\arraybackslash\hsize=.22\hsize}X}
\newcolumntype{s}{>{\hsize=.45\hsize}X}


%%%%%%%%%%%%%%%%%%%%%%%%%%%%%%%%
% INI creacio subsubsubsecion
%%%%%%%%%%%%%%%%%%%%%%%%%%%%%%%%
\usepackage{titlesec}

\let\oldparagraph\paragraph
\renewcommand{\paragraph}[1]{\oldparagraph{\hspace{30pt}#1}}


\titleclass{\subsubsubsection}{straight}[\subsection]

\newcounter{subsubsubsection}[subsubsection]
\renewcommand\thesubsubsubsection{\thesubsubsection.\arabic{subsubsubsection}}
\renewcommand\theparagraph{\thesubsubsubsection.\arabic{paragraph}} % optional; useful if paragraphs are to be numbered

\titleformat{\subsubsubsection}
  {\normalfont\normalsize\bfseries}{\thesubsubsubsection}{1em}{}
\titlespacing*{\subsubsubsection}
{0pt}{3.25ex plus 1ex minus .2ex}{1.5ex plus .2ex}


\makeatletter
\renewcommand\paragraph{\@startsection{paragraph}{5}{10pt}%
  {3.25ex \@plus1ex \@minus.2ex}%
  {-1em}%
  {\normalfont\normalsize\bfseries}}
\renewcommand\subparagraph{\@startsection{subparagraph}{6}{\parindent}%
  {3.25ex \@plus1ex \@minus .2ex}%
  {-1em}%
  {\normalfont\normalsize\bfseries}}
\def\toclevel@subsubsubsection{4}
\def\toclevel@paragraph{5}
\def\toclevel@paragraph{6}
\def\l@subsubsubsection{\@dottedtocline{4}{7em}{4em}}
\def\l@paragraph{\@dottedtocline{5}{10em}{5em}}
\def\l@subparagraph{\@dottedtocline{6}{14em}{6em}}
\makeatother

\setcounter{secnumdepth}{4}
\setcounter{tocdepth}{3}
%%%%%%%%%%%%%%%%%%%%%%%%%%%%%%%%
% FIN creacio subsubsubtitle
%%%%%%%%%%%%%%%%%%%%%%%%%%%%%%%%


%%%%%%%%%%%%%%%%%%%%%%%%%%%%%%%%
% INI creacio equation
%%%%%%%%%%%%%%%%%%%%%%%%%%%%%%%%
\DeclareFloatingEnvironment[fileext=loeq,placement={!ht},name=Equació]{equationfloat}
%%%%%%%%%%%%%%%%%%%%%%%%%%%%%%%%
% FIN creacio equation
%%%%%%%%%%%%%%%%%%%%%%%%%%%%%%%%

\titleformat{\section}[block]
  {\clearpage\normalfont\Large\bfseries}
  {\thesection}{1em}{}


\renewcommand\refname{Bibliografia}


\linespread{1.25}


\makeatletter
\renewenvironment{abstract}{%
    \if@twocolumn
      \section*{\abstractname}%
    \else 
      \begin{center}%
        {\bfseries \Large\itshape\abstractname\vspace{\z@}}%  % Título en negrita, 14pt y cursiva
      \end{center}%
      \quotation
    \fi}
    {\if@twocolumn\else\endquotation\fi}
\makeatother

% \fancypagestyle{main}{
%     \fancyhf{}
%     \lhead{\small Aplicació Web per a la Visualització i Comparació
% de Trajectòries d’Algorismes d’Optimització}
%     % \rhead{\includegraphics[width=1.2cm]{resources/Icons/logo-fib.png}}
%     \cfoot{\thepage}
%     \setlength{\footskip}{35pt}
% }

\newcommand{\q}[1]{``#1''}

\begin{document}

\hypersetup{
    pdfborder=0 0 0,
}

%% PORTADA I TABLE OF CONTENTS %%

%%%%                                   %%%%
%         INDEX DE LA MEMORIA             %
%%%%                                   %%%%

% \includepdf[pages=-]{portada_fib.pdf}
\input{Extra/GEP}
%\input{Extra/portada}
% \cleardoublepage

% \pagestyle{main}
% \setcounter{page}{11}
\include{toc}
% \cleardoublepage
\input{Contexto y alcance/main}

\input{Planificación temporal/main}

\input{Gestión económica/main}

\input{Informe de sostenibilidad/main}

% \cleardoublepage

% \cleardoublepage


% \input{Annex/main}


\clearpage
% Estilo bibliográfico IEEE
\bibliographystyle{IEEEtran}
\renewcommand{\refname}{Referencias}
\bibliography{references.bib}

%%%%                                   %%%%
%         FINAL DE LA MEMORIA             %
%%%%                                   %%%%

\end{document}


%% VERSIÓ INCIAL MEMORIA %%


\documentclass[a4paper, 11pt, twoside]{article}
\usepackage{wrapfig}
% \usepackage{abstract}
\usepackage{graphicx} % Required for inserting images
\usepackage[a4paper, left=1.15in, right=1.15in, top=1.5in, bottom=1.2in]{geometry}
\usepackage{float}
\usepackage{array}
\usepackage{booktabs}
\usepackage[T1]{fontenc}
\usepackage{amsmath}
\usepackage{tabularx}
\usepackage{graphicx} % Required for inserting images
\usepackage{tikz, xcolor}
\usepackage{hyperref}
\usepackage{cite}
\usepackage{pdfpages}
\usepackage{pdflscape}
\usepackage{afterpage}
\usepackage{textpos}
\usepackage{fancyhdr}
\usepackage{caption}
\usepackage{lmodern,textcomp}
\usepackage[catalan]{babel}
\usepackage{lipsum}
\usepackage{newfloat}
\usepackage{listings}
\usepackage{placeins}
\usepackage[ruled]{algorithm2e}
\usepackage{tcolorbox}
% Estilo de citas IEEE sencillo: usar paquete cite para comprimir y ordenar referencias numéricas
\usepackage{cite}
\usepackage[toc,page]{appendix}
\usepackage{forest}
\usepackage{pdflscape}
\usepackage{graphicx}
\usepackage{geometry}
\usepackage{quoting,xparse}
\usepackage{svg}



% \NewDocumentCommand{\bywhom}{m}{% the Bourbaki trick
%   {\nobreak\hfill\penalty50\hskip1em\null\nobreak
%    \hfill\mbox{\normalfont(#1)}%
%    \parfillskip=0pt \finalhyphendemerits=0 \par}%
% }

% AIXÒ SERVEIX PER A MODIFICAR UN PROBLEMA AMB ELS CARACTERS --> " + PARAULA 

% \NewDocumentEnvironment{pquotation}{m}
%   {\begin{quoting}[
%      indentfirst=true,
%      leftmargin=\parindent,
%      rightmargin=\parindent]\itshape}
%   {\bywhom{#1}
  
%   \end{quoting}}

% AIXÒ PER INSERIR FRAGMENTS DE CODI

\lstnewenvironment{code}[1][]%
{
   \noindent
   \minipage{\linewidth}
   \vspace{0.5\baselineskip}
   \lstset{
     basicstyle=\ttfamily\footnotesize,
     frame=single,
     breaklines=true,    % Permitir que las líneas se rompan automáticamente
     breakatwhitespace=true, % Permitir la ruptura en los espacios en blanco
     #1
   }
}
{\endminipage}

\lstset{
    language=Java,
    basicstyle=\ttfamily\small,  % Define el tipo de letra y tamaño del código
    keywordstyle=\color{blue},   % Define el color de las palabras clave de Java
    stringstyle=\color{red},     % Define el color de las cadenas de texto
    commentstyle=\color{green},  % Define el color de los comentarios
    morecomment=[s][\color{gray}]{/**}{*/},  % Estilo para comentarios de bloque
    breaklines=true,  % Configura el quiebre automático de línea
    numbers=left,  % Posición de los números de línea
    numberstyle=\tiny\color{gray},  % Estilo para los números de línea
    showspaces=false,  % No mostrar espacios con underscores
    showstringspaces=false,  % No mostrar espacios dentro de strings
    frame=single,  % Agregar un marco alrededor del código
    captionpos=b,  % b = bottom, t = top
}

\definecolor{eclipseStrings}{RGB}{42,0.0,255}
\definecolor{eclipseKeywords}{RGB}{127,0,85}
\colorlet{numb}{magenta!60!black}

\lstdefinelanguage{json}{
    basicstyle=\ttfamily\small,
    showstringspaces=false,
    breaklines=true,
    frame=single,
    literate=
     *{0}{{{\color{numb}0}}}{1}
      {1}{{{\color{numb}1}}}{1}
      {2}{{{\color{numb}2}}}{1}
      {3}{{{\color{numb}3}}}{1}
      {4}{{{\color{numb}4}}}{1}
      {5}{{{\color{numb}5}}}{1}
      {6}{{{\color{numb}6}}}{1}
      {7}{{{\color{numb}7}}}{1}
      {8}{{{\color{numb}8}}}{1}
      {9}{{{\color{numb}9}}}{1}
      {:}{{{\color{punct}{:}}}}{1}
      {,}{{{\color{punct}{,}}}}{1}
      {\{}{{{\color{delim}{\{}}}}{1}
      {\}}{{{\color{delim}{\}}}}}{1}
      {[}{{{\color{delim}{[}}}}{1}
      {]}{{{\color{delim}{]}}}}{1},
    morestring=[b]"
}

\lstdefinelanguage{JavaScript}{
  morekeywords=[1]{break, continue, delete, else, for, function, if, in,
    new, return, this, typeof, var, void, while, with},
  % Literals, primitive types, and reference types.
  morekeywords=[2]{false, null, true, boolean, number, undefined,
    Array, Boolean, Date, Math, Number, String, Object},
  % Built-ins.
  morekeywords=[3]{eval, parseInt, parseFloat, escape, unescape},
  sensitive,
  morecomment=[s]{/*}{*/},
  morecomment=[l]//,
  morecomment=[s]{/**}{*/}, % JavaDoc style comments
  morestring=[b]',
  morestring=[b]"
}[keywords, comments, strings]

\lstdefinelanguage{svelte}{
    morekeywords={async, await, then, if, each, in, on, let, const, function, import, export, from, with, typeof, new, var, return},
    sensitive=true,
    morecomment=[l]{//},
    morecomment=[s]{/*}{*/},
    morecomment=[s]{<!--}{-->},
    morestring=[b]',
    morestring=[b]",
    morestring=[b]`,
    ndkeywords={import, export, return, if, for, while, function, switch, case, break, var, let, const, async, await, then, catch, try, typeof, new, class, extends, implements, interface, constructor, public, private, protected, static, get, set},
    alsoletter={:},
    moredelim=**[is][\color{orange}]{\@\{}{\@\}},
    moredelim=**[is][\color{cyan}]{\$\{}{\}},
    keywordstyle=\color{blue}\bfseries,
    commentstyle=\color{gray}\ttfamily,
    stringstyle=\color{red}\ttfamily,
    identifierstyle=\color{black},
    breaklines=true,
    basicstyle=\ttfamily\small,
    columns=fullflexible,
    showstringspaces=false,
    tabsize=2,
    frame=single,
    captionpos=b,
}


\lstdefinestyle{java}{
    language=Java,
    basicstyle=\ttfamily\small,
    keywordstyle=\color{blue},
    stringstyle=\color{red},
    commentstyle=\color{gray},
    morecomment=[s][\color{gray}]{/**}{*/},
    breaklines=true,
    numbers=left,
    numberstyle=\tiny\color{gray},
    showspaces=false,
    showstringspaces=false,
    frame=single,
    captionpos=b
}

\usepackage{color}
\definecolor{numb}{rgb}{0.58,0,0.82}
\definecolor{punct}{rgb}{0.3,0.3,0.3}
\definecolor{delim}{rgb}{0.5,0.5,0.5}
\definecolor{string}{rgb}{0.6,0.1,0.1}

\lstdefinestyle{json}{
    language=json,
    basicstyle=\ttfamily\small,
    keywordstyle=\color{blue},
    stringstyle=\color{string},
    commentstyle=\color{green},
    morecomment=[s][\color{gray}]{/**}{*/},
    breaklines=true,
    numbers=left,
    numberstyle=\tiny\color{gray},
    showspaces=false,
    showstringspaces=false,
    frame=single,
    captionpos=b
}

\lstdefinelanguage{drl}{
    basicstyle=\ttfamily\small,
    keywordstyle=\color{blue}\bfseries,
    stringstyle=\color{red},
    commentstyle=\color{green},
    morecomment=[s][\color{green}]{/*}{*/},
    morecomment=[l][\color{green}]{//},
    morestring=[b]",
    keywords={rule, when, then, end, salience, eval, update, retract},
    sensitive=true,
    breaklines=true,
    numbers=left,
    numberstyle=\tiny\color{gray},
    showspaces=false,
    showstringspaces=false,
    frame=single,
    captionpos=b
}

\lstdefinestyle{drl}{
    language=drl,
    basicstyle=\ttfamily\small,
    keywordstyle=\color{blue}\bfseries,
    stringstyle=\color{red},
    commentstyle=\color{green},
    breaklines=true,
    numbers=left,
    numberstyle=\tiny\color{gray},
    showspaces=false,
    showstringspaces=false,
    frame=single,
    captionpos=b
}



\renewcommand{\lstlistingname}{Codi}
\renewcommand{\lstlistlistingname}{Índex de Fragments de Codis}
\renewcommand{\appendixpagename}{Annexos}
\renewcommand{\appendixtocname}{Annexos}

\renewcommand*{\listalgorithmcfname}{Llistat de pseudocodi}
\renewcommand*{\algorithmcfname}{Pseudocodi}
\renewcommand*{\algorithmautorefname}{pseudocodi}

\definecolor{darkgreen}{rgb}{0.0, 0.3, 0.0}
\definecolor{lightgreen}{rgb}{0.9, 0.95, 0.9}

\newcommand{\todo}[1]{
    \begin{tcolorbox}[colback=lightgreen, colframe=darkgreen, title=TODO]
        #1
    \end{tcolorbox}
}


% \usepackage{indentfirst} 
% \setlength{\parindent}{1cm}


\usepackage{parskip}
% \setlength{\parindent}{10pt}

% \setlength{\headheight}{15.6pt}
\setlength{\headheight}{16.05263pt}




\vspace{5mm}

\newcolumntype{u}{>{\centering\arraybackslash\hsize=.22\hsize}X}
\newcolumntype{s}{>{\hsize=.45\hsize}X}


%%%%%%%%%%%%%%%%%%%%%%%%%%%%%%%%
% INI creacio subsubsubsecion
%%%%%%%%%%%%%%%%%%%%%%%%%%%%%%%%
\usepackage{titlesec}

\let\oldparagraph\paragraph
\renewcommand{\paragraph}[1]{\oldparagraph{\hspace{30pt}#1}}


\titleclass{\subsubsubsection}{straight}[\subsection]

\newcounter{subsubsubsection}[subsubsection]
\renewcommand\thesubsubsubsection{\thesubsubsection.\arabic{subsubsubsection}}
\renewcommand\theparagraph{\thesubsubsubsection.\arabic{paragraph}} % optional; useful if paragraphs are to be numbered

\titleformat{\subsubsubsection}
  {\normalfont\normalsize\bfseries}{\thesubsubsubsection}{1em}{}
\titlespacing*{\subsubsubsection}
{0pt}{3.25ex plus 1ex minus .2ex}{1.5ex plus .2ex}


\makeatletter
\renewcommand\paragraph{\@startsection{paragraph}{5}{10pt}%
  {3.25ex \@plus1ex \@minus.2ex}%
  {-1em}%
  {\normalfont\normalsize\bfseries}}
\renewcommand\subparagraph{\@startsection{subparagraph}{6}{\parindent}%
  {3.25ex \@plus1ex \@minus .2ex}%
  {-1em}%
  {\normalfont\normalsize\bfseries}}
\def\toclevel@subsubsubsection{4}
\def\toclevel@paragraph{5}
\def\toclevel@paragraph{6}
\def\l@subsubsubsection{\@dottedtocline{4}{7em}{4em}}
\def\l@paragraph{\@dottedtocline{5}{10em}{5em}}
\def\l@subparagraph{\@dottedtocline{6}{14em}{6em}}
\makeatother

\setcounter{secnumdepth}{4}
\setcounter{tocdepth}{3}
%%%%%%%%%%%%%%%%%%%%%%%%%%%%%%%%
% FIN creacio subsubsubtitle
%%%%%%%%%%%%%%%%%%%%%%%%%%%%%%%%


%%%%%%%%%%%%%%%%%%%%%%%%%%%%%%%%
% INI creacio equation
%%%%%%%%%%%%%%%%%%%%%%%%%%%%%%%%
\DeclareFloatingEnvironment[fileext=loeq,placement={!ht},name=Equació]{equationfloat}
%%%%%%%%%%%%%%%%%%%%%%%%%%%%%%%%
% FIN creacio equation
%%%%%%%%%%%%%%%%%%%%%%%%%%%%%%%%

\titleformat{\section}[block]
  {\clearpage\normalfont\Large\bfseries}
  {\thesection}{1em}{}


\renewcommand\refname{Bibliografia}


\linespread{1.25}


\makeatletter
\renewenvironment{abstract}{%
    \if@twocolumn
      \section*{\abstractname}%
    \else 
      \begin{center}%
        {\bfseries \Large\itshape\abstractname\vspace{\z@}}%  % Título en negrita, 14pt y cursiva
      \end{center}%
      \quotation
    \fi}
    {\if@twocolumn\else\endquotation\fi}
\makeatother

% \fancypagestyle{main}{
%     \fancyhf{}
%     \lhead{\small Aplicació Web per a la Visualització i Comparació
% de Trajectòries d’Algorismes d’Optimització}
%     % \rhead{\includegraphics[width=1.2cm]{resources/Icons/logo-fib.png}}
%     \cfoot{\thepage}
%     \setlength{\footskip}{35pt}
% }

\newcommand{\q}[1]{``#1''}

\begin{document}

\hypersetup{
    pdfborder=0 0 0,
}

%% PORTADA I TABLE OF CONTENTS %%

%%%%                                   %%%%
%         INDEX DE LA MEMORIA             %
%%%%                                   %%%%

% \includepdf[pages=-]{portada_fib.pdf}
\input{Extra/GEP}
%\input{Extra/portada}
% \cleardoublepage

% \pagestyle{main}
% \setcounter{page}{11}
\include{toc}
% \cleardoublepage
\input{Contexto y alcance/main}

\input{Planificación temporal/main}

\input{Gestión económica/main}

\input{Informe de sostenibilidad/main}

% \cleardoublepage

% \cleardoublepage


% \input{Annex/main}


\clearpage
% Estilo bibliográfico IEEE
\bibliographystyle{IEEEtran}
\renewcommand{\refname}{Referencias}
\bibliography{references.bib}

%%%%                                   %%%%
%         FINAL DE LA MEMORIA             %
%%%%                                   %%%%

\end{document}


%% VERSIÓ INCIAL MEMORIA %%


\section{Alcance}

\input{Contexto y alcance/Alcance/Objetivo principal}

\input{Contexto y alcance/Alcance/Objetivos secundarios}

\input{Contexto y alcance/Alcance/Requerimientos funcionales}

\input{Contexto y alcance/Alcance/Requerimientos no funcionales}

\input{Contexto y alcance/Alcance/Posibles obstaculos y riesgos}
\documentclass[a4paper, 11pt, twoside]{article}
\usepackage{wrapfig}
% \usepackage{abstract}
\usepackage{graphicx} % Required for inserting images
\usepackage[a4paper, left=1.15in, right=1.15in, top=1.5in, bottom=1.2in]{geometry}
\usepackage{float}
\usepackage{array}
\usepackage{booktabs}
\usepackage[T1]{fontenc}
\usepackage{amsmath}
\usepackage{tabularx}
\usepackage{graphicx} % Required for inserting images
\usepackage{tikz, xcolor}
\usepackage{hyperref}
\usepackage{cite}
\usepackage{pdfpages}
\usepackage{pdflscape}
\usepackage{afterpage}
\usepackage{textpos}
\usepackage{fancyhdr}
\usepackage{caption}
\usepackage{lmodern,textcomp}
\usepackage[catalan]{babel}
\usepackage{lipsum}
\usepackage{newfloat}
\usepackage{listings}
\usepackage{placeins}
\usepackage[ruled]{algorithm2e}
\usepackage{tcolorbox}
% Estilo de citas IEEE sencillo: usar paquete cite para comprimir y ordenar referencias numéricas
\usepackage{cite}
\usepackage[toc,page]{appendix}
\usepackage{forest}
\usepackage{pdflscape}
\usepackage{graphicx}
\usepackage{geometry}
\usepackage{quoting,xparse}
\usepackage{svg}



% \NewDocumentCommand{\bywhom}{m}{% the Bourbaki trick
%   {\nobreak\hfill\penalty50\hskip1em\null\nobreak
%    \hfill\mbox{\normalfont(#1)}%
%    \parfillskip=0pt \finalhyphendemerits=0 \par}%
% }

% AIXÒ SERVEIX PER A MODIFICAR UN PROBLEMA AMB ELS CARACTERS --> " + PARAULA 

% \NewDocumentEnvironment{pquotation}{m}
%   {\begin{quoting}[
%      indentfirst=true,
%      leftmargin=\parindent,
%      rightmargin=\parindent]\itshape}
%   {\bywhom{#1}
  
%   \end{quoting}}

% AIXÒ PER INSERIR FRAGMENTS DE CODI

\lstnewenvironment{code}[1][]%
{
   \noindent
   \minipage{\linewidth}
   \vspace{0.5\baselineskip}
   \lstset{
     basicstyle=\ttfamily\footnotesize,
     frame=single,
     breaklines=true,    % Permitir que las líneas se rompan automáticamente
     breakatwhitespace=true, % Permitir la ruptura en los espacios en blanco
     #1
   }
}
{\endminipage}

\lstset{
    language=Java,
    basicstyle=\ttfamily\small,  % Define el tipo de letra y tamaño del código
    keywordstyle=\color{blue},   % Define el color de las palabras clave de Java
    stringstyle=\color{red},     % Define el color de las cadenas de texto
    commentstyle=\color{green},  % Define el color de los comentarios
    morecomment=[s][\color{gray}]{/**}{*/},  % Estilo para comentarios de bloque
    breaklines=true,  % Configura el quiebre automático de línea
    numbers=left,  % Posición de los números de línea
    numberstyle=\tiny\color{gray},  % Estilo para los números de línea
    showspaces=false,  % No mostrar espacios con underscores
    showstringspaces=false,  % No mostrar espacios dentro de strings
    frame=single,  % Agregar un marco alrededor del código
    captionpos=b,  % b = bottom, t = top
}

\definecolor{eclipseStrings}{RGB}{42,0.0,255}
\definecolor{eclipseKeywords}{RGB}{127,0,85}
\colorlet{numb}{magenta!60!black}

\lstdefinelanguage{json}{
    basicstyle=\ttfamily\small,
    showstringspaces=false,
    breaklines=true,
    frame=single,
    literate=
     *{0}{{{\color{numb}0}}}{1}
      {1}{{{\color{numb}1}}}{1}
      {2}{{{\color{numb}2}}}{1}
      {3}{{{\color{numb}3}}}{1}
      {4}{{{\color{numb}4}}}{1}
      {5}{{{\color{numb}5}}}{1}
      {6}{{{\color{numb}6}}}{1}
      {7}{{{\color{numb}7}}}{1}
      {8}{{{\color{numb}8}}}{1}
      {9}{{{\color{numb}9}}}{1}
      {:}{{{\color{punct}{:}}}}{1}
      {,}{{{\color{punct}{,}}}}{1}
      {\{}{{{\color{delim}{\{}}}}{1}
      {\}}{{{\color{delim}{\}}}}}{1}
      {[}{{{\color{delim}{[}}}}{1}
      {]}{{{\color{delim}{]}}}}{1},
    morestring=[b]"
}

\lstdefinelanguage{JavaScript}{
  morekeywords=[1]{break, continue, delete, else, for, function, if, in,
    new, return, this, typeof, var, void, while, with},
  % Literals, primitive types, and reference types.
  morekeywords=[2]{false, null, true, boolean, number, undefined,
    Array, Boolean, Date, Math, Number, String, Object},
  % Built-ins.
  morekeywords=[3]{eval, parseInt, parseFloat, escape, unescape},
  sensitive,
  morecomment=[s]{/*}{*/},
  morecomment=[l]//,
  morecomment=[s]{/**}{*/}, % JavaDoc style comments
  morestring=[b]',
  morestring=[b]"
}[keywords, comments, strings]

\lstdefinelanguage{svelte}{
    morekeywords={async, await, then, if, each, in, on, let, const, function, import, export, from, with, typeof, new, var, return},
    sensitive=true,
    morecomment=[l]{//},
    morecomment=[s]{/*}{*/},
    morecomment=[s]{<!--}{-->},
    morestring=[b]',
    morestring=[b]",
    morestring=[b]`,
    ndkeywords={import, export, return, if, for, while, function, switch, case, break, var, let, const, async, await, then, catch, try, typeof, new, class, extends, implements, interface, constructor, public, private, protected, static, get, set},
    alsoletter={:},
    moredelim=**[is][\color{orange}]{\@\{}{\@\}},
    moredelim=**[is][\color{cyan}]{\$\{}{\}},
    keywordstyle=\color{blue}\bfseries,
    commentstyle=\color{gray}\ttfamily,
    stringstyle=\color{red}\ttfamily,
    identifierstyle=\color{black},
    breaklines=true,
    basicstyle=\ttfamily\small,
    columns=fullflexible,
    showstringspaces=false,
    tabsize=2,
    frame=single,
    captionpos=b,
}


\lstdefinestyle{java}{
    language=Java,
    basicstyle=\ttfamily\small,
    keywordstyle=\color{blue},
    stringstyle=\color{red},
    commentstyle=\color{gray},
    morecomment=[s][\color{gray}]{/**}{*/},
    breaklines=true,
    numbers=left,
    numberstyle=\tiny\color{gray},
    showspaces=false,
    showstringspaces=false,
    frame=single,
    captionpos=b
}

\usepackage{color}
\definecolor{numb}{rgb}{0.58,0,0.82}
\definecolor{punct}{rgb}{0.3,0.3,0.3}
\definecolor{delim}{rgb}{0.5,0.5,0.5}
\definecolor{string}{rgb}{0.6,0.1,0.1}

\lstdefinestyle{json}{
    language=json,
    basicstyle=\ttfamily\small,
    keywordstyle=\color{blue},
    stringstyle=\color{string},
    commentstyle=\color{green},
    morecomment=[s][\color{gray}]{/**}{*/},
    breaklines=true,
    numbers=left,
    numberstyle=\tiny\color{gray},
    showspaces=false,
    showstringspaces=false,
    frame=single,
    captionpos=b
}

\lstdefinelanguage{drl}{
    basicstyle=\ttfamily\small,
    keywordstyle=\color{blue}\bfseries,
    stringstyle=\color{red},
    commentstyle=\color{green},
    morecomment=[s][\color{green}]{/*}{*/},
    morecomment=[l][\color{green}]{//},
    morestring=[b]",
    keywords={rule, when, then, end, salience, eval, update, retract},
    sensitive=true,
    breaklines=true,
    numbers=left,
    numberstyle=\tiny\color{gray},
    showspaces=false,
    showstringspaces=false,
    frame=single,
    captionpos=b
}

\lstdefinestyle{drl}{
    language=drl,
    basicstyle=\ttfamily\small,
    keywordstyle=\color{blue}\bfseries,
    stringstyle=\color{red},
    commentstyle=\color{green},
    breaklines=true,
    numbers=left,
    numberstyle=\tiny\color{gray},
    showspaces=false,
    showstringspaces=false,
    frame=single,
    captionpos=b
}



\renewcommand{\lstlistingname}{Codi}
\renewcommand{\lstlistlistingname}{Índex de Fragments de Codis}
\renewcommand{\appendixpagename}{Annexos}
\renewcommand{\appendixtocname}{Annexos}

\renewcommand*{\listalgorithmcfname}{Llistat de pseudocodi}
\renewcommand*{\algorithmcfname}{Pseudocodi}
\renewcommand*{\algorithmautorefname}{pseudocodi}

\definecolor{darkgreen}{rgb}{0.0, 0.3, 0.0}
\definecolor{lightgreen}{rgb}{0.9, 0.95, 0.9}

\newcommand{\todo}[1]{
    \begin{tcolorbox}[colback=lightgreen, colframe=darkgreen, title=TODO]
        #1
    \end{tcolorbox}
}


% \usepackage{indentfirst} 
% \setlength{\parindent}{1cm}


\usepackage{parskip}
% \setlength{\parindent}{10pt}

% \setlength{\headheight}{15.6pt}
\setlength{\headheight}{16.05263pt}




\vspace{5mm}

\newcolumntype{u}{>{\centering\arraybackslash\hsize=.22\hsize}X}
\newcolumntype{s}{>{\hsize=.45\hsize}X}


%%%%%%%%%%%%%%%%%%%%%%%%%%%%%%%%
% INI creacio subsubsubsecion
%%%%%%%%%%%%%%%%%%%%%%%%%%%%%%%%
\usepackage{titlesec}

\let\oldparagraph\paragraph
\renewcommand{\paragraph}[1]{\oldparagraph{\hspace{30pt}#1}}


\titleclass{\subsubsubsection}{straight}[\subsection]

\newcounter{subsubsubsection}[subsubsection]
\renewcommand\thesubsubsubsection{\thesubsubsection.\arabic{subsubsubsection}}
\renewcommand\theparagraph{\thesubsubsubsection.\arabic{paragraph}} % optional; useful if paragraphs are to be numbered

\titleformat{\subsubsubsection}
  {\normalfont\normalsize\bfseries}{\thesubsubsubsection}{1em}{}
\titlespacing*{\subsubsubsection}
{0pt}{3.25ex plus 1ex minus .2ex}{1.5ex plus .2ex}


\makeatletter
\renewcommand\paragraph{\@startsection{paragraph}{5}{10pt}%
  {3.25ex \@plus1ex \@minus.2ex}%
  {-1em}%
  {\normalfont\normalsize\bfseries}}
\renewcommand\subparagraph{\@startsection{subparagraph}{6}{\parindent}%
  {3.25ex \@plus1ex \@minus .2ex}%
  {-1em}%
  {\normalfont\normalsize\bfseries}}
\def\toclevel@subsubsubsection{4}
\def\toclevel@paragraph{5}
\def\toclevel@paragraph{6}
\def\l@subsubsubsection{\@dottedtocline{4}{7em}{4em}}
\def\l@paragraph{\@dottedtocline{5}{10em}{5em}}
\def\l@subparagraph{\@dottedtocline{6}{14em}{6em}}
\makeatother

\setcounter{secnumdepth}{4}
\setcounter{tocdepth}{3}
%%%%%%%%%%%%%%%%%%%%%%%%%%%%%%%%
% FIN creacio subsubsubtitle
%%%%%%%%%%%%%%%%%%%%%%%%%%%%%%%%


%%%%%%%%%%%%%%%%%%%%%%%%%%%%%%%%
% INI creacio equation
%%%%%%%%%%%%%%%%%%%%%%%%%%%%%%%%
\DeclareFloatingEnvironment[fileext=loeq,placement={!ht},name=Equació]{equationfloat}
%%%%%%%%%%%%%%%%%%%%%%%%%%%%%%%%
% FIN creacio equation
%%%%%%%%%%%%%%%%%%%%%%%%%%%%%%%%

\titleformat{\section}[block]
  {\clearpage\normalfont\Large\bfseries}
  {\thesection}{1em}{}


\renewcommand\refname{Bibliografia}


\linespread{1.25}


\makeatletter
\renewenvironment{abstract}{%
    \if@twocolumn
      \section*{\abstractname}%
    \else 
      \begin{center}%
        {\bfseries \Large\itshape\abstractname\vspace{\z@}}%  % Título en negrita, 14pt y cursiva
      \end{center}%
      \quotation
    \fi}
    {\if@twocolumn\else\endquotation\fi}
\makeatother

% \fancypagestyle{main}{
%     \fancyhf{}
%     \lhead{\small Aplicació Web per a la Visualització i Comparació
% de Trajectòries d’Algorismes d’Optimització}
%     % \rhead{\includegraphics[width=1.2cm]{resources/Icons/logo-fib.png}}
%     \cfoot{\thepage}
%     \setlength{\footskip}{35pt}
% }

\newcommand{\q}[1]{``#1''}

\begin{document}

\hypersetup{
    pdfborder=0 0 0,
}

%% PORTADA I TABLE OF CONTENTS %%

%%%%                                   %%%%
%         INDEX DE LA MEMORIA             %
%%%%                                   %%%%

% \includepdf[pages=-]{portada_fib.pdf}
\input{Extra/GEP}
%\input{Extra/portada}
% \cleardoublepage

% \pagestyle{main}
% \setcounter{page}{11}
\include{toc}
% \cleardoublepage
\input{Contexto y alcance/main}

\input{Planificación temporal/main}

\input{Gestión económica/main}

\input{Informe de sostenibilidad/main}

% \cleardoublepage

% \cleardoublepage


% \input{Annex/main}


\clearpage
% Estilo bibliográfico IEEE
\bibliographystyle{IEEEtran}
\renewcommand{\refname}{Referencias}
\bibliography{references.bib}

%%%%                                   %%%%
%         FINAL DE LA MEMORIA             %
%%%%                                   %%%%

\end{document}


%% VERSIÓ INCIAL MEMORIA %%



\documentclass[a4paper, 11pt, twoside]{article}
\usepackage{wrapfig}
% \usepackage{abstract}
\usepackage{graphicx} % Required for inserting images
\usepackage[a4paper, left=1.15in, right=1.15in, top=1.5in, bottom=1.2in]{geometry}
\usepackage{float}
\usepackage{array}
\usepackage{booktabs}
\usepackage[T1]{fontenc}
\usepackage{amsmath}
\usepackage{tabularx}
\usepackage{graphicx} % Required for inserting images
\usepackage{tikz, xcolor}
\usepackage{hyperref}
\usepackage{cite}
\usepackage{pdfpages}
\usepackage{pdflscape}
\usepackage{afterpage}
\usepackage{textpos}
\usepackage{fancyhdr}
\usepackage{caption}
\usepackage{lmodern,textcomp}
\usepackage[catalan]{babel}
\usepackage{lipsum}
\usepackage{newfloat}
\usepackage{listings}
\usepackage{placeins}
\usepackage[ruled]{algorithm2e}
\usepackage{tcolorbox}
% Estilo de citas IEEE sencillo: usar paquete cite para comprimir y ordenar referencias numéricas
\usepackage{cite}
\usepackage[toc,page]{appendix}
\usepackage{forest}
\usepackage{pdflscape}
\usepackage{graphicx}
\usepackage{geometry}
\usepackage{quoting,xparse}
\usepackage{svg}



% \NewDocumentCommand{\bywhom}{m}{% the Bourbaki trick
%   {\nobreak\hfill\penalty50\hskip1em\null\nobreak
%    \hfill\mbox{\normalfont(#1)}%
%    \parfillskip=0pt \finalhyphendemerits=0 \par}%
% }

% AIXÒ SERVEIX PER A MODIFICAR UN PROBLEMA AMB ELS CARACTERS --> " + PARAULA 

% \NewDocumentEnvironment{pquotation}{m}
%   {\begin{quoting}[
%      indentfirst=true,
%      leftmargin=\parindent,
%      rightmargin=\parindent]\itshape}
%   {\bywhom{#1}
  
%   \end{quoting}}

% AIXÒ PER INSERIR FRAGMENTS DE CODI

\lstnewenvironment{code}[1][]%
{
   \noindent
   \minipage{\linewidth}
   \vspace{0.5\baselineskip}
   \lstset{
     basicstyle=\ttfamily\footnotesize,
     frame=single,
     breaklines=true,    % Permitir que las líneas se rompan automáticamente
     breakatwhitespace=true, % Permitir la ruptura en los espacios en blanco
     #1
   }
}
{\endminipage}

\lstset{
    language=Java,
    basicstyle=\ttfamily\small,  % Define el tipo de letra y tamaño del código
    keywordstyle=\color{blue},   % Define el color de las palabras clave de Java
    stringstyle=\color{red},     % Define el color de las cadenas de texto
    commentstyle=\color{green},  % Define el color de los comentarios
    morecomment=[s][\color{gray}]{/**}{*/},  % Estilo para comentarios de bloque
    breaklines=true,  % Configura el quiebre automático de línea
    numbers=left,  % Posición de los números de línea
    numberstyle=\tiny\color{gray},  % Estilo para los números de línea
    showspaces=false,  % No mostrar espacios con underscores
    showstringspaces=false,  % No mostrar espacios dentro de strings
    frame=single,  % Agregar un marco alrededor del código
    captionpos=b,  % b = bottom, t = top
}

\definecolor{eclipseStrings}{RGB}{42,0.0,255}
\definecolor{eclipseKeywords}{RGB}{127,0,85}
\colorlet{numb}{magenta!60!black}

\lstdefinelanguage{json}{
    basicstyle=\ttfamily\small,
    showstringspaces=false,
    breaklines=true,
    frame=single,
    literate=
     *{0}{{{\color{numb}0}}}{1}
      {1}{{{\color{numb}1}}}{1}
      {2}{{{\color{numb}2}}}{1}
      {3}{{{\color{numb}3}}}{1}
      {4}{{{\color{numb}4}}}{1}
      {5}{{{\color{numb}5}}}{1}
      {6}{{{\color{numb}6}}}{1}
      {7}{{{\color{numb}7}}}{1}
      {8}{{{\color{numb}8}}}{1}
      {9}{{{\color{numb}9}}}{1}
      {:}{{{\color{punct}{:}}}}{1}
      {,}{{{\color{punct}{,}}}}{1}
      {\{}{{{\color{delim}{\{}}}}{1}
      {\}}{{{\color{delim}{\}}}}}{1}
      {[}{{{\color{delim}{[}}}}{1}
      {]}{{{\color{delim}{]}}}}{1},
    morestring=[b]"
}

\lstdefinelanguage{JavaScript}{
  morekeywords=[1]{break, continue, delete, else, for, function, if, in,
    new, return, this, typeof, var, void, while, with},
  % Literals, primitive types, and reference types.
  morekeywords=[2]{false, null, true, boolean, number, undefined,
    Array, Boolean, Date, Math, Number, String, Object},
  % Built-ins.
  morekeywords=[3]{eval, parseInt, parseFloat, escape, unescape},
  sensitive,
  morecomment=[s]{/*}{*/},
  morecomment=[l]//,
  morecomment=[s]{/**}{*/}, % JavaDoc style comments
  morestring=[b]',
  morestring=[b]"
}[keywords, comments, strings]

\lstdefinelanguage{svelte}{
    morekeywords={async, await, then, if, each, in, on, let, const, function, import, export, from, with, typeof, new, var, return},
    sensitive=true,
    morecomment=[l]{//},
    morecomment=[s]{/*}{*/},
    morecomment=[s]{<!--}{-->},
    morestring=[b]',
    morestring=[b]",
    morestring=[b]`,
    ndkeywords={import, export, return, if, for, while, function, switch, case, break, var, let, const, async, await, then, catch, try, typeof, new, class, extends, implements, interface, constructor, public, private, protected, static, get, set},
    alsoletter={:},
    moredelim=**[is][\color{orange}]{\@\{}{\@\}},
    moredelim=**[is][\color{cyan}]{\$\{}{\}},
    keywordstyle=\color{blue}\bfseries,
    commentstyle=\color{gray}\ttfamily,
    stringstyle=\color{red}\ttfamily,
    identifierstyle=\color{black},
    breaklines=true,
    basicstyle=\ttfamily\small,
    columns=fullflexible,
    showstringspaces=false,
    tabsize=2,
    frame=single,
    captionpos=b,
}


\lstdefinestyle{java}{
    language=Java,
    basicstyle=\ttfamily\small,
    keywordstyle=\color{blue},
    stringstyle=\color{red},
    commentstyle=\color{gray},
    morecomment=[s][\color{gray}]{/**}{*/},
    breaklines=true,
    numbers=left,
    numberstyle=\tiny\color{gray},
    showspaces=false,
    showstringspaces=false,
    frame=single,
    captionpos=b
}

\usepackage{color}
\definecolor{numb}{rgb}{0.58,0,0.82}
\definecolor{punct}{rgb}{0.3,0.3,0.3}
\definecolor{delim}{rgb}{0.5,0.5,0.5}
\definecolor{string}{rgb}{0.6,0.1,0.1}

\lstdefinestyle{json}{
    language=json,
    basicstyle=\ttfamily\small,
    keywordstyle=\color{blue},
    stringstyle=\color{string},
    commentstyle=\color{green},
    morecomment=[s][\color{gray}]{/**}{*/},
    breaklines=true,
    numbers=left,
    numberstyle=\tiny\color{gray},
    showspaces=false,
    showstringspaces=false,
    frame=single,
    captionpos=b
}

\lstdefinelanguage{drl}{
    basicstyle=\ttfamily\small,
    keywordstyle=\color{blue}\bfseries,
    stringstyle=\color{red},
    commentstyle=\color{green},
    morecomment=[s][\color{green}]{/*}{*/},
    morecomment=[l][\color{green}]{//},
    morestring=[b]",
    keywords={rule, when, then, end, salience, eval, update, retract},
    sensitive=true,
    breaklines=true,
    numbers=left,
    numberstyle=\tiny\color{gray},
    showspaces=false,
    showstringspaces=false,
    frame=single,
    captionpos=b
}

\lstdefinestyle{drl}{
    language=drl,
    basicstyle=\ttfamily\small,
    keywordstyle=\color{blue}\bfseries,
    stringstyle=\color{red},
    commentstyle=\color{green},
    breaklines=true,
    numbers=left,
    numberstyle=\tiny\color{gray},
    showspaces=false,
    showstringspaces=false,
    frame=single,
    captionpos=b
}



\renewcommand{\lstlistingname}{Codi}
\renewcommand{\lstlistlistingname}{Índex de Fragments de Codis}
\renewcommand{\appendixpagename}{Annexos}
\renewcommand{\appendixtocname}{Annexos}

\renewcommand*{\listalgorithmcfname}{Llistat de pseudocodi}
\renewcommand*{\algorithmcfname}{Pseudocodi}
\renewcommand*{\algorithmautorefname}{pseudocodi}

\definecolor{darkgreen}{rgb}{0.0, 0.3, 0.0}
\definecolor{lightgreen}{rgb}{0.9, 0.95, 0.9}

\newcommand{\todo}[1]{
    \begin{tcolorbox}[colback=lightgreen, colframe=darkgreen, title=TODO]
        #1
    \end{tcolorbox}
}


% \usepackage{indentfirst} 
% \setlength{\parindent}{1cm}


\usepackage{parskip}
% \setlength{\parindent}{10pt}

% \setlength{\headheight}{15.6pt}
\setlength{\headheight}{16.05263pt}




\vspace{5mm}

\newcolumntype{u}{>{\centering\arraybackslash\hsize=.22\hsize}X}
\newcolumntype{s}{>{\hsize=.45\hsize}X}


%%%%%%%%%%%%%%%%%%%%%%%%%%%%%%%%
% INI creacio subsubsubsecion
%%%%%%%%%%%%%%%%%%%%%%%%%%%%%%%%
\usepackage{titlesec}

\let\oldparagraph\paragraph
\renewcommand{\paragraph}[1]{\oldparagraph{\hspace{30pt}#1}}


\titleclass{\subsubsubsection}{straight}[\subsection]

\newcounter{subsubsubsection}[subsubsection]
\renewcommand\thesubsubsubsection{\thesubsubsection.\arabic{subsubsubsection}}
\renewcommand\theparagraph{\thesubsubsubsection.\arabic{paragraph}} % optional; useful if paragraphs are to be numbered

\titleformat{\subsubsubsection}
  {\normalfont\normalsize\bfseries}{\thesubsubsubsection}{1em}{}
\titlespacing*{\subsubsubsection}
{0pt}{3.25ex plus 1ex minus .2ex}{1.5ex plus .2ex}


\makeatletter
\renewcommand\paragraph{\@startsection{paragraph}{5}{10pt}%
  {3.25ex \@plus1ex \@minus.2ex}%
  {-1em}%
  {\normalfont\normalsize\bfseries}}
\renewcommand\subparagraph{\@startsection{subparagraph}{6}{\parindent}%
  {3.25ex \@plus1ex \@minus .2ex}%
  {-1em}%
  {\normalfont\normalsize\bfseries}}
\def\toclevel@subsubsubsection{4}
\def\toclevel@paragraph{5}
\def\toclevel@paragraph{6}
\def\l@subsubsubsection{\@dottedtocline{4}{7em}{4em}}
\def\l@paragraph{\@dottedtocline{5}{10em}{5em}}
\def\l@subparagraph{\@dottedtocline{6}{14em}{6em}}
\makeatother

\setcounter{secnumdepth}{4}
\setcounter{tocdepth}{3}
%%%%%%%%%%%%%%%%%%%%%%%%%%%%%%%%
% FIN creacio subsubsubtitle
%%%%%%%%%%%%%%%%%%%%%%%%%%%%%%%%


%%%%%%%%%%%%%%%%%%%%%%%%%%%%%%%%
% INI creacio equation
%%%%%%%%%%%%%%%%%%%%%%%%%%%%%%%%
\DeclareFloatingEnvironment[fileext=loeq,placement={!ht},name=Equació]{equationfloat}
%%%%%%%%%%%%%%%%%%%%%%%%%%%%%%%%
% FIN creacio equation
%%%%%%%%%%%%%%%%%%%%%%%%%%%%%%%%

\titleformat{\section}[block]
  {\clearpage\normalfont\Large\bfseries}
  {\thesection}{1em}{}


\renewcommand\refname{Bibliografia}


\linespread{1.25}


\makeatletter
\renewenvironment{abstract}{%
    \if@twocolumn
      \section*{\abstractname}%
    \else 
      \begin{center}%
        {\bfseries \Large\itshape\abstractname\vspace{\z@}}%  % Título en negrita, 14pt y cursiva
      \end{center}%
      \quotation
    \fi}
    {\if@twocolumn\else\endquotation\fi}
\makeatother

% \fancypagestyle{main}{
%     \fancyhf{}
%     \lhead{\small Aplicació Web per a la Visualització i Comparació
% de Trajectòries d’Algorismes d’Optimització}
%     % \rhead{\includegraphics[width=1.2cm]{resources/Icons/logo-fib.png}}
%     \cfoot{\thepage}
%     \setlength{\footskip}{35pt}
% }

\newcommand{\q}[1]{``#1''}

\begin{document}

\hypersetup{
    pdfborder=0 0 0,
}

%% PORTADA I TABLE OF CONTENTS %%

%%%%                                   %%%%
%         INDEX DE LA MEMORIA             %
%%%%                                   %%%%

% \includepdf[pages=-]{portada_fib.pdf}
%%% Portada
\begin{titlepage}
\begin{center}
    \includegraphics[width=1\linewidth]{extra/figures/logo-upc-fib.png}
    \vspace{30mm}

    \LARGE{\textsc{\textbf{{INTEGRACIÓN Y AUTOMATIZACIÓN DE COSTES DE INFRAESTRUCTURA HPC EN ENTORNOS DE INVESTIGACIÓN}}}} \\
    \vspace{25mm}
      
        
    \textbf{\Large{\textsc{PAU NAVARRO ÁLVAREZ}}} \\
    \vspace{20mm}
        
    
    
    \large{\textbf{Director/a} ROBERTO RIVEIRO INSUA (FUNDACIO INSTITUT DE RECERCA BIOMEDICA (IRB BARCELONA))} \\
    \large{\textbf{Ponente:} ALEX PAJUELO GONZÁLEZ (Departamento de Arquitectura de Computadores)} \\
    \large{\textbf{Titulación:} Grado en Ingeniería Informática (Tecnologías de la información)} \\
    
    \large{\textbf{Lliurament 2:} }Planificación temporal \\
    \large{\textbf{Tutor GEP:} CAROLINA CONSOLIDACIÓN SEGURA (Departament d'Organització d'Empreses)} \\
    \large{\textbf{Fecha:}} 29 de Septiembre de 2025\\

    \textbf{\large{Memoria del trabajo de fin de grado}} \\
    \vspace{3mm}
    \textbf{\large{Facultat d'Informàtica de Barcelona (FIB)}} \\
    \vspace{3mm}
    \textbf{\large{Universitat Politècnica de Catalunya (UPC) - BarcelonaTech}} \\

\end{center}
\end{titlepage}
%\input{Extra/portada}
% \cleardoublepage

% \pagestyle{main}
% \setcounter{page}{11}
\renewcommand{\contentsname}{Continguts}

\tableofcontents
\clearpage
\listoffigures
\begingroup
    \let\clearpage\relax
    \listoftables
\endgroup
\clearpage

\newpage
% \cleardoublepage
\input{Contexto y alcance/Contextualización/main}
\input{Contexto y alcance/Justificación/main}
\input{Contexto y alcance/Alcance/main}
\input{Contexto y alcance/Metodología y herramientas de seguimiento/main}

\documentclass[a4paper, 11pt, twoside]{article}
\usepackage{wrapfig}
% \usepackage{abstract}
\usepackage{graphicx} % Required for inserting images
\usepackage[a4paper, left=1.15in, right=1.15in, top=1.5in, bottom=1.2in]{geometry}
\usepackage{float}
\usepackage{array}
\usepackage{booktabs}
\usepackage[T1]{fontenc}
\usepackage{amsmath}
\usepackage{tabularx}
\usepackage{graphicx} % Required for inserting images
\usepackage{tikz, xcolor}
\usepackage{hyperref}
\usepackage{cite}
\usepackage{pdfpages}
\usepackage{pdflscape}
\usepackage{afterpage}
\usepackage{textpos}
\usepackage{fancyhdr}
\usepackage{caption}
\usepackage{lmodern,textcomp}
\usepackage[catalan]{babel}
\usepackage{lipsum}
\usepackage{newfloat}
\usepackage{listings}
\usepackage{placeins}
\usepackage[ruled]{algorithm2e}
\usepackage{tcolorbox}
% Estilo de citas IEEE sencillo: usar paquete cite para comprimir y ordenar referencias numéricas
\usepackage{cite}
\usepackage[toc,page]{appendix}
\usepackage{forest}
\usepackage{pdflscape}
\usepackage{graphicx}
\usepackage{geometry}
\usepackage{quoting,xparse}
\usepackage{svg}



% \NewDocumentCommand{\bywhom}{m}{% the Bourbaki trick
%   {\nobreak\hfill\penalty50\hskip1em\null\nobreak
%    \hfill\mbox{\normalfont(#1)}%
%    \parfillskip=0pt \finalhyphendemerits=0 \par}%
% }

% AIXÒ SERVEIX PER A MODIFICAR UN PROBLEMA AMB ELS CARACTERS --> " + PARAULA 

% \NewDocumentEnvironment{pquotation}{m}
%   {\begin{quoting}[
%      indentfirst=true,
%      leftmargin=\parindent,
%      rightmargin=\parindent]\itshape}
%   {\bywhom{#1}
  
%   \end{quoting}}

% AIXÒ PER INSERIR FRAGMENTS DE CODI

\lstnewenvironment{code}[1][]%
{
   \noindent
   \minipage{\linewidth}
   \vspace{0.5\baselineskip}
   \lstset{
     basicstyle=\ttfamily\footnotesize,
     frame=single,
     breaklines=true,    % Permitir que las líneas se rompan automáticamente
     breakatwhitespace=true, % Permitir la ruptura en los espacios en blanco
     #1
   }
}
{\endminipage}

\lstset{
    language=Java,
    basicstyle=\ttfamily\small,  % Define el tipo de letra y tamaño del código
    keywordstyle=\color{blue},   % Define el color de las palabras clave de Java
    stringstyle=\color{red},     % Define el color de las cadenas de texto
    commentstyle=\color{green},  % Define el color de los comentarios
    morecomment=[s][\color{gray}]{/**}{*/},  % Estilo para comentarios de bloque
    breaklines=true,  % Configura el quiebre automático de línea
    numbers=left,  % Posición de los números de línea
    numberstyle=\tiny\color{gray},  % Estilo para los números de línea
    showspaces=false,  % No mostrar espacios con underscores
    showstringspaces=false,  % No mostrar espacios dentro de strings
    frame=single,  % Agregar un marco alrededor del código
    captionpos=b,  % b = bottom, t = top
}

\definecolor{eclipseStrings}{RGB}{42,0.0,255}
\definecolor{eclipseKeywords}{RGB}{127,0,85}
\colorlet{numb}{magenta!60!black}

\lstdefinelanguage{json}{
    basicstyle=\ttfamily\small,
    showstringspaces=false,
    breaklines=true,
    frame=single,
    literate=
     *{0}{{{\color{numb}0}}}{1}
      {1}{{{\color{numb}1}}}{1}
      {2}{{{\color{numb}2}}}{1}
      {3}{{{\color{numb}3}}}{1}
      {4}{{{\color{numb}4}}}{1}
      {5}{{{\color{numb}5}}}{1}
      {6}{{{\color{numb}6}}}{1}
      {7}{{{\color{numb}7}}}{1}
      {8}{{{\color{numb}8}}}{1}
      {9}{{{\color{numb}9}}}{1}
      {:}{{{\color{punct}{:}}}}{1}
      {,}{{{\color{punct}{,}}}}{1}
      {\{}{{{\color{delim}{\{}}}}{1}
      {\}}{{{\color{delim}{\}}}}}{1}
      {[}{{{\color{delim}{[}}}}{1}
      {]}{{{\color{delim}{]}}}}{1},
    morestring=[b]"
}

\lstdefinelanguage{JavaScript}{
  morekeywords=[1]{break, continue, delete, else, for, function, if, in,
    new, return, this, typeof, var, void, while, with},
  % Literals, primitive types, and reference types.
  morekeywords=[2]{false, null, true, boolean, number, undefined,
    Array, Boolean, Date, Math, Number, String, Object},
  % Built-ins.
  morekeywords=[3]{eval, parseInt, parseFloat, escape, unescape},
  sensitive,
  morecomment=[s]{/*}{*/},
  morecomment=[l]//,
  morecomment=[s]{/**}{*/}, % JavaDoc style comments
  morestring=[b]',
  morestring=[b]"
}[keywords, comments, strings]

\lstdefinelanguage{svelte}{
    morekeywords={async, await, then, if, each, in, on, let, const, function, import, export, from, with, typeof, new, var, return},
    sensitive=true,
    morecomment=[l]{//},
    morecomment=[s]{/*}{*/},
    morecomment=[s]{<!--}{-->},
    morestring=[b]',
    morestring=[b]",
    morestring=[b]`,
    ndkeywords={import, export, return, if, for, while, function, switch, case, break, var, let, const, async, await, then, catch, try, typeof, new, class, extends, implements, interface, constructor, public, private, protected, static, get, set},
    alsoletter={:},
    moredelim=**[is][\color{orange}]{\@\{}{\@\}},
    moredelim=**[is][\color{cyan}]{\$\{}{\}},
    keywordstyle=\color{blue}\bfseries,
    commentstyle=\color{gray}\ttfamily,
    stringstyle=\color{red}\ttfamily,
    identifierstyle=\color{black},
    breaklines=true,
    basicstyle=\ttfamily\small,
    columns=fullflexible,
    showstringspaces=false,
    tabsize=2,
    frame=single,
    captionpos=b,
}


\lstdefinestyle{java}{
    language=Java,
    basicstyle=\ttfamily\small,
    keywordstyle=\color{blue},
    stringstyle=\color{red},
    commentstyle=\color{gray},
    morecomment=[s][\color{gray}]{/**}{*/},
    breaklines=true,
    numbers=left,
    numberstyle=\tiny\color{gray},
    showspaces=false,
    showstringspaces=false,
    frame=single,
    captionpos=b
}

\usepackage{color}
\definecolor{numb}{rgb}{0.58,0,0.82}
\definecolor{punct}{rgb}{0.3,0.3,0.3}
\definecolor{delim}{rgb}{0.5,0.5,0.5}
\definecolor{string}{rgb}{0.6,0.1,0.1}

\lstdefinestyle{json}{
    language=json,
    basicstyle=\ttfamily\small,
    keywordstyle=\color{blue},
    stringstyle=\color{string},
    commentstyle=\color{green},
    morecomment=[s][\color{gray}]{/**}{*/},
    breaklines=true,
    numbers=left,
    numberstyle=\tiny\color{gray},
    showspaces=false,
    showstringspaces=false,
    frame=single,
    captionpos=b
}

\lstdefinelanguage{drl}{
    basicstyle=\ttfamily\small,
    keywordstyle=\color{blue}\bfseries,
    stringstyle=\color{red},
    commentstyle=\color{green},
    morecomment=[s][\color{green}]{/*}{*/},
    morecomment=[l][\color{green}]{//},
    morestring=[b]",
    keywords={rule, when, then, end, salience, eval, update, retract},
    sensitive=true,
    breaklines=true,
    numbers=left,
    numberstyle=\tiny\color{gray},
    showspaces=false,
    showstringspaces=false,
    frame=single,
    captionpos=b
}

\lstdefinestyle{drl}{
    language=drl,
    basicstyle=\ttfamily\small,
    keywordstyle=\color{blue}\bfseries,
    stringstyle=\color{red},
    commentstyle=\color{green},
    breaklines=true,
    numbers=left,
    numberstyle=\tiny\color{gray},
    showspaces=false,
    showstringspaces=false,
    frame=single,
    captionpos=b
}



\renewcommand{\lstlistingname}{Codi}
\renewcommand{\lstlistlistingname}{Índex de Fragments de Codis}
\renewcommand{\appendixpagename}{Annexos}
\renewcommand{\appendixtocname}{Annexos}

\renewcommand*{\listalgorithmcfname}{Llistat de pseudocodi}
\renewcommand*{\algorithmcfname}{Pseudocodi}
\renewcommand*{\algorithmautorefname}{pseudocodi}

\definecolor{darkgreen}{rgb}{0.0, 0.3, 0.0}
\definecolor{lightgreen}{rgb}{0.9, 0.95, 0.9}

\newcommand{\todo}[1]{
    \begin{tcolorbox}[colback=lightgreen, colframe=darkgreen, title=TODO]
        #1
    \end{tcolorbox}
}


% \usepackage{indentfirst} 
% \setlength{\parindent}{1cm}


\usepackage{parskip}
% \setlength{\parindent}{10pt}

% \setlength{\headheight}{15.6pt}
\setlength{\headheight}{16.05263pt}




\vspace{5mm}

\newcolumntype{u}{>{\centering\arraybackslash\hsize=.22\hsize}X}
\newcolumntype{s}{>{\hsize=.45\hsize}X}


%%%%%%%%%%%%%%%%%%%%%%%%%%%%%%%%
% INI creacio subsubsubsecion
%%%%%%%%%%%%%%%%%%%%%%%%%%%%%%%%
\usepackage{titlesec}

\let\oldparagraph\paragraph
\renewcommand{\paragraph}[1]{\oldparagraph{\hspace{30pt}#1}}


\titleclass{\subsubsubsection}{straight}[\subsection]

\newcounter{subsubsubsection}[subsubsection]
\renewcommand\thesubsubsubsection{\thesubsubsection.\arabic{subsubsubsection}}
\renewcommand\theparagraph{\thesubsubsubsection.\arabic{paragraph}} % optional; useful if paragraphs are to be numbered

\titleformat{\subsubsubsection}
  {\normalfont\normalsize\bfseries}{\thesubsubsubsection}{1em}{}
\titlespacing*{\subsubsubsection}
{0pt}{3.25ex plus 1ex minus .2ex}{1.5ex plus .2ex}


\makeatletter
\renewcommand\paragraph{\@startsection{paragraph}{5}{10pt}%
  {3.25ex \@plus1ex \@minus.2ex}%
  {-1em}%
  {\normalfont\normalsize\bfseries}}
\renewcommand\subparagraph{\@startsection{subparagraph}{6}{\parindent}%
  {3.25ex \@plus1ex \@minus .2ex}%
  {-1em}%
  {\normalfont\normalsize\bfseries}}
\def\toclevel@subsubsubsection{4}
\def\toclevel@paragraph{5}
\def\toclevel@paragraph{6}
\def\l@subsubsubsection{\@dottedtocline{4}{7em}{4em}}
\def\l@paragraph{\@dottedtocline{5}{10em}{5em}}
\def\l@subparagraph{\@dottedtocline{6}{14em}{6em}}
\makeatother

\setcounter{secnumdepth}{4}
\setcounter{tocdepth}{3}
%%%%%%%%%%%%%%%%%%%%%%%%%%%%%%%%
% FIN creacio subsubsubtitle
%%%%%%%%%%%%%%%%%%%%%%%%%%%%%%%%


%%%%%%%%%%%%%%%%%%%%%%%%%%%%%%%%
% INI creacio equation
%%%%%%%%%%%%%%%%%%%%%%%%%%%%%%%%
\DeclareFloatingEnvironment[fileext=loeq,placement={!ht},name=Equació]{equationfloat}
%%%%%%%%%%%%%%%%%%%%%%%%%%%%%%%%
% FIN creacio equation
%%%%%%%%%%%%%%%%%%%%%%%%%%%%%%%%

\titleformat{\section}[block]
  {\clearpage\normalfont\Large\bfseries}
  {\thesection}{1em}{}


\renewcommand\refname{Bibliografia}


\linespread{1.25}


\makeatletter
\renewenvironment{abstract}{%
    \if@twocolumn
      \section*{\abstractname}%
    \else 
      \begin{center}%
        {\bfseries \Large\itshape\abstractname\vspace{\z@}}%  % Título en negrita, 14pt y cursiva
      \end{center}%
      \quotation
    \fi}
    {\if@twocolumn\else\endquotation\fi}
\makeatother

% \fancypagestyle{main}{
%     \fancyhf{}
%     \lhead{\small Aplicació Web per a la Visualització i Comparació
% de Trajectòries d’Algorismes d’Optimització}
%     % \rhead{\includegraphics[width=1.2cm]{resources/Icons/logo-fib.png}}
%     \cfoot{\thepage}
%     \setlength{\footskip}{35pt}
% }

\newcommand{\q}[1]{``#1''}

\begin{document}

\hypersetup{
    pdfborder=0 0 0,
}

%% PORTADA I TABLE OF CONTENTS %%

%%%%                                   %%%%
%         INDEX DE LA MEMORIA             %
%%%%                                   %%%%

% \includepdf[pages=-]{portada_fib.pdf}
\input{Extra/GEP}
%\input{Extra/portada}
% \cleardoublepage

% \pagestyle{main}
% \setcounter{page}{11}
\include{toc}
% \cleardoublepage
\input{Contexto y alcance/main}

\input{Planificación temporal/main}

\input{Gestión económica/main}

\input{Informe de sostenibilidad/main}

% \cleardoublepage

% \cleardoublepage


% \input{Annex/main}


\clearpage
% Estilo bibliográfico IEEE
\bibliographystyle{IEEEtran}
\renewcommand{\refname}{Referencias}
\bibliography{references.bib}

%%%%                                   %%%%
%         FINAL DE LA MEMORIA             %
%%%%                                   %%%%

\end{document}


%% VERSIÓ INCIAL MEMORIA %%



\section{Gestión económica}

La gestión económica de un proyecto es fundamental para asegurar su viabilidad y exito a largo plazo. Ya hemos visto en secciones anteriores, los otros dos vertices 
del triángulo de la gestión de proyectos: el alcance y el tiempo. En esta sección, nos centraremos en cómo gestionar eficazmente los recursos financieros para garantizar
que el proyecto se mantenga dentro del presupuesto y se logren los objetivos establecidos. 

\input{Gestión económica/Recursos humanos.tex}

\input{Gestión económica/Recursos Hardware.tex}

\input{Gestión económica/Recursos Software.tex}

\input{Gestión económica/Contingencias.tex}

\section{Informe de sostenibilidad}

\input{Informe de sostenibilidad/Autoevaluación.tex}


\input{Informe de sostenibilidad/Dimensión ambiental.tex}

\input{Informe de sostenibilidad/Dimensión económica.tex}

\input{Informe de sostenibilidad/Dimensión social.tex}


% \cleardoublepage

% \cleardoublepage


% \documentclass[a4paper, 11pt, twoside]{article}
\usepackage{wrapfig}
% \usepackage{abstract}
\usepackage{graphicx} % Required for inserting images
\usepackage[a4paper, left=1.15in, right=1.15in, top=1.5in, bottom=1.2in]{geometry}
\usepackage{float}
\usepackage{array}
\usepackage{booktabs}
\usepackage[T1]{fontenc}
\usepackage{amsmath}
\usepackage{tabularx}
\usepackage{graphicx} % Required for inserting images
\usepackage{tikz, xcolor}
\usepackage{hyperref}
\usepackage{cite}
\usepackage{pdfpages}
\usepackage{pdflscape}
\usepackage{afterpage}
\usepackage{textpos}
\usepackage{fancyhdr}
\usepackage{caption}
\usepackage{lmodern,textcomp}
\usepackage[catalan]{babel}
\usepackage{lipsum}
\usepackage{newfloat}
\usepackage{listings}
\usepackage{placeins}
\usepackage[ruled]{algorithm2e}
\usepackage{tcolorbox}
% Estilo de citas IEEE sencillo: usar paquete cite para comprimir y ordenar referencias numéricas
\usepackage{cite}
\usepackage[toc,page]{appendix}
\usepackage{forest}
\usepackage{pdflscape}
\usepackage{graphicx}
\usepackage{geometry}
\usepackage{quoting,xparse}
\usepackage{svg}



% \NewDocumentCommand{\bywhom}{m}{% the Bourbaki trick
%   {\nobreak\hfill\penalty50\hskip1em\null\nobreak
%    \hfill\mbox{\normalfont(#1)}%
%    \parfillskip=0pt \finalhyphendemerits=0 \par}%
% }

% AIXÒ SERVEIX PER A MODIFICAR UN PROBLEMA AMB ELS CARACTERS --> " + PARAULA 

% \NewDocumentEnvironment{pquotation}{m}
%   {\begin{quoting}[
%      indentfirst=true,
%      leftmargin=\parindent,
%      rightmargin=\parindent]\itshape}
%   {\bywhom{#1}
  
%   \end{quoting}}

% AIXÒ PER INSERIR FRAGMENTS DE CODI

\lstnewenvironment{code}[1][]%
{
   \noindent
   \minipage{\linewidth}
   \vspace{0.5\baselineskip}
   \lstset{
     basicstyle=\ttfamily\footnotesize,
     frame=single,
     breaklines=true,    % Permitir que las líneas se rompan automáticamente
     breakatwhitespace=true, % Permitir la ruptura en los espacios en blanco
     #1
   }
}
{\endminipage}

\lstset{
    language=Java,
    basicstyle=\ttfamily\small,  % Define el tipo de letra y tamaño del código
    keywordstyle=\color{blue},   % Define el color de las palabras clave de Java
    stringstyle=\color{red},     % Define el color de las cadenas de texto
    commentstyle=\color{green},  % Define el color de los comentarios
    morecomment=[s][\color{gray}]{/**}{*/},  % Estilo para comentarios de bloque
    breaklines=true,  % Configura el quiebre automático de línea
    numbers=left,  % Posición de los números de línea
    numberstyle=\tiny\color{gray},  % Estilo para los números de línea
    showspaces=false,  % No mostrar espacios con underscores
    showstringspaces=false,  % No mostrar espacios dentro de strings
    frame=single,  % Agregar un marco alrededor del código
    captionpos=b,  % b = bottom, t = top
}

\definecolor{eclipseStrings}{RGB}{42,0.0,255}
\definecolor{eclipseKeywords}{RGB}{127,0,85}
\colorlet{numb}{magenta!60!black}

\lstdefinelanguage{json}{
    basicstyle=\ttfamily\small,
    showstringspaces=false,
    breaklines=true,
    frame=single,
    literate=
     *{0}{{{\color{numb}0}}}{1}
      {1}{{{\color{numb}1}}}{1}
      {2}{{{\color{numb}2}}}{1}
      {3}{{{\color{numb}3}}}{1}
      {4}{{{\color{numb}4}}}{1}
      {5}{{{\color{numb}5}}}{1}
      {6}{{{\color{numb}6}}}{1}
      {7}{{{\color{numb}7}}}{1}
      {8}{{{\color{numb}8}}}{1}
      {9}{{{\color{numb}9}}}{1}
      {:}{{{\color{punct}{:}}}}{1}
      {,}{{{\color{punct}{,}}}}{1}
      {\{}{{{\color{delim}{\{}}}}{1}
      {\}}{{{\color{delim}{\}}}}}{1}
      {[}{{{\color{delim}{[}}}}{1}
      {]}{{{\color{delim}{]}}}}{1},
    morestring=[b]"
}

\lstdefinelanguage{JavaScript}{
  morekeywords=[1]{break, continue, delete, else, for, function, if, in,
    new, return, this, typeof, var, void, while, with},
  % Literals, primitive types, and reference types.
  morekeywords=[2]{false, null, true, boolean, number, undefined,
    Array, Boolean, Date, Math, Number, String, Object},
  % Built-ins.
  morekeywords=[3]{eval, parseInt, parseFloat, escape, unescape},
  sensitive,
  morecomment=[s]{/*}{*/},
  morecomment=[l]//,
  morecomment=[s]{/**}{*/}, % JavaDoc style comments
  morestring=[b]',
  morestring=[b]"
}[keywords, comments, strings]

\lstdefinelanguage{svelte}{
    morekeywords={async, await, then, if, each, in, on, let, const, function, import, export, from, with, typeof, new, var, return},
    sensitive=true,
    morecomment=[l]{//},
    morecomment=[s]{/*}{*/},
    morecomment=[s]{<!--}{-->},
    morestring=[b]',
    morestring=[b]",
    morestring=[b]`,
    ndkeywords={import, export, return, if, for, while, function, switch, case, break, var, let, const, async, await, then, catch, try, typeof, new, class, extends, implements, interface, constructor, public, private, protected, static, get, set},
    alsoletter={:},
    moredelim=**[is][\color{orange}]{\@\{}{\@\}},
    moredelim=**[is][\color{cyan}]{\$\{}{\}},
    keywordstyle=\color{blue}\bfseries,
    commentstyle=\color{gray}\ttfamily,
    stringstyle=\color{red}\ttfamily,
    identifierstyle=\color{black},
    breaklines=true,
    basicstyle=\ttfamily\small,
    columns=fullflexible,
    showstringspaces=false,
    tabsize=2,
    frame=single,
    captionpos=b,
}


\lstdefinestyle{java}{
    language=Java,
    basicstyle=\ttfamily\small,
    keywordstyle=\color{blue},
    stringstyle=\color{red},
    commentstyle=\color{gray},
    morecomment=[s][\color{gray}]{/**}{*/},
    breaklines=true,
    numbers=left,
    numberstyle=\tiny\color{gray},
    showspaces=false,
    showstringspaces=false,
    frame=single,
    captionpos=b
}

\usepackage{color}
\definecolor{numb}{rgb}{0.58,0,0.82}
\definecolor{punct}{rgb}{0.3,0.3,0.3}
\definecolor{delim}{rgb}{0.5,0.5,0.5}
\definecolor{string}{rgb}{0.6,0.1,0.1}

\lstdefinestyle{json}{
    language=json,
    basicstyle=\ttfamily\small,
    keywordstyle=\color{blue},
    stringstyle=\color{string},
    commentstyle=\color{green},
    morecomment=[s][\color{gray}]{/**}{*/},
    breaklines=true,
    numbers=left,
    numberstyle=\tiny\color{gray},
    showspaces=false,
    showstringspaces=false,
    frame=single,
    captionpos=b
}

\lstdefinelanguage{drl}{
    basicstyle=\ttfamily\small,
    keywordstyle=\color{blue}\bfseries,
    stringstyle=\color{red},
    commentstyle=\color{green},
    morecomment=[s][\color{green}]{/*}{*/},
    morecomment=[l][\color{green}]{//},
    morestring=[b]",
    keywords={rule, when, then, end, salience, eval, update, retract},
    sensitive=true,
    breaklines=true,
    numbers=left,
    numberstyle=\tiny\color{gray},
    showspaces=false,
    showstringspaces=false,
    frame=single,
    captionpos=b
}

\lstdefinestyle{drl}{
    language=drl,
    basicstyle=\ttfamily\small,
    keywordstyle=\color{blue}\bfseries,
    stringstyle=\color{red},
    commentstyle=\color{green},
    breaklines=true,
    numbers=left,
    numberstyle=\tiny\color{gray},
    showspaces=false,
    showstringspaces=false,
    frame=single,
    captionpos=b
}



\renewcommand{\lstlistingname}{Codi}
\renewcommand{\lstlistlistingname}{Índex de Fragments de Codis}
\renewcommand{\appendixpagename}{Annexos}
\renewcommand{\appendixtocname}{Annexos}

\renewcommand*{\listalgorithmcfname}{Llistat de pseudocodi}
\renewcommand*{\algorithmcfname}{Pseudocodi}
\renewcommand*{\algorithmautorefname}{pseudocodi}

\definecolor{darkgreen}{rgb}{0.0, 0.3, 0.0}
\definecolor{lightgreen}{rgb}{0.9, 0.95, 0.9}

\newcommand{\todo}[1]{
    \begin{tcolorbox}[colback=lightgreen, colframe=darkgreen, title=TODO]
        #1
    \end{tcolorbox}
}


% \usepackage{indentfirst} 
% \setlength{\parindent}{1cm}


\usepackage{parskip}
% \setlength{\parindent}{10pt}

% \setlength{\headheight}{15.6pt}
\setlength{\headheight}{16.05263pt}




\vspace{5mm}

\newcolumntype{u}{>{\centering\arraybackslash\hsize=.22\hsize}X}
\newcolumntype{s}{>{\hsize=.45\hsize}X}


%%%%%%%%%%%%%%%%%%%%%%%%%%%%%%%%
% INI creacio subsubsubsecion
%%%%%%%%%%%%%%%%%%%%%%%%%%%%%%%%
\usepackage{titlesec}

\let\oldparagraph\paragraph
\renewcommand{\paragraph}[1]{\oldparagraph{\hspace{30pt}#1}}


\titleclass{\subsubsubsection}{straight}[\subsection]

\newcounter{subsubsubsection}[subsubsection]
\renewcommand\thesubsubsubsection{\thesubsubsection.\arabic{subsubsubsection}}
\renewcommand\theparagraph{\thesubsubsubsection.\arabic{paragraph}} % optional; useful if paragraphs are to be numbered

\titleformat{\subsubsubsection}
  {\normalfont\normalsize\bfseries}{\thesubsubsubsection}{1em}{}
\titlespacing*{\subsubsubsection}
{0pt}{3.25ex plus 1ex minus .2ex}{1.5ex plus .2ex}


\makeatletter
\renewcommand\paragraph{\@startsection{paragraph}{5}{10pt}%
  {3.25ex \@plus1ex \@minus.2ex}%
  {-1em}%
  {\normalfont\normalsize\bfseries}}
\renewcommand\subparagraph{\@startsection{subparagraph}{6}{\parindent}%
  {3.25ex \@plus1ex \@minus .2ex}%
  {-1em}%
  {\normalfont\normalsize\bfseries}}
\def\toclevel@subsubsubsection{4}
\def\toclevel@paragraph{5}
\def\toclevel@paragraph{6}
\def\l@subsubsubsection{\@dottedtocline{4}{7em}{4em}}
\def\l@paragraph{\@dottedtocline{5}{10em}{5em}}
\def\l@subparagraph{\@dottedtocline{6}{14em}{6em}}
\makeatother

\setcounter{secnumdepth}{4}
\setcounter{tocdepth}{3}
%%%%%%%%%%%%%%%%%%%%%%%%%%%%%%%%
% FIN creacio subsubsubtitle
%%%%%%%%%%%%%%%%%%%%%%%%%%%%%%%%


%%%%%%%%%%%%%%%%%%%%%%%%%%%%%%%%
% INI creacio equation
%%%%%%%%%%%%%%%%%%%%%%%%%%%%%%%%
\DeclareFloatingEnvironment[fileext=loeq,placement={!ht},name=Equació]{equationfloat}
%%%%%%%%%%%%%%%%%%%%%%%%%%%%%%%%
% FIN creacio equation
%%%%%%%%%%%%%%%%%%%%%%%%%%%%%%%%

\titleformat{\section}[block]
  {\clearpage\normalfont\Large\bfseries}
  {\thesection}{1em}{}


\renewcommand\refname{Bibliografia}


\linespread{1.25}


\makeatletter
\renewenvironment{abstract}{%
    \if@twocolumn
      \section*{\abstractname}%
    \else 
      \begin{center}%
        {\bfseries \Large\itshape\abstractname\vspace{\z@}}%  % Título en negrita, 14pt y cursiva
      \end{center}%
      \quotation
    \fi}
    {\if@twocolumn\else\endquotation\fi}
\makeatother

% \fancypagestyle{main}{
%     \fancyhf{}
%     \lhead{\small Aplicació Web per a la Visualització i Comparació
% de Trajectòries d’Algorismes d’Optimització}
%     % \rhead{\includegraphics[width=1.2cm]{resources/Icons/logo-fib.png}}
%     \cfoot{\thepage}
%     \setlength{\footskip}{35pt}
% }

\newcommand{\q}[1]{``#1''}

\begin{document}

\hypersetup{
    pdfborder=0 0 0,
}

%% PORTADA I TABLE OF CONTENTS %%

%%%%                                   %%%%
%         INDEX DE LA MEMORIA             %
%%%%                                   %%%%

% \includepdf[pages=-]{portada_fib.pdf}
\input{Extra/GEP}
%\input{Extra/portada}
% \cleardoublepage

% \pagestyle{main}
% \setcounter{page}{11}
\include{toc}
% \cleardoublepage
\input{Contexto y alcance/main}

\input{Planificación temporal/main}

\input{Gestión económica/main}

\input{Informe de sostenibilidad/main}

% \cleardoublepage

% \cleardoublepage


% \input{Annex/main}


\clearpage
% Estilo bibliográfico IEEE
\bibliographystyle{IEEEtran}
\renewcommand{\refname}{Referencias}
\bibliography{references.bib}

%%%%                                   %%%%
%         FINAL DE LA MEMORIA             %
%%%%                                   %%%%

\end{document}


%% VERSIÓ INCIAL MEMORIA %%




\clearpage
% Estilo bibliográfico IEEE
\bibliographystyle{IEEEtran}
\renewcommand{\refname}{Referencias}
\bibliography{references.bib}

%%%%                                   %%%%
%         FINAL DE LA MEMORIA             %
%%%%                                   %%%%

\end{document}


%% VERSIÓ INCIAL MEMORIA %%



\section{Gestión económica}

La gestión económica de un proyecto es fundamental para asegurar su viabilidad y exito a largo plazo. Ya hemos visto en secciones anteriores, los otros dos vertices 
del triángulo de la gestión de proyectos: el alcance y el tiempo. En esta sección, nos centraremos en cómo gestionar eficazmente los recursos financieros para garantizar
que el proyecto se mantenga dentro del presupuesto y se logren los objetivos establecidos. 

\input{Gestión económica/Recursos humanos.tex}

\subsection{Recursos hardware}
Los gastos de los recursos hardware hacen referencia al coste de todo los equipos físicos necesarios para el desarrollo del proyecto.
En este caso, se ha considerado el uso del mismo equipo par los diferentes roles desempeñados, que contiene el material 
necesario para que cada uno pueda realizar sus funciones. El equipo utilizado es el siguiente:

\begin{table}[H]
    \centering
    \begin{tabular}{|c|c|c|c|c|}
        \hline
        \textbf{Cantidad} & \textbf{Nombre} & \textbf{Coste}\\
        \hline
        3 & hp elitebook 640 14 inch g11 notebook pc & 859,00 € \\
        3 & hp 125 Wired Keyboard PERP & 27,10 € \\
        3 & hp Wired Mouse 1000 & 5,99 € \\
        \hline
    \end{tabular}
    \caption{Resumen del tiempo aproximado para cada tarea. Elaboración propia.}
    \label{tab:estimaciones}
\end{table}


\subsection{Recursos software}
La naturaleza de la organización como entidad sin ánimo de lucro impulsa una estrategia 
rigurosa de optimización de costes, priorizando el uso de soluciones de código abierto 
(\textit{Open Source}) siempre que sea posible. Este enfoque no solo minimiza el gasto 
financiero, sino que también fomenta la transparencia y la adaptabilidad técnica del proyecto.

Para la gestión de recursos de computación y planificación de tareas, se ha optado por 
Slurm, una solución de código abierto ampliamente reconocida en entornos de alto 
rendimiento. En cuanto a la infraestructura colaborativa, la empresa dispone de un acuerdo 
para el uso de Google Workspace, lo que permite gestionar la comunicación y la 
documentación bajo una filosofía de costes reducidos o nulos que se alinea con la estrategia 
\textit{Open Source}.

Para la infraestructura de datos, se utiliza Google Cloud en su plan Estándar, con 
almacenamiento y análisis en BigQuery On-Demand. Según la estimación de cargas del 
proyecto, se prevé un volumen aproximado de:
\begin{itemize}
    \item \textbf{Almacenamiento}: 100--200 GB de datos históricos generados por Slurm, lo que implica un coste mensual aproximado de 4 a 8€ (0,02 €/GB/mes).
    \item \textbf{Consultas}: unos 2 TB procesadas al mes, lo que corresponde a 10 € mensuales (5 €/TB procesado).
\end{itemize}

En total, se estima un coste global de unos \textbf{20–25 € mensuales}, es decir, alrededor de \textbf{80–100 €} para los 4 meses de duración del proyecto. Se ha imputado un presupuesto fijo de \textbf{100 €} para cubrir holguras en picos de almacenamiento o consultas.

Finalmente, la gestión de tareas se realiza a través de \textbf{Jira}, donde se aprovecha una reducción de precio substancial por el estado non-profit de la empresa.

\begin{table}[H]
    \begin{center}
    \begin{tabular}{|c|c|c|c|}
        \hline
        \textbf{Nombre del recurso} & \textbf{Coste} & \textbf{Vida útil} & \textbf{Coste Horario}\\
        \hline
        Google Cloud (BigQuery On-Demand, 4 meses) & 100,00 € & 539 h & 0,186 € \\
        Jira (3 usuarios, coste reducido) & 50,00 € & 539 h & 0,093 € \\
        Slurm (Open Source) & 0,00 € & 539 h & 0,000 € \\
        Google Workspace (Acuerdo non-profit) & 0,00 € & 539 h & 0,000 € \\
        \hline
        \multicolumn{3}{|l|}{\textbf{Total horario} (asumiendo los 3 roles)} & \textbf{0,279 €} \\
        \hline
        \multicolumn{3}{|l|}{\textbf{Total del proyecto} (asumiendo los 3 roles)} & \textbf{150,00 €} \\
        \hline
    \end{tabular}
    \caption{Coste horario de los recursos software indirectos. Elaboración propia.}
    \label{tab:recursos_software}
    \end{center}
\end{table}


\subsection{Contingencias}
Existe una gran posibilidad de que durante la ejecución del proyecto surjan imprevistos que 
requieran recursos adicionales no contemplados en el presupuesto inicial. Es por eso que 
para mitigar el riesgo que implican estos imprevistos, se ha decidido asignar un porcentaje 
del presupuesto total a contingencias. El porcentaje habitual de contingencia para un proyecto 
de desarrollo de software (o sector tecnológico en general) suele oscilar entre el 10\% y el 20\% 
del presupuesto total. Se ha optado por asignar un \textbf{15\%} del presupuesto total a contingencias, 
lo que nos porporciona los siguientes valores:


\begin{table}[H]
    \begin{center}
    \begin{tabular}{|c|c|c|c|}
        \hline
        \textbf{Recurso} & \textbf{Importe} & \textbf{Contingencias} \\ 
        \hline
        Recursos humanos & 12.176,58 € & 1.826,49 € \\
        Recursos materiales directos & 2.676,00 € & 401,40 € \\
        Recursos materiales indirectos  & 6.447,00 € & 967,05 € \\
        \hline
        \multicolumn{2}{|l|}{\textbf{Total}}  & \textbf{3.194,94 €} \\
        \hline
    \end{tabular}
    \caption{Tabla de contingencias. Elaboración propia.}
    \label{tab:calculo_contingencias}
    \end{center}
\end{table}

\section{Informe de sostenibilidad}

\input{Informe de sostenibilidad/Autoevaluación.tex}


\input{Informe de sostenibilidad/Dimensión ambiental.tex}

\subsection{Dimensión económica}

\textbf{¿Has estimado el coste de la realización del proyecto (recursos humanos
y materiales)?}

Sí, se ha realizado una estimación del coste del proyecto, una parte fundamental del proyecto.
Dicha estimación se encuentra en los apartados 6.x en la Gestión económica del proyecto.

\textbf{¿Cómo se resuelve actualmente el problema que quieres abordar (estado del arte)? 
¿En qué mejorará económicamente tu solución a las existentes?}

Siguiendo con la respuesta del apartado anterorior no existe una solución actual al problema a abordar.
Pero se espera que la solución propuesta mejore económicamente las cuentas de la institución. Al optimizar
los costes de los recursos y asignarlos de manera eficiente, se espera que la institución pueda reaprovechar
recursos y así poder comprar equipamiento nuevo que mejore la calidad del servicio ofrecido.

\input{Informe de sostenibilidad/Dimensión social.tex}


% \cleardoublepage

% \cleardoublepage


% \documentclass[a4paper, 11pt, twoside]{article}
\usepackage{wrapfig}
% \usepackage{abstract}
\usepackage{graphicx} % Required for inserting images
\usepackage[a4paper, left=1.15in, right=1.15in, top=1.5in, bottom=1.2in]{geometry}
\usepackage{float}
\usepackage{array}
\usepackage{booktabs}
\usepackage[T1]{fontenc}
\usepackage{amsmath}
\usepackage{tabularx}
\usepackage{graphicx} % Required for inserting images
\usepackage{tikz, xcolor}
\usepackage{hyperref}
\usepackage{cite}
\usepackage{pdfpages}
\usepackage{pdflscape}
\usepackage{afterpage}
\usepackage{textpos}
\usepackage{fancyhdr}
\usepackage{caption}
\usepackage{lmodern,textcomp}
\usepackage[catalan]{babel}
\usepackage{lipsum}
\usepackage{newfloat}
\usepackage{listings}
\usepackage{placeins}
\usepackage[ruled]{algorithm2e}
\usepackage{tcolorbox}
% Estilo de citas IEEE sencillo: usar paquete cite para comprimir y ordenar referencias numéricas
\usepackage{cite}
\usepackage[toc,page]{appendix}
\usepackage{forest}
\usepackage{pdflscape}
\usepackage{graphicx}
\usepackage{geometry}
\usepackage{quoting,xparse}
\usepackage{svg}



% \NewDocumentCommand{\bywhom}{m}{% the Bourbaki trick
%   {\nobreak\hfill\penalty50\hskip1em\null\nobreak
%    \hfill\mbox{\normalfont(#1)}%
%    \parfillskip=0pt \finalhyphendemerits=0 \par}%
% }

% AIXÒ SERVEIX PER A MODIFICAR UN PROBLEMA AMB ELS CARACTERS --> " + PARAULA 

% \NewDocumentEnvironment{pquotation}{m}
%   {\begin{quoting}[
%      indentfirst=true,
%      leftmargin=\parindent,
%      rightmargin=\parindent]\itshape}
%   {\bywhom{#1}
  
%   \end{quoting}}

% AIXÒ PER INSERIR FRAGMENTS DE CODI

\lstnewenvironment{code}[1][]%
{
   \noindent
   \minipage{\linewidth}
   \vspace{0.5\baselineskip}
   \lstset{
     basicstyle=\ttfamily\footnotesize,
     frame=single,
     breaklines=true,    % Permitir que las líneas se rompan automáticamente
     breakatwhitespace=true, % Permitir la ruptura en los espacios en blanco
     #1
   }
}
{\endminipage}

\lstset{
    language=Java,
    basicstyle=\ttfamily\small,  % Define el tipo de letra y tamaño del código
    keywordstyle=\color{blue},   % Define el color de las palabras clave de Java
    stringstyle=\color{red},     % Define el color de las cadenas de texto
    commentstyle=\color{green},  % Define el color de los comentarios
    morecomment=[s][\color{gray}]{/**}{*/},  % Estilo para comentarios de bloque
    breaklines=true,  % Configura el quiebre automático de línea
    numbers=left,  % Posición de los números de línea
    numberstyle=\tiny\color{gray},  % Estilo para los números de línea
    showspaces=false,  % No mostrar espacios con underscores
    showstringspaces=false,  % No mostrar espacios dentro de strings
    frame=single,  % Agregar un marco alrededor del código
    captionpos=b,  % b = bottom, t = top
}

\definecolor{eclipseStrings}{RGB}{42,0.0,255}
\definecolor{eclipseKeywords}{RGB}{127,0,85}
\colorlet{numb}{magenta!60!black}

\lstdefinelanguage{json}{
    basicstyle=\ttfamily\small,
    showstringspaces=false,
    breaklines=true,
    frame=single,
    literate=
     *{0}{{{\color{numb}0}}}{1}
      {1}{{{\color{numb}1}}}{1}
      {2}{{{\color{numb}2}}}{1}
      {3}{{{\color{numb}3}}}{1}
      {4}{{{\color{numb}4}}}{1}
      {5}{{{\color{numb}5}}}{1}
      {6}{{{\color{numb}6}}}{1}
      {7}{{{\color{numb}7}}}{1}
      {8}{{{\color{numb}8}}}{1}
      {9}{{{\color{numb}9}}}{1}
      {:}{{{\color{punct}{:}}}}{1}
      {,}{{{\color{punct}{,}}}}{1}
      {\{}{{{\color{delim}{\{}}}}{1}
      {\}}{{{\color{delim}{\}}}}}{1}
      {[}{{{\color{delim}{[}}}}{1}
      {]}{{{\color{delim}{]}}}}{1},
    morestring=[b]"
}

\lstdefinelanguage{JavaScript}{
  morekeywords=[1]{break, continue, delete, else, for, function, if, in,
    new, return, this, typeof, var, void, while, with},
  % Literals, primitive types, and reference types.
  morekeywords=[2]{false, null, true, boolean, number, undefined,
    Array, Boolean, Date, Math, Number, String, Object},
  % Built-ins.
  morekeywords=[3]{eval, parseInt, parseFloat, escape, unescape},
  sensitive,
  morecomment=[s]{/*}{*/},
  morecomment=[l]//,
  morecomment=[s]{/**}{*/}, % JavaDoc style comments
  morestring=[b]',
  morestring=[b]"
}[keywords, comments, strings]

\lstdefinelanguage{svelte}{
    morekeywords={async, await, then, if, each, in, on, let, const, function, import, export, from, with, typeof, new, var, return},
    sensitive=true,
    morecomment=[l]{//},
    morecomment=[s]{/*}{*/},
    morecomment=[s]{<!--}{-->},
    morestring=[b]',
    morestring=[b]",
    morestring=[b]`,
    ndkeywords={import, export, return, if, for, while, function, switch, case, break, var, let, const, async, await, then, catch, try, typeof, new, class, extends, implements, interface, constructor, public, private, protected, static, get, set},
    alsoletter={:},
    moredelim=**[is][\color{orange}]{\@\{}{\@\}},
    moredelim=**[is][\color{cyan}]{\$\{}{\}},
    keywordstyle=\color{blue}\bfseries,
    commentstyle=\color{gray}\ttfamily,
    stringstyle=\color{red}\ttfamily,
    identifierstyle=\color{black},
    breaklines=true,
    basicstyle=\ttfamily\small,
    columns=fullflexible,
    showstringspaces=false,
    tabsize=2,
    frame=single,
    captionpos=b,
}


\lstdefinestyle{java}{
    language=Java,
    basicstyle=\ttfamily\small,
    keywordstyle=\color{blue},
    stringstyle=\color{red},
    commentstyle=\color{gray},
    morecomment=[s][\color{gray}]{/**}{*/},
    breaklines=true,
    numbers=left,
    numberstyle=\tiny\color{gray},
    showspaces=false,
    showstringspaces=false,
    frame=single,
    captionpos=b
}

\usepackage{color}
\definecolor{numb}{rgb}{0.58,0,0.82}
\definecolor{punct}{rgb}{0.3,0.3,0.3}
\definecolor{delim}{rgb}{0.5,0.5,0.5}
\definecolor{string}{rgb}{0.6,0.1,0.1}

\lstdefinestyle{json}{
    language=json,
    basicstyle=\ttfamily\small,
    keywordstyle=\color{blue},
    stringstyle=\color{string},
    commentstyle=\color{green},
    morecomment=[s][\color{gray}]{/**}{*/},
    breaklines=true,
    numbers=left,
    numberstyle=\tiny\color{gray},
    showspaces=false,
    showstringspaces=false,
    frame=single,
    captionpos=b
}

\lstdefinelanguage{drl}{
    basicstyle=\ttfamily\small,
    keywordstyle=\color{blue}\bfseries,
    stringstyle=\color{red},
    commentstyle=\color{green},
    morecomment=[s][\color{green}]{/*}{*/},
    morecomment=[l][\color{green}]{//},
    morestring=[b]",
    keywords={rule, when, then, end, salience, eval, update, retract},
    sensitive=true,
    breaklines=true,
    numbers=left,
    numberstyle=\tiny\color{gray},
    showspaces=false,
    showstringspaces=false,
    frame=single,
    captionpos=b
}

\lstdefinestyle{drl}{
    language=drl,
    basicstyle=\ttfamily\small,
    keywordstyle=\color{blue}\bfseries,
    stringstyle=\color{red},
    commentstyle=\color{green},
    breaklines=true,
    numbers=left,
    numberstyle=\tiny\color{gray},
    showspaces=false,
    showstringspaces=false,
    frame=single,
    captionpos=b
}



\renewcommand{\lstlistingname}{Codi}
\renewcommand{\lstlistlistingname}{Índex de Fragments de Codis}
\renewcommand{\appendixpagename}{Annexos}
\renewcommand{\appendixtocname}{Annexos}

\renewcommand*{\listalgorithmcfname}{Llistat de pseudocodi}
\renewcommand*{\algorithmcfname}{Pseudocodi}
\renewcommand*{\algorithmautorefname}{pseudocodi}

\definecolor{darkgreen}{rgb}{0.0, 0.3, 0.0}
\definecolor{lightgreen}{rgb}{0.9, 0.95, 0.9}

\newcommand{\todo}[1]{
    \begin{tcolorbox}[colback=lightgreen, colframe=darkgreen, title=TODO]
        #1
    \end{tcolorbox}
}


% \usepackage{indentfirst} 
% \setlength{\parindent}{1cm}


\usepackage{parskip}
% \setlength{\parindent}{10pt}

% \setlength{\headheight}{15.6pt}
\setlength{\headheight}{16.05263pt}




\vspace{5mm}

\newcolumntype{u}{>{\centering\arraybackslash\hsize=.22\hsize}X}
\newcolumntype{s}{>{\hsize=.45\hsize}X}


%%%%%%%%%%%%%%%%%%%%%%%%%%%%%%%%
% INI creacio subsubsubsecion
%%%%%%%%%%%%%%%%%%%%%%%%%%%%%%%%
\usepackage{titlesec}

\let\oldparagraph\paragraph
\renewcommand{\paragraph}[1]{\oldparagraph{\hspace{30pt}#1}}


\titleclass{\subsubsubsection}{straight}[\subsection]

\newcounter{subsubsubsection}[subsubsection]
\renewcommand\thesubsubsubsection{\thesubsubsection.\arabic{subsubsubsection}}
\renewcommand\theparagraph{\thesubsubsubsection.\arabic{paragraph}} % optional; useful if paragraphs are to be numbered

\titleformat{\subsubsubsection}
  {\normalfont\normalsize\bfseries}{\thesubsubsubsection}{1em}{}
\titlespacing*{\subsubsubsection}
{0pt}{3.25ex plus 1ex minus .2ex}{1.5ex plus .2ex}


\makeatletter
\renewcommand\paragraph{\@startsection{paragraph}{5}{10pt}%
  {3.25ex \@plus1ex \@minus.2ex}%
  {-1em}%
  {\normalfont\normalsize\bfseries}}
\renewcommand\subparagraph{\@startsection{subparagraph}{6}{\parindent}%
  {3.25ex \@plus1ex \@minus .2ex}%
  {-1em}%
  {\normalfont\normalsize\bfseries}}
\def\toclevel@subsubsubsection{4}
\def\toclevel@paragraph{5}
\def\toclevel@paragraph{6}
\def\l@subsubsubsection{\@dottedtocline{4}{7em}{4em}}
\def\l@paragraph{\@dottedtocline{5}{10em}{5em}}
\def\l@subparagraph{\@dottedtocline{6}{14em}{6em}}
\makeatother

\setcounter{secnumdepth}{4}
\setcounter{tocdepth}{3}
%%%%%%%%%%%%%%%%%%%%%%%%%%%%%%%%
% FIN creacio subsubsubtitle
%%%%%%%%%%%%%%%%%%%%%%%%%%%%%%%%


%%%%%%%%%%%%%%%%%%%%%%%%%%%%%%%%
% INI creacio equation
%%%%%%%%%%%%%%%%%%%%%%%%%%%%%%%%
\DeclareFloatingEnvironment[fileext=loeq,placement={!ht},name=Equació]{equationfloat}
%%%%%%%%%%%%%%%%%%%%%%%%%%%%%%%%
% FIN creacio equation
%%%%%%%%%%%%%%%%%%%%%%%%%%%%%%%%

\titleformat{\section}[block]
  {\clearpage\normalfont\Large\bfseries}
  {\thesection}{1em}{}


\renewcommand\refname{Bibliografia}


\linespread{1.25}


\makeatletter
\renewenvironment{abstract}{%
    \if@twocolumn
      \section*{\abstractname}%
    \else 
      \begin{center}%
        {\bfseries \Large\itshape\abstractname\vspace{\z@}}%  % Título en negrita, 14pt y cursiva
      \end{center}%
      \quotation
    \fi}
    {\if@twocolumn\else\endquotation\fi}
\makeatother

% \fancypagestyle{main}{
%     \fancyhf{}
%     \lhead{\small Aplicació Web per a la Visualització i Comparació
% de Trajectòries d’Algorismes d’Optimització}
%     % \rhead{\includegraphics[width=1.2cm]{resources/Icons/logo-fib.png}}
%     \cfoot{\thepage}
%     \setlength{\footskip}{35pt}
% }

\newcommand{\q}[1]{``#1''}

\begin{document}

\hypersetup{
    pdfborder=0 0 0,
}

%% PORTADA I TABLE OF CONTENTS %%

%%%%                                   %%%%
%         INDEX DE LA MEMORIA             %
%%%%                                   %%%%

% \includepdf[pages=-]{portada_fib.pdf}
%%% Portada
\begin{titlepage}
\begin{center}
    \includegraphics[width=1\linewidth]{extra/figures/logo-upc-fib.png}
    \vspace{30mm}

    \LARGE{\textsc{\textbf{{INTEGRACIÓN Y AUTOMATIZACIÓN DE COSTES DE INFRAESTRUCTURA HPC EN ENTORNOS DE INVESTIGACIÓN}}}} \\
    \vspace{25mm}
      
        
    \textbf{\Large{\textsc{PAU NAVARRO ÁLVAREZ}}} \\
    \vspace{20mm}
        
    
    
    \large{\textbf{Director/a} ROBERTO RIVEIRO INSUA (FUNDACIO INSTITUT DE RECERCA BIOMEDICA (IRB BARCELONA))} \\
    \large{\textbf{Ponente:} ALEX PAJUELO GONZÁLEZ (Departamento de Arquitectura de Computadores)} \\
    \large{\textbf{Titulación:} Grado en Ingeniería Informática (Tecnologías de la información)} \\
    
    \large{\textbf{Lliurament 2:} }Planificación temporal \\
    \large{\textbf{Tutor GEP:} CAROLINA CONSOLIDACIÓN SEGURA (Departament d'Organització d'Empreses)} \\
    \large{\textbf{Fecha:}} 29 de Septiembre de 2025\\

    \textbf{\large{Memoria del trabajo de fin de grado}} \\
    \vspace{3mm}
    \textbf{\large{Facultat d'Informàtica de Barcelona (FIB)}} \\
    \vspace{3mm}
    \textbf{\large{Universitat Politècnica de Catalunya (UPC) - BarcelonaTech}} \\

\end{center}
\end{titlepage}
%\input{Extra/portada}
% \cleardoublepage

% \pagestyle{main}
% \setcounter{page}{11}
\renewcommand{\contentsname}{Continguts}

\tableofcontents
\clearpage
\listoffigures
\begingroup
    \let\clearpage\relax
    \listoftables
\endgroup
\clearpage

\newpage
% \cleardoublepage
\input{Contexto y alcance/Contextualización/main}
\input{Contexto y alcance/Justificación/main}
\input{Contexto y alcance/Alcance/main}
\input{Contexto y alcance/Metodología y herramientas de seguimiento/main}

\documentclass[a4paper, 11pt, twoside]{article}
\usepackage{wrapfig}
% \usepackage{abstract}
\usepackage{graphicx} % Required for inserting images
\usepackage[a4paper, left=1.15in, right=1.15in, top=1.5in, bottom=1.2in]{geometry}
\usepackage{float}
\usepackage{array}
\usepackage{booktabs}
\usepackage[T1]{fontenc}
\usepackage{amsmath}
\usepackage{tabularx}
\usepackage{graphicx} % Required for inserting images
\usepackage{tikz, xcolor}
\usepackage{hyperref}
\usepackage{cite}
\usepackage{pdfpages}
\usepackage{pdflscape}
\usepackage{afterpage}
\usepackage{textpos}
\usepackage{fancyhdr}
\usepackage{caption}
\usepackage{lmodern,textcomp}
\usepackage[catalan]{babel}
\usepackage{lipsum}
\usepackage{newfloat}
\usepackage{listings}
\usepackage{placeins}
\usepackage[ruled]{algorithm2e}
\usepackage{tcolorbox}
% Estilo de citas IEEE sencillo: usar paquete cite para comprimir y ordenar referencias numéricas
\usepackage{cite}
\usepackage[toc,page]{appendix}
\usepackage{forest}
\usepackage{pdflscape}
\usepackage{graphicx}
\usepackage{geometry}
\usepackage{quoting,xparse}
\usepackage{svg}



% \NewDocumentCommand{\bywhom}{m}{% the Bourbaki trick
%   {\nobreak\hfill\penalty50\hskip1em\null\nobreak
%    \hfill\mbox{\normalfont(#1)}%
%    \parfillskip=0pt \finalhyphendemerits=0 \par}%
% }

% AIXÒ SERVEIX PER A MODIFICAR UN PROBLEMA AMB ELS CARACTERS --> " + PARAULA 

% \NewDocumentEnvironment{pquotation}{m}
%   {\begin{quoting}[
%      indentfirst=true,
%      leftmargin=\parindent,
%      rightmargin=\parindent]\itshape}
%   {\bywhom{#1}
  
%   \end{quoting}}

% AIXÒ PER INSERIR FRAGMENTS DE CODI

\lstnewenvironment{code}[1][]%
{
   \noindent
   \minipage{\linewidth}
   \vspace{0.5\baselineskip}
   \lstset{
     basicstyle=\ttfamily\footnotesize,
     frame=single,
     breaklines=true,    % Permitir que las líneas se rompan automáticamente
     breakatwhitespace=true, % Permitir la ruptura en los espacios en blanco
     #1
   }
}
{\endminipage}

\lstset{
    language=Java,
    basicstyle=\ttfamily\small,  % Define el tipo de letra y tamaño del código
    keywordstyle=\color{blue},   % Define el color de las palabras clave de Java
    stringstyle=\color{red},     % Define el color de las cadenas de texto
    commentstyle=\color{green},  % Define el color de los comentarios
    morecomment=[s][\color{gray}]{/**}{*/},  % Estilo para comentarios de bloque
    breaklines=true,  % Configura el quiebre automático de línea
    numbers=left,  % Posición de los números de línea
    numberstyle=\tiny\color{gray},  % Estilo para los números de línea
    showspaces=false,  % No mostrar espacios con underscores
    showstringspaces=false,  % No mostrar espacios dentro de strings
    frame=single,  % Agregar un marco alrededor del código
    captionpos=b,  % b = bottom, t = top
}

\definecolor{eclipseStrings}{RGB}{42,0.0,255}
\definecolor{eclipseKeywords}{RGB}{127,0,85}
\colorlet{numb}{magenta!60!black}

\lstdefinelanguage{json}{
    basicstyle=\ttfamily\small,
    showstringspaces=false,
    breaklines=true,
    frame=single,
    literate=
     *{0}{{{\color{numb}0}}}{1}
      {1}{{{\color{numb}1}}}{1}
      {2}{{{\color{numb}2}}}{1}
      {3}{{{\color{numb}3}}}{1}
      {4}{{{\color{numb}4}}}{1}
      {5}{{{\color{numb}5}}}{1}
      {6}{{{\color{numb}6}}}{1}
      {7}{{{\color{numb}7}}}{1}
      {8}{{{\color{numb}8}}}{1}
      {9}{{{\color{numb}9}}}{1}
      {:}{{{\color{punct}{:}}}}{1}
      {,}{{{\color{punct}{,}}}}{1}
      {\{}{{{\color{delim}{\{}}}}{1}
      {\}}{{{\color{delim}{\}}}}}{1}
      {[}{{{\color{delim}{[}}}}{1}
      {]}{{{\color{delim}{]}}}}{1},
    morestring=[b]"
}

\lstdefinelanguage{JavaScript}{
  morekeywords=[1]{break, continue, delete, else, for, function, if, in,
    new, return, this, typeof, var, void, while, with},
  % Literals, primitive types, and reference types.
  morekeywords=[2]{false, null, true, boolean, number, undefined,
    Array, Boolean, Date, Math, Number, String, Object},
  % Built-ins.
  morekeywords=[3]{eval, parseInt, parseFloat, escape, unescape},
  sensitive,
  morecomment=[s]{/*}{*/},
  morecomment=[l]//,
  morecomment=[s]{/**}{*/}, % JavaDoc style comments
  morestring=[b]',
  morestring=[b]"
}[keywords, comments, strings]

\lstdefinelanguage{svelte}{
    morekeywords={async, await, then, if, each, in, on, let, const, function, import, export, from, with, typeof, new, var, return},
    sensitive=true,
    morecomment=[l]{//},
    morecomment=[s]{/*}{*/},
    morecomment=[s]{<!--}{-->},
    morestring=[b]',
    morestring=[b]",
    morestring=[b]`,
    ndkeywords={import, export, return, if, for, while, function, switch, case, break, var, let, const, async, await, then, catch, try, typeof, new, class, extends, implements, interface, constructor, public, private, protected, static, get, set},
    alsoletter={:},
    moredelim=**[is][\color{orange}]{\@\{}{\@\}},
    moredelim=**[is][\color{cyan}]{\$\{}{\}},
    keywordstyle=\color{blue}\bfseries,
    commentstyle=\color{gray}\ttfamily,
    stringstyle=\color{red}\ttfamily,
    identifierstyle=\color{black},
    breaklines=true,
    basicstyle=\ttfamily\small,
    columns=fullflexible,
    showstringspaces=false,
    tabsize=2,
    frame=single,
    captionpos=b,
}


\lstdefinestyle{java}{
    language=Java,
    basicstyle=\ttfamily\small,
    keywordstyle=\color{blue},
    stringstyle=\color{red},
    commentstyle=\color{gray},
    morecomment=[s][\color{gray}]{/**}{*/},
    breaklines=true,
    numbers=left,
    numberstyle=\tiny\color{gray},
    showspaces=false,
    showstringspaces=false,
    frame=single,
    captionpos=b
}

\usepackage{color}
\definecolor{numb}{rgb}{0.58,0,0.82}
\definecolor{punct}{rgb}{0.3,0.3,0.3}
\definecolor{delim}{rgb}{0.5,0.5,0.5}
\definecolor{string}{rgb}{0.6,0.1,0.1}

\lstdefinestyle{json}{
    language=json,
    basicstyle=\ttfamily\small,
    keywordstyle=\color{blue},
    stringstyle=\color{string},
    commentstyle=\color{green},
    morecomment=[s][\color{gray}]{/**}{*/},
    breaklines=true,
    numbers=left,
    numberstyle=\tiny\color{gray},
    showspaces=false,
    showstringspaces=false,
    frame=single,
    captionpos=b
}

\lstdefinelanguage{drl}{
    basicstyle=\ttfamily\small,
    keywordstyle=\color{blue}\bfseries,
    stringstyle=\color{red},
    commentstyle=\color{green},
    morecomment=[s][\color{green}]{/*}{*/},
    morecomment=[l][\color{green}]{//},
    morestring=[b]",
    keywords={rule, when, then, end, salience, eval, update, retract},
    sensitive=true,
    breaklines=true,
    numbers=left,
    numberstyle=\tiny\color{gray},
    showspaces=false,
    showstringspaces=false,
    frame=single,
    captionpos=b
}

\lstdefinestyle{drl}{
    language=drl,
    basicstyle=\ttfamily\small,
    keywordstyle=\color{blue}\bfseries,
    stringstyle=\color{red},
    commentstyle=\color{green},
    breaklines=true,
    numbers=left,
    numberstyle=\tiny\color{gray},
    showspaces=false,
    showstringspaces=false,
    frame=single,
    captionpos=b
}



\renewcommand{\lstlistingname}{Codi}
\renewcommand{\lstlistlistingname}{Índex de Fragments de Codis}
\renewcommand{\appendixpagename}{Annexos}
\renewcommand{\appendixtocname}{Annexos}

\renewcommand*{\listalgorithmcfname}{Llistat de pseudocodi}
\renewcommand*{\algorithmcfname}{Pseudocodi}
\renewcommand*{\algorithmautorefname}{pseudocodi}

\definecolor{darkgreen}{rgb}{0.0, 0.3, 0.0}
\definecolor{lightgreen}{rgb}{0.9, 0.95, 0.9}

\newcommand{\todo}[1]{
    \begin{tcolorbox}[colback=lightgreen, colframe=darkgreen, title=TODO]
        #1
    \end{tcolorbox}
}


% \usepackage{indentfirst} 
% \setlength{\parindent}{1cm}


\usepackage{parskip}
% \setlength{\parindent}{10pt}

% \setlength{\headheight}{15.6pt}
\setlength{\headheight}{16.05263pt}




\vspace{5mm}

\newcolumntype{u}{>{\centering\arraybackslash\hsize=.22\hsize}X}
\newcolumntype{s}{>{\hsize=.45\hsize}X}


%%%%%%%%%%%%%%%%%%%%%%%%%%%%%%%%
% INI creacio subsubsubsecion
%%%%%%%%%%%%%%%%%%%%%%%%%%%%%%%%
\usepackage{titlesec}

\let\oldparagraph\paragraph
\renewcommand{\paragraph}[1]{\oldparagraph{\hspace{30pt}#1}}


\titleclass{\subsubsubsection}{straight}[\subsection]

\newcounter{subsubsubsection}[subsubsection]
\renewcommand\thesubsubsubsection{\thesubsubsection.\arabic{subsubsubsection}}
\renewcommand\theparagraph{\thesubsubsubsection.\arabic{paragraph}} % optional; useful if paragraphs are to be numbered

\titleformat{\subsubsubsection}
  {\normalfont\normalsize\bfseries}{\thesubsubsubsection}{1em}{}
\titlespacing*{\subsubsubsection}
{0pt}{3.25ex plus 1ex minus .2ex}{1.5ex plus .2ex}


\makeatletter
\renewcommand\paragraph{\@startsection{paragraph}{5}{10pt}%
  {3.25ex \@plus1ex \@minus.2ex}%
  {-1em}%
  {\normalfont\normalsize\bfseries}}
\renewcommand\subparagraph{\@startsection{subparagraph}{6}{\parindent}%
  {3.25ex \@plus1ex \@minus .2ex}%
  {-1em}%
  {\normalfont\normalsize\bfseries}}
\def\toclevel@subsubsubsection{4}
\def\toclevel@paragraph{5}
\def\toclevel@paragraph{6}
\def\l@subsubsubsection{\@dottedtocline{4}{7em}{4em}}
\def\l@paragraph{\@dottedtocline{5}{10em}{5em}}
\def\l@subparagraph{\@dottedtocline{6}{14em}{6em}}
\makeatother

\setcounter{secnumdepth}{4}
\setcounter{tocdepth}{3}
%%%%%%%%%%%%%%%%%%%%%%%%%%%%%%%%
% FIN creacio subsubsubtitle
%%%%%%%%%%%%%%%%%%%%%%%%%%%%%%%%


%%%%%%%%%%%%%%%%%%%%%%%%%%%%%%%%
% INI creacio equation
%%%%%%%%%%%%%%%%%%%%%%%%%%%%%%%%
\DeclareFloatingEnvironment[fileext=loeq,placement={!ht},name=Equació]{equationfloat}
%%%%%%%%%%%%%%%%%%%%%%%%%%%%%%%%
% FIN creacio equation
%%%%%%%%%%%%%%%%%%%%%%%%%%%%%%%%

\titleformat{\section}[block]
  {\clearpage\normalfont\Large\bfseries}
  {\thesection}{1em}{}


\renewcommand\refname{Bibliografia}


\linespread{1.25}


\makeatletter
\renewenvironment{abstract}{%
    \if@twocolumn
      \section*{\abstractname}%
    \else 
      \begin{center}%
        {\bfseries \Large\itshape\abstractname\vspace{\z@}}%  % Título en negrita, 14pt y cursiva
      \end{center}%
      \quotation
    \fi}
    {\if@twocolumn\else\endquotation\fi}
\makeatother

% \fancypagestyle{main}{
%     \fancyhf{}
%     \lhead{\small Aplicació Web per a la Visualització i Comparació
% de Trajectòries d’Algorismes d’Optimització}
%     % \rhead{\includegraphics[width=1.2cm]{resources/Icons/logo-fib.png}}
%     \cfoot{\thepage}
%     \setlength{\footskip}{35pt}
% }

\newcommand{\q}[1]{``#1''}

\begin{document}

\hypersetup{
    pdfborder=0 0 0,
}

%% PORTADA I TABLE OF CONTENTS %%

%%%%                                   %%%%
%         INDEX DE LA MEMORIA             %
%%%%                                   %%%%

% \includepdf[pages=-]{portada_fib.pdf}
\input{Extra/GEP}
%\input{Extra/portada}
% \cleardoublepage

% \pagestyle{main}
% \setcounter{page}{11}
\include{toc}
% \cleardoublepage
\input{Contexto y alcance/main}

\input{Planificación temporal/main}

\input{Gestión económica/main}

\input{Informe de sostenibilidad/main}

% \cleardoublepage

% \cleardoublepage


% \input{Annex/main}


\clearpage
% Estilo bibliográfico IEEE
\bibliographystyle{IEEEtran}
\renewcommand{\refname}{Referencias}
\bibliography{references.bib}

%%%%                                   %%%%
%         FINAL DE LA MEMORIA             %
%%%%                                   %%%%

\end{document}


%% VERSIÓ INCIAL MEMORIA %%



\section{Gestión económica}

La gestión económica de un proyecto es fundamental para asegurar su viabilidad y exito a largo plazo. Ya hemos visto en secciones anteriores, los otros dos vertices 
del triángulo de la gestión de proyectos: el alcance y el tiempo. En esta sección, nos centraremos en cómo gestionar eficazmente los recursos financieros para garantizar
que el proyecto se mantenga dentro del presupuesto y se logren los objetivos establecidos. 

\input{Gestión económica/Recursos humanos.tex}

\input{Gestión económica/Recursos Hardware.tex}

\input{Gestión económica/Recursos Software.tex}

\input{Gestión económica/Contingencias.tex}

\section{Informe de sostenibilidad}

\input{Informe de sostenibilidad/Autoevaluación.tex}


\input{Informe de sostenibilidad/Dimensión ambiental.tex}

\input{Informe de sostenibilidad/Dimensión económica.tex}

\input{Informe de sostenibilidad/Dimensión social.tex}


% \cleardoublepage

% \cleardoublepage


% \documentclass[a4paper, 11pt, twoside]{article}
\usepackage{wrapfig}
% \usepackage{abstract}
\usepackage{graphicx} % Required for inserting images
\usepackage[a4paper, left=1.15in, right=1.15in, top=1.5in, bottom=1.2in]{geometry}
\usepackage{float}
\usepackage{array}
\usepackage{booktabs}
\usepackage[T1]{fontenc}
\usepackage{amsmath}
\usepackage{tabularx}
\usepackage{graphicx} % Required for inserting images
\usepackage{tikz, xcolor}
\usepackage{hyperref}
\usepackage{cite}
\usepackage{pdfpages}
\usepackage{pdflscape}
\usepackage{afterpage}
\usepackage{textpos}
\usepackage{fancyhdr}
\usepackage{caption}
\usepackage{lmodern,textcomp}
\usepackage[catalan]{babel}
\usepackage{lipsum}
\usepackage{newfloat}
\usepackage{listings}
\usepackage{placeins}
\usepackage[ruled]{algorithm2e}
\usepackage{tcolorbox}
% Estilo de citas IEEE sencillo: usar paquete cite para comprimir y ordenar referencias numéricas
\usepackage{cite}
\usepackage[toc,page]{appendix}
\usepackage{forest}
\usepackage{pdflscape}
\usepackage{graphicx}
\usepackage{geometry}
\usepackage{quoting,xparse}
\usepackage{svg}



% \NewDocumentCommand{\bywhom}{m}{% the Bourbaki trick
%   {\nobreak\hfill\penalty50\hskip1em\null\nobreak
%    \hfill\mbox{\normalfont(#1)}%
%    \parfillskip=0pt \finalhyphendemerits=0 \par}%
% }

% AIXÒ SERVEIX PER A MODIFICAR UN PROBLEMA AMB ELS CARACTERS --> " + PARAULA 

% \NewDocumentEnvironment{pquotation}{m}
%   {\begin{quoting}[
%      indentfirst=true,
%      leftmargin=\parindent,
%      rightmargin=\parindent]\itshape}
%   {\bywhom{#1}
  
%   \end{quoting}}

% AIXÒ PER INSERIR FRAGMENTS DE CODI

\lstnewenvironment{code}[1][]%
{
   \noindent
   \minipage{\linewidth}
   \vspace{0.5\baselineskip}
   \lstset{
     basicstyle=\ttfamily\footnotesize,
     frame=single,
     breaklines=true,    % Permitir que las líneas se rompan automáticamente
     breakatwhitespace=true, % Permitir la ruptura en los espacios en blanco
     #1
   }
}
{\endminipage}

\lstset{
    language=Java,
    basicstyle=\ttfamily\small,  % Define el tipo de letra y tamaño del código
    keywordstyle=\color{blue},   % Define el color de las palabras clave de Java
    stringstyle=\color{red},     % Define el color de las cadenas de texto
    commentstyle=\color{green},  % Define el color de los comentarios
    morecomment=[s][\color{gray}]{/**}{*/},  % Estilo para comentarios de bloque
    breaklines=true,  % Configura el quiebre automático de línea
    numbers=left,  % Posición de los números de línea
    numberstyle=\tiny\color{gray},  % Estilo para los números de línea
    showspaces=false,  % No mostrar espacios con underscores
    showstringspaces=false,  % No mostrar espacios dentro de strings
    frame=single,  % Agregar un marco alrededor del código
    captionpos=b,  % b = bottom, t = top
}

\definecolor{eclipseStrings}{RGB}{42,0.0,255}
\definecolor{eclipseKeywords}{RGB}{127,0,85}
\colorlet{numb}{magenta!60!black}

\lstdefinelanguage{json}{
    basicstyle=\ttfamily\small,
    showstringspaces=false,
    breaklines=true,
    frame=single,
    literate=
     *{0}{{{\color{numb}0}}}{1}
      {1}{{{\color{numb}1}}}{1}
      {2}{{{\color{numb}2}}}{1}
      {3}{{{\color{numb}3}}}{1}
      {4}{{{\color{numb}4}}}{1}
      {5}{{{\color{numb}5}}}{1}
      {6}{{{\color{numb}6}}}{1}
      {7}{{{\color{numb}7}}}{1}
      {8}{{{\color{numb}8}}}{1}
      {9}{{{\color{numb}9}}}{1}
      {:}{{{\color{punct}{:}}}}{1}
      {,}{{{\color{punct}{,}}}}{1}
      {\{}{{{\color{delim}{\{}}}}{1}
      {\}}{{{\color{delim}{\}}}}}{1}
      {[}{{{\color{delim}{[}}}}{1}
      {]}{{{\color{delim}{]}}}}{1},
    morestring=[b]"
}

\lstdefinelanguage{JavaScript}{
  morekeywords=[1]{break, continue, delete, else, for, function, if, in,
    new, return, this, typeof, var, void, while, with},
  % Literals, primitive types, and reference types.
  morekeywords=[2]{false, null, true, boolean, number, undefined,
    Array, Boolean, Date, Math, Number, String, Object},
  % Built-ins.
  morekeywords=[3]{eval, parseInt, parseFloat, escape, unescape},
  sensitive,
  morecomment=[s]{/*}{*/},
  morecomment=[l]//,
  morecomment=[s]{/**}{*/}, % JavaDoc style comments
  morestring=[b]',
  morestring=[b]"
}[keywords, comments, strings]

\lstdefinelanguage{svelte}{
    morekeywords={async, await, then, if, each, in, on, let, const, function, import, export, from, with, typeof, new, var, return},
    sensitive=true,
    morecomment=[l]{//},
    morecomment=[s]{/*}{*/},
    morecomment=[s]{<!--}{-->},
    morestring=[b]',
    morestring=[b]",
    morestring=[b]`,
    ndkeywords={import, export, return, if, for, while, function, switch, case, break, var, let, const, async, await, then, catch, try, typeof, new, class, extends, implements, interface, constructor, public, private, protected, static, get, set},
    alsoletter={:},
    moredelim=**[is][\color{orange}]{\@\{}{\@\}},
    moredelim=**[is][\color{cyan}]{\$\{}{\}},
    keywordstyle=\color{blue}\bfseries,
    commentstyle=\color{gray}\ttfamily,
    stringstyle=\color{red}\ttfamily,
    identifierstyle=\color{black},
    breaklines=true,
    basicstyle=\ttfamily\small,
    columns=fullflexible,
    showstringspaces=false,
    tabsize=2,
    frame=single,
    captionpos=b,
}


\lstdefinestyle{java}{
    language=Java,
    basicstyle=\ttfamily\small,
    keywordstyle=\color{blue},
    stringstyle=\color{red},
    commentstyle=\color{gray},
    morecomment=[s][\color{gray}]{/**}{*/},
    breaklines=true,
    numbers=left,
    numberstyle=\tiny\color{gray},
    showspaces=false,
    showstringspaces=false,
    frame=single,
    captionpos=b
}

\usepackage{color}
\definecolor{numb}{rgb}{0.58,0,0.82}
\definecolor{punct}{rgb}{0.3,0.3,0.3}
\definecolor{delim}{rgb}{0.5,0.5,0.5}
\definecolor{string}{rgb}{0.6,0.1,0.1}

\lstdefinestyle{json}{
    language=json,
    basicstyle=\ttfamily\small,
    keywordstyle=\color{blue},
    stringstyle=\color{string},
    commentstyle=\color{green},
    morecomment=[s][\color{gray}]{/**}{*/},
    breaklines=true,
    numbers=left,
    numberstyle=\tiny\color{gray},
    showspaces=false,
    showstringspaces=false,
    frame=single,
    captionpos=b
}

\lstdefinelanguage{drl}{
    basicstyle=\ttfamily\small,
    keywordstyle=\color{blue}\bfseries,
    stringstyle=\color{red},
    commentstyle=\color{green},
    morecomment=[s][\color{green}]{/*}{*/},
    morecomment=[l][\color{green}]{//},
    morestring=[b]",
    keywords={rule, when, then, end, salience, eval, update, retract},
    sensitive=true,
    breaklines=true,
    numbers=left,
    numberstyle=\tiny\color{gray},
    showspaces=false,
    showstringspaces=false,
    frame=single,
    captionpos=b
}

\lstdefinestyle{drl}{
    language=drl,
    basicstyle=\ttfamily\small,
    keywordstyle=\color{blue}\bfseries,
    stringstyle=\color{red},
    commentstyle=\color{green},
    breaklines=true,
    numbers=left,
    numberstyle=\tiny\color{gray},
    showspaces=false,
    showstringspaces=false,
    frame=single,
    captionpos=b
}



\renewcommand{\lstlistingname}{Codi}
\renewcommand{\lstlistlistingname}{Índex de Fragments de Codis}
\renewcommand{\appendixpagename}{Annexos}
\renewcommand{\appendixtocname}{Annexos}

\renewcommand*{\listalgorithmcfname}{Llistat de pseudocodi}
\renewcommand*{\algorithmcfname}{Pseudocodi}
\renewcommand*{\algorithmautorefname}{pseudocodi}

\definecolor{darkgreen}{rgb}{0.0, 0.3, 0.0}
\definecolor{lightgreen}{rgb}{0.9, 0.95, 0.9}

\newcommand{\todo}[1]{
    \begin{tcolorbox}[colback=lightgreen, colframe=darkgreen, title=TODO]
        #1
    \end{tcolorbox}
}


% \usepackage{indentfirst} 
% \setlength{\parindent}{1cm}


\usepackage{parskip}
% \setlength{\parindent}{10pt}

% \setlength{\headheight}{15.6pt}
\setlength{\headheight}{16.05263pt}




\vspace{5mm}

\newcolumntype{u}{>{\centering\arraybackslash\hsize=.22\hsize}X}
\newcolumntype{s}{>{\hsize=.45\hsize}X}


%%%%%%%%%%%%%%%%%%%%%%%%%%%%%%%%
% INI creacio subsubsubsecion
%%%%%%%%%%%%%%%%%%%%%%%%%%%%%%%%
\usepackage{titlesec}

\let\oldparagraph\paragraph
\renewcommand{\paragraph}[1]{\oldparagraph{\hspace{30pt}#1}}


\titleclass{\subsubsubsection}{straight}[\subsection]

\newcounter{subsubsubsection}[subsubsection]
\renewcommand\thesubsubsubsection{\thesubsubsection.\arabic{subsubsubsection}}
\renewcommand\theparagraph{\thesubsubsubsection.\arabic{paragraph}} % optional; useful if paragraphs are to be numbered

\titleformat{\subsubsubsection}
  {\normalfont\normalsize\bfseries}{\thesubsubsubsection}{1em}{}
\titlespacing*{\subsubsubsection}
{0pt}{3.25ex plus 1ex minus .2ex}{1.5ex plus .2ex}


\makeatletter
\renewcommand\paragraph{\@startsection{paragraph}{5}{10pt}%
  {3.25ex \@plus1ex \@minus.2ex}%
  {-1em}%
  {\normalfont\normalsize\bfseries}}
\renewcommand\subparagraph{\@startsection{subparagraph}{6}{\parindent}%
  {3.25ex \@plus1ex \@minus .2ex}%
  {-1em}%
  {\normalfont\normalsize\bfseries}}
\def\toclevel@subsubsubsection{4}
\def\toclevel@paragraph{5}
\def\toclevel@paragraph{6}
\def\l@subsubsubsection{\@dottedtocline{4}{7em}{4em}}
\def\l@paragraph{\@dottedtocline{5}{10em}{5em}}
\def\l@subparagraph{\@dottedtocline{6}{14em}{6em}}
\makeatother

\setcounter{secnumdepth}{4}
\setcounter{tocdepth}{3}
%%%%%%%%%%%%%%%%%%%%%%%%%%%%%%%%
% FIN creacio subsubsubtitle
%%%%%%%%%%%%%%%%%%%%%%%%%%%%%%%%


%%%%%%%%%%%%%%%%%%%%%%%%%%%%%%%%
% INI creacio equation
%%%%%%%%%%%%%%%%%%%%%%%%%%%%%%%%
\DeclareFloatingEnvironment[fileext=loeq,placement={!ht},name=Equació]{equationfloat}
%%%%%%%%%%%%%%%%%%%%%%%%%%%%%%%%
% FIN creacio equation
%%%%%%%%%%%%%%%%%%%%%%%%%%%%%%%%

\titleformat{\section}[block]
  {\clearpage\normalfont\Large\bfseries}
  {\thesection}{1em}{}


\renewcommand\refname{Bibliografia}


\linespread{1.25}


\makeatletter
\renewenvironment{abstract}{%
    \if@twocolumn
      \section*{\abstractname}%
    \else 
      \begin{center}%
        {\bfseries \Large\itshape\abstractname\vspace{\z@}}%  % Título en negrita, 14pt y cursiva
      \end{center}%
      \quotation
    \fi}
    {\if@twocolumn\else\endquotation\fi}
\makeatother

% \fancypagestyle{main}{
%     \fancyhf{}
%     \lhead{\small Aplicació Web per a la Visualització i Comparació
% de Trajectòries d’Algorismes d’Optimització}
%     % \rhead{\includegraphics[width=1.2cm]{resources/Icons/logo-fib.png}}
%     \cfoot{\thepage}
%     \setlength{\footskip}{35pt}
% }

\newcommand{\q}[1]{``#1''}

\begin{document}

\hypersetup{
    pdfborder=0 0 0,
}

%% PORTADA I TABLE OF CONTENTS %%

%%%%                                   %%%%
%         INDEX DE LA MEMORIA             %
%%%%                                   %%%%

% \includepdf[pages=-]{portada_fib.pdf}
\input{Extra/GEP}
%\input{Extra/portada}
% \cleardoublepage

% \pagestyle{main}
% \setcounter{page}{11}
\include{toc}
% \cleardoublepage
\input{Contexto y alcance/main}

\input{Planificación temporal/main}

\input{Gestión económica/main}

\input{Informe de sostenibilidad/main}

% \cleardoublepage

% \cleardoublepage


% \input{Annex/main}


\clearpage
% Estilo bibliográfico IEEE
\bibliographystyle{IEEEtran}
\renewcommand{\refname}{Referencias}
\bibliography{references.bib}

%%%%                                   %%%%
%         FINAL DE LA MEMORIA             %
%%%%                                   %%%%

\end{document}


%% VERSIÓ INCIAL MEMORIA %%




\clearpage
% Estilo bibliográfico IEEE
\bibliographystyle{IEEEtran}
\renewcommand{\refname}{Referencias}
\bibliography{references.bib}

%%%%                                   %%%%
%         FINAL DE LA MEMORIA             %
%%%%                                   %%%%

\end{document}


%% VERSIÓ INCIAL MEMORIA %%




\clearpage
% Estilo bibliográfico IEEE
\bibliographystyle{IEEEtran}
\renewcommand{\refname}{Referencias}
\bibliography{references.bib}

%%%%                                   %%%%
%         FINAL DE LA MEMORIA             %
%%%%                                   %%%%

\end{document}


%% VERSIÓ INCIAL MEMORIA %%




\clearpage
\bibliographystyle{plainnat}
% \bibliography{references}


%%%%                                   %%%%
%         FINAL DE LA MEMORIA             %
%%%%                                   %%%%

\end{document}


%% VERSIÓ INCIAL MEMORIA %%

